\documentclass[11pt, a4paper]{ctexart}
\usepackage{../common/ncd-style}

\hypersetup{
  pdftitle={新概念丹麦语 5 · Tanke og Perspektiv},
  pdfauthor={Jaehaerys},
  pdfsubject={Ny Koncept Dansk 5}
}

\begin{document}

% =====================================================================
% Book Cover Page
% =====================================================================
\begin{titlepage}
  \centering
  \vspace*{2.5cm}
  {\color{NCDRed}\rule{\linewidth}{2.5pt}}\\[1.2cm]
  {\fontsize{34}{40}\selectfont\bfseries\color{NCDNavy} 新概念丹麦语 5}\\[0.5cm]
  {\huge\itshape\color{NCDTeal} Ny Koncept Dansk 5 · Tanke og Perspektiv}\\[0.8cm]
  {\large\bfseries\color{NCDGray} 精研篇（全 16 课完整版） \quad|\quad 对标 NCE 4 \& DU3 Modul 5 (PD3) $\rightarrow$ Modul 6 (Studieprøven, B2 $\rightarrow$ C1)}\\[1.2cm]
  {\color{NCDRed}\rule{\linewidth}{2.5pt}}\\[2cm]

  \begin{tcolorbox}[
    colback=NCDBgBlue,
    colframe=NCDNavy,
    width=0.90\linewidth,
    arc=4pt,
    boxrule=0.8pt,
    left=14pt, right=14pt, top=12pt, bottom=12pt
  ]
    \textbf{\color{NCDNavy}本册学习目标：}实现从“熟练公共辩论”到“数据评述、反讽识读、哲学伦理思辨与学术综述写作”的 B2 $\rightarrow$ C1 跨越（每课 450--650 词）。全册以 16 课高密度精读收官《新概念丹麦语》104 课体系，专攻 \textbf{PD3 / Studieprøven 统计图表三步分析法}、\textbf{丹麦式低调反讽（\textit{Underdrivelse / Litote}）与社论高阶让步}、\textbf{引述动词立场光谱（\textit{konstatere / hævde / påstå / slå fast}）}、\textbf{三层以上主从复合句与文化不可译词辨析（\textit{dannelse / tryghed / fællesskab}）} 以及 \textbf{多源观点综合论述（\textit{Syntese}）与双语哲思散文写作}。
  \end{tcolorbox}

  \vfill
  {\large\color{NCDGray} 编写体例：导读 $\rightarrow$ 精研双语篇与语音 $\rightarrow$ 屈折词汇表 $\rightarrow$ 高阶修辞与句法精讲 $\rightarrow$ 句型操练 $\rightarrow$ PD3 / Studieprøven 实战与详解}
\end{titlepage}

% Table of Contents
\tableofcontents
\clearpage

% =====================================================================
% Unit 1
% =====================================================================
\ncdunit{1}{Data, Tendenser og Analyse}{数据洞察与趋势分析（直击 PD3 图表）}{Lektion 1--4：统计图表全套高阶表达 · PD3 图表写作与口试黄金三步法 · 悖论与动态同步结构 · 社会学宏观分析语汇}

\ncdlesson{1}{Demografiens stille revolution: Når samfundet ældes}{人口结构的静默革命：当社会老龄化}
{DU3 Modul 5（PD3, B2）$\rightarrow$ Modul 6（Studieprøven, C1）｜对标 NCE 4 \textit{Statistisk analyse \& Demografisk udvikling}}
{\textbf{统计图表全套高阶表达}：占比与结构（\textit{udgøre en andel på}）· 递增与峰值（\textit{stige støt til, toppe i}）· 拐点与趋平（\textit{flade ud, aftage}）· 差距与超越（\textit{en markant forskel mellem A og B, overstige}）}
{欢迎来到《新概念丹麦语 5 · Tanke og Perspektiv》！前面四册带你完成了从日常交流到 B2- 公共辩论的积累，而第五册将正式打通从 B2 迈向 C1（PD3 与 Studieprøven）的学术思辨之路。在 PD3 写作与口试中，读懂统计图表只是第一步；真正拉开分数差距的，在于能否摆脱“数据上升了、下降了”的平面复述，转而运用精准的\textbf{占比、峰值、拐点与差距对比动词}，揭示数字背后的社会结构变迁。}

% ---------------------------------------------------------------------
\begin{foerduleaeser}
带着下面两个宏观问题阅读数据表、Tekst A 和 Tekst B：
\begin{enumerate}
  \item \textit{Hvornår overstiger andelen af borgere på 65 år og derover ifølge tabellen og Tekst A andelen af 0--17-årige, og hvordan udvikler væksten i ældregruppen sig efter 2035?}（根据数据表和 Tekst A，65 岁及以上人口占比将在何时首次超过 0--17 岁人群？2035 年之后老年人口的增长态势又会发生什么变化？）
  \item \textit{Hvorfor advarer cheføkonomen i Tekst B mod at betragte befolkningens aldring som en uundgåelig krise, og hvilke to politiske greb peger hun på?}（在 Tekst B 中，首席经济学家为什么警告不要把人口老龄化视为不可避免的危机？她指出了哪两项政策抓手？）
\end{enumerate}
\end{foerduleaeser}

% ---------------------------------------------------------------------
\ncdsubsec{二、Tekster（人口数据评论与政策分析双语篇）}

\subsection*{Tekst A：\textit{Danmarks alderspyramide vender på hovedet}（人口统计与经济分析专栏 · 约 370 词）}

\begin{tipbox}{语篇说明}
本篇围绕下方的丹麦人口结构演变与预测表展开分析。阅读时请紧盯两类高阶表达：\textbf{\color{NCDNavy}深蓝色标记的占比、递增与差距对比结构}（\textit{udgøre en andel på, stige støt til, en markant forskel mellem}），以及\textbf{\color{NCDRed}深红色标记的超越、峰值与趋平拐点动词}（\textit{overstige, toppe i, flade ud, aftage}）。
\end{tipbox}

\begin{tipbox}{数据图表参考：Danmarks befolkningsudvikling efter aldersgruppe (2005--2045, illustrativ prognose)}
\renewcommand{\arraystretch}{1.22}
{\small
\begin{tabularx}{\linewidth}{@{} >{\bfseries\raggedright\arraybackslash}p{4.8cm} >{\centering\arraybackslash}X >{\centering\arraybackslash}X >{\centering\arraybackslash}X >{\centering\arraybackslash}X >{\centering\arraybackslash}X @{}}
\toprule
Aldersgruppe / Nøgletal & 2005 & 2015 & 2025 & 2035 (prog.) & 2045 (prog.) \\
\midrule
0--17 år (andel af befolkningen) & 22,3 \% & 20,6 \% & 19,1 \% & 18,4 \% & 18,2 \% \\
18--64 år (erhvervsaktiv alder) & 62,6 \% & 60,8 \% & 59,7 \% & 57,5 \% & 56,6 \% \\
65+ år (andel af befolkningen) & 15,1 \% & 18,6 \% & 21,2 \% & 24,1 \% & 25,2 \% \\
80+ år (andel af befolkningen) & 4,1 \% & 4,3 \% & 5,6 \% & 7,4 \% & 8,5 \% \\
\midrule
Gennemsnitlig levealder (år) & 77,9 & 80,6 & 82,1 & 83,4 & 84,2 \\
\bottomrule
\end{tabularx}
}
\end{tipbox}

Når man betragter Danmarks demografiske udvikling fra 2005 frem mod 2045, træder et tydeligt mønster frem: Det danske samfund ældes i et tempo, som grundlæggende forandrer balancen mellem generationerne. Mens den gennemsnitlige levealder i perioden \da{stiger støt til} godt 84 år i 2045, falder fødselstallene samtidig til et historisk lavt niveau. Resultatet er ikke blot flere ældre borgere, men en gennemgribende forskydning af hele befolkningens alderssammensætning.

Et blik på fordelingen mellem unge og ældre illustrerer dette skift med stor tydelighed. I 2005 \da{udgjorde} børn og unge mellem 0 og 17 år \da{en andel på} 22,3 procent af befolkningen, mens borgere på 65 år og derover kun udgjorde 15,1 procent. Der var altså \da{en markant forskel mellem} de to aldersgrupper på 7,2 procentpoint. Tyve år senere er billedet imidlertid vendt på hovedet: I 2025 når ældregruppen op på 21,2 procent og \dabold{overstiger} dermed for første gang i nyere tid de yngste årgange, som samtidig er faldet til 19,1 procent.

Endnu mere opsigtsvækkende er udviklingen i væksttempoet frem mod midten af århundredet. Andelen af borgere på 65 år og derover vokser særlig kraftigt mellem 2015 og 2035, hvor tiårsvæksten \dabold{topper i} perioden frem mod 2035 med næsten tre procentpoint på blot et årti. Derefter begynder kurven gradvist at \dabold{flade ud}, så stigningen fra 2035 til 2045 blot udgør 1,1 procentpoint. Parallelle tendenser ses i den yngste gruppe, hvor faldet ligeledes \dabold{aftager} mærkbart og stort set stagnerer omkring 18,2 procent i 2045.

Bagsiden af denne udvikling rammer den erhvervsaktive del af befolkningen mellem 18 og 64 år, hvis andel reduceres fra 62,6 procent i 2005 til blot 56,6 procent i 2045. Samtidig mere end fordobles gruppen af borgere over 80 år fra 4,1 til 8,5 procent. For velfærdsstaten betyder det, at færre hænder i arbejde skal finansiere både pensioner, et mere specialiseret sundhedsvæsen og en langt mere ressourcekrævende ældrepleje.

Samlet set peger tallene på en langsigtet udfordring for Danmarks forsørgerbyrde. Selvom væksten blandt de 65-årige på sigt stabiliseres, fortsætter presset på både arbejdsmarkedet og de kommunale budgetter, fordi den ældste del af befolkningen bliver ved med at vokse. Dermed bliver spørgsmålet om, hvordan man fastholder høj beskæftigelse og sikrer kvaliteten i kernevelfærden, en af de mest afgørende politiske opgaver i de kommende årtier.

\begin{transbox}
\small
\textbf{\color{NCDTeal}[Tekst A 中文精译]}\\
当我们审视丹麦从 2005 年到 2045 年的人口结构演变时，一个清晰的规律显现出来：丹麦社会正以一种从根本上改变代际平衡的速度老龄化。在这一时期内，平均预期寿命稳步攀升至 2045 年的 84 岁以上，而出生率却同时跌至历史低位。其结果不仅是老年公民数量增多，更是整个人口年龄构成的深刻转移。

只要看一眼年轻群体与老年群体的分布比例，就能非常清楚地说明这一转变。2005 年，0 至 17 岁的儿童及青少年占总人口的 22.3\%，而 65 岁及以上公民仅占 15.1\%。也就是说，这两个年龄组之间存在着高达 7.2 个百分点的显著差距。然而二十年后，局面却完全倒转：到了 2025 年，老年群体占比达到 21.2\%，从而在近代历史上首次超过了同时期已降至 19.1\% 的最年轻群体。

更加引人注目的是直至本世纪中叶的增速演变趋势。65 岁及以上人口占比在 2015 年至 2035 年间增长尤为迅猛，其中十年增幅在迈向 2035 年的阶段达到峰值，短短十年内就攀升了近三个百分点。此后，曲线开始逐渐趋平，使得 2035 年至 2045 年间的增幅仅为 1.1 个百分点。类似的趋势也出现在最年轻群体中：其降幅同样明显放缓，到 2045 年基本上在 18.2\% 左右停滞下来。

这一趋势的另一面则冲击着 18 至 64 岁的劳动年龄人口——其占比从 2005 年的 62.6\% 缩减至 2045 年的仅仅 56.6\%。与此同时，80 岁以上群体占比翻了一番多，从 4.1\% 升至 8.5\%。对福利国家而言，这意味着更少在职劳动者必须共同支撑养老金支出、更加专科化的医疗体系以及对资源需求高得多的老年照护。

总体而言，这些数字指向了丹麦抚养比面临的一项长期挑战。尽管 65 岁以上群体的整体增速从长远来看会趋于平稳，但劳动力市场和各市镇财政预算承受的压力仍在持续，因为最高龄的那部分人口仍在不断壮大。由此，如何在未来数十年间保持高就业率并守住核心福利服务的质量，正成为最关键的政治课题之一。
\end{transbox}

\subsection*{Tekst B：\textit{Kommentar: Demografi er ikke skæbne --- det er et politisk valg}（政策分析短评 · 约 130 词）}

\begin{emailbox}{政策评论｜Af Sofie Lindberg, cheføkonom i Tænketanken Fremtid}
Tallene for Danmarks aldring skaber ofte dramatiske overskrifter om et kommende kollaps i velfærdsstaten. Men en højere levealder bør først og fremmest ses som en sundhedsmæssig succes. Demografi er nemlig ikke en fastlåst skæbne, hvis samfundets rammer tilpasses i tide.

Når prognoserne viser en stigende forsørgerbyrde, glemmer debatten tit, at dagens 68-årige generelt er langt sundere end tidligere generationer. Hvis flere seniorer får lyst og mulighed for at blive på arbejdsmarkedet gennem fleksible vilkår, kan en gradvist højere pensionsalder afbøde meget af presset på de offentlige finanser.

Samtidig må vi undgå, at politikkerne udelukkende handler om at få folk til at arbejde længere. Løsningen afhænger lige så meget af investeringer i ny teknologi og højere produktivitet i ældreplejen, så færre medarbejdere kan levere høj kvalitet uden at brænde ud.
\end{emailbox}

\textbf{Tekst B 阅读核对（Rigtigt / Forkert / Kort svar，答案见第六部分）：}
\begin{enumerate}
  \item \textit{Sofie mener, at det stigende antal ældre uundgåeligt vil føre til et velfærdskollaps.} \underline{\hspace{2.2cm}}
  \item \textit{Hvilken positiv forskel fremhæver Sofie hos dagens 68-årige sammenlignet med tidligere generationer?} \underline{\hspace{2.2cm}}
  \item \textit{Hvilke to ting ud over længere arbejdsliv bør samfundet ifølge Sofie investere i for at sikre ældreplejen?} \underline{\hspace{2.2cm}}
\end{enumerate}

\begin{tipbox}{本课发音与语调要点（Udtale \& Prosodi）}
\begin{itemize}
  \item \textbf{图表与人口学核心词重音与静音字母}：\textit{demografi} \ipa{demoɡʁɑˈfiːˀ}（重音在最末音节 \textit{-fi}，带喉塞音）；\textit{en andel} \ipa{ˈanˌd̥eːˀl}；\textit{at overstige} \ipa{ˈʌʊ̯ɐˌsd̥iːˀə}（\textit{g} 在词中不发音）；\textit{støt} \ipa{ˈsd̥øt}（短促元音 \ipa{ø}）。
  \item \textbf{复合词与硬读音提醒}：\textit{erhvervsaktiv} \ipa{ɛɐ̯ˈvɛɐ̯ˀvsˌɑɡtˢiʊ̯ˀ}（词首 \textit{er-} 不重读，\textit{h} 在 \textit{hv} 前不发音！）；\textit{en forsørgerbyrde} \ipa{fʌˈsɶɐ̯ˀɐˌbyɐ̯d̥ə}（重音在第二音节）；\textit{at flade ud} \ipa{ˈflæːðə ˈuðˀ}（含两个软 \textit{d}）。
  \item \textbf{丹麦语小数与百分点读法}：注意丹麦语小数点写作逗号，朗读时读作 \textit{komma}，如 \textit{21,2 procent} 读作 \textit{enogtyve komma to procent}；而区分差距时务必清晰重读 \textit{procentpoint} \ipa{pʰʁoˈsɛnˀdˌpʰɔɪ̯nd}（注意末尾 \textit{-t} 在 \textit{procent} 中往往同化弱化）。
\end{itemize}
\end{tipbox}

\begin{transbox}
\small
\textbf{\color{NCDTeal}[Tekst B 中文对照]}\\
有关丹麦人口老龄化的数据，往往会催生出诸如“福利国家即将崩溃”之类的夸张标题。但预期寿命的延长，首先应当被视作公共健康领域的一项成功。只要社会制度框架能够及时调整，人口结构就绝非注定无法改变的宿命。

当各项预测显示抚养比不断攀升时，公共辩论常常忽略了一个事实：今天的 68 岁老人普遍比以往任何几代人都更加健康。如果通过更灵活的工作条件让更多资深员工既有意愿又有机会留在劳动力市场，循序渐进地提高退休年龄便能够缓冲公共财政面临的大部分压力。

与此同时，我们必须避免将政策仅仅局限于让人们延长工作年限。解决方案同样取决于向养老照护领域的新技术与更高生产力进行投资，使人数更少的员工也能在避免职业倦怠的前提下提供高质量的服务。
\end{transbox}

% ---------------------------------------------------------------------
\ncdsubsec{三、Gloser（核心生词表）}

\begin{tipbox}{词表重点}
图表动词的完成时助动词区分极其关键：\da{at udgøre}（占据/构成）和 \da{at overstige}（超过）均为及物动词，完成时用 \textbf{har}（\textit{har udgjort / har oversteget}）；而表示曲线自身变化的不及物动词 \da{at aftage}（减弱）和 \da{at flade ud}（趋平）以及反身性过程动词 \da{at ældes}（变老），完成时一律用 \textbf{er}（\textit{er aftaget / er fladet ud / er ældet}）。此外，注意区分 \textit{en procent}（百分比）和 \textit{et procentpoint}（百分点，单复同形）。
\end{tipbox}

\begin{glosertable}
名词 & en andel & \textit{(andelen, andele, andelene)} & 份额，占比 \\
名词 & et procentpoint & \textit{(procentpointet, procentpoint, procentpointene)} & 百分点 \\
名词 & en forskel & \textit{(forskellen, forskelle, forskellene)} & 差异，差距 \\
名词 & en levealder & \textit{(levealderen, levealdre, levealdrene)} & 预期寿命，平均寿命 \\
名词 & en prognose & \textit{(prognosen, prognoser, prognoserne)} & 预测，趋势展望 \\
名词 & en aldersgruppe & \textit{(aldersgruppen, aldersgrupper, aldersgrupperne)} & 年龄组，年龄段 \\
名词 & en forsørgerbyrde & \textit{(forsørgerbyrden, forsørgerbyrder, forsørgerbyrderne)} & 抚养比，供养负担 \\
名词 & en tilbagetrækning & \textit{(tilbagetrækningen, tilbagetrækninger, tilbagetrækningerne)} & 退休，退出职场 \\
名词 & en produktivitet & \textit{(produktiviteten, kun ental)} & 生产力，劳动效率 \\
\midrule
动词 [har] & at udgøre & \textit{(udgør, udgjorde, har udgjort)} & 构成，占（……比例） \\
动词 [har] & at overstige & \textit{(overstiger, oversteg, har oversteget)} & 超过，超出（数值） \\
动词 [har] & at toppe & \textit{(topper, toppede, har toppet)} & 达到峰值，登顶 \\
动词 [har] & at afbøde & \textit{(afbøder, afbødede, har afbødet)} & 缓冲，减轻（冲击） \\
动词 [er] & at flade ud & \textit{(flader ud, fladede ud, er fladet ud)} & 趋平，增速放缓 \\
动词 [er] & at aftage & \textit{(aftager, aftog, er aftaget)} & 减弱，放缓，递减 \\
动词 [er] & at ældes & \textit{(ældes, ældedes, er ældet)} & 变老，老龄化 \\
\midrule
形容词 & markant & \textit{(markant, markante)} & 显著的，明显的 \\
形容词 & erhvervsaktiv & \textit{(erhvervsaktivt, erhvervsaktive)} & 处于工作年龄的，在职的 \\
\end{glosertable}

% ---------------------------------------------------------------------
\ncdsubsec{四、Grammatik \& Stilistik（图表描述核心语法与修辞精讲）}

\begin{grammatikbox}{语法点 1：PD3 / Studieprøven 统计图表四大核心表达模块}
在考场写作或口语面对柱状图、折线图与统计表格时，按“占比 $\rightarrow$ 趋势与峰值 $\rightarrow$ 拐点趋平 $\rightarrow$ 比较超越”四个模块调度精确动词，可立刻取代单调的 \textit{stiger / falder}：

\vspace{0.3em}
\renewcommand{\arraystretch}{1.24}
{\small
\begin{tabularx}{\linewidth}{@{} >{\bfseries\raggedright\arraybackslash}p{2.7cm} >{\raggedright\arraybackslash}p{4.8cm} >{\raggedright\arraybackslash}X @{}}
\toprule
\textbf{数据功能模块} & \textbf{核心高阶短语} & \textbf{标准书面示范句} \\
\midrule
\da{1. 占比与构成} &
\textit{udgøre en andel på X \%}\newline
\textit{svare til X \% af det samlede antal} &
\textit{I 2025 \da{udgør} de ældre \da{en andel på} 21,2 procent af befolkningen.} \\
\midrule
\da{2. 持续递增与峰值} &
\textit{stige støt til X \%}\newline
\textit{toppe i [år/periode] med X} &
\textit{Væksten \dabold{topper i} perioden frem mod 2035 med knap tre procentpoint.} \\
\midrule
\dabold{3. 拐点趋平与减速} &
\textit{flade ud omkring X \%}\newline
\textit{aftage mærkbart / stagnere} &
\textit{Efter 2035 begynder kurven gradvist at \dabold{flade ud}, og faldet \dabold{aftager}.} \\
\midrule
\dabold{4. 差距与超越} &
\textit{en markant forskel mellem A og B}\newline
\textit{overstige / mere end fordobles} &
\textit{Der er \da{en markant forskel mellem} de to grupper, før ældregruppen \dabold{overstiger} de unge.} \\
\bottomrule
\end{tabularx}
}
\end{grammatikbox}

\begin{grammatikbox}{语法点 2：图表描述四大考场高危排雷（中/英母语迁移纠偏）}
\begin{enumerate}
  \item \textbf{\textit{procent}（百分比）vs \textit{procentpoint}（百分点）绝对不能混用！}：从 15,1 \% 升至 24,1 \%，差额是 \da{9 procentpoint}（9 个百分点），绝不是“9 procent”（因为 24,1 比 15,1 相对增长了约 60\%，而非 9\%！）。表示绝对百分点差额时，务必使用 \textit{et fald / en stigning på X procentpoint}。
  \item \textbf{介词 \textit{på}（幅度和数值）vs \textit{til}（终点）vs \textit{med}（伴随幅度）}：
    \begin{itemize}
      \item \textit{Andelen steg \dabold{fra} 15 \% \dabold{til} 21 \%}（终点用 \textit{til}）。
      \item \textit{En stigning \dabold{på} 6 procentpoint}（名词后接变化幅度必须用 \textit{på}）。
      \item \textit{Andelen er steget \dabold{med} 6 procentpoint}（动词 \textit{stige / falde} 后接变化量常用 \textit{med}，或用 \textit{til} 接最终数值）。
    \end{itemize}
  \item \textbf{英语小数点陷阱与动词完成时助动词 \textit{er / har}}：丹麦语小数必须写\textbf{逗号}（\ding{51}\ \textit{21,2 \%}，\ding{55}\ \textit{*21.2 \%}）；变化类动词在完成时中用 \dabold{er}：\ding{51}\ \textit{Andelen \dabold{er steget} støt, mens kurven nu \dabold{er fladet ud}.}（\ding{55}\ \textit{*Andelen har steget støt.}）。
  \item \textbf{前置时间状语必触发 V2 倒装，且不加英语逗号}：\ding{51}\ \textit{Mellem 2035 og 2045 \dabold{flader} væksten gradvist ud.}（绝不能写 \ding{55}\ \textit{*Mellem 2035 og 2045, væksten flader ud.}）。
\end{enumerate}
\end{grammatikbox}

% ---------------------------------------------------------------------
\ncdsubsec{五、Mønsterøvelser（统计图表句型与语序操练）}

\textbf{练习 A：图表高阶词汇与介词精准填空（每空填入最准确的 1 个丹麦语词）}
\begin{enumerate}
  \item I 2005 udgjorde børn og unge en \underline{\hspace{2.2cm}} på 22,3 procent af hele befolkningen.
  \item Den gennemsnitlige levealder er steget \underline{\hspace{2.2cm}} til godt 84 år i prognosen.
  \item Andelen af ældre på 65+ vil i 2025 for første gang \underline{\hspace{2.2cm}} gruppen af 0--17-årige.
  \item Der er en markant forskel \underline{\hspace{2.2cm}} unge og ældre i den offentlige debat om pension.
  \item Efter ti år med hastig vækst begynder kurven endelig at \underline{\hspace{2.2cm}} ud efter 2035.
  \item Faldet i de yngste årgange \underline{\hspace{2.2cm}} mærkbart, når vi nærmer os år 2045.
  \item Stigningen fra 15,1 til 24,1 procent svarer til en fremgang på 9 \underline{\hspace{2.2cm}}.
  \item Væksttempoet i ældregruppen \underline{\hspace{2.2cm}} i årtiet frem mod 2035, før det stabiliseres.
\end{enumerate}

\textbf{练习 B：图表表达升级与典型考场病句修正（按括号要求改写或纠错）}
\begin{enumerate}
  \item 【纠错】修正助动词、英语小数点与变化介词：\textit{Fra 2005 til 2025 har andelen af ældre steget til 21.2 procent, hvilket er en stigning af 6,1 procent.} $\rightarrow$
  \item 【纠错】修正句首时间短语逗号与 V2 倒装：\textit{Mellem 2035 og 2045, væksten flader langsomt ud.} $\rightarrow$
  \item 【纠错】修正 \textit{mellem / af} 搭配错误：\textit{Der er en stor forskel imellem den offentlige og den private sektor på omkring tolv procent.}（请改写为精确书面对比句）$\rightarrow$
  \item 【升级】用 \textit{udgøre en andel på} 替换平淡表达：\textit{65-årige og ældre er 21,2 procent af alle danskere i 2025.} $\rightarrow$
  \item 【升级】用 \textit{overstige} 合并对比两项数据：\textit{I 2025 er der 21,2 procent ældre, men kun 19,1 procent børn og unge.} $\rightarrow$
  \item 【升级】用 \textit{toppe i ... hvorefter ... flade ud} 重组句子：\textit{Tilvæksten er størst i 2035, og bagefter bliver den næsten den samme.} $\rightarrow$
\end{enumerate}

\textbf{练习 C：中译丹（灵活运用图表占比、峰值、拐点与政策应对表达）}
\begin{enumerate}
  \item 2005 年，老年人口占总人口的 15.1\%。
  \item 两个年龄组之间存在着显著的差距。
  \item 2025 年，老年人口占比将超过儿童及青少年。
  \item 在 2035 年达到峰值后，该增长曲线将逐渐趋平。
  \item 随着平均寿命稳步上升，社会面临着越来越高的抚养比。
  \item 更高的退休年龄和新技术投资可以缓冲公共财政面临的压力。
\end{enumerate}

% ---------------------------------------------------------------------
\ncdsubsec{六、PD3 / Studieprøven 实战微任务与参考答案}

\begin{emailbox}{PD3 Skrivning / Mundtlig Del 实战微任务（Statistikanalyse · mindst 200--250 ord）}
\textbf{Situation}: Du skal til eksamen analysere tabellen fra Tekst A om Danmarks befolkningsudvikling fra 2005 til 2045.\\
\textbf{Opgave}: Skriv en sammenhængende analyse på \textbf{mindst 200--250 ord}, hvor du:
\begin{enumerate}
  \item Beskriver de væsentligste tendenser i tabellen ved hjælp af præcise udtryk for andel, stigning, højdepunkt, udfladning og forskelle mellem aldersgrupper.
  \item Forklarer kort, hvilke konsekvenser udviklingen kan få for velfærdsstaten og arbejdsmarkedet.
  \item Vurderer med inddragelse af Tekst B, hvordan samfundet bedst kan håndtere den demografiske forandring.
\end{enumerate}
\end{emailbox}

\begin{answerbox}
\textbf{Før du læser 参考回答}
\begin{enumerate}
  \item \textit{Ifølge tabellen og Tekst A overstiger andelen af borgere på 65 år og derover andelen af 0--17-årige for første gang i 2025 (21,2 \% mod 19,1 \%). Væksten i ældregruppen topper frem mod 2035, hvorefter kurven gradvist flader ud mellem 2035 og 2045.}
  \item \textit{Cheføkonomen advarer mod at se aldring som en uundgåelig krise, fordi højere levealder først og fremmest er en sundhedsmæssig succes. Hun peger på to politiske greb: fleksible vilkår med en gradvist højere pensionsalder for sunde seniorer samt investeringer i ny teknologi og højere produktivitet i ældreplejen.}
\end{enumerate}

\textbf{Tekst B 阅读核对答案}
\begin{enumerate}
  \item \da{Forkert}（她明确指出人口老龄化并非注定带来崩溃的宿命：\textit{Demografi er nemlig ikke en fastlåst skæbne}）
  \item \da{At dagens 68-årige generelt er langt sundere end tidligere generationer.}（今天的 68 岁人群普遍比前几代人健康得多）
  \item \da{Ny teknologi og højere produktivitet i ældreplejen.}（在老年照护领域投资新技术并提高劳动生产力）
\end{enumerate}

\textbf{练习 A 答案}
\begin{enumerate}
  \item \da{andel}（固定结构：\textit{udgøre en andel på...} 占……的比例）
  \item \da{støt}（副词搭配：\textit{stige støt til...} 稳步上升至……）
  \item \da{overstige}（助动词 \textit{vil} 后接动词原形 \textit{overstige}，表示“超过”）
  \item \da{mellem}（对比两组事物固定搭配：\textit{en forskel mellem A og B}）
  \item \da{flade}（固定短语动词：\textit{flade ud} 趋平、放缓）
  \item \da{aftager}（现在时：\textit{aftage} 减弱、缩小幅度）
  \item \da{procentpoint}（绝对百分点差额必须用单复同形的 \textit{procentpoint}）
  \item \da{topper}（现在时：\textit{toppe i [periode]} 在某阶段达到峰值）
\end{enumerate}

\textbf{练习 B 答案（参考）}
\begin{enumerate}
  \item \textit{Fra 2005 til 2025 er andelen af ældre steget til 21,2 procent, hvilket er en stigning på 6,1 procentpoint.}（\textit{stige} 用助动词 \textit{er}，小数用逗号，名词后幅度用 \textit{på} 且单位为 \textit{procentpoint}）
  \item \textit{Mellem 2035 og 2045 flader væksten langsomt ud.}（去掉句首状语后的英语式逗号，保持 V2 动词 \textit{flader} 紧跟第二位）
  \item \textit{Der er en markant forskel på omkring tolv procentpoint mellem den offentlige og den private sektor.}（使用 \textit{en markant forskel på X procentpoint mellem A og B} 使表述更严谨）
  \item \textit{I 2025 udgør borgere på 65 år og derover en andel på 21,2 procent af alle danskere.}（用 \textit{udgøre en andel på} 替换平淡系动词 \textit{er}）
  \item \textit{I 2025 overstiger andelen af ældre med 21,2 procent for første gang andelen af børn og unge på 19,1 procent.}（以动词 \textit{overstige} 一句话精准概括曲线交叉点）
  \item \textit{Tilvæksten topper i 2035, hvorefter kurven gradvist flader ud.}（利用 \textit{toppe i + hvorefter + flade ud} 形成高级句法衔接）
\end{enumerate}

\textbf{练习 C 答案}
\begin{enumerate}
  \item \textit{I 2005 udgjorde ældregruppen en andel på 15,1 procent af hele befolkningen.}
  \item \textit{Der findes en markant forskel mellem de to aldersgrupper.}
  \item \textit{I 2025 overstiger andelen af ældre antallet af børn og unge.}
  \item \textit{Efter at væksten topper i 2035, flader kurven gradvist ud.}
  \item \textit{Når levealderen stiger støt, står samfundet over for en voksende forsørgerbyrde.}
  \item \textit{En højere pensionsalder og investeringer i ny teknologi kan afbøde presset på økonomien.}
\end{enumerate}

\textbf{PD3 / Studieprøven 高分范文（约 240 词，不含标题）}\\
\textit{\textbf{Når balancen mellem generationer skifter: Udfordringer og løsninger}}

\textit{Tabellen over Danmarks befolkningsudvikling fra 2005 til 2045 afslører en markant demografisk forskydning mellem generationerne. Mens den gennemsnitlige levealder i perioden \da{stiger støt til} 84,2 år, falder andelen af børn og unge mellem 0 og 17 år fra 22,3 til 18,2 procent, hvorefter dette fald aftager mod slutningen af perioden. Samtidig sker der en afgørende krydsning i 2025: Her \da{udgør} borgere på 65 år og derover \da{en andel på} 21,2 procent og \dabold{overstiger} dermed for første gang de yngste årgange på 19,1 procent. Væksten i ældregruppen \dabold{topper i} tiåret frem mod 2035 med knap tre procentpoint, hvorefter stigningen gradvist \dabold{flader ud} til blot 1,1 procentpoint frem mod 2045. Desuden er der \da{en markant forskel mellem} 2005 og 2045 i den erhvervsaktive gruppe, som skrumper med seks procentpoint.}

\textit{Denne udvikling får væsentlige konsekvenser for både arbejdsmarkedet og velfærdsstaten. Når færre borgere befinder sig i den erhvervsaktive alder, mens andelen af 80+-årige mere end fordobles fra 4,1 til 8,5 procent, øges presset markant på sundhedsvæsenet, ældreplejen og de kommunale budgetter. En hastigt voksende forsørgerbyrde skaber samtidig risiko for akut mangel på hænder i kernevelfærden.}

\textit{Alligevel er der ingen grund til ren dommedagspessimisme. Som cheføkonomen fremhæver, afspejler stigende levealder først og fremmest et sundere og stærkere samfund. Hvis Danmark formår både at tilbyde fleksible arbejdsvilkår, der fastholder sunde seniorer længere i arbejde, og samtidig investere målrettet i ny velfærdsteknologi og højere produktivitet, kan samfundet \da{afbøde} de økonomiske udfordringer og sikre et holdbart velfærdssamfund fremover.}

\small\textbf{名师点评}：这篇约 240 词的范文严格契合 PD3 / Studieprøven 图表题评分标准。第一段精准调动 \textit{stiger støt til, aftager, udgør en andel på, overstiger, topper i, flader ud, en markant forskel mellem} 等全部核心词组，并准确使用了 \textit{procentpoint} 与逗号小数；第二段提炼“劳动力萎缩与照护压力”两大制度后果；第三段顺畅融入 Tekst B 视角完成辩证收束，结构平衡、语言凝练，稳居 C1/12 分水准。
\end{answerbox}

\ncdlesson{2}{Deltidsarbejde og kønsroller på det danske arbejdsmarked}{丹麦劳动力市场上的兼职工作与性别角色}
{DU3 Modul 5（PD3）核心模块｜对标 \textit{Skriftlig fremstilling (Diagramanalyse) \& Mundtlig eksamen Del 2}}
{\textbf{PD3 图表写作与口试黄金三步法}：\ding{172} 宏观概括数据趋势（\textit{overordnet set, kendetegnes ved}）$\rightarrow$ \ding{173} 剖析社会与制度成因（\textit{forklares delvist ved, hænger tæt sammen med}）$\rightarrow$ \ding{174} 评估未来影响与对策（\textit{på sigt indebærer det, modvirke ulighed}）}
{面对 PD3 笔试和口试的统计图表题，中国考生最容易犯的失误是“只报数字、没有结构”，把五分钟变成枯燥的读数流水账。真正的考场高分答案遵循经典的\textbf{“黄金三步法”}：先用一句话定调宏观走势并锁定核心反差，再从劳动力市场结构与文化惯性剖析多重成因，最后落脚到个人与社会的长远影响及政策对策。本课通过丹麦男女兼职就业率这一典型议题，带你完整演练从 B2 到 C1 的结构化分析框架。}

% ---------------------------------------------------------------------
\begin{foerduleaeser}
带着下面两个宏观问题阅读数据表、Tekst A 和 Tekst B：
\begin{enumerate}
  \item \textit{Hvordan anvender Tekst A tretrinsmodellen til både at sammenfatte udviklingen i deltidsarbejde fra 2010 til 2024 og forklare årsagerne bag forskellen mellem kønnene?}（Tekst A 是如何运用“三步法”，既概括 2010 至 2024 年丹麦兼职工作的整体走势，又剖析男女差异背后的多重成因的？）
  \item \textit{Hvorfor advarer fagforeningsrepræsentanten Mette i Tekst B mod at presse ansatte i plejesektoren over på fuldtid uden at forbedre arbejdsmiljøet først?}（在 Tekst B 中，工会代表梅特为什么警告不要在未改善工作环境前强推护理行业员工转全职？）
\end{enumerate}
\end{foerduleaeser}

% ---------------------------------------------------------------------
\ncdsubsec{二、Tekster（三步法图表范文与工会回应双语篇）}

\subsection*{Tekst A：\textit{Ligestilling i tal: Hvem arbejder på deltid i Danmark?}（PD3 三步法图表深度评论 · 约 355 词）}

\begin{tipbox}{语篇说明}
本篇以丹麦就业与性别统计数据为蓝本，完整示范 PD3 图表评论的“三步架构”。阅读时请留意\textbf{\color{NCDNavy}深蓝色标记的宏观概括与结构描述语}（Trin 1），以及\textbf{\color{NCDRed}深红色标记的因果剖析、长远后果与对策路标}（Trin 2 \& Trin 3）。
\end{tipbox}

\begin{tipbox}{数据图表参考：Andel beskæftigede på deltid efter køn og sektor (Danmark, \% af beskæftigede)}
\renewcommand{\arraystretch}{1.22}
{\small
\begin{tabularx}{\linewidth}{@{} >{\bfseries\raggedright\arraybackslash}p{5.5cm} >{\centering\arraybackslash}X >{\centering\arraybackslash}X >{\centering\arraybackslash}X >{\centering\arraybackslash}X @{}}
\toprule
Kategori & 2010 & 2018 & 2024 & Ændring (2010--2024) \\
\midrule
Kvinder i alt på deltid & 35,4 \% & 32,1 \% & 30,2 \% & -5,2 pct.point \\
Mænd i alt på deltid & 13,2 \% & 15,4 \% & 17,1 \% & +3,9 pct.point \\
Kønsgab (kvinder minus mænd) & 22,2 pct.point & 16,7 pct.point & 13,1 pct.point & -9,1 pct.point \\
\midrule
Offentlig pleje og omsorg (deltid) & 48,0 \% & 44,5 \% & 41,8 \% & -6,2 pct.point \\
Privat industri og IT (deltid) & 9,5 \% & 10,8 \% & 11,6 \% & +2,1 pct.point \\
\bottomrule
\end{tabularx}
}
\end{tipbox}

\da{Overordnet set kendetegnes} udviklingen i dansk deltidsarbejde fra 2010 til 2024 ved en gradvis udligning mellem kønnene, selvom der fortsat består en markant skævhed. Mens andelen af kvinder på deltid er faldet fra 35,4 til 30,2 procent, er mændenes andel steget fra 13,2 til 17,1 procent. Dermed er kønsgabet skrumpet fra 22,2 til 13,1 procentpoint. Samtidig \da{afspejler tallene en tydelig kontrast} mellem brancher: I den offentlige pleje- og omsorgssektor arbejder hele 41,8 procent på deltid i 2024, hvorimod tallet i privat industri og IT blot ligger på 11,6 procent.

Denne fordeling \dabold{forklares delvist ved} det stærkt kønsopdelte danske arbejdsmarked. Kvinder udgør fortsat det store flertal i kommunale velfærdsfag som sundhed, børnepasning og ældrepleje. Her planlægges vagtplaner ofte omkring 30 eller 32 timer om ugen, hvilket gør deltid til en indbygget norm i sektoren. Omvendt domineres industrien og IT-branchen stadig af mænd samt af en arbejdskultur, hvor fuldtidsstillinger tages for givet som standard.

Herudover \dabold{hænger uligheden tæt sammen med} kulturelle kønsroller og fordelingen af omsorgsopgaver i familielivet. Når småbørnsfamilier skal få hverdagen til at hænge sammen, vælger kvinder statistisk set oftere at gå ned i tid for at tage sig af børn og hjem. Ganske vist har nye regler om øremærket barsel skubbet flere fædre mod større deltagelse i familieansvaret, men traditionelle mønstre i arbejdslivet forandrer sig langt langsommere end lovgivningen.

\dabold{På længere sigt indebærer dette} betydelige konsekvenser for både den enkelte og samfundet som helhed. For mange kvinder fører mange års deltidsarbejde til en lavere livstidsindkomst, færre lederposter og en mærkbart mindre pensionsopsparing i alderdommen. For den danske velfærdsstat betyder de mange deltidsstillinger samtidig, at kommuner og hospitaler går glip af tiltrængte arbejdstimer midt i en periode med alvorlig mangel på arbejdskraft.

\dabold{For at modvirke denne skævhed} peger flere eksperter og arbejdsgivere på et behov for både bedre normeringer og en lovfæstet ret til fuldtid i det offentlige. Hvis flere fuldtidsstillinger kombineres med mere fleksible vagtplaner og en kultur, hvor også mænd i højere grad tager del i omsorgsopgaverne, kan udligningen fortsætte. Samtidig bør politikere huske, at frivillig deltid fortsat må være en legitim mulighed for borgere, der bevidst prioriterer fritid og trivsel over højere løn.

\begin{transbox}
\small
\textbf{\color{NCDTeal}[Tekst A 中文精译]}\\
总体而言，2010 年至 2024 年间丹麦兼职工作的发展趋势以两性之间的逐步趋同为特征，尽管显著的失衡依然存在。女性从事兼职的比例从 35.4\% 降至 30.2\%，而男性的比例则从 13.2\% 升至 17.1\%。由此，性别差距从 22.2 个百分点缩小到 13.1 个百分点。与此同时，数据也折射出行业之间的鲜明反差：2024 年公共护理与照料部门的兼职比例高达 41.8\%，而私营工业和信息技术行业这一数字仅为 11.6\%。

这一分布状况部分可以由高度按性别分流的丹麦劳动力市场来解释。女性在市政福利行业——如医疗、托幼和养老照护领域——仍然占据绝大多数。在这些行业里，排班表往往围绕每周 30 或 32 小时来设计，使得兼职成为行业内部固有的常态。相反，工业和 IT 行业仍由男性主导，并伴随着一种将全职岗位视为理所当然的职场文化。

除此之外，这种不平等还与文化性别角色以及家庭内部照料任务的分配紧密相连。当有年幼子女的家庭需要维持日常运转时，从统计上看女性仍然更频繁地选择减少工时来照顾孩子和家庭。诚然，关于专享育儿假的新规已经推动更多父亲更多地分担家庭责任，但职场生活中的传统模式变化得比法律条文慢得多。

从长远来看，这意味着对个人和整个社会都将产生重大后果。对许多女性而言，多年的兼职工作导致终身收入偏低、晋升管理岗位减少以及晚年养老金储蓄明显缩水。而对丹麦福利国家而言，大量兼职岗位同时意味着市镇机构和医院在劳动力严重短缺的时期白白错失了亟需的工作小时数。

为了扭转这种失衡，多位专家和雇主指出，公共部门既需要更充裕的人员配置标准，也需要落实法定转为全职的权利。如果能将更多全职岗位与更灵活的排班制度结合起来，并培育一种让男性同样更多参与照料事务的职场文化，这种差距缩小的趋势就能延续下去。与此同时，决策者应当记住，对于那些有意识地将休闲和身心福祉置于高薪之上的公民来说，自愿兼职依然必须保留为一个正当的选择。
\end{transbox}

\subsection*{Tekst B：\textit{Fagforeningssvar: ”Gør ikke fuldtid til en tvangstrøje”}（工会回应短评 · 约 140 词）}

\begin{emailbox}{工会立场回应｜Af Mette Holm, næstformand i Velfærdsforbundet}
Tallene for deltidsarbejde i plejesektoren vækker ofte bekymring blandt politikere, som ønsker, at alle ansatte hurtigt skal gå op på 37 timer. Det mål lyder enkelt på papiret, men i praksis risikerer man at overse årsagen til, at så mange sosu-assistenter og pædagoger faktisk vælger deltid.

For vores medlemmer handler deltid sjældent om manglende engagement. Tværtimod bruger mange en 30-timers arbejdsuge som et nødvendigt værn mod fysisk og psykisk nedslidning i en hverdag præget af højt arbejdspres, tunge løft og skiftende vagter. Hvis man blot presser medarbejdere over på fuldtid uden bedre normeringer, vil sygefraværet sandsynligvis stige markant.

Vi bakker fuldt op om retten til fuldtid for dem, der ønsker højere løn og pension. Men vejen frem går ikke gennem tvang --- den går gennem et bedre arbejdsmiljø, hvor en fuldtidsstilling faktisk er til at holde til gennem et helt arbejdsliv.
\end{emailbox}

\textbf{Tekst B 阅读核对（Rigtigt / Forkert / Kort svar，答案见第六部分）：}
\begin{enumerate}
  \item \textit{Mette mener, at ansatte i plejesektoren vælger deltid, fordi de mangler engagement i deres arbejde.} \underline{\hspace{2.2cm}}
  \item \textit{Hvad risikerer der at ske ifølge Mette, hvis man presser medarbejdere på fuldtid uden bedre normeringer?} \underline{\hspace{2.2cm}}
  \item \textit{Hvad skal der til ifølge Mette, før flere vælger fuldtid frivilligt?} \underline{\hspace{2.2cm}}
\end{enumerate}

\begin{tipbox}{本课发音与语调要点（Udtale \& Prosodi）}
\begin{itemize}
  \item \textbf{劳动市场与性别议题高频词}：\textit{et arbejdsmarked} \ipa{ˈɑːˌbɑɪ̯ˀd̥sˌmɑːɡəð}（复合词主重音在第一音节，中间 \textit{d} 弱化，词尾含软 \textit{d}）；\textit{beskæftiget} \ipa{beˈsɡɛfd̥iəð}（重音在第二音节）；\textit{kønsopdelt} \ipa{ˈkʰønˀsˌʌb̥d̥eːˀld}（含两处喉塞音）。
  \item \textbf{抽象动作与社会学术语}：\textit{en udligning} \ipa{ˈuðˌliːˀneŋ}（软 \textit{d} 接长音 \textit{i}，\textit{g} 在 \textit{-lign-} 中不发音！）；\textit{en nedslidning} \ipa{ˈneðˌsliðˀneŋ}（两个软 \textit{d}）；\textit{at kendetegne} \ipa{ˈkʰɛnəˌtˢɑɪ̯ˀnə}（\textit{eg} 读双元音 \ipa{ɑɪ̯ˀ}）。
  \item \textbf{PD3 口试三步路标语音停顿}：在口试叙述图表时，每切换一个步骤（\textit{Overordnet set... / Dette forklares ved... / På længere sigt...}），务必将句首路标词重读并做约半秒的清晰停顿，能立刻让考官听出你的逻辑框架。
\end{itemize}
\end{tipbox}

\begin{transbox}
\small
\textbf{\color{NCDTeal}[Tekst B 中文对照]}\\
护理行业的兼职统计数字常常令政治人物感到担忧，他们希望所有员工都能迅速提高到每周 37 小时。这一目标在纸面上听起来很简单，但在实践中人们却面临忽略一个核心原因的风险——那就是为何如此多的护理助理和幼教人员实际上会主动选择兼职。

对我们的工会会员来说，选择兼职极少是因为缺乏敬业精神。恰恰相反，许多人在高强度工作压力、繁重体力搬抬和频繁倒班的日常环境中，把每周 30 小时的工作制当作抵御身心劳损的必要防线。如果不在改善人员配比的情况下强行迫使员工转成全职，病假缺勤率很可能会显著上升。

我们完全支持让那些希望获得更高薪水和养老金的员工享有全职权利。但未来的出路绝不是依靠强制手段——而在于打造更好的工作环境，让一个全职岗位真正能够支撑人走完整个职业生涯。
\end{transbox}

% ---------------------------------------------------------------------
\ncdsubsec{三、Gloser（核心生词表）}

\begin{tipbox}{词表重点}
注意剖析原因时的动词方向性：\da{at kendetegne}（以……为特征，常用被动 \textit{kendetegnes ved}）和 \da{at forklare}（解释，常用被动 \textit{forklares ved / med}）在图表陈述中极高频；而 \da{at indebære}（意味着、蕴含后果）和 \da{at modvirke}（抵消、抑制负面趋势）则是 Step 3 评估影响与对策的核心高阶动词。注意 \textit{arbejdspres} 和 \textit{deltidsarbejde} 均为不可数中性名词（\textit{kun ental}）。
\end{tipbox}

\begin{glosertable}
名词 & et arbejdsmarked & \textit{(arbejdsmarkedet, arbejdsmarkeder, arbejdsmarkederne)} & 劳动力市场 \\
名词 & et deltidsarbejde & \textit{(deltidsarbejdet, kun ental)} & 兼职工作 \\
名词 & en udligning & \textit{(udligningen, udligninger, udligningerne)} & 趋同，缩小差距 \\
名词 & en skævhed & \textit{(skævheden, skævheder, skævhederne)} & 失衡，倾斜偏差 \\
名词 & en kønsrolle & \textit{(kønsrollen, kønsroller, kønsrollerne)} & 性别角色 \\
名词 & en vagtplan & \textit{(vagtplanen, vagtplaner, vagtplanerne)} & 排班表，值班表 \\
名词 & en pensionsopsparing & \textit{(pensionsopsparingen, pensionsopsparinger, pensionsopsparingerne)} & 养老金储蓄 \\
名词 & en normering & \textit{(normeringen, normeringer, normeringerne)} & 人员配置标准，编制 \\
名词 & et arbejdspres & \textit{(arbejdspresset, kun ental)} & 工作压力 \\
名词 & en nedslidning & \textit{(nedslidningen, nedslidninger, nedslidningerne)} & 身心劳损，职业损耗 \\
\midrule
动词 [har] & at kendetegne & \textit{(kendetegner, kendetegnede, har kendetegnet)} & 以……为特征，标志着 \\
动词 [har] & at indebære & \textit{(indebærer, indebar, har indebåret)} & 意味着，带来（后果） \\
动词 [har] & at modvirke & \textit{(modvirker, modvirkede, har modvirket)} & 抵消，抑制，扭转 \\
动词 [har] & at afspejle & \textit{(afspejler, afspejlede, har afspejlet)} & 反映，折射出 \\
动词 [har] & at tilskrive & \textit{(tilskriver, tilskrev, har tilskrevet)} & 把……归因于 \\
\midrule
形容词 & kønsopdelt & \textit{(kønsopdelt, kønsopdelte)} & 按性别分流的，性别隔离的 \\
形容词 & overordnet & \textit{(overordnet, overordnede)} & 宏观的，总体的 \\
形容词 & frivillig & \textit{(frivilligt, frivillige)} & 自愿的 \\
\end{glosertable}

% ---------------------------------------------------------------------
\ncdsubsec{四、Grammatik \& Stilistik（PD3 三步法骨架与因果剖析语法精讲）}

\begin{grammatikbox}{语法点 1：PD3 图表写作与口试“黄金三步法”全景模板}
无论是 PD3 笔试中的数据分析短文，还是口试 Del 2 面对考官抽到的统计图，只要严格按以下三步展开，就能确保内容饱满、逻辑不偏题且语域达到 B2 $\rightarrow$ C1：

\vspace{0.3em}
\renewcommand{\arraystretch}{1.24}
{\small
\begin{tabularx}{\linewidth}{@{} >{\bfseries\raggedright\arraybackslash}p{3.0cm} >{\raggedright\arraybackslash}p{4.6cm} >{\raggedright\arraybackslash}X @{}}
\toprule
\textbf{三步结构} & \textbf{核心论证任务} & \textbf{考场高分句型公式} \\
\midrule
\da{Trin 1：宏观概括}\newline （约占 30\% 篇幅） &
不急着念单个数字，先概括总趋势 + 点明最突出的行业或群体反差 &
\textit{\da{Overordnet set kendetegnes} udviklingen \da{ved}...}\newline
\textit{Tallene \da{afspejler en tydelig kontrast} mellem A og B.} \\
\midrule
\dabold{Trin 2：多维成因}\newline （约占 35\% 篇幅） &
从\textbf{制度/市场结构}与\textbf{文化/生活惯性}两个维度解释“为什么” &
\textit{Denne forskel \dabold{forklares delvist ved}...}\newline
\textit{Herudover \dabold{hænger} skævheden \dabold{tæt sammen med}...} \\
\midrule
\dabold{Trin 3：影响与对策}\newline （约占 35\% 篇幅） &
评估个人与社会长远代价，提出平衡且有条件的政策干预措施 &
\textit{\dabold{På længere sigt indebærer dette}, at...}\newline
\textit{\dabold{For at modvirke} denne ulighed bør man...} \\
\bottomrule
\end{tabularx}
}
\end{grammatikbox}

\begin{grammatikbox}{语法点 2：超越 \textit{fordi} 的四种 C1 归因动词与考场避坑铁律}
在 C1 学术和政论表达中，一直重复 \textit{Det er fordi...} 会显得语体单薄。学会用被动与中性因果动词能把论证推高一个台阶：
\begin{enumerate}
  \item \textbf{\da{forklares ved / med} vs \da{skyldes}}：
    \begin{itemize}
      \item \textit{Skævheden \dabold{forklares delvist ved} det kønsopdelte arbejdsmarked.}（部分可由……解释——注意主动是 \textit{forklare noget med/ved}）。
      \item \textit{Udviklingen \dabold{skyldes} både strukturelle og kulturelle forhold.}（源于/归因于——\textbf{铁律：\textit{skyldes} 自身已是被动形态，后面绝不能加 \textit{af}！}\ding{55}\ \textit{*Udviklingen skyldes af normeringer.}）。
    \end{itemize}
  \item \textbf{\da{kan tilskrives}（可归功/归因于）与 \da{hænger sammen med}（与……相关联）}：\textit{Tilbagegangen \dabold{kan tilskrives} et stigende arbejdspres.}（注意名词短语直接接在 \textit{tilskrives} 后，同样不加 \textit{af}）。
  \item \textbf{区分从句连词 \textit{fordi} 与介词短语 \textit{på grund af}}：\textit{fordi} 后面必须接“主语 + 动词”完整从句，而 \textit{på grund af} 后接名词短语：\ding{51}\ \textit{Mange vælger deltid \dabold{på grund af} højt arbejdspres.}（绝不写 \ding{55}\ \textit{*fordi højt arbejdspres}）。
  \item \textbf{句首评价副词 \textit{Overordnet set} 与目的不定式 \textit{For at modvirke...} 的 V2 倒装}：放在句首时，后面主句动词必须占据第二位：\ding{51}\ \textit{Overordnet set \dabold{viser} tallene en klar udligning.} / \textit{For at modvirke ulighed \dabold{bør} kommunen styrke normeringerne.}
\end{enumerate}
\end{grammatikbox}

% ---------------------------------------------------------------------
\ncdsubsec{五、Mønsterøvelser（三步法句型与归因操练）}

\textbf{练习 A：三步法标志词与因果介词精准填空（每空填入最准确的 1 个丹麦语词）}
\begin{enumerate}
  \item \underline{\hspace{2.2cm}} set viser statistikken fra 2010 til 2024 en tydelig udligning mellem kønnene.
  \item Udviklingen kendetegnes \underline{\hspace{2.2cm}} et fald i kvinders deltid og en stigning blandt mænd.
  \item Forskellen mellem brancher forklares delvist \underline{\hspace{2.2cm}} det stærkt kønsopdelte arbejdsmarked.
  \item Valget af deltid hænger tæt sammen \underline{\hspace{2.2cm}} både vagtplaner og kulturelle kønsroller.
  \item Den høje andel af deltid i plejen \underline{\hspace{2.2cm}} i høj grad et tungt arbejdspres i hverdagen. (skylde，现在被动形态)
  \item På længere sigt \underline{\hspace{2.2cm}} mange års deltid en lavere pensionsopsparing for den enkelte. (indebære，现在时)
  \item For at \underline{\hspace{2.2cm}} skævheden på arbejdsmarkedet kræves der bedre normeringer i daginstitutioner.
  \item Tallene \underline{\hspace{2.2cm}} en markant kontrast mellem offentlig omsorg og privat industri. (afspejle，现在时)
\end{enumerate}

\newpage
\textbf{练习 B：图表评论句型升级与语法纠错（按括号要求改写或纠错）}
\begin{enumerate}
  \item 【纠错】修正 \textit{skyldes} 的多余介词与英语式逗号：\textit{Overordnet set, forskellen skyldes af traditionelle kønsroller.} $\rightarrow$
  \item 【纠错】修正 \textit{fordi} 与名词短语误用：\textit{Mange pædagoger arbejder på deltid fordi det store arbejdspres.} $\rightarrow$
  \item 【纠错】修正目的状语前置时的主句 V2 倒装：\textit{For at modvirke nedslidning kommunerne må forbedre arbejdsmiljøet.} $\rightarrow$
  \item 【升级】用 \textit{kendetegnes ved} 升级 Step 1 宏观概括句：\textit{Fra 2010 til 2024 er udviklingen sådan, at kønsgabet bliver mindre.} $\rightarrow$
  \item 【升级】用 \textit{forklares delvist ved} 改写 Step 2 归因句：\textit{Det er også på grund af vagtplanerne i plejen, at så mange arbejder deltid.} $\rightarrow$
  \item 【升级】用 \textit{På længere sigt indebærer dette, at...} 改写 Step 3 后果句：\textit{Hvis folk arbejder deltid i mange år, får de mindre i pension senere.} $\rightarrow$
\end{enumerate}

\textbf{练习 C：中译丹（灵活运用 PD3 三步法核心连接词与词汇）}
\begin{enumerate}
  \item 总体而言，2024 年男女之间的差距已明显缩小。
  \item 这些数字清晰地折射出公共部门与私营部门之间的反差。
  \item 这一失衡部分可以用按性别分流的劳动力市场来解释。
  \item 选择兼职与家庭内部的照料任务分配紧密相连。
  \item 从长远来看，多年的兼职工作意味着更低的养老金储蓄。
  \item 为了扭转这种不平等，政府既需要改善人员配置，也需要创造更灵活的排班表。
\end{enumerate}

% ---------------------------------------------------------------------
\ncdsubsec{六、PD3 / Studieprøven 实战微任务与参考答案}

\begin{emailbox}{PD3 Skrivning / Mundtlig Del 实战微任务（Tretrinsanalyse · mindst 200--250 ord）}
\textbf{Situation}: Til PD3-eksamen får du udleveret tabellen fra Tekst A om deltidsarbejde efter køn og sektor fra 2010 til 2024.\\
\textbf{Opgave}: Skriv et struktureret analyseindlæg på \textbf{mindst 200--250 ord}, hvor du bruger \textbf{tretrinsmodellen}:
\begin{enumerate}
  \item \textbf{Trin 1 (Overblik)}: Sammenfat hovedtendenserne og den største kontrast mellem sektorerne.
  \item \textbf{Trin 2 (Årsager)}: Forklar både de strukturelle forhold og de kulturelle kønsroller bag tallene.
  \item \textbf{Trin 3 (Konsekvenser \& løsninger)}: Vurder de langsigtede konsekvenser og diskutér, hvordan man kan skabe flere fuldtidsstillinger uden at øge nedslidningen (inddrag gerne synspunktet fra Tekst B).
\end{enumerate}
\end{emailbox}

\begin{answerbox}
\textbf{Før du læser 参考回答}
\begin{enumerate}
  \item \textit{Tekst A bruger tretrinsmodellen ved først at beskrive den overordnede udligning af kønsgabet (fra 22,2 til 13,1 procentpoint) og kontrasten mellem pleje og industri, derefter forklare fordelingen via det kønsopdelte arbejdsmarked og omsorgsroller, og til sidst vurdere konsekvenserne for pension og velfærd samt pege på fleksibel fuldtid.}
  \item \textit{Mette advarer mod at presse ansatte over på fuldtid, fordi mange i plejesektoren vælger deltid som et værn mod fysisk og psykisk nedslidning. Uden bedre normeringer og et sundere arbejdsmiljø risikerer man blot højere sygefravær.}
\end{enumerate}

\textbf{Tekst B 阅读核对答案}
\begin{enumerate}
  \item \da{Forkert}（她强调选择兼职极少是因为缺乏敬业精神，而是为了抵御劳损）
  \item \da{At sygefraværet sandsynligvis vil stige markant.}（病假缺勤率很可能会大幅上升）
  \item \da{Et bedre arbejdsmiljø og bedre normeringer.}（更好的工作环境和更充裕的人员配置）
\end{enumerate}

\textbf{练习 A 答案}
\begin{enumerate}
  \item \da{Overordnet}（固定评价开篇短语：\textit{Overordnet set} 总体而言）
  \item \da{ved}（固定搭配：\textit{kendetegnes ved} 以……为特征）
  \item \da{ved} / \da{med}（固定被动归因搭配：\textit{forklares delvist ved / med}）
  \item \da{med}（固定动词短语：\textit{hænge tæt sammen med} 与……紧密相关）
  \item \da{skyldes}（注意该词本身即被动含义，直接接原因名词短语，不加 \textit{af}）
  \item \da{indebærer}（现在时动词：“意味着、带来后果”）
  \item \da{modvirke}（动词不定式：“抵消、扭转”负面趋势）
  \item \da{afspejler}（现在时动词：“反映、折射出”）
\end{enumerate}

\textbf{练习 B 答案（参考）}
\begin{enumerate}
  \item \textit{Overordnet set skyldes forskellen traditionelle kønsroller.}（去掉英语式逗号并遵循 V2 倒装，同时删去 \textit{skyldes} 后多余的介词 \textit{af}）
  \item \textit{Mange pædagoger arbejder på deltid på grund af det store arbejdspres.}（名词短语前必须用介词短语 \textit{på grund af}，而非连词 \textit{fordi}）
  \item \textit{For at modvirke nedslidning må kommunerne forbedre arbejdsmiljøet.}（目的不定式放句首时，主句发生 V2 倒装，情态动词 \textit{må} 置于主语 \textit{kommunerne} 前）
  \item \textit{Fra 2010 til 2024 kendetegnes udviklingen ved et gradvist mindre kønsgab.}（用 \textit{kendetegnes ved} 将口语化描述升格为学术概括句）
  \item \textit{Dette forklares delvist ved vagtplanerne i plejesektoren.}（用 C1 被动因果句式取代冗长的分裂句）
  \item \textit{På længere sigt indebærer mange års deltid en mindre pensionsopsparing.}（将 \textit{hvis} 条件句凝练为“时间视角 + \textit{indebære} + 抽象名词短语”）
\end{enumerate}

\textbf{练习 C 答案}
\begin{enumerate}
  \item \textit{Overordnet set er forskellen mellem kønnene mindsket tydeligt i 2024.}
  \item \textit{Tallene afspejler en klar kontrast mellem den offentlige og den private sektor.}
  \item \textit{Denne skævhed forklares delvist ved et kønsopdelt arbejdsmarked.}
  \item \textit{Valget af deltid hænger tæt sammen med fordelingen af omsorgsopgaver i familien.}
  \item \textit{På længere sigt indebærer mange års deltid en lavere pensionsopsparing.}
  \item \textit{For at modvirke uligheden kræves både bedre normeringer og mere fleksible vagtplaner.}
\end{enumerate}

\textbf{PD3 / Studieprøven 高分范文（约 235 词，不含标题）}\\
\textit{\textbf{Vejen mod reel ligestilling på arbejdsmarkedet kræver mere end fuldtid}}

\textit{\da{Overordnet set kendetegnes} udviklingen i dansk deltidsarbejde mellem 2010 og 2024 \da{ved} en langsom udligning mellem kønnene. Andelen af kvinder på deltid er faldet med 5,2 procentpoint fra 35,4 til 30,2 procent, mens mænds andel samtidig er steget fra 13,2 til 17,1 procent, hvilket har mindsket det samlede kønsgab til 13,1 procentpoint. Alligevel \da{afspejler tallene} fortsat en skarp kontrast mellem den offentlige plejesektor med hele 41,8 procent på deltid og privat industri og IT med blot 11,6 procent.}

\textit{Denne vedvarende skævhed \dabold{forklares delvist ved} vores kønsopdelte arbejdsmarked, hvor kommunale velfærdsfag som børnepasning og ældrepleje med skiftende vagtplaner primært beskæftiger kvinder. Her bliver deltidsstillinger ofte bygget ind i selve driften som standard. Desuden \dabold{hænger tendensen tæt sammen med} traditionelle kønsroller i familien, fordi kvinder statistisk set stadig oftere reducerer deres arbejdstid af hensyn til børn og hverdagsomsorg, selv efter indførelsen af øremærket barsel.}

\textit{\dabold{På længere sigt indebærer} udbredt deltid betydelige konsekvenser for både den enkelte og samfundet: Kvinder får lavere livstidsindkomst og en mærkbart mindre pensionsopsparing, mens velfærdsstaten mister tiltrængte arbejdstimer i en tid med personalemangel. \dabold{For at modvirke} denne ulighed bør kommunerne indføre en tydelig ret til fuldtid kombineret med mere fleksible vagtplaner. Samtidig har fagforeningen en afgørende pointe: Fuldtid må aldrig blive en tvang. Hvis flere ansatte i plejen skal vælge 37 timer frivilligt, kræver det først bedre normeringer og et sundere arbejdsmiljø, så ingen slides ned før pensionsalderen.}

\small\textbf{名师点评}：这篇约 235 词的回答堪称 PD3 三步法标杆：首段用 \textit{Overordnet set kendetegnes... ved} 与 \textit{afspejler} 抓出趋势与行业反差；中段以 \textit{forklares delvist ved} 和 \textit{hænger tæt sammen med} 切入结构与文化双重视角；末段结合 \textit{På længere sigt indebærer} 与 Tekst B 的工会关怀完成辩证对策，无一废句，直达考试高分带。
\end{answerbox}

\ncdlesson{3}{Urbaniseringens pris: Hvem har råd til at bo i byen?}{城市化的代价：谁住得起大城市？}
{DU3 Modul 5（PD3）$\rightarrow$ Modul 6（Studieprøven, B2 $\rightarrow$ C1）｜对标 NCE 4 \textit{Samfundsanalyse \& Boligpolitik}}
{\textbf{悖论与动态同步结构}：反讽/悖论导入（\textit{paradoksalt nok}）· 动态同步（\textit{i takt med at...}）· 正反馈螺旋（\textit{en selvforstærkende effekt}）· 比例递进（\textit{jo mere ..., desto ...}）· 代价取舍（\textit{på bekostning af}）}
{在谈论哥本哈根或奥胡斯的高房价时，B1 水平的表达往往停留在“房子太贵、工资不够”。而要写出打动 Studieprøven 考官的社会学深度，你必须学会描绘事物之间的\textbf{动态联动与制度悖论}：城市越宜居、越绿色，反而越把普通劳动者挤出市中心；通勤距离越远，社会分化就越剧烈。本课将带你掌握 \textbf{\textit{paradoksalt nok}}、\textbf{\textit{i takt med at}} 与 \textbf{\textit{jo ..., desto ...}} 的严密句法，精准刻画城市化背后的自我强化循环。}

% ---------------------------------------------------------------------
\begin{foerduleaeser}
带着下面两个宏观问题阅读 Tekst A 和 Tekst B：
\begin{enumerate}
  \item \textit{Hvilket paradoks opstår der ifølge Tekst A, når storbyer som København bliver grønnere og mere attraktive, og hvordan skaber boligmarkedet en selvforstærkende effekt?}（根据 Tekst A，当哥本哈根等大城市变得更加绿色宜居时，会产生什么悖论？房地产市场又是如何制造“自我强化效应”的？）
  \item \textit{Hvordan adskiller Henriks og Sofies forslag til at løse boligpresset sig fra hinanden i byrådsdebatten i Tekst B?}（在 Tekst B 的市政议会辩论中，亨里克和索菲针对缓解住房压力的方案有何分歧，又在何处达成妥协？）
\end{enumerate}
\end{foerduleaeser}

% ---------------------------------------------------------------------
\ncdsubsec{二、Tekster（城市社会学分析与市政对谈双语篇）}

\subsection*{Tekst A：\textit{Når attraktive byer prissætter deres egne borgere ud}（城市社会学分析专栏 · 约 345 词）}

\begin{tipbox}{语篇说明}
本篇剖析丹麦大城市人口集聚背后的住房危机与通勤不平等。阅读时请特别关注\textbf{\color{NCDNavy}深蓝色标记的动态同步与比例递进句型}（\textit{i takt med at..., jo mere..., desto...}），以及\textbf{\color{NCDRed}深红色标记的悖论、反馈循环与代价取舍表达}（\textit{paradoksalt nok, en selvforstærkende effekt, på bekostning af}）。
\end{tipbox}

I mere end to årtier har Danmark oplevet en markant urbanisering, hvor både studerende, børnefamilier og vidensvirksomheder søger mod København og Aarhus. Ved første øjekast ligner udviklingen en ubetinget succeshistorie med grønne cykelstier, rene havnebade, moderne metroforbindelser og et blomstrende kulturliv. Men \da{i takt med at} storbyerne er blevet stadig mere attraktive og klimavenlige, er boligpriserne stukket af fra helt almindelige lønninger, så adgangen til byen nu afhænger af formue frem for behov.

Her opstår der et slående bypolitisk dilemma. \dabold{Paradoksalt nok} er det netop de vellykkede investeringer i byrum, kollektiv trafik og grønt miljø, der nu truer byens sociale mangfoldighed. Når nedslidte industrikvarterer omdannes til eftertragtede boligområder langs havnen, stiger grundværdierne eksplosivt, og sygeplejersker, pædagoger samt politibetjente har simpelthen ikke længere råd til at bo tæt på de borgere og arbejdspladser, de servicerer i hverdagen.

På boligmarkedet sætter denne udvikling hurtigt \dabold{en selvforstærkende effekt} i gang. \da{Jo mere} eksklusivt et kvarter fremstår, \da{desto} større bliver interessen fra velstillede købere og private pensionskasser, hvilket igen presser kvadratmeterpriserne endnu højere op. Samtidig prioriterer kommuner og udviklere ofte dyre ejerboliger eller luksuriøse lejeboliger, som giver hurtige indtægter og stærke skatteydere, \dabold{på bekostning af} billige almene boliger med en husleje til at betale.

Konsekvensen mærkes langt uden for kommunegrænsen i form af voksende pendlerulighed. Unge familier og ansatte i servicefag fortrænges gradvist til oplandskommuner, hvor boligerne stadig er til at betale. \da{Jo længere} væk fra arbejdspladsen disse borgere presses, \da{desto} flere timer og kroner sluger den daglige transport gennem morgenens trafikpropper. Dermed opstår der ikke blot voksende økonomiske forskelle, men også en skæv fordeling af fritid, energi og familietid mellem centrum og opland.

Hvis storbyerne ikke skal ende som lukkede enklaver for de mest velhavende, kræves der en langt mere målrettet bolig- og planpolitik. En aktiv kommunal jordpolitik, faste kvoter for almene boliger ved nybyggeri samt hurtige togforbindelser til omegnen kan bremse segregeringen og sikre blandede kvarterer. Kun ved at tænke grøn vækst og social balance sammen kan byen bevare plads til alle de mennesker, der får dens skoler, hospitaler og gader til at fungere.

\begin{transbox}
\small
\textbf{\color{NCDTeal}[Tekst A 中文精译]}\\
二十多年来，丹麦经历了显著的城市化进程，大学生、有孩子的家庭以及知识密集型企业纷纷向哥本哈根和奥胡斯集聚。乍看之下，这一发展宛如一个毫无瑕疵的成功故事：绿意盎然的自行车道、洁净的港湾泳池、现代化的地铁线路以及繁荣的文化生活。然而，随着大城市变得愈发具有吸引力且更加气候友好，房价却早已远远甩开了普通人的薪资水平，以至于能否进入城市生活如今更多取决于资产而非实际需求。

一个引人注目的城市政策困境由此浮现。令人深感悖论的是，恰恰是那些对城市空间、公共交通和绿色生态环境的成功投资，如今正在威胁着城市原本的社会多样性。当破旧的老工业区沿着港湾被改造成备受追捧的住宅区时，土地价值呈爆发式飙升，护士、幼教老师以及警察等群体干脆再也负担不起住在自己日常所服务的市民与工作岗位附近。

在房地产市场上，这种演变迅速启动了一种自我强化效应。一个街区显得越高端稀缺，富裕买家和私营养老基金的兴趣就越浓厚，而这反过来又将每平方米房价推得更高。与此同时，市镇政府和开发商往往优先建设能带来快速收益和高收入纳税人的高价产权房或奢华租赁公寓，却以牺牲租金可负担的平价社会福利住房为代价。

这一后果以日益扩大的通勤不平等形式，在市镇边界之外的广大地区被深切感受到。年轻家庭和服务业从业者逐渐被挤向房价尚可负担的周边卫星市镇。这些市民被推离工作地点越远，穿过早高峰拥堵车流的日常通勤吞噬的时间和金钱就越多。由此产生的不仅是不断扩大的经济差距，还有市中心与外围地区之间在闲暇时间、精力与家庭陪伴上的失衡分配。

如果不想让大城市最终沦为只服务于最富裕阶层的封闭飞地，就必须推行更有针对性的住房与城市规划政策。积极的市政土地政策、新建项目中固定的社会福利住房配额，以及通往周边地区的快捷轨道交通，都能够遏制社会隔离并保障混合街区。唯有将绿色增长与社会平衡统筹考虑，城市才能为所有维系其学校、医院和街道运转的人留出容身之地。
\end{transbox}

\subsection*{Tekst B：\textit{Debat fra byrådssalen: Vækst og blandede bydele}（市政议会政策对谈 · 约 155 词）}

\textbf{阅读任务（理解与核对）：}\textit{Læs debatten mellem finansordfører Henrik og byplanlægger Sofie, og besvar spørgsmålene nedenfor.}

\begin{emailbox}{市政辩论对谈｜Henrik (Finansordfører) \& Sofie (Byplanlægger)}
\textbf{Henrik:} Vi ønsker jo alle en levende by, men jeg er bekymret for forslaget om at tvinge udviklere til at bygge 30 procent almene boliger i alle nye projekter. \da{Jo strammere} krav kommunen stiller til private investorer, \da{desto} færre boliger bliver der i sidste ende bygget --- og så stiger priserne bare endnu mere på grund af boligmangel.

\textbf{Sofie:} Den risiko skal vi selvfølgelig tage alvorligt, Henrik. Men \dabold{paradoksalt nok} hjælper et stort nybyggeri af dyre lejligheder ikke en sygeplejerske, hvis huslejen ligger over 16.000 kroner. Når vi bygger udelukkende efter markedets højeste afkast, sker det \dabold{på bekostning af} byens sammenhængskraft.

\textbf{Henrik:} Det anerkender jeg delvist. Måske ligger kompromiset i en hurtigere kommunal sagsbehandling, så vi udbygger både private og almene projekter i takt med befolkningstilvæksten --- især omkring nye S-togsstationer i forstæderne.

\textbf{Sofie:} Et rigtig godt bud! Kombinerer vi blandede ejerformer i centrum med bedre regional transport, kan vi nemlig bryde den \dabold{selvforstærkende effekt} uden at bremse byens udvikling.
\end{emailbox}

\textbf{Tekst B 阅读核对（Rigtigt / Forkert / Kort svar，答案见第六部分）：}
\begin{enumerate}
  \item \textit{Henrik frygter, at strenge krav om 30 procent almene boliger vil føre til færre nybyggede boliger i alt.} \underline{\hspace{2.2cm}}
  \item \textit{Hvorfor mener Sofie ikke, at mange nye, dyre lejligheder hjælper almindelige lønmodtagere som sygeplejersker?} \underline{\hspace{2.2cm}}
  \item \textit{Hvilket konkret kompromis foreslår Henrik for at sikre både flere boliger og hurtigere udbygning?} \underline{\hspace{2.2cm}}
\end{enumerate}

\begin{tipbox}{本课发音与语调要点（Udtale \& Prosodi）}
\begin{itemize}
  \item \textbf{城市化与住房政策词汇}：\textit{urbanisering} \ipa{uɐ̯bɑniˈseːˀɐeŋ}（重音在 \textit{-se-} 且带喉塞音，\textit{-ering} 读 \ipa{ˈeːˀɐeŋ}）；\textit{en kvadratmeterpris} \ipa{kvɑˈd̥ʁɑːd̥ˌmeːˀdɐˌpʰʁiːˀs}；\textit{almen} \ipa{ˈalˌmeːˀn}（常见复数 \textit{almene boliger} \ipa{ˈalˌmeːˀnə ˈb̥oːliɐ}，注意 \textit{bolig} 中 \textit{g} 不发音！）。
  \item \textbf{动态结构词的静音字母与 \textit{stød}}：\textit{paradoksalt} \ipa{pʰɑʁɑd̥ʌɡˈsæːˀld}（重音在最末音节）；\textit{selvforstærkende} \ipa{ˈsɛlˀfʌˌsd̥ɛɐ̯ɡənə}（\textit{selv} 末尾 \textit{v} 常省略不读！）；\textit{på bekostning af} \ipa{pʰɔ beˈkʰʌsd̥neŋ ˈɑʊ̯ˀ}（\textit{af} 读 \ipa{ɑʊ̯ˀ}）。
  \item \textbf{比例比较句语调起伏}：朗读 \textit{Jo mere X, desto mere Y} 时，前半句在 \textit{Jo mere} 处微扬，到逗号前略作悬停；后半句在 \textit{desto} 后的比较级重音陡然落地，凸显因果递增张力。
\end{itemize}
\end{tipbox}

\begin{transbox}
\small
\textbf{\color{NCDTeal}[Tekst B 中文对照]}\\
\textbf{亨里克（财政发言人）：}我们当然都希望拥有一座充满活力的城市，但我对在所有新项目中强制要求开发商建设 30\% 社会福利住房的提案感到担忧。市政府对私营投资者施加的要求越严苛，最终建成的住房总数就会越少——而这样一来，房价只会因为房源短缺而涨得更凶。

\textbf{索菲（城市规划师）：}我们当然应当认真对待这种风险，亨里克。但具有讽刺意味的悖论在于，如果月租超过一万六千克朗，哪怕新建再多高价公寓也帮不到一名护士。当我们完全追逐市场最高回报率来建房时，其代价往往就是城市的社会凝聚力。

\textbf{亨里克：}这一点我部分认可。也许折中方案在于加快市政审批流程，使我们能够随着人口增长同步推进私营项目与社会福利住房项目——尤其是在郊区的新城铁车站周边。

\textbf{索菲：}这是一个非常好的建议！只要我们把市中心的多元产权混合模式与更完善的区域公共交通结合起来，就能在不扼杀城市发展的前提下打破这种自我强化的恶性循环。
\end{transbox}

% ---------------------------------------------------------------------
\ncdsubsec{三、Gloser（核心生词表）}

\begin{tipbox}{词表重点}
丹麦社会特有的 \da{almen bolig}（复数常说 \textit{almene boliger}）并非中文语境下简陋的“救济房”，而是由非营利住房协会（\textit{boligselskaber}）运营、面向广大普通工薪阶层的平价福利住房，是理解丹麦城市政策辩论的关键概念！此外，动词 \da{at prissætte}（定价）与英文短语 \textit{price out} 对应为 \textit{at prissætte nogen ud}（因高价将某群体挤出）。
\end{tipbox}

\begin{glosertable}
名词 & en urbanisering & \textit{(urbaniseringen, kun ental)} & 城市化，都市化 \\
名词 & et boligmarked & \textit{(boligmarkedet, boligmarkeder, boligmarkederne)} & 房地产市场，住房市场 \\
名词 & en kvadratmeterpris & \textit{(kvadratmeterprisen, kvadratmeterpriser, kvadratmeterpriserne)} & 每平方米房价 \\
名词 & en mangfoldighed & \textit{(mangfoldigheden, kun ental)} & 多样性，多元共存 \\
名词 & en byfornyelse & \textit{(byfornyelsen, byfornyelser, byfornyelserne)} & 城市更新，旧区改造 \\
名词 & en oplandskommune & \textit{(oplandskommunen, oplandskommuner, oplandskommunerne)} & 周边市镇，外围卫星城 \\
名词 & en segregering & \textit{(segregeringen, segregeringer, segregeringerne)} & 社会隔离，阶层分区居住 \\
名词 & en ejerform & \textit{(ejerformen, ejerformer, ejerformerne)} & 住房产权形式 \\
名词 & en sagsbehandling & \textit{(sagsbehandlingen, sagsbehandlinger, sagsbehandlingerne)} & 行政审批，案件处理 \\
名词 & en effekt & \textit{(effekten, effekter, effekterne)} & 效应，效果 \\
\midrule
动词 [har] & at prissætte & \textit{(prissætter, prissatte, har prissat)} & 给……定价（\textit{ud} 挤出） \\
动词 [har] & at fortrænge & \textit{(fortrænger, fortrængte, har fortrængt)} & 排挤，取代，驱离 \\
动词 [har] & at presse & \textit{(presser, pressede, har presset)} & 施压，推高，挤迫 \\
动词 [har] & at udbygge & \textit{(udbygger, udbyggede, har udbygget)} & 扩建，拓展完善 \\
动词 [har] & at bremse & \textit{(bremser, bremsede, har bremset)} & 抑制，放缓，刹住 \\
\midrule
形容词 & eftertragtet & \textit{(eftertragtet, eftertragtede)} & 抢手的，备受追捧的 \\
形容词 & almen & \textit{(alment, almene)} & 公益性的，非营利福利性质的 \\
副词短语 & paradoksalt nok & \textit{(不变，句首评价副词)} & 令人悖论的是，具有讽刺意味的是 \\
\end{glosertable}

% ---------------------------------------------------------------------
\ncdsubsec{四、Grammatik \& Stilistik（动态同步、比例递进与悖论句法精讲）}

\begin{grammatikbox}{语法点 1：动态同步、比例递进与悖论修辞总览矩阵}
在 C1 级别的社会评论与城市分析文章中，静态列举原因往往不够立体；用好下面五组句式，才能让论述呈现“随时间同步演化、相互加速”的动态张力：

\vspace{0.3em}
\renewcommand{\arraystretch}{1.24}
{\small
\begin{tabularx}{\linewidth}{@{} >{\bfseries\raggedright\arraybackslash}p{2.9cm} >{\raggedright\arraybackslash}p{4.6cm} >{\raggedright\arraybackslash}X @{}}
\toprule
\textbf{修辞与逻辑功能} & \textbf{句型结构公式} & \textbf{高阶社会评论例句} \\
\midrule
\dabold{1. 悖论揭示} &
\textit{\dabold{Paradoksalt nok} + V2 倒装} &
\textit{\dabold{Paradoksalt nok} fører byfornyelse ofte til mindre social mangfoldighed.} \\
\midrule
\da{2. 动态同步推进} &
\textit{\da{I takt med at} [主+动]..., [动+主]...}\newline
（或 \textit{i takt med} + 名词） &
\textit{\da{I takt med at} byen bliver grønnere, stiger boligpriserne hurtigt.} \\
\midrule
\da{3. 比例同步递进} &
\textit{\da{Jo} + 比较级 [从句语序], \da{desto} + 比较级 [V2 倒装]} &
\textit{\da{Jo mere} attraktiv bydelen bliver, \da{desto} højere presses huslejen op.} \\
\midrule
\dabold{4. 正反馈恶性循环} &
\textit{sætte \dabold{en selvforstærkende effekt} i gang / bryde...} &
\textit{Høje priser sætter \dabold{en selvforstærkende effekt} i gang på boligmarkedet.} \\
\midrule
\dabold{5. 代价与利益取舍} &
\textit{\dabold{på bekostning af} + 名词短语} &
\textit{Luksusbyggeri sker ofte \dabold{på bekostning af} billige almene boliger.} \\
\bottomrule
\end{tabularx}
}
\end{grammatikbox}

\begin{grammatikbox}{语法点 2：\textit{Jo ..., desto ...} 语序铁律与三大考场陷阱}
\begin{enumerate}
  \item \textbf{\da{Jo ..., desto ...} 的核心秘诀：“前从句不倒装，后主句必倒装”！}
    \begin{itemize}
      \item 前半句由 \textit{jo} 引导，属于\textbf{从属从句}，内部保持“主语 +（\textit{ikke}）+ 动词”正常语序；当“\textit{jo + 比较级 + 名词}”本身充当从句主语时，后面必须插入形式主语占位词 \dabold{der}：\\
      \ding{51}\ \textit{Jo flere familier \dabold{der} flytter ud, \da{desto} længere bliver pendlertiden.}
      \item 后半句由 \textit{desto} 引导，属于\textbf{主句}，而“\textit{desto + 比较级（+ 被修饰名词）}”整个词组占领句首第一顺位，因此后面\textbf{必须立刻接变位动词}（V2 倒装）：\\
      \ding{51}\ \textit{Jo dyrere boligerne bliver, \da{desto sværere får} unge lønmodtagere ved at blive.}（绝不可写成 \ding{55}\ \textit{*...desto unge får sværere ved...}）。
    \end{itemize}
  \item \textbf{\textit{Paradoksalt nok} 与 \textit{I takt med at...} 不加逗号，且强触发 V2 倒装}：\\
    \ding{51}\ \textit{Paradoksalt nok \dabold{har} sygeplejersker ikke råd til at bo tæt på hospitalet.}（英语常写 \textit{Paradoxically, ...}，但丹麦语\textbf{绝不能在句首副词后打逗号}！）。
  \item \textbf{\textit{på bekostning af} 的固定形态}：表示“以……为代价”，永远用单数无冠词形式 \textit{på bekostning af} + 名词短语；不要误加定冠词写成 \ding{55}\ \textit{*på bekostningen af} 或改用介词 \ding{55}\ \textit{*på bekostning for}。
\end{enumerate}
\end{grammatikbox}

% ---------------------------------------------------------------------
\ncdsubsec{五、Mønsterøvelser（悖论与动态结构操练）}

\textbf{练习 A：动态结构与核心搭配精准填空（每空填入最准确的 1 个丹麦语词）}
\begin{enumerate}
  \item I \underline{\hspace{2.2cm}} med at storbyen tiltrækker flere studerende, vokser presset på boligmarkedet.
  \item \underline{\hspace{2.2cm}} nok fører grøn byfornyelse ofte til, at almindelige familier presses ud af kvarteret.
  \item Jo højere kvadratmeterprisen stiger, \underline{\hspace{2.2cm}} sværere bliver det for unge at købe lejlighed.
  \item Jo flere pendlere \underline{\hspace{2.2cm}} bosætter sig i oplandet, desto større bliver trængslen på motorvejen.
  \item Eksklusive byggerier prioriteres desværre ofte på \underline{\hspace{2.2cm}} af billige almene boliger.
  \item Når velstillede købere dominerer et område, sætter det en \underline{\hspace{2.2cm}} effekt i gang.
  \item Kommunen ønsker at \underline{\hspace{2.2cm}} segregeringen gennem faste kvoter for blandede ejerformer.
  \item Hurtigere kommunal \underline{\hspace{2.2cm}} kan gøre det lettere at bygge flere betalelige boliger i tide.
\end{enumerate}

\textbf{练习 B：动态比例句语序纠错与句型升级（修正语病或按要求重构）}
\begin{enumerate}
  \item 【纠错】修正英语逗号与 V2 倒装错误：\textit{Paradoksalt nok, manglen på billige boliger bliver værre i grønne bydele.} $\rightarrow$
  \item 【纠错】修正 \textit{jo...desto...} 的主从句语序与占位词：\textit{Jo flere dyre lejligheder bygges i centrum, desto færre lærere kan bo nær skolen.} $\rightarrow$
  \item 【纠错】修正介词短语错用：\textit{Byen vokser hurtigt på bekostningen for den sociale mangfoldighed.} $\rightarrow$
  \item 【升级】用 \textit{I takt med at...} 合并两句：\textit{Byerne investerer i nye metrostationer. Samtidig stiger huslejen hurtigt.} $\rightarrow$
  \item 【升级】用 \textit{Jo ..., desto ...} 重写因果递进句：\textit{Hvis borgere flytter længere væk fra jobbet, bruger de mere tid på daglig transport.} $\rightarrow$
  \item 【升级】用 \textit{sætte en selvforstærkende effekt i gang} 浓缩说明：\textit{Når et kvarter bliver populært, stiger priserne, og det gør kvarteret endnu mere eksklusivt.} $\rightarrow$
\end{enumerate}

\textbf{练习 C：中译丹（活用悖论、同步演化与比例递进句式）}
\begin{enumerate}
  \item 随着大城市变得越来越宜居，普通家庭往往负担不起市中心的房租。
  \item 令人悖论的是，成功的老旧街区改造有时会削弱城市的社会多样性。
  \item 市中心建得越多高价公寓，外围市镇承受的通勤压力就越大。
  \item 奢华住房项目常常是以牺牲平价福利住房为代价推进的。
  \item 当富裕买家接管一个街区时，就会启动一种自我强化效应。
  \item 城市规划必须将经济增长与公共交通结合起来，才能遏制社会隔离。
\end{enumerate}

% ---------------------------------------------------------------------
\ncdsubsec{六、PD3 / Studieprøven 实战微任务与参考答案}

\begin{emailbox}{Studieprøven / PD3 Debatindlæg 实战微任务（Samfundsanalyse · mindst 220--250 ord）}
\textbf{Situation}: I en avisdebat om boligpolitikken i de danske storbyer skriver en økonom, at markedet alene bør bestemme, hvad der bygges, fordi flere nye boliger automatisk vil sænke priserne for alle.\\
\textbf{Opgave}: Skriv et analytisk debatindlæg på \textbf{mindst 220--250 ord}, hvor du diskuterer urbaniseringens pris i storbyerne. Du skal:
\begin{enumerate}
  \item Analysere paradokset omkring stigende attraktivitet og sociale udfordringer i centrum.
  \item Forklare den selvforstærkende effekt og konsekvenserne for pendlerne i oplandskommunerne.
  \item Tage stilling til økonomens påstand og foreslå konkrete politiske tiltag, der kan sikre blandede bydele (brug gerne \textit{paradoksalt nok, i takt med at, jo ..., desto ...} og \textit{på bekostning af}).
\end{enumerate}
\end{emailbox}

\begin{answerbox}
\textbf{Før du læser 参考回答}
\begin{enumerate}
  \item \textit{Paradokset består i, at netop de vellykkede grønne og kulturelle investeringer får grundværdierne til at stige så kraftigt, at pædagoger og sygeplejersker prissættes ud af byen. Boligmarkedet skaber en selvforstærkende effekt ved, at eksklusive kvarterer tiltrækker flere velstillede købere og investorer, hvilket presser kvadratmeterpriserne endnu højere op.}
  \item \textit{Henrik frygter, at strenge kvoter på 30 procent almene boliger vil bremse privat nybyggeri og skabe boligmangel, mens Sofie mener, at dyrt byggeri alene ikke hjælper almindelige lønmodtagere. Til sidst mødes de om hurtigere sagsbehandling, stationsnær udbygning og blandede ejerformer.}
\end{enumerate}

\textbf{Tekst B 阅读核对答案}
\begin{enumerate}
  \item \da{Rigtigt}（他明确担忧强制配额会吓退私人投资者，导致建房总数减少）
  \item \da{Fordi huslejen ligger over 16.000 kroner, hvilket almindelige lønmodtagere ikke kan betale.}
  \item \da{En hurtigere kommunal sagsbehandling kombineret med udbygning omkring nye S-togsstationer i forstæderne.}
\end{enumerate}

\textbf{练习 A 答案}
\begin{enumerate}
  \item \da{takt}（固定动态同步连接短语：\textit{i takt med at...} 随着……同步发展）
  \item \da{Paradoksalt}（句首评价副词短语：\textit{Paradoksalt nok} 令人悖论的是）
  \item \da{desto}（比例比较结构配对：\textit{jo ..., desto ...}）
  \item \da{der}（当 \textit{jo flere + 名词} 在从句中作主语时，必须接占位关系词 \textit{der}）
  \item \da{bekostning}（固定搭配：\textit{på bekostning af} 以牺牲……为代价）
  \item \da{selvforstærkende}（固定学术搭配：\textit{en selvforstærkende effekt} 自我强化效应）
  \item \da{bremse} / \da{modvirke}（动词不定式：遏制/抑制社会隔离）
  \item \da{sagsbehandling}（名词：市政行政审批流程）
\end{enumerate}

\textbf{练习 B 答案（参考）}
\begin{enumerate}
  \item \textit{Paradoksalt nok bliver manglen på billige boliger værre i grønne bydele.}（删去英语式逗号，并保证变位动词 \textit{bliver} 紧跟在 \textit{Paradoksalt nok} 之后）
  \item \textit{Jo flere dyre lejligheder der bygges i centrum, desto færre lærere kan bo nær skolen.}（前句加入主语占位词 \textit{der}；后句将 \textit{desto færre lærere} 视为句首成分，紧接倒装动词 \textit{kan}）
  \item \textit{Byen vokser hurtigt på bekostning af den sociale mangfoldighed.}（固定为单数无冠词 \textit{på bekostning af}）
  \item \textit{I takt med at byerne investerer i nye metrostationer, stiger huslejen hurtigt.}（以 \textit{I takt med at...} 体现发展同步性，并注意主句 V2 倒装 \textit{stiger huslejen}）
  \item \textit{Jo længere væk fra jobbet borgerne flytter, desto mere tid bruger de på daglig transport.}（严格遵守“前从句不倒装、后主句 \textit{bruger de} 倒装”）
  \item \textit{Når et kvarter bliver populært, sætter prisstigningerne en selvforstærkende effekt i gang.}（用核心词组精炼概括正反馈螺旋）
\end{enumerate}

\textbf{练习 C 答案}
\begin{enumerate}
  \item \textit{I takt med at storbyerne bliver mere attraktive, har almindelige familier ikke råd til huslejen i centrum.}
  \item \textit{Paradoksalt nok kan en vellykket byfornyelse svække byens sociale mangfoldighed.}
  \item \textit{Jo flere dyre lejligheder der bygges i centrum, desto større bliver pendlerpresset i oplandskommunerne.}
  \item \textit{Luksusprojekter gennemføres ofte på bekostning af billige almene boliger.}
  \item \textit{Når velstillede købere overtager et kvarter, sættes en selvforstærkende effekt i gang.}
  \item \textit{Byplanlægningen skal kombinere vækst med offentlig transport for at bremse segregeringen.}
\end{enumerate}

\textbf{PD3 / Studieprøven 高分范文（约 240 词，不含标题）}\\
\textit{\textbf{Markedet alene skaber ikke en by for alle}}

\textit{Det lyder umiddelbart besnærende, når en økonom hævder, at et helt frit boligmarked automatisk vil dæmpe prispresset gennem øget nybyggeri. Erfaringerne fra København og Aarhus tegner imidlertid et langt mere nuanceret og bekymrende billede. \da{I takt med at} storbyerne gennem de seneste årtier er blevet grønnere, tryggere og kulturelt mere levende, er kvadratmeterpriserne steget langt hurtigere end almindelige lønninger.}

\textit{\dabold{Paradoksalt nok} risikerer selve byens succes dermed at undergrave dens sociale mangfoldighed. Når private investorer udelukkende bygger efter det højeste afkast, opføres der primært dyre ejer- og lejeboliger \dabold{på bekostning af} billige almene boliger. Det sætter \dabold{en selvforstærkende effekt} i gang på boligmarkedet: \da{Jo mere} eksklusivt et kvarter bliver, \da{desto} færre pædagoger, sygeplejersker og politibetjente har råd til at blive boende tæt på deres arbejdsplads og de borgere, de servicerer.}

\textit{Resultatet er en stigende pendlerulighed mellem bykernen og omegnen. \da{Jo længere} væk fra arbejdspladsen almindelige lønmodtagere fortrænges, \da{desto} flere timer og ressourcer mister familierne på daglig transport i myldretiden. Et rent frit marked løser altså ikke denne skæve udvikling, fordi efterspørgslen efter attraktive adresser bliver ved med at presse huslejen opad.}

\textit{Derfor har kommunerne brug for en aktiv og langsigtet planpolitik. Ved at kombinere hurtigere sagsbehandling med klare krav til blandede ejerformer, faste kvoter for almene boliger og stærke S-togsforbindelser til oplandskommunerne kan byerne både rumme fortsat vækst og sikre, at almindelige borgere også i fremtiden har råd til at bo i byen.}

\small\textbf{名师点评}：这篇约 240 词的评论文章精准回应了经济学家观点并剖析了城市化代价。全篇高频且无误地驾驭了本课五大支柱：\textit{I takt med at...}（同步推移）、\textit{Paradoksalt nok}（首句倒装揭示悖论）、\textit{på bekostning af}（利益权衡）、\textit{en selvforstærkende effekt}（机制提炼）以及两处高难度的 \textit{Jo ..., desto ...} 句式，兼具社会学洞察与 C1 句法掌控力。
\end{answerbox}

\ncdlesson{4}{Fritidsvaner i forandring: Fra foreningsliv til individuel træning}{变迁中的休闲：从社团生活到个人健身}
{DU3 Modul 5（PD3）$\rightarrow$ Modul 6（Studieprøven, C1）｜Unit 1 社会学宏观视野总汇演}
{\textbf{社会学抽象语汇与现象升华}：核心社科概念（\textit{individualisering, sammenhængskraft, kommercialisering, uforpligtende fællesskaber, institutionaliseret}）· 由表及里的结构推论（\textit{set i et bredere perspektiv, vidne om, i lyset af, et skift fra ... mod ...}）}
{在 Studieprøven 写作和深度访谈中，仅仅停留在“现在去健身房的人多了、参加传统俱乐部的人少了”这种具体生活观察，很难拿到 C1 高分。考官期待看到的是\textbf{“社会学升华能力”}：把日常习惯的变迁抽象为个人主义化、商业化、低约束社群与社会凝聚力之间的张力，并将其与前三课学过的人口老龄化、双职工时间挤压和城市化通勤连接起来。本课是 Unit 1 的收官综合演练！}

% ---------------------------------------------------------------------
\begin{foerduleaeser}
带着下面两个宏观问题阅读 Tekst A 和 Tekst B：
\begin{enumerate}
  \item \textit{Hvordan fortolker Tekst A overgangen fra traditionelt foreningsliv til kommercielle fitnesscentre, og hvordan hænger dette skift sammen med tendenserne fra lektion 1--3?}（Tekst A 如何解读从传统社团生活向商业健身房的过渡？这一转变又如何与第 1 至第 3 课探讨的宏观趋势相互关联？）
  \item \textit{Hvorfor mener kultursociologen Kasper i Tekst B ikke, at unge danskere har mistet interessen for fællesskab, selvom de vælger det traditionelle foreningsliv fra?}（在 Tekst B 中，为什么文化社会学家卡斯珀认为丹麦年轻人即便放弃传统社团，也并没有丧失对群体归属感的兴趣？）
\end{enumerate}
\end{foerduleaeser}

% ---------------------------------------------------------------------
\ncdsubsec{二、Tekster（社会学趋势分析与学界回应双语篇）}

\subsection*{Tekst A：\textit{Medlemskort eller fitness-app? Det danske foreningsliv under pres}（社会学文化分析专栏 · 约 360 词）}

\begin{tipbox}{语篇说明}
本篇以丹麦人休闲与运动习惯的变迁为切口，示范如何从日常细节提炼宏观社会学判断。阅读时请重点标出\textbf{\color{NCDNavy}深蓝色标记的社会学核心抽象语汇}（\textit{individualisering, sammenhængskraft, kommercialisering, uforpligtende fællesskaber, institutionaliseret}），以及\textbf{\color{NCDRed}深红色标记的宏观推论与单元串联路标}（\textit{set i et bredere perspektiv, i lyset af, vidner om}）。
\end{tipbox}

Hver aften året rundt træder tusindvis af danskere ind i døgnåbne fitnesscentre med hovedtelefoner på, mens lokale idrætsforeninger kæmper hårdt for at rekruttere nye trænere og frivillige til bestyrelsen. Ved første øjekast ligner det blot en uskyldig ændring i vores motionsvaner og fritidsforbrug. \dabold{Set i et bredere sociologisk perspektiv} \dabold{vidner udviklingen imidlertid om} et dybere \dabold{skift fra} kollektiv forpligtelse \dabold{mod} individuel optimering i det moderne hverdagsliv.

I mere end et århundrede har det \da{institutionaliserede} foreningsliv udgjort en grundpille i den danske \da{sammenhængskraft}. I den lokale fodboldklub, spejdergruppe eller gymnastikforening mødtes borgere på tværs af generationer og sociale skel, ikke blot som betalende kunder, men som ligeværdige medlemmer med fælles ansvar. På den årlige generalforsamling og i det praktiske arbejde omkring halbestyrelsen lærte både unge og ældre hverdagsdemokrati og gensidig respekt i praksis.

De seneste årtier har imidlertid bragt en hastig \da{kommercialisering} og \da{individualisering} af fritiden. Moderne fitnesskæder, digitale løbe-apps og fleksible holdtilbud er skræddersyet præcis til den enkeltes stramme kalender. Når man køber et privat medlemskab, kan man træne klokken seks om morgenen eller ni om aftenen uden at skylde nogen en forklaring, hvis man springer næste uges træning over på grund af arbejde.

Denne forandring kan kun forstås fuldt ud \dabold{i lyset af} de strukturelle pres, som blev belyst tidligere i dette kapitel. Når dobbeltarbejdende familier kæmper med skiftende vagtplaner, når ældre borgere forventes at blive længere på arbejdsmarkedet, og når pendlere i voksende storbyer bruger stadig flere timer på transport, bliver faste mødetider og frivillige vagter svære at passe ind. I stedet søger mange mod \da{uforpligtende fællesskaber} --- såsom åbne løbegrupper eller midlertidige træningshold --- hvor man kan nyde socialt samvær uden langvarige pligter.

Spørgsmålet er blot, hvad samfundet mister, når fritiden først og fremmest organiseres som en privat vare. Kommercielle tilbud skaber sjældent den dybe, gensidige tillid, som opstår, når naboer løfter en fælles opgave over tid. Samtidig risikerer dyrere private abonnementer at skabe nye sociale skel mellem familier med råd til personlig træning og udsatte børn, som tidligere fandt trygge rammer i den lokale forening. Hvis civilsamfundet skal bevare sin styrke i et aldrende og travlt samfund, må foreningerne både forny deres organisering og fastholde værdien af forpligtende fællesskab.

\begin{transbox}
\small
\textbf{\color{NCDTeal}[Tekst A 中文精译]}\\
一年四季的每个傍晚，成千上万名丹麦人戴着耳机走进全天候营业的健身房，而地方体育社团却在苦苦招募志愿教练和理事会成员。乍一看，这似乎只是我们运动习惯与休闲消费上的无伤大雅的变化。然而，放在更宏阔的社会学视野下来看，这一演变折射出一种在现代日常生活中从集体责任走向个人自我优化的深层转向。

一个多世纪以来，高度制度化的社团生活一直构成丹麦社会凝聚力的基石。在当地的足球俱乐部、童子军团体或体操协会里，跨越不同代际和社会阶层的市民汇聚一堂——他们不仅是付费消费者，更是承担共同责任的平等成员。在一年一度的会员大会以及场馆理事会的日常实务中，年轻人和长辈都在实践中学会了日常民主与彼此尊重。

然而，近几十年却带来了休闲领域的迅速商业化与个人主义化。现代连锁健身房、数字跑步应用以及灵活的团课产品，都是为个人紧张的日程表精准量身定制的。只要购买一张私人会员卡，人们就可以在早晨六点或晚上九点前去锻炼，倘若因为工作原因跳过下周的训练，也不必向任何人交代。

唯有放在本单元前文所剖析的结构性压力背景下，才能充分理解这一变化。当双职工家庭在频繁倒班的工作制中疲于奔命，当资深公民被期待留在职场更长时间，且不断扩张的大城市通勤者在交通上耗费越来越多小时，固定的训练时段与志愿值班便越来越难以嵌入日程。取而代之的是，许多人转向低约束性的社群——例如开放跑团或临时训练小组——在那里人们既能享受社交陪伴，又无需承担长期义务。

问题在于，当休闲生活主要作为一种私人市场商品来组织时，社会究竟会失去什么。商业化服务很少能够孕育出那种邻里长期共同扛起一项任务时所诞生的深层相互信任。与此同时，更加昂贵的私人会籍费还可能制造新的社会鸿沟：一边是负担得起私教课程的家庭，另一边则是原本能够在地方社团中找到安心归属的弱势儿童。如果公民社会要在一个日趋老龄化且步履匆匆的社会中守住自身力量，各类社团就既需要革新组织方式，也必须守护好“有责任担当的共同体”这一核心价值。
\end{transbox}

\subsection*{Tekst B：\textit{Modperspektiv: Nye fællesskaber er ikke nødvendigvis svagere}（社会学评论回应 · 约 140 词）}

\begin{emailbox}{报刊学术短评｜Af Kasper Nørgaard, kultursociolog}
Bekymringen for det danske foreningslivs undergang overser ofte, at menneskers behov for fællesskab slet ikke er forsvundet --- det har blot ændret form. Mange unge danskere fravælger ikke samarbejde af rendyrket egoisme, men fordi ældre foreningsstrukturer i praksis opleves som tunge, ufleksible og alt for institutionaliserede.

I både storbyer og mindre provinsbyer ser vi i dag en eksplosion af nye, mere smidige samværsformer: frivillige løbeklubber uden kontingent, lokale byhaver samt ugentlige fællesspisninger i gårdlauget. Disse uforpligtende fællesskaber har lave tærskler, så også nye tilflyttere, enlige og travle studerende let kan deltage og bidrage aktivt, når hverdagen tillader det.

Individualisering betyder altså ikke nødvendigvis social kulde eller ligegyldighed. Hvis traditionelle foreninger tør møde borgerne med mere fleksible roller, digitale løsninger og færre formelle møder, kan nyt og gammelt civilsamfund sagtens supplere hinanden og fortsat styrke Danmarks sammenhængskraft.
\end{emailbox}

\textbf{Tekst B 阅读核对（Rigtigt / Forkert / Kort svar，答案见第六部分）：}
\begin{enumerate}
  \item \textit{Kasper mener, at unge danskere vælger traditionelle foreninger fra, fordi de slet ikke ønsker fællesskab længere.} \underline{\hspace{2.2cm}}
  \item \textit{Hvilke tre eksempler på nye, fleksible samværsformer nævner Kasper i sin kommentar?} \underline{\hspace{2.2cm}}
  \item \textit{Hvad bør foreninger ændre ifølge Kasper for at passe bedre til moderne borgere?} \underline{\hspace{2.2cm}}
\end{enumerate}

\begin{tipbox}{本课发音与语调要点（Udtale \& Prosodi）}
\begin{itemize}
  \item \textbf{拉丁/法语词源社会学长词重音铁律}：凡以 \textit{-isering} 结尾的抽象过程名词，主重音一律落在倒数第三音节的长元音 \textit{-se-} 上（带喉塞音）：\textit{individualisering} \ipa{end̥ivid̥uɑliˈseːˀɐeŋ}；\textit{kommercialisering} \ipa{kʰʌmɛɐ̯siɑliˈseːˀɐeŋ}；同理，形容词 \textit{institutionaliseret} \ipa{ensd̥itˢuʃonɑliˈseːˀɐð} 主重音也在 \textit{-se-} 上。
  \item \textbf{核心文化复合词}：\textit{sammenhængskraft} \ipa{ˈsɑmənˌhɛŋˀsˌkʰʁɑfd}（重音在首音节，\textit{hængs} 带喉塞音）；\textit{uforpligtende} \ipa{uˌfʌˈpʰleɡd̥ənə}（注意 \textit{-ende} 词尾的 \textit{d} 完全不发音）；\textit{en tærskel} \ipa{ˈtˢɛɐ̯sɡəl}。
  \item \textbf{动词不规则读音提醒}：\textit{at forny}（现在时 \textit{fornyr} \ipa{fʌˈnyːˀɐ}，重音在 \textit{ny}，注意拼写不加 \textit{e}！）。
\end{itemize}
\end{tipbox}

\begin{transbox}
\small
\textbf{\color{NCDTeal}[Tekst B 中文对照]}\\
对丹麦社团生活衰落的担忧往往忽视了一点：人类对群体归属感的基本需求根本没有消失——它只是改变了形式。许多丹麦年轻人并非出于纯粹的自私才放弃传统社团合作，而是因为旧有的社团架构在实践中显得过于沉重、僵化且高度制度化。

无论在大城市还是中小城镇，我们如今都看到新型、更灵活的共处方式正在井喷：无需会费的志愿跑团、社区共享菜园以及邻里庭院里每周一次的共餐活动。这些低约束社群的参与门槛很低，让新迁入的居民、独居人士和忙碌的大学生在日常时间允许时都能轻松加入并积极贡献力量。

因此，个人主义化并不必然意味着人情冷漠或漠不关心。只要传统社团敢于用更灵活的角色分工、数字化工具以及更少的正式会议来迎合市民需求，新兴与传统的公民社会完全能够相辅相成，继续巩固丹麦的社会凝聚力。
\end{transbox}

% ---------------------------------------------------------------------
\ncdsubsec{三、Gloser（核心生词表）}

\begin{tipbox}{词表重点}
本课重点攻克\textbf{C1 社会学抽象名词后缀规律}：以 \textit{-isering} 结尾的词（\textit{individualisering, kommercialisering, urbanisering}）表示社会转型过程，全部为 \textbf{en} 性且通常只有单数（\textit{kun ental}）；而以 \textit{-skab} 结尾的抽象或集合概念（\textit{et fællesskab, et medlemskab}）多为 \textbf{et} 性！另外务必牢记动词 \da{at forny} 的现在时是 \dabold{fornyr}（无 \textit{-er}），以及 \textit{tærskel} 在定冠词和复数中会脱落 \textit{e}（\textit{tærsklen, tærskler}）。
\end{tipbox}

\begin{glosertable}
名词 & en individualisering & \textit{(individualiseringen, kun ental)} & 个人主义化，个体化 \\
名词 & en sammenhængskraft & \textit{(sammenhængskraften, kun ental)} & 社会凝聚力 \\
名词 & en kommercialisering & \textit{(kommercialiseringen, kun ental)} & 商业化 \\
名词 & et fællesskab & \textit{(fællesskabet, fællesskaber, fællesskaberne)} & 共同体，社群，群体归属 \\
名词 & et foreningsliv & \textit{(foreningslivet, kun ental)} & 社团生活，协会文化 \\
名词 & en frivillighed & \textit{(frivilligheden, kun ental)} & 志愿精神，志愿参与 \\
名词 & en generalforsamling & \textit{(generalforsamlingen, generalforsamlinger, generalforsamlingerne)} & 会员大会，年度股东/成员大会 \\
名词 & en tærskel & \textit{(tærsklen, tærskler, tærsklerne)} & 门槛，界限 \\
名词 & et skel & \textit{(skellet, skel, skellene)} & 鸿沟，分界线 \\
名词 & en forpligtelse & \textit{(forpligtelsen, forpligtelser, forpligtelserne)} & 义务，责任约束 \\
名词 & et medlemskab & \textit{(medlemskabet, medlemskaber, medlemskaberne)} & 会员资格，会籍 \\
\midrule
动词 [har] & at vidne (om) & \textit{(vidner, vidnede, har vidnet)} & 见证，折射出，表明 \\
动词 [har] & at rekruttere & \textit{(rekrutterer, rekrutterede, har rekrutteret)} & 招募，吸纳 \\
动词 [har] & at undergrave & \textit{(undergraver, undergravede, har undergravet)} & 侵蚀，从根本上削弱 \\
动词 [har] & at forny & \textit{(fornyr, fornyede, har fornyet)} & 革新，更新换代 \\
\midrule
形容词 & uforpligtende & \textit{(uforpligtende, 不变)} & 无约束的，轻松自愿的 \\
形容词 & institutionaliseret & \textit{(institutionaliseret, institutionaliserede)} & 制度化的，体制化建制的 \\
形容词 & kommerciel & \textit{(kommercielt, kommercielle)} & 商业性的，营利性的 \\
\end{glosertable}

% ---------------------------------------------------------------------
\ncdsubsec{四、Grammatik \& Stilistik（社会学升华框架与抽象名词化精讲）}

\begin{grammatikbox}{语法点 1：从生活细节升华到社会学判断的四层推进矩阵}
Studieprøven 与 PD3 高分段最看重“以小见大”的论述力：不要停留在“大家都在用健身 App”，而要用下面四个语言阶梯把现象提炼成结构分析：

\vspace{0.3em}
\renewcommand{\arraystretch}{1.24}
{\small
\begin{tabularx}{\linewidth}{@{} >{\bfseries\raggedright\arraybackslash}p{2.9cm} >{\raggedright\arraybackslash}p{4.6cm} >{\raggedright\arraybackslash}X @{}}
\toprule
\textbf{分析阶梯} & \textbf{核心学术句型} & \textbf{C1 社会学升华示范句} \\
\midrule
\da{1. 现象折射本质} &
\textit{Ved første øjekast ligner det...}\newline
\textit{Set i et bredere perspektiv \dabold{vidner} det \dabold{om}...} &
\textit{Set i et bredere perspektiv \dabold{vidner} de nye træningsvaner \dabold{om} en stigende individualisering.} \\
\midrule
\da{2. 演变方向定性} &
\textit{et skift \da{fra} A \da{mod} B}\newline
\textit{en kommercialisering af [领域]} &
\textit{Vi ser et tydeligt skift \da{fra} kollektiv forpligtelse \da{mod} fleksible, kommercielle tilbud.} \\
\midrule
\dabold{3. 关联宏观背景} &
\textit{\dabold{I lyset af} [宏观趋势]...}\newline
（串联 L1 老龄化 / L2 工时 / L3 通勤） &
\textit{\dabold{I lyset af} længere pendlertider bliver faste vagter i foreningen svære at passe ind.} \\
\midrule
\dabold{4. 评估社会张力} &
\textit{på tværs af generationer og sociale skel}\newline
\textit{undergrave sammenhængskraften} &
\textit{Dyrere private medlemskaber kan skabe nye skel og \dabold{undergrave} den lokale \da{sammenhængskraft}.} \\
\bottomrule
\end{tabularx}
}
\end{grammatikbox}

\begin{grammatikbox}{语法点 2：抽象形容词变形与社会学常见语病排查}
\begin{enumerate}
  \item \textbf{以 \textit{-ende} 结尾的分词形容词不变形 vs 以 \textit{-et} 结尾的过去分词形容词要变 \textit{-ede}}：
    \begin{itemize}
      \item \da{uforpligtende}（现在分词作定语）在任何性别、单复数和定冠词前都\textbf{保持不变}：\textit{et uforpligtende møde / uforpligtende fællesskaber}。
      \item \da{institutionaliseret}（过去分词作定语）在通性/中性单数不定式作 \textit{-et}，但在\textbf{复数或定冠词后必须变为 \dabold{-ede}}：\ding{51}\ \textit{det institutionaliserede foreningsliv / alt for institutionaliserede strukturer}（切忌漏掉变化写成 \ding{55}\ \textit{*institutionaliseret strukturer}）。
    \end{itemize}
  \item \textbf{动词作主语升级为抽象名词作主语（Nominalstil）}：与其写口语化的 \textit{Når fritiden bliver mere kommerciel, mister vi tillid.}，不如提升为 \da{Kommercialiseringen af fritiden risikerer at svække den gensidige tillid.}——这正是 B2 与 C1 的文体分水岭！
  \item \textbf{高频介词固定搭配铁律}：
    \begin{itemize}
      \item \da{på tværs af} generationer og sociale skel（跨代际与社会阶层——绝不照英语直译为 \ding{55}\ \textit{*på kryds af}）。
      \item \da{i lyset af} de strukturelle pres（鉴于/放在……的背景下——注意定冠词 \textit{lyset} 不能省略，绝不可写 \ding{55}\ \textit{*i lys af}）。
      \item \da{at vidne om} + 名词短语（折射出/证明了——不可漏掉介词 \textit{om}）。
    \end{itemize}
\end{enumerate}
\end{grammatikbox}

% ---------------------------------------------------------------------
\ncdsubsec{五、Mønsterøvelser（社会学词汇、升华与综合操练）}

\textbf{练习 A：社会学分析核心词与搭配精准填空（每空填入最准确的 1 个丹麦语词）}
\begin{enumerate}
  \item Set i et bredere sociologisk perspektiv \underline{\hspace{2.2cm}} fitnessbølgen om en dybere samfundsændring.
  \item De nye motionsvaner afspejler et tydeligt \underline{\hspace{2.2cm}} fra kollektive pligter mod individuel fleksibilitet.
  \item I \underline{\hspace{2.2cm}} af længere pendlertider har mange travle familier svært ved at nå faste møder.
  \item I den lokale sportsklub mødes borgerne på \underline{\hspace{2.2cm}} af alder, indkomst og uddannelse.
  \item Mange unge søger i dag \underline{\hspace{2.2cm}} fællesskaber, hvor man kan deltage uden faste bestyrelsesposter.
  \item Hvis alle trækker sig ind i private tilbud, kan det på sigt undergrave byens sociale \underline{\hspace{2.2cm}}.
  \item Foreningerne skal \underline{\hspace{2.2cm}} deres arbejdsformer, hvis de vil tiltrække travle tilflyttere.
  \item De gamle strukturer opleves af mange unge som alt for \underline{\hspace{2.2cm}} og præget af tunge regler.
\end{enumerate}

\textbf{练习 B：抽象句式升华与语法排雷（将口语描述升级为 C1 社会学评论或修正错句）}
\begin{enumerate}
  \item 【纠错】修正形容词定冠词变形与介词漏用：\textit{Det institutionaliseret foreningsliv vidner en stærk kultur på kryds af generationer.} $\rightarrow$
  \item 【纠错】修正固定短语冠词与动词变位：\textit{I lys af travle hverdage fornyer lokale klubber nu deres frivillighed.} $\rightarrow$
  \item 【纠错】修正 \textit{på tværs af} 与复数搭配：\textit{De nye løbeklubber har lave tærskel og åbner for fællesskab tværs af sociale skel.} $\rightarrow$
  \item 【升级】用抽象名词 \textit{Kommercialiseringen af fritiden} 改写：\textit{Når private firmaer tjener flere penge på vores fritid, bliver det sværere for fattige børn at være med.} $\rightarrow$
  \item 【升级】用 \textit{Set i et bredere perspektiv vidner... om...} 升华日常观察：\textit{Flere folk træner alene med musik i ørerne, og det viser, at samfundet er blevet mere individuelt.} $\rightarrow$
  \item 【升级】串联 Unit 1：用 \textit{I lyset af + V2 倒装} 合并原因与结果：\textit{Befolkningen ældes og folk pendler langt, så frivillige foreninger mangler hjælpere.} $\rightarrow$
\end{enumerate}

\textbf{练习 C：中译丹（灵活运用本课社会学语汇与宏观分析句式）}
\begin{enumerate}
  \item 从更广阔的视角来看，健身热潮折射出日益增强的个人主义化。
  \item 传统社团在跨越代际和社会鸿沟方面一直发挥着重要作用。
  \item 鉴于大城市漫长的通勤时间，许多人更偏爱低约束性的社群。
  \item 休闲领域的商业化可能在长远上削弱地方的社会凝聚力。
  \item 许多年轻人认为传统社团结构过于制度化且门槛过高。
  \item 如果传统协会能够革新自身组织形式，它们依然能够吸引志愿参与者。
\end{enumerate}

% ---------------------------------------------------------------------
\ncdsubsec{六、PD3 / Studieprøven 实战微任务与参考答案}

\begin{emailbox}{Unit 1 收官综合任务｜Studieprøven / PD3 Debatindlæg（mindst 250 ord）}
\textbf{Situation}: I en avisartikel fremføres det, at det danske foreningsliv er døende, fordi moderne borgere kun tænker på sig selv og deres egen træning.\\
\textbf{Opgave}: Skriv et reflekteret og sociologisk funderet debatindlæg på \textbf{mindst 250 ord}, hvor du nuancerer påstanden. Du skal:
\begin{enumerate}
  \item Analysere skiftet fra traditionelt foreningsliv til individuel og kommerciel træning ved hjælp af præcise sociologiske begreber (\textit{individualisering, kommercialisering, uforpligtende fællesskaber, sammenhængskraft}).
  \item Forklare forandringen \textbf{i lyset af} mindst to af de bredere samfundstendenser fra Unit 1 (f.eks. en aldrende befolkning, tidspres på arbejdsmarkedet eller længere pendling i storbyerne).
  \item Vurdere med inspiration fra Tekst B, hvordan både nye og traditionelle fællesskaber kan styrke civilsamfundet fremover.
\end{enumerate}
\end{emailbox}

\begin{answerbox}
\textbf{Før du læser 参考回答}
\begin{enumerate}
  \item \textit{Tekst A fortolker skiftet som et udtryk for individualisering og kommercialisering af fritiden, hvor kollektiv forpligtelse erstattes af fleksibel selvoptimering. Samtidig kobles denne udvikling til lektion 1--3 ved, at travle dobbeltarbejdende familier og længere storby-pendling gør faste foreningsmøder svære at passe ind.}
  \item \textit{Kasper mener ikke, at unge danskere har mistet interessen for fællesskab, fordi de stadig søger sociale aktiviteter som gratis løbeklubber, byhaver og fællesspisninger. De vælger blot traditionelle foreninger fra, fordi de oplever strukturerne som tunge og alt for institutionaliserede.}
\end{enumerate}

\textbf{Tekst B 阅读核对答案}
\begin{enumerate}
  \item \da{Forkert}（他明确指出年轻人对群体归属感的需求并未消失，只是不再喜欢高度制度化的旧模式）
  \item \da{Frivillige løbeklubber uden kontingent, lokale byhaver og ugentlige fællesspisninger i gårdlauget.}
  \item \da{At møde borgerne med mere fleksible roller og færre formelle møder.}（提供更灵活的角色分工并减少繁琐正式会议）
\end{enumerate}

\textbf{练习 A 答案}
\begin{enumerate}
  \item \da{vidner}（固定搭配：\textit{vidne om} 折射出、表明）
  \item \da{skift}（社会演变表述：\textit{et skift fra A mod B} 从 A 向 B 的转向）
  \item \da{lyset}（固定学术介词结构：\textit{i lyset af} 鉴于、在……背景下）
  \item \da{tværs}（固定搭配：\textit{på tværs af} 跨越……界限）
  \item \da{uforpligtende}（不变形分词形容词：“低约束的、轻松自愿的”）
  \item \da{sammenhængskraft}（丹麦社会学高频核心词：“社会凝聚力”）
  \item \da{forny}（动词原形：革新、更新）
  \item \da{institutionaliserede}（复数名词 \textit{strukturer} 修饰语，过去分词形容词须加 \textit{-ede}）
\end{enumerate}

\textbf{练习 B 答案（参考）}
\begin{enumerate}
  \item \textit{Det institutionaliserede foreningsliv vidner om en stærk kultur på tværs af generationer.}（定冠词后分词形容词加 \textit{-ede}；补齐 \textit{vidne om} 的介词 \textit{om}；将英语式 \textit{på kryds af} 改为 \textit{på tværs af}）
  \item \textit{I lyset af travle hverdage fornyr lokale klubber nu deres frivillighed.}（必须使用定冠词形态 \textit{I lyset af}；动词 \textit{at forny} 的现在时拼写为 \textit{fornyr}）
  \item \textit{De nye løbeklubber har lave tærskler og åbner for fællesskab på tværs af sociale skel.}（可数名词用复数 \textit{lave tærskler}；完整介词短语为 \textit{på tværs af}）
  \item \textit{Kommercialiseringen af fritiden risikerer at skabe nye sociale skel for udsatte børn.}（以抽象名词短语作主语取代 \textit{når} 口语条件句，大幅提升语体）
  \item \textit{Set i et bredere perspektiv vidner de individuelle træningsvaner om en stigende individualisering i samfundet.}（使用固定升华框架把具体观察转化为社会趋势判断）
  \item \textit{I lyset af en aldrende befolkning og længere pendling kæmper frivillige foreninger for at rekruttere hjælpere.}（将前置条件融合为 \textit{I lyset af...} 并确保主句动词 \textit{kæmper} 遵循 V2 倒装）
\end{enumerate}

\textbf{练习 C 答案}
\begin{enumerate}
  \item \textit{Set i et bredere perspektiv vidner fitnessbølgen om en stigende individualisering.}
  \item \textit{Traditionelle foreninger har spillet en vigtig rolle på tværs af generationer og sociale skel.}
  \item \textit{I lyset af lange pendlertider i storbyerne foretrækker mange uforpligtende fællesskaber.}
  \item \textit{Kommercialiseringen af fritiden kan på længere sigt undergrave den lokale sammenhængskraft.}
  \item \textit{Mange unge oplever traditionelle foreningsstrukturer som alt for institutionaliserede og med høje tærskler.}
  \item \textit{Hvis traditionelle foreninger fornyr deres organisering, kan de stadig tiltrække frivillige deltagere.}
\end{enumerate}

\textbf{PD3 / Studieprøven 高分范文（约 265 词，不含标题）}\\
\textit{\textbf{Foreningslivet dør ikke --- det skifter ham i en travl tid}}

\textit{Påstanden om, at det danske foreningsliv er døende på grund af ren egoisme, rammer forbi skiven. \dabold{Set i et bredere sociologisk perspektiv vidner} de tomme bestyrelsesposter og fyldte fitnesscentre nemlig ikke om manglende socialt behov hos borgerne, men om et strukturelt \dabold{skift fra} traditionelle, faste pligter \dabold{mod} fleksible og kommercielle tilbud, der passer langt bedre til en fragmenteret hverdag.}

\textit{Gennem generationer har det \da{institutionaliserede} foreningsliv styrket Danmarks \da{sammenhængskraft} ved at bringe naboer sammen \da{på tværs af} alder, indkomst og sociale skel. Men \dabold{i lyset af} nutidens øvrige samfundstendenser er rammerne ændret markant. Når dobbeltarbejdende familier kæmper med skiftende vagtplaner, når ældre borgere forventes at arbejde længere på arbejdsmarkedet, og når høje boligpriser presser lønmodtagere ud i oplandskommunerne med længere daglig pendling, bliver faste aftentræninger og årlige generalforsamlinger svære at passe ind. Den stigende \da{individualisering} og \da{kommercialisering} af fritiden er derfor ofte et praktisk svar på et konkret tidspres i hverdagen frem for manglende interesse for andre mennesker.}

\textit{Alligevel indebærer udviklingen en reel samfundsmæssig risiko. Hvis dyre private abonnementer erstatter billige klubber, kan det skabe nye sociale skel mellem velstillede og udsatte familier samt på sigt \dabold{undergrave} den gensidige tillid i lokalområdet. Samtidig viser uformelle løbeklubber, lokale byhaver og ugentlige fællesspisninger, at danskerne stadig aktivt søger \da{uforpligtende fællesskaber} med lave tærskler og hurtig adgang.}

\textit{Løsningen er derfor hverken at skælde individet ud eller at opgive den klassiske foreningsmodel. Hvis lokale foreninger formår at \da{forny} sig med mere fleksible frivillige roller, digitale løsninger og færre tunge møder, kan både nye og traditionelle fællesskaber fortsat supplere hinanden og bære det danske civilsamfund trygt ind i fremtiden.}

\small\textbf{名师点评}：作为 Unit 1 收官大作，这篇约 265 词的议论文示范了最高水准的跨课综合能力！它没有孤立地谈健身房现象，而是用 \textit{i lyset af} 巧妙串联起第一课的老龄化延迟退休、第二课的排班工时压力与第三课的高房价通勤困境，并将 \textit{institutionaliserede, sammenhængskraft, individualisering, kommercialisering, uforpligtende fællesskaber} 与正反辩证完美融合，堪称 Studieprøven C1 范本。
\end{answerbox}


% =====================================================================
% Unit 2
% =====================================================================
\ncdunit{2}{Retorik, Ironi og Mediekritik}{修辞、反讽与媒体评论}{Lektion 5--8：丹麦式反讽与低调陈述（Underdrivelse/Litote）· 报刊社论（Leder）高阶让步 · 引述动词态度光谱 · 多重主从复合长句节奏}

\ncdlesson{5}{Dansk ironi og kunsten at underdrive}{丹麦式反讽与低调委婉的艺术}
{DU3 Modul 5 (PD3) $\rightarrow$ Modul 6 (Studieprøven, B2 $\rightarrow$ C1)｜对标 NCE 4 \textit{Retorik, Ironi \& Pragmatik}}
{\textbf{丹麦语修辞精髓：低调陈述与反语（\textit{Underdrivelse / Litote}）}（\textit{ikke helt dårligt, en smule udfordrende, det kan vist næppe kaldes en triumf, ikke ligefrem uproblematisk}）· 否定微调副词（\textit{næppe, ligefrem, knap nok, langt fra}）的语用强度与语序}
{从 B2 迈向 C1，语法正确只是入场券，真正区分“说丹麦语的人”和“懂丹麦社会密码的人”，是对\textbf{低调陈述（\textit{underdrivelse}）与反语（\textit{litote}）}的精准识读与运用。在丹麦职场、报刊评论与日常交谈中，最高级别的赞美往往藏在“不算太糟”里，而最严厉的批评却常常披着“略具挑战性”的外衣。吃透本课，你不仅能听出丹麦人温和语气下的锋芒，还能在 PD3 与 Studieprøven 评论写作中展现出高级的分寸感与冷峻幽默！}

% ---------------------------------------------------------------------
\begin{foerduleaeser}
带着下面两个宏观问题阅读 Tekst A 和 Tekst B：
\begin{enumerate}
  \item \textit{Hvorfor kan danskernes brug af litoter og underdrivelser skabe misforståelser for udlændinge på arbejdspladsen ifølge Tekst A?}（根据 Tekst A，为什么丹麦人对否定反语和低调陈述的使用会在职场上给外国人造成误解？）
  \item \textit{Hvordan ser klummeskribenten i Tekst B på underdrivelsens rolle, når den flytter fra middagsselskabet ind i politik og ledelse?}（在 Tekst B 中，专栏作者如何看待低调陈述从私人餐桌进入政治与管理领域后的角色变化？）
\end{enumerate}
\end{foerduleaeser}

% ---------------------------------------------------------------------
\ncdsubsec{二、Tekster（文化修辞随笔与辛辣专栏回应双语篇）}

\subsection*{Tekst A：\textit{Når ”ikke helt dårligt” betyder fremragende}（文化与语用随笔 · 约 360 词）}

\begin{tipbox}{语篇说明}
本篇深度剖析丹麦公共交流与职场文化中最具代表性的修辞机制——\textit{underdrivelse} 与 \textit{litote}。阅读时请留意\textbf{\color{NCDNavy}深蓝色加粗的肯定性低调褒奖}（通过双否或弱化表达由衷认可），以及\textbf{\color{NCDRed}深红色加粗的委婉批评与冷峻反讽}（表面克制，实则指向严重问题）。
\end{tipbox}

Efter tre måneder i en dansk ingeniørvirksomhed stod den nyansatte projektleder Jonas over for sin første store evaluering. Projektet var blevet afleveret til tiden, budgettet holdt, og kunden var mere end tilfreds med resultatet. Da afdelingschefen bagefter kiggede forbi hans skrivebord og sagde, at indsatsen da var \da{”ikke helt dårlig”}, sank Jonas' hjerte en smule. Først senere forklarede en kollega ham, at i en dansk kontekst er en sådan bemærkning faktisk \da{langt fra ringe} --- ja, ofte svarer den til en varm anerkendelse af et stykke fremragende arbejde.

Denne særegne forkærlighed for at nedtone både ros og kritik kaldes i retorikken for \textit{underdrivelse}, og når den formes gennem en benægtelse af det modsatte, taler man om en \textit{litote}. Hvor andre kulturer gerne griber til store superlativer som ”fantastisk” eller ”katastrofalt”, foretrækker dansk dagligsprog og debatprosa en besindig mellemvej. At kalde en vellykket koncert for \da{”ikke uden kvalitet”} eller en hård vinter for \da{”lidt frisk i vinden”} skaber en uformel distance, hvor man demonstrerer, at man ikke tager sig selv alt for højtideligt. Samtidig afspejler tonen et lighedsideal, hvor prangende begejstring let kan virke påtrængende og skabe unødig afstand mellem mennesker.

Bagsiden viser sig imidlertid, når underdrivelsen bruges til at pakke alvorlige problemer ind. Når en rapport efter et mislykket it-projekt beskriver forløbet som \dabold{”en smule udfordrende”}, eller når en kommentator tørt bemærker, at en ministers håndtering af sagen \dabold{”vist næppe kan kaldes en triumf”}, ligger der en skarp brod gemt under den rolige overflade. Tilsvarende betyder sætningen ”forslaget er \dabold{ikke ligefrem uproblematisk}” på godt dansk, at forslaget rummer alvorlige mangler og bør tænkes grundigt om, før det vedtages.

For udlændinge, der lærer sproget, opstår misforståelserne næsten uundgåeligt, hvis man tager ordene alt for bogstaveligt. Den, der hører, at en opgave har \dabold{”et vist forbedringspotentiale”}, eller at ”vi nok lige skal kigge på et par detaljer”, tror måske, at der blot mangler en hurtig finpudsning, mens chefen i virkeligheden efterlyser en gennemgribende omarbejdelse. Kunsten består derfor ikke i blot at lære gloserne, men i at afkode den sociale kontekst, pauseføringen og det lille skæve smil, der røber ironien. Først når man mestrer dette skjulte register, kan man virkelig læse mellem linjerne i Danmark.

\begin{transbox}
\small
\textbf{\color{NCDTeal}[Tekst A 中文精译]}\\
在一家丹麦工程公司工作了三个月后，新入职的项目经理约纳斯迎来了自己的第一次大型评估。项目按时交付了，预算没有超支，客户对结果也感到格外满意。然而当部门主管随后走到他的办公桌前，说这次的表现“倒也不算太糟”时，约纳斯的心却沉了一下。直到后来，一位同事才向他解释说，在丹麦语境下，这样一句评价其实“远谈不上差劲”——甚至往往就相当于对一份出色工作的热烈肯定。

这种既习惯给赞美降温、也习惯给批评留白的独特偏好，在修辞学上被称为“低调陈述（underdrivelse）”；而当它通过否定反面来表达肯定时，人们便称之为“反语/曲言（litote）”。在别的文化或许习惯动用“无与伦比”或“灾难性的”这类最高级形容词的地方，丹麦的日常口语与议论散文却偏爱一条克制的中间道路。把一场成功的音乐会称作“并非全无水准”，或者把严酷的寒冬形容为“风稍微有点提神”，能够营造出一种不拘礼节的距离感，借此表明说话者并没有把自己看得太了不起。与此同时，这种语调也折射出一种平等主义理想——在那里，张扬夸张的赞叹很容易显得唐突，并在人与人之间制造不必要的距离。

然而，当低调陈述被用来包裹严重问题时，它的另一面便显露出来。当一份报告在某次失败的 IT 项目之后将整个过程描述为“有那么一点点挑战性”，或者当一位评论员冷冷地指出，某位大臣对事件的处理“恐怕很难称之为一场凯旋”时，平静的表面之下其实藏着尖锐的刺。同样地，句子“这项提案算不上完全没有问题”换作直白大白话，意思就是该提案存在严重缺陷，在获得通过前理应彻底重新考虑。

对于正在学习丹麦语的外国人而言，如果把字面意思太当真，误解几乎不可避免。听到自己的方案“具备一定的改进空间”，或者“我们大概还得顺便看一下几个细节”时，对方也许以为只需要快速润色一下，而上司实际上期待的是彻底重写。因此，真正的功夫不在于单纯背下生词，而在于解码社会语境、停顿节奏以及那抹泄露反讽意味的微妙微笑。唯有当你驾驭了这层隐蔽的语域，你才算真正读懂了丹麦社会的弦外之音。
\end{transbox}

\subsection*{Tekst B：\textit{Klumme: Når ironien bliver et tågeslør}（锐评专栏回应 · 约 145 词）}

\textbf{阅读任务（理解与核对）：}\textit{Læs kulturkommentatoren Sofies skarpe svar på Tekst A, og besvar spørgsmålene nedenfor.}

\begin{emailbox}{Klumme i dagbladet｜Fra: Sofie Kragh (kulturkommentator) \quad Emne: Sv: Kunsten at sige alt ved at sige næsten ingenting}
Det er selvfølgelig charmerende, når danskere smiler beskedent og kalder en personlig succes for \da{”ikke helt tosset”}. Men når den samme retorik flytter ind i bestyrelseslokaler og ministerier, bliver charmen hurtigt \dabold{en smule anstrengende}.

Alt for ofte fungerer den hyggelige underdrivelse nemlig som et bekvemt tågeslør, der skjuler reelle konflikter og svigt. Når et milliardprojekt overskrider tidsplanen med tre år, og den ansvarlige kalder det ”en mindre justering”, eller når en åbenlys krise i velfærden omtales som \dabold{”ikke ligefrem optimal”}, handler det sjældent om harmløs selvironi. Det handler om ansvarsforflygtigelse.

Ironi kræver modtagere, der kender koden, og modige afsendere, der tør mene noget i offentligheden. Uden tydelighed risikerer vi, at høfligheden overdøver kritikken, og at vælgerne mister tålmodigheden. Derfor kunne dansk offentlig debat vist godt tåle lidt færre litoter og lidt mere klar tale --- selvom det naturligvis ville være \dabold{en kende ubehageligt} for vores nationale selvforståelse.
\end{emailbox}

\textbf{Tekst B 阅读核对（Rigtigt / Forkert / Kort svar，答案见第六部分）：}
\begin{enumerate}
  \item \textit{Sofie mener, at underdrivelse altid er charmerende, også når den bruges af ministre og chefer.} \\\underline{\hspace{2.2cm}}
  \item \textit{Hvad kan en hyggelig underdrivelse ifølge Sofie i virkeligheden bruges til at skjule i politik og ledelse?} \underline{\hspace{2.2cm}}
  \item \textit{Hvad efterlyser Sofie til sidst i den danske offentlige debat, selvom det kan føles en kende ubehageligt?} \underline{\hspace{2.2cm}}
\end{enumerate}

\begin{tipbox}{本课发音与语调要点（Udtale \& Prosodi）}
\begin{itemize}
  \item \textbf{反讽与否定词关键读音}：\textit{en litote} \ipa{liˈtˢoːdə}（重音在第二音节）；\textit{en ironi} \ipa{iʁoˈniːˀ}（重音在词尾带 \textit{stød}）；\textit{næppe} \ipa{ˈnɛbə}；\textit{ligefrem} \ipa{ˈliːəˌfʁɑmˀ}（\textit{lige} 中的 \textit{g} 不发音，末尾 \textit{-frem} 带 \textit{stød}）；\textit{uproblematisk} \ipa{upʰʁob̥ləˈmæːˀd̥isɡ̊}。
  \item \textbf{极易吞音的长复合词}：\textit{selvironi} \ipa{ˈsɛlʔiʁoˌniːˀ}（\textit{selv} 结尾的 \textit{v} 常弱化或脱落）；\textit{ansvarsforflygtigelse} \ipa{ˈansvɑːsˌfʌˈflyɡ̊d̥iəlsə}（词中 \textit{g} 在 \textit{-ig-} 中不发音）；\textit{forbedringspotentiale} \ipa{fʌˈb̥ɛðˀʁeŋspʰotˢɛnˌɕæːlə}（含软 d + \textit{stød}，词尾 \textit{-tiale} 读作 \ipa{ɕæːlə}）。
  \item \textbf{反讽语调的“降调平冷感”}：丹麦语的低调反语从不高声夸张！朗读 \textit{det kan vist næppe kaldes en triumf} 或 \textit{ikke ligefrem optimal} 时，整体音高偏低、节奏平稳略慢，仅在 \textit{næppe} 或 \textit{ligefrem} 上轻点重音，留半拍沉默，反讽效果才最地道。
\end{itemize}
\end{tipbox}

\begin{transbox}
\small
\textbf{\color{NCDTeal}[Tekst B 中文对照]}\\
\textbf{日报专栏（索菲·克拉格，文化评论员｜回应：靠几乎什么都不说来表达一切的艺术）：}当丹麦人谦虚地微笑着，把一项个人成就称作“倒也还不算太蠢”时，这当然挺迷人的。然而一旦同一套修辞搬进董事会会议室和政府各部，这份魅力很快就会变得“有那么一点吃力”。

因为这种惬意的低调陈述实在太频繁地充当了一层方便的迷雾，用来遮掩真实的冲突与失职。当一个耗资十亿克朗的项目比原定时间表延期三年，而负责人却称之为“一次小幅调整”，或者当公共福利中一场显而易见的危机被轻描淡写地称为“算不上完全理想”时，这很少关乎无害的自嘲幽默。这关乎逃避责任。

反讽需要懂得密码的接收者，也需要敢于在公共空间持有立场的发送者。如果没有清晰度，我们就会面临礼貌压倒批评、选民最终失去耐心的风险。正因如此，丹麦的公共讨论恐怕倒是可以承受少一点曲言反语、多一点直白坦诚——尽管这对我们的民族自我认知来说，当然会有“那么一丁点不自在”。
\end{transbox}

% ---------------------------------------------------------------------
\ncdsubsec{三、Gloser（核心生词表）}

\begin{tipbox}{词表重点}
动词 \textit{at underdrive}（低调陈述）与 \textit{at overdrive}（夸张）同属强变化动词，过去式与过去分词完全平行：\textit{underdriver, underdrev, har underdrevet}。名词 \textit{et tågeslør}（迷雾、烟幕）与 \textit{et forbrug/forhold} 类似，单复同形；注意 \textit{næppe}（几乎不/恐怕不）本身已含否定意义，切勿再与 \textit{ikke} 叠用。
\end{tipbox}

\begin{glosertable}
名词 & en underdrivelse & \textit{(underdrivelsen, underdrivelser, underdrivelserne)} & 低调陈述，轻描淡写 \\
名词 & en litote & \textit{(litoten, litoter, litoterne)} & 反语，曲言法（否定反面表肯定） \\
名词 & en selvironi & \textit{(selvironien, kun ental)} & 自嘲，自我反讽 \\
名词 & en anerkendelse & \textit{(anerkendelsen, anerkendelser, anerkendelserne)} & 认可，赞赏，承认 \\
名词 & et lighedsideal & \textit{(lighedsidealet, lighedsidealer, lighedsidealerne)} & 平等主义理想 \\
名词 & et forbedringspotentiale & \textit{(forbedringspotentialet, forbedringspotentialer, forbedringspotentialerne)} & 改进空间，提升潜力 \\
名词 & en omarbejdelse & \textit{(omarbejdelsen, omarbejdelser, omarbejdelserne)} & 彻底修改，重写，改编 \\
名词 & et tågeslør & \textit{(tågesløret, tågeslør, tågeslørene)} & 迷雾，烟幕，掩人耳目的说辞 \\
名词 & en ansvarsforflygtigelse & \textit{(ansvarsforflygtigelsen, kun ental)} & 推卸责任，逃避担当 \\
名词 & en selvforståelse & \textit{(selvforståelsen, selvforståelser, selvforståelserne)} & 自我认知，自我定位 \\
\midrule
动词 [har] & at underdrive & \textit{(underdriver, underdrev, har underdrevet)} & 轻描淡写，低调表达 \\
动词 [har] & at nedtone & \textit{(nedtoner, nedtonede, har nedtonet)} & 淡化，给……降温 \\
动词 [har] & at afkode & \textit{(afkoder, afkodede, har afkodet)} & 解码，识破（潜台词） \\
动词 [har] & at overdøve & \textit{(overdøver, overdøvede, har overdøvet)} & 压倒，盖过（声音/批评） \\
\midrule
形容词 & besindig & \textit{(besindigt, besindige)} & 冷静克制的，沉稳的 \\
形容词 & påtrængende & \textit{(påtrængende, 不变)} & 唐突冒昧的，强加于人的 \\
副词 & næppe & \textit{(副词，不变)} & 几乎不，恐怕未必 \\
副词 & ligefrem & \textit{(副词，不变)} & 简直，算得上，直接地 \\
\end{glosertable}

% ---------------------------------------------------------------------
\ncdsubsec{四、Grammatik \& Stilistik（丹麦式反讽、Litote 光谱与语序精讲）}

\begin{grammatikbox}{修辞点 1：丹麦语低调陈述（\textit{Underdrivelse}）与反语（\textit{Litote}）实战光谱}
在 PD3 与 Studieprøven 阅读理解（\textit{Læsning}）中，能否读出作者字面之下的真实立场往往决定成败；而在专栏或辩论写作中恰当点缀一两处 \textit{litote}，则能立即赋予文章纯正的丹麦笔锋：

\vspace{0.3em}
\renewcommand{\arraystretch}{1.24}
{\small
\begin{tabularx}{\linewidth}{@{} >{\bfseries\raggedright\arraybackslash}p{4.3cm} >{\raggedright\arraybackslash}p{4.5cm} >{\raggedright\arraybackslash}X @{}}
\toprule
\textbf{丹麦语地道表达} & \textbf{字面意思 vs 真实语用力度} & \textbf{典型语境与修辞功能} \\
\midrule
\da{ikke helt dårligt} /\newline \da{langt fra ringe} & 字面：不算太差 / 远非低劣\newline 实意：\textbf{相当出色、值得肯定} & 谦虚肯定他人成果或自谦成功，避免最高级带来的浮夸感。 \\
\midrule
\da{ikke uden kvalitet} /\newline \da{ikke helt tosset} & 字面：并非毫无水准\newline 实意：\textbf{很有水准、颇具亮点} & 书评、影评或文化闲谈中带点克制笑意的赞赏。 \\
\midrule
\dabold{en smule udfordrende} /\newline \dabold{en kende ubehageligt} & 字面：有一点挑战 / 略微不适\newline 实意：\textbf{相当棘手、极为尴尬难堪} & 职场报告、官场公关中用极软量词包装硬伤或危机。 \\
\midrule
\dabold{ikke ligefrem uproblematisk} /\newline \dabold{ikke helt optimalt} & 字面：算不上毫无问题\newline 实意：\textbf{问题重重、存在严重缺陷} & 批判性社论或专家质疑中最常见的高级委婉否定。 \\
\midrule
\dabold{kan vist næppe kaldes en triumf} & 字面：恐怕很难称作一场凯旋\newline 实意：\textbf{简直是一场彻头彻尾的失败} & 典型冷幽默反讽（搭配 \textit{vist + næppe} 双重降调，讽刺力拉满）。 \\
\bottomrule
\end{tabularx}
}
\end{grammatikbox}

\begin{grammatikbox}{语法点 2：低调微调副词（\textit{næppe, ligefrem, langt fra, knap nok}）的句法铁律}
\begin{enumerate}
  \item \textbf{\dabold{næppe} 自带完全否定力，严禁再加 \textit{ikke}}：\\
    \textit{næppe} 意为“几乎不、恐怕未必”。在语法上它与 \textit{ikke / aldrig} 处于同一互斥槽位（中场否定位置），同时出现会导致逻辑崩溃或严重语病：\\
    \ding{51}\ \textit{Forslaget kan \dabold{næppe} kaldes gennemtænkt.}（\ding{55}\ \textit{Forslaget kan næppe ikke kaldes...}）
  \item \textbf{\dabold{ikke ligefrem} 作为一个整体否定强化模块移动}：\\
    主句中放在变位动词之后（或助动词与实义动词之间）；从句中则必须整体站在变位动词\textbf{之前}：\\
    主句：\textit{Planen er \dabold{ikke ligefrem} uproblematisk.} / 从句：\textit{Han indrømmede, at planen \dabold{ikke ligefrem} var uproblematisk.}
  \item \textbf{复数或长段弱化量词 \da{en smule / en kende / et vist} 的搭配分工}：\\
    \ding{172} \da{en smule} 与 \da{en kende}（书面略带雅致反讽感）均作副词性程度修饰语，后接形容词原级，形容词保持单数/主语一致性（如 \textit{situationen er en kende ubehagelig}）；\\
    \ding{173} \da{et vist / en vis}（一定的、某种程度的）作形容词定语，必须随中心名词性数变化（\textit{en vis skepsis}，但 \textit{et vist forbedringspotentiale}）。
  \item \textbf{否定词置于句首触发 V2 倒装的修辞重音}：\\
    把 \textit{Langt fra} 或 \textit{Næppe} 推到句首用于反讽强调时，后面主句同样触发 V2 倒装：\textit{\dabold{Næppe} havde ministeren afsluttet sin tale, før kritikken begyndte.}
\end{enumerate}
\end{grammatikbox}

\begin{grammatikbox}{语用与写作指南（Stilistik）：何时“留白”，何时“直说”？}
\begin{itemize}
  \item \textbf{识破“礼貌陷阱”三大高频词}：如果丹麦主管说 \textit{”Det ser spændende ud, men...”}（这看起来挺有意思，不过……）、\textit{”Det kan vi godt kigge nærmere på”}（我们可以再仔细看看），或者评价方案 \textit{”har et vist forbedringspotentiale”}，千万别当成已经通过！这通常意味着需要实质性返工。
  \item \textbf{PD3 / Studieprøven 写作红线}：在正式学术综述（\textit{Referat}）或统计图表描述中，必须用明确精确的动词（\textit{falde markant, udgøre en tredjedel}），切忌用 \textit{en smule} 模糊客观数据；只有在\textbf{文化评论、社论或辩论投稿的结语/反驳句}中，使用一句克制的 \textit{ikke ligefrem uproblematisk} 或 \textit{det kan vist næppe kaldes...}，才能达到画龙点睛的高分效果。
  \item \textbf{本课必背文化习语与搭配}：\da{at læse mellem linjerne}（读懂弦外之音）；\da{at tage noget for pålydende}（照字面全盘当真）；\da{at pakke kritikken ind}（把批评委婉包裹起来）；\da{på godt dansk}（用直白通俗的丹麦话来说）。
\end{itemize}
\end{grammatikbox}

% ---------------------------------------------------------------------
\ncdsubsec{五、Mønsterøvelser（修辞识读与句型升级操练）}

\textbf{练习 A：低调陈述与否定反语高频搭配填空（每空填入最准确的 1 个丹麦语词）}
\begin{enumerate}
  \item Efter den vellykkede fremlæggelse sagde chefen tørt, at resultatet da var ikke \underline{\hspace{2.2cm}} dårligt.
  \item Når et milliardprojekt overskrider budgettet med halvdelen, kan det vist \underline{\hspace{2.2cm}} kaldes en triumf.
  \item Kommentatoren påpegede, at ministerens nye reformforslag ikke \underline{\hspace{2.2cm}} er uproblematisk.
  \item Projektet var langt \underline{\hspace{2.2cm}} ringe; tværtimod høstede det stor anerkendelse i branchen.
  \item For udlændinge kan danskernes skjulte koder virke en \underline{\hspace{2.2cm}} udfordrende at afkode i starten.
  \item Man bør aldrig tage en dansk underdrivelse helt \underline{\hspace{2.2cm}} pålydende uden at se på konteksten.
  \item Rapporten antyder diplomatisk, at den nuværende løsning rummer et \underline{\hspace{2.2cm}} forbedringspotentiale.
  \item Det ville naturligvis være en \underline{\hspace{2.2cm}} ubehageligt, hvis politikerne pludselig skulle tale helt uden omsvøb.
\end{enumerate}

\textbf{练习 B：低调反语纠错与修辞升级（修正语序与否定错误，或将直白句升级为地道丹麦式反语）}
\begin{enumerate}
  \item 【否定叠加纠错】\textit{Den nye aftale kan vist næppe ikke løse hele krisen på arbejdsmarkedet.} $\rightarrow$
  \item 【从句语序纠错】\textit{Hun understregede, at forslaget er ikke ligefrem uproblematisk for kommunerne.} $\rightarrow$
  \item 【词性一致纠错】\textit{Direktøren mente, at afdelingen havde en vist forbedringspotentiale på det digitale område.} $\rightarrow$
  \item 【低调褒奖升级】用 litote 改写过满的赞美：\textit{Koncerten var helt fantastisk, og musikerne spillede fremragende.}（用 \textit{ikke uden kvalitet} 与 \textit{langt fra ringe} 改写）$\rightarrow$
  \item 【冷峻反讽升级】把直白批评升级为丹麦专栏式反语：\textit{Ministerens plan var en total fiasko.}（用 \textit{kan vist næppe kaldes en triumf} 改写）$\rightarrow$
  \item 【职场委婉批评升级】把生硬否定改写为丹麦主管式反馈：\textit{Dit udkast har mange fejl og er ikke klart til mødet.}（用 \textit{en smule udfordrende} 或 \textit{et vist forbedringspotentiale} 改写）$\rightarrow$
\end{enumerate}

\textbf{练习 C：中译丹（精准运用 Litote、低调陈述与语用动词撰写 C1 句子）}
\begin{enumerate}
  \item 在丹麦，说一个结果“不算太糟”往往意味着高度认可。
  \item 当局将该项目的延误轻描淡写地称为“一次小幅调整”。
  \item 这个提案算不上毫无问题，因为它忽视了高昂的成本。
  \item 如果只按照字面意思理解这种反讽，人们就很容易误判批评的严重性。
  \item 我恐怕很难将这场仓促达成的政治妥协称之为一次胜利。
  \item 当低调陈述沦为遮盖失职的烟幕时，它便不再迷人。
\end{enumerate}

% ---------------------------------------------------------------------
\ncdsubsec{六、PD3 / Studieprøven 实战微任务与参考答案}

\begin{emailbox}{PD3 / Studieprøven 实战写作任务（Klumme / Debatindlæg · mindst 200 ord）}
\textbf{Opgave}: I Danmark betragtes brugen af underdrivelser og ironi ofte som en charmerende del af kulturen, der skaber ligeværd og en afslappet tone. Kritikere indvender derimod, at denne stil kan skabe forvirring på internationale arbejdspladser og fungere som ansvarsforflygtigelse i den politiske debat.\\[3pt]
Skriv et nuanceret \textbf{debatindlæg eller en klumme på 215--260 ord} (mindst 200 ord), hvor du diskuterer fordele og ulemper ved den danske kunst at underdrive. Inddrag gerne \textbf{mindst tre konkrete eksempler på litoter eller underdrivelser} i din tekst, og tag stilling til, hvornår klar tale er nødvendig.
\end{emailbox}

\begin{answerbox}
\textbf{Før du læser 参考回答}
\begin{enumerate}
  \item \textit{Fordi udlændinge ofte tager underdrivelser alt for bogstaveligt. Hvis en chef kalder en opgave ”ikke helt dårlig”, kan man tro, at chefen er utilfreds, selvom det betyder ros; omvendt kan ”et vist forbedringspotentiale” dække over krav om en fuldstændig omarbejdelse.}
  \item \textit{Hun mener, at underdrivelsen i politik og ledelse mister sin charme og i stedet fungerer som et tågeslør, der skjuler reelle konflikter, svigt og ansvarsforflygtigelse.}
\end{enumerate}

\textbf{Tekst B 阅读核对答案}
\begin{enumerate}
  \item \da{Forkert}（她认为在董事会会议室和政府各部里，这种魅力很快就会变得“有点吃力”，甚至成为逃避责任的烟幕）
  \item \da{Reelle konflikter, svigt og ansvarsforflygtigelse.}（真实的冲突、失职与推卸责任）
  \item \da{Lidt færre litoter og lidt mere klar tale.}（少一点曲言反语，多一点直白坦诚）
\end{enumerate}

\textbf{练习 A 答案}
\begin{enumerate}
  \item \da{helt}（高频低调褒奖固定搭配：\textit{ikke helt dårligt} 不算太糟，实为不错）
  \item \da{næppe}（冷峻反讽搭配：\textit{vist næppe kaldes en triumf} 恐怕很难称作胜利）
  \item \da{ligefrem}（强化委婉否定：\textit{ikke ligefrem uproblematisk} 算不上完全没问题）
  \item \da{fra}（反语赞赏搭配：\textit{langt fra ringe} 远谈不上差劲）
  \item \da{smule}（程度弱化短语：\textit{en smule udfordrende} 有点棘手/具挑战性）
  \item \da{for}（固定搭配：\textit{at tage noget for pålydende} 照字面意思全盘当真）
  \item \da{vist}（中性名词 \textit{et forbedringspotentiale} 前须用中性定语 \textit{et vist}）
  \item \da{kende}（书面风格度弱化副词：\textit{en kende ubehageligt} 有那么一丁点不自在）
\end{enumerate}

\newpage
\textbf{练习 B 答案（参考）}
\begin{enumerate}
  \item \textit{Den nye aftale kan vist næppe løse hele krisen på arbejdsmarkedet.}（\textit{næppe} 本身已表否定，须删去多余的 \textit{ikke}）
  \item \textit{Hun understregede, at forslaget ikke ligefrem er uproblematisk for kommunerne.}（在 \textit{at} 引导的从句中，中场否定模块 \textit{ikke ligefrem} 须置于变位动词 \textit{er} 之前）
  \item \textit{Direktøren mente, at afdelingen havde et vist forbedringspotentiale på det digitale område.}（中心词 \textit{potentiale} 为中性 \textit{et}-名词，故须配中性词形 \textit{et vist}）
  \item \textit{Koncerten var bestemt ikke uden kvalitet, og musikerne leverede en præstation, der var langt fra ringe.}（用双否曲言法替代浮夸的最高级，更贴合丹麦乐评格调）
  \item \textit{Ministerens plan kan vist næppe kaldes en triumf.}（借 \textit{vist næppe} 将直接抨击转为冷静尖锐的专栏式反讽）
  \item \textit{Dit udkast er en smule udfordrende at gå ind i endnu og rummer nok et vist forbedringspotentiale før mødet.}（利用 \textit{en smule} 与 \textit{et vist} 委婉传达大修需求，同时保留同事体面）
\end{enumerate}

\textbf{练习 C 答案}
\begin{enumerate}
  \item \textit{I Danmark betyder det ofte stor anerkendelse, når man kalder et resultat for ikke helt dårligt.}
  \item \textit{Myndighederne nedtonede projektets forsinkelse som en mindre justering.}
  \item \textit{Dette forslag er ikke ligefrem uproblematisk, fordi det overser de høje omkostninger.}
  \item \textit{Hvis man tager denne ironi for pålydende, undervurderer man let alvoren i kritikken.}
  \item \textit{Dette hastige politiske kompromis kan vist næppe kaldes en sejr.}
  \item \textit{Når underdrivelsen bliver et tågeslør over svigt, mister den sin charme.}
\end{enumerate}

\textbf{PD3 / Studieprøven 高分范文（234 词，不含标题）}\\
\textit{\textbf{Mellem lunt smil og klart sprog: Hvad vi vinder og taber ved underdrivelsen}}

\textit{Når en dansker beskriver en stor triumf som \da{”ikke helt dårlig”}, smiler udlændinge ofte over den særlige nordiske beskedenhed. \da{Langt fra alle} kulturer dyrker kunsten at nedtone sin egen succes på den måde, og i dagligdagen har denne form for selvironi en vigtig social funktion. Den skaber en uformel tone, dæmper egoer og minder os om, at ingen bør tage sig selv alt for højtideligt.}

\textit{Alligevel er denne charmerende vane \dabold{ikke ligefrem uproblematisk}, når den flytter ind i moderne, internationale virksomheder eller i den politiske debat. For nye kolleger fra udlandet kan en bemærkning om, at et projekt har ”et vist forbedringspotentiale”, let tages for pålydende som en lille detalje, selvom chefen i virkeligheden ønsker en ny retning. Dermed risikerer høfligheden at skabe unødig forvirring i stedet for tryghed.}

\textit{Endnu mere alvorligt bliver det, når magthavere bruger lune underdrivelser som et praktisk tågeslør. Hvis en alvorlig fejl i sundhedsvæsenet eller et budgetskred bliver omtalt som \dabold{”en smule udfordrende”}, svækker sproget den demokratiske ansvarlighed. En offentlig skandale \dabold{kan vist næppe kaldes} ”en mindre justering”, uden at borgerne mister tilliden til afsenderen.}

\textit{Løsningen er selvfølgelig ikke at afskaffe den danske ironi --- et samfund helt uden underspillet humor ville være \dabold{en kende trist}. Vi bør derimod skelne skarpere mellem hyggelige pauser og alvorlige beslutninger: Ved kaffemaskinen kan vi trygt lade litoterne blomstre, men når ansvaret skal placeres, skylder vi hinanden at kalde en fejl for en fejl.}

\small\textbf{名师点评：}这篇 234 词的 C1 评论文章完美展现了“知反讽而不滥用反讽”的大家风范。结构上采用四段辩证推进：肯定文化温度 $\rightarrow$ 剖析职场误读 $\rightarrow$ 批判政治责任模糊 $\rightarrow$ 提出情境区分对策。文中自然嵌入 \textit{”ikke helt dårlig”, Langt fra alle, ikke ligefrem uproblematisk, ”et vist forbedringspotentiale”, ”en smule udfordrende”, kan vist næppe kaldes}，甚至在收尾点睛一笔 \textit{en kende trist}（“那倒是有那么一丁点可悲”），精准戳中丹麦评委的审美心弦！
\end{answerbox}

\ncdlesson{6}{Lederen i søndagsavisen: Velfærd eller skattelettelser?}{周日版社论：福利优先还是减税优先？}
{DU3 Modul 5 (PD3) $\rightarrow$ Modul 6 (Studieprøven, B2 $\rightarrow$ C1)｜对标 NCE 4 \textit{Avisleder \& Avanceret Indrømmelse}}
{\textbf{报刊社论（\textit{Leder}）高阶让步句式}（\textit{Hvorvidt..., står hen i det uvisse / Uanset hvordan man vender og drejer det... / Bevares, men spørgsmålet forbliver... / Det ændrer imidlertid ikke ved den kendsgerning, at...}）· 主语从句与让步从句的 V2 倒装节奏}
{翻开丹麦各大主流报纸（如 \textit{Politiken}、\textit{Berlingske} 或 \textit{Weekendavisen}）的周日版社论（\textit{søndagsleder}），你会发现最顶级的政治评论从不简单喊口号，而是善用\textbf{“先承认对方半步，再以无可辩驳的事实将论点拉回核心”}的高阶让步句式。本课聚焦丹麦政坛永恒的财政大辩论——福利扩张还是减税激励——带你吃透 C1 社论最犀利的四大让步骨架，彻底摆脱低阶写作中单调重复的 \textit{selvom} 与 \textit{men}！}

% ---------------------------------------------------------------------
\begin{foerduleaeser}
带着下面两个宏观问题阅读 Tekst A 和 Tekst B：
\begin{enumerate}
  \item \textit{Hvorfor mener søndagslederen i Tekst A, at et ekstra økonomisk råderum ikke automatisk løser velfærdsstatens langsigtede udfordringer?}（根据 Tekst A 中的周日版社论，为什么额外的财政盈余空间并不能自动解决福利国家的长期挑战？）
  \item \textit{Hvilken indvending rejser cheføkonomen i Tekst B mod lederens måde at stille velfærd og skattelettelser op over for hinanden på?}（在 Tekst B 中，首席经济学家对社论将福利与减税对立起来的方式提出了怎样的异议？）
\end{enumerate}
\end{foerduleaeser}

% ---------------------------------------------------------------------
\ncdsubsec{二、Tekster（周日版报纸社论与经济学家反驳回应双语篇）}

\subsection*{Tekst A：\textit{Søndagens leder: Den svære kunst at prioritere}（报刊社论 · 约 360 词）}

\begin{tipbox}{语篇说明}
本篇采用地道丹麦大报社论（\textit{Leder}）的权威笔法撰写，探讨丹麦年度财政预算案（\textit{Finansloven}）谈判的核心分歧。阅读时请特别跟踪\textbf{\color{NCDNavy}深蓝色加粗的高阶让步与悬置句式}（\textit{Hvorvidt..., står hen i det uvisse / Bevares, men...}），以及\textbf{\color{NCDRed}深红色加粗的强势拉回与事实锚定框架}（\textit{Uanset hvordan man vender og drejer det... / Det ændrer imidlertid ikke ved den kendsgerning, at...}）。
\end{tipbox}

Hvert efterår, når forhandlingerne om næste års finanslov skydes i gang på Christiansborg, genopstår den samme klassiske strid i dansk politik. På den ene side kræver fagforeninger og flere partier massive investeringer i sygehuse, folkeskoler og ældrepleje for at sikre velfærdsstatens kvalitet og nedbringe ventelisterne. På den anden side argumenterer borgerlige partier og erhvervsorganisationer for markante skattelettelser, der skal belønne arbejdsindsatsen, øge beskæftigelsen og styrke Danmarks internationale konkurrenceevne.

\da{Hvorvidt} regeringen i år formår at samle et bredt flertal bag en langsigtet aftale, \da{står endnu hen i det uvisse}. De seneste måneder har budt på opjusteringer af det økonomiske råderum, hvilket naturligt har vakt store forhåbninger hos næsten alle politiske partier og interesseorganisationer. \dabold{Det ændrer imidlertid ikke ved den kendsgerning, at} det demografiske pres vokser år for år. Når stadig flere ældre lever længere, mens færre unge træder ind på arbejdsmarkedet, ædes de ekstra milliarder hurtigt op, blot for at fastholde det nuværende serviceniveau i hjemmeplejen og på hospitalerne.

\da{Bevares}, skattelettelser i bunden kan gøre det en smule mere attraktivt at tage et ekstra vagtskifte eller gå fra deltid til fuldtid. \da{Men spørgsmålet forbliver}, om en skattereform alene kan skaffe de tusindvis af sygeplejersker og pædagoger, som landets kommuner mangler allerede i dag. Omvendt må man også stille det ubehagelige modspørgsmål til velfærdens varmeste forsvarere: Hvad nytter flere penge i statskassen, hvis et højt skattetryk på sigt dæmper lysten til at starte nye virksomheder og skabe den velstand, som hele den danske model hviler på?

\dabold{Uanset hvordan man vender og drejer det}, kommer Danmark ikke uden om en ærlig samtale om prioritering. Den politiske fristelse til at love både lavere skat, grønne investeringer og mere velfærd på én gang er forståelig i en tid med skærpet vælgerkonkurrence. Alligevel risikerer kortsigtet valgflæsk at udhule tilliden til den samfundskontrakt, som gennem generationer har bundet borgere og stat sammen.

En moden offentlig debat kræver derfor mod til at vælge til og fra. Hvis fremtidens velfærd skal være økonomisk bæredygtig, må politikerne turde træffe upopulære beslutninger frem for at skubbe regningen videre til næste valgperiode. Det er først, når man anerkender, at enhver krone kun kan bruges én gang, at det politiske kompromis bliver reelt ansvarligt.

\begin{transbox}
\small
\textbf{\color{NCDTeal}[Tekst A 中文精译]}\\
每年秋天，当克里斯蒂安堡宫启动下一年度财政预算案谈判时，丹麦政坛总会重演同一场经典争论。一边是工会与多个政党，要求向医院、公立中小学和养老护理投入巨额资金，以保障福利国家的服务品质并缩短候诊名单；另一边则是中右翼政党与工商业界组织，主张推出力度可观的减税措施，以奖赏劳动付出、提升就业率并增强丹麦的国际竞争力。

政府今年究竟能否在长期协议背后凝聚起广泛多数，目前仍处于未知之数。近几个月来，政府接连上调了财政可用空间预期，这自然点燃了几乎所有政党与利益团体的热烈期待。然而，这丝毫改变不了一个冷峻的事实：人口结构带来的压力正在逐年加剧。当越来越多的老人寿命延长、而进入劳动力市场的年轻人却日益减少时，多出来的数十亿克朗仅仅为了维持居家护理与医院现有的服务水准就会被迅速吞噬殆尽。

诚然，在低收入档实施减税，或许能让多加一个班或从兼职转为全职工作显得稍具吸引力。但核心问题依然存在：单靠一场税制改革，能否填补全国各地市镇眼下就急缺的数千名护士和幼教人员？反过来，人们也必须向福利制度最热情的捍卫者提出那个令人不适的反问：如果高企的税负长远来看抑制了人们创办新企业的意愿、削弱了整个丹麦模式所赖以运转的财富创造力，那么国库里再多钱又有何用？

不管怎么左思右想、从哪个角度权衡，丹麦都无法回避一场关于优先次序的坦诚对话。在选民争夺日益激烈的当下，政治人物同时许诺“既减税、又投入绿色转型、再扩大福利”的诱惑诚然可以理解。尽管如此，短视的竞选空头支票极有可能掏空民众对社会契约的信任——正是这份契约，历代以来将公民与国家紧密联结在一起。

因此，一场成熟的公共讨论需要敢于取舍的勇气。倘若未来的福利制度要具备财政可持续性，政治家们就必须敢于作出不受欢迎的抉择，而不是把账单推给下一个选举周期。唯有当人们承认每一克朗都只能花一次时，政治妥协才真正称得上负责任。
\end{transbox}

\subsection*{Tekst B：\textit{Modsvar: Vækst og velfærd er ikke hinandens fjender}（经济学家反驳短评 · 约 155 词）}

\textbf{阅读任务（理解与核对）：}\textit{Læs cheføkonomen Thomas Lindbergs modsvar til søndagens leder, og besvar spørgsmålene nedenfor.}

\begin{emailbox}{Modsvar til søndagslederen｜Fra: Thomas Lindberg (cheføkonom) \quad Emne: En falsk modsætning mellem skattelettelser og velfærd}
Søndagens leder efterlyser mod til at prioritere, og det er svært at være uenig i. Men selve præmissen hviler på en falsk modsætning mellem velfærd og skattelettelser, som alt for ofte låser den politiske debat fast.

\da{Bevares}, en ukritisk gavebod uden dækning vil selvfølgelig skade økonomien og skabe inflation på sigt. \da{Men spørgsmålet forbliver}, hvordan vi overhovedet skaffer de hænder og de skattekroner, som fremtidens sygehuse kræver. Målrettede lettelser på skatten på arbejde får nemlig flere til at lægge ekstra timer og blive længere på arbejdsmarkedet, hvilket direkte øger både arbejdsudbuddet og produktiviteten.

\dabold{Uanset hvordan man vender og drejer det}, kan en regering ikke fordele en velstand, som ikke først er blevet skabt af dygtige medarbejdere og private virksomheder. \dabold{Det ændrer imidlertid ikke ved den kendsgerning, at} offentlige investeringer i uddannelse og infrastruktur også styrker erhvervslivet. Derfor handler finansloven ikke om enten velfærd eller skattelettelser, men om kloge reformer, der får kagen til at vokse.
\end{emailbox}

\textbf{Tekst B 阅读核对（Rigtigt / Forkert / Kort svar，答案见第六部分）：}
\begin{enumerate}
  \item \textit{Thomas er fuldstændig uenig med lederen i, at politikerne bør have mod til at prioritere.} \underline{\hspace{2.2cm}}
  \item \textit{Hvilken konkret effekt har målrettede skattelettelser på arbejde ifølge Thomas?} \underline{\hspace{2.2cm}}
  \item \textit{Hvad bør finansloven ifølge Thomas egentlig handle om i stedet for et ”enten-eller”?} \underline{\hspace{2.2cm}}
\end{enumerate}

\begin{tipbox}{本课发音与语调要点（Udtale \& Prosodi）}
\begin{itemize}
  \item \textbf{财经与社论高频词读音陷阱}：\textit{en finanslov} \ipa{fiˈnansˌlʌʊ̯}（\textit{lov} 读 \ipa{lʌʊ̯}）；\textit{et råderum} \ipa{ˈʁɔːðəˌʁʌmˀ}（软 d，末尾 \textit{-rum} 带 \textit{stød}）；\textit{en kendsgerning} \ipa{ˈkʰɛnˀsˌɡɛɐ̯neŋ}（\textit{kends-} 中 \textit{d} 不发音且带 \textit{stød}！）；\textit{et arbejdsudbud} \ipa{ˈɑːb̥ɑjˀd̥sˌuðb̥uð}（两个软 d）；\textit{hvorvidt} \ipa{vɒˈvid}（\textit{h} 与尾 \textit{d} 均不发音）。
  \item \textbf{社论语气词重音}：\textit{Bevares} \ipa{b̥eˈvɑːˀɐs}（重音在第二音节，带 \textit{stød}，后接短停顿）；\textit{uvis} \ipa{uˈvis}（重音在 \textit{-vis}）；\textit{uanset} \ipa{ˈuːænˌseːˀd̥}（重音落在否定前缀 \textit{u-} 上！）。
  \item \textbf{社论长让步句的“升降调弧线”}：朗读 \textit{Hvorvidt regeringen formår..., står hen i det uvisse} 时，从句末尾微扬表示“话未说完”，逗号停顿后主句动词 \textit{står} 稳重下沉，落到 \textit{uvisse} 收尾，传递大报评论员的笃定权威感。
\end{itemize}
\end{tipbox}

\begin{transbox}
\small
\textbf{\color{NCDTeal}[Tekst B 中文对照]}\\
\textbf{对周日社论的回应（托马斯·林德伯格，首席经济学家｜减税与福利之间是一个伪对立）：}周日版社论呼吁政治人物拿出魄力来排定优先次序，这一点很难令人反对。但文章本身的前提，却建立在福利与减税之间的一个伪对立之上——这种对立往往将政治讨论死死锁住。

诚然，毫无资金保障、不加节制的“大撒币”长远来看当然会损害经济并引发通胀。但问题依然在于，我们究竟如何才能筹集到未来医院运转所必需的人手与税收克朗。因为有针对性地降低劳动所得税，会促使更多人多加班、更长久地留在劳动力市场上，从而直接提升劳动力供给与整体生产率。

无论你怎么权衡推敲，一届政府都无法去分配尚未被优秀劳动者与私营企业创造出来的财富。不过，这同样改变不了一个事实：政府在教育与基础设施上的公共投资也会反过来增强工商业活力。因此，财政预算案并不关乎“要么福利、要么减税”的二选一，而关乎能够把蛋糕做大的明智改革。
\end{transbox}

% ---------------------------------------------------------------------
\ncdsubsec{三、Gloser（核心生词表）}

\begin{tipbox}{词表重点}
注意几个经济政论关键名词的屈折规律：\textit{et arbejdsudbud}（劳动力供给）变定指时双写 \textit{d}（\textit{arbejdsudbuddet}）；\textit{en finanslov} 的复数是 \textit{finanslove}；强变化复合动词 \textit{at fastholde}（维持/坚持）遵循 \textit{holde} 的变化（\textit{fastholdt, har fastholdt}），而 \textit{at fordele}（分配）为第 2 类弱变化（\textit{fordelte, har fordelt}）。
\end{tipbox}

\begin{glosertable}
名词 & en leder & \textit{(lederen, ledere, lederne)} & 报刊社论；领导者 \\
名词 & en finanslov & \textit{(finansloven, finanslove, finanslovene)} & 国家年度财政预算案 \\
名词 & et råderum & \textit{(råderummet, kun ental)} & 财政可用空间，回旋余地 \\
名词 & en skattelettelse & \textit{(skattelettelsen, skattelettelser, skattelettelserne)} & 减税，税收减免 \\
名词 & et skattetryk & \textit{(skattetrykket, kun ental)} & 税负，整体税收负担 \\
名词 & et arbejdsudbud & \textit{(arbejdsudbuddet, kun ental)} & 劳动力供给 \\
名词 & en kendsgerning & \textit{(kendsgerningen, kendsgerninger, kendsgerningerne)} & 客观事实，既成事实 \\
名词 & en modsætning & \textit{(modsætningen, modsætninger, modsætningerne)} & 对立，矛盾 \\
名词 & en velstand & \textit{(velstanden, kun ental)} & 繁荣，富裕程度 \\
名词 & et valgflæsk & \textit{(valgflæsket, kun ental)} & 竞选空头支票，讨好选民的承诺 \\
\midrule
动词 [har] & at prioritere & \textit{(prioriterer, prioriterede, har prioriteret)} & 排定优先次序，取舍 \\
动词 [har] & at fastholde & \textit{(fastholder, fastholdt, har fastholdt)} & 维持，保持，坚持 \\
动词 [har] & at fordele & \textit{(fordeler, fordelte, har fordelt)} & 分配，分摊 \\
动词 [har] & at rokke (ved) & \textit{(rokker, rokkede, har rokket)} & 动摇，改变（事实/立场） \\
动词 [har] & at udhule & \textit{(udhuler, udhulede, har udhulet)} & 掏空，逐渐削弱 \\
\midrule
形容词 & uvis & \textit{(uvist, uvisse)} & 不确定的，未知的 \\
形容词 & kortsigtet & \textit{(kortsigtet, kortsigtede)} & 短视的，短期的 \\
连词 & uanset & \textit{(从属连词/介词，不变)} & 无论，不管（后接 \textit{hv-} 词或名词） \\
\end{glosertable}

% ---------------------------------------------------------------------
\ncdsubsec{四、Grammatik \& Stilistik（报刊社论高阶让步句法精讲）}

\begin{grammatikbox}{修辞点 1：从 B1 到 C1——让步反驳句式的三级进阶矩阵}
在 PD3 与 Studieprøven 写作（尤其是评论文与社论回复）中，如何优雅地“承认对手有理之处，但不被对手带偏结论”是核心评分指标：

\vspace{0.3em}
\renewcommand{\arraystretch}{1.24}
{\small
\begin{tabularx}{\linewidth}{@{} >{\bfseries\raggedright\arraybackslash}p{2.2cm} >{\raggedright\arraybackslash}p{5.5cm} >{\raggedright\arraybackslash}X @{}}
\toprule
\textbf{能力层级} & \textbf{让步—拉回骨架公式} & \textbf{语用效果与适用场景} \\
\midrule
B1 基础级 & \textit{Selvom A er rigtigt, mener jeg B.} & 略显单薄平铺，难以支撑大报社论的辩证张力。 \\
\midrule
B2 过渡级\newline （NCD 4） & \textit{Ganske vist..., men ikke desto mindre...} & 清晰工整，适合标准读者辩论投稿（\textit{Debatindlæg}）。 \\
\midrule
\da{C1 社论级 1}\newline （悬置争议） & \textit{\da{Hvorvidt} [从句], \da{står hen i det uvisse}.} & 展现老练评论员的审慎，不妄下断语，先标定不确定性再深挖结构问题。 \\
\midrule
\da{C1 社论级 2}\newline （退一步反问） & \textit{\da{Bevares}, [让步句]. \da{Men spørgsmålet forbliver}, om/hvordan...} & 先以“诚然/好吧”大方认账，随即把焦点锁死在对方回避的关键追问上。 \\
\midrule
\dabold{C1 社论级 3}\newline （总括全盘） & \textit{\dabold{Uanset hvordan man vender og drejer det}, [V2 主句]...} & 无论从哪个角度推敲都绕不开的硬核结论，气势浑厚。 \\
\midrule
\dabold{C1 社论级 4}\newline （锚定铁律） & \textit{\dabold{Det ændrer / rokker imidlertid ikke ved den kendsgerning, at}...} & 无论前文让步多少，最终用“然而这改变/动摇不了一个铁的事实……”一锤定音。 \\
\bottomrule
\end{tabularx}
}
\end{grammatikbox}

\begin{grammatikbox}{语法点 2：四大 C1 社论句式的语序、介词与主谓铁律}
\begin{enumerate}
  \item \textbf{\da{Hvorvidt..., står hen i det uvisse}：整个从句打包充当主语！}\\
    由 \textit{Hvorvidt} 引导的间接疑问从句位于句首时，整个从句作为一个“单一抽象主语（第 1 顺位）”，后面直接接单数主句变位动词 \da{står}；同时 \textit{Hvorvidt} 从句内部绝不倒装：\\
    \ding{51}\ \textit{\da{Hvorvidt} forslaget \dabold{faktisk virker}, \da{står} stadig hen i det uvisse.}（\ding{55}\ \textit{...hvorvidt virker forslaget...} 或在逗号后加多余的 \textit{det}：除非采用左移外置，社论书面体绝不重复 \textit{det}！）
  \item \textbf{\dabold{Uanset hvordan man vender og drejer det,} + 必触发 V2 倒装！}\\
    由 \textit{uanset hvordan / uanset om / uanset hvor meget} 引导的让步从句整体占据主句前场，紧随逗号之后必须是\textbf{主句变位动词}（V2）：\\
    \ding{51}\ \textit{Uanset hvordan man vender og drejer det, \dabold{kræver} velfærden reformer.}（\ding{55}\ \textit{...det, velfærden kræver reformer.}）
  \item \textbf{\dabold{Det ændrer / rokker imidlertid ikke ved den kendsgerning, at...} 的介词锁定}\\
    动词 \textit{at ændre}（此处表“对……有所改变”）与 \textit{at rokke}（动摇）在此搭配介词 \dabold{ved}！不可遗漏 \textit{ved}，也不可用 \textit{på/af}；副词 \textit{imidlertid} 与 \textit{ikke} 依次站在主句动词之后，后面的同位语从句由 \textit{at} 引导，保持正常从句语序：\\
    \ding{51}\ \textit{Det ændrer imidlertid ikke \dabold{ved} den kendsgerning, at udgifterne \dabold{fortsat stiger}.}
  \item \textbf{\da{Bevares,} 的语体与标点位置}\\
    \textit{Bevares} 原意近乎“姑且算如此吧 / 诚然”，作句首独立的让步插话语，后面打逗号，随后接正常语序的让步句（注意：\textit{Bevares} 单独加逗号时，其后的主句\textbf{不倒装}！例如：\textit{Bevares, skattelettelser \da{kan} give et lille løft.}）。
\end{enumerate}
\end{grammatikbox}

\begin{grammatikbox}{社论写作排雷与高频政经搭配（Stilistik）}
\begin{itemize}
  \item \textbf{区分 \textit{at råde over}（拥有/支配）与 \textit{at råde til}（建议）}：财政新闻里说“政府拥有更大的财政空间”，用 \da{råde over}（如 \textit{Staten råder over et større råderum}）；而说“建议政府推迟减税”，才用 \da{råde ... til}。
  \item \textbf{避免“假二选一”的社论套话}：当你想批评辩论双方把复杂问题过度简化时，地道表达是 \da{en falsk modsætning}（伪对立）或 \da{et forsimplet enten-eller}（过于简化的非此即彼）。
  \item \textbf{本课社论高频修辞动词与短语}：\da{at skyde forhandlingerne i gang}（启动谈判）；\da{at komme uden om noget}（回避掉某事）；\da{at vælge til og fra}（做出取舍）；\da{at skubbe regningen videre}（把账单推给未来）；\da{at hvile på et spinkelt grundlag}（建立在薄弱基础之上）。
\end{itemize}
\end{grammatikbox}

% ---------------------------------------------------------------------
\ncdsubsec{五、Mønsterøvelser（社论让步句式与语序操练）}

\textbf{练习 A：高阶社论让步结构与政经搭配填空（每空填入最准确的 1 个丹麦语词）}
\begin{enumerate}
  \item \underline{\hspace{2.2cm}} politikerne kan nå frem til et bredt forlig inden jul, er endnu usikkert.
  \item Hvad der sker med reformen efter næste folketingsvalg, står fortsat hen i det \underline{\hspace{2.2cm}}.
  \item \underline{\hspace{2.2cm}} hvordan man vender og drejer det, må enhver krone tjenes, før den kan bruges.
  \item Uanset hvordan man vender og \underline{\hspace{2.2cm}} det, vokser behovet for arbejdskraft i plejesektoren.
  \item \underline{\hspace{2.2cm}}, en lille skattelettelse kan motivere nogle, men løser den manglen på sygeplejersker?
  \item Bevares, forslaget lyder sympatisk, men spørgsmålet \underline{\hspace{2.2cm}}, hvordan det skal finansieres.
  \item Det ændrer imidlertid ikke ved den \underline{\hspace{2.2cm}}, at antallet af ældre borgere stiger markant.
  \item Nye tal for væksten rokker ikke \underline{\hspace{2.2cm}} den kendsgerning, at kommunerne mangler pædagoger.
\end{enumerate}

\textbf{练习 B：社论语序纠错与句式升格（修正从句语序、介词及 V2 错误，或升级为 C1 社论句式）}
\begin{enumerate}
  \item 【主语从句语序纠错】\textit{Hvorvidt vil regeringen prioritere skattelettelser i år, står endnu hen i det uvisse.} $\rightarrow$
  \item 【让步从句后 V2 纠错】\textit{Uanset hvordan man vender og drejer det, vi kan ikke undgå svære prioriteringer.} $\rightarrow$
  \item 【介词与从句语序纠错】\textit{Det ændrer imidlertid ikke på den kendsgerning, at stiger udgifterne til sundhed hvert år.} $\rightarrow$
  \item 【B1 升格为 C1 社论悬置句】用 \textit{Hvorvidt..., står hen i det uvisse} 升级：\textit{Vi ved ikke endnu, om flere arbejdstimer kan betale hele regningen.} $\rightarrow$
  \item 【B1 升格为 C1 强让步反问】用 \textit{Bevares, ... Men spørgsmålet forbliver, ...} 升级：\textit{Selvom kortere ventelister er gode, ved vi stadig ikke, hvor pengene skal komme fra.} $\rightarrow$
  \item 【B2 升格为 C1 事实锚定】用 \textit{Det ændrer imidlertid ikke ved den kendsgerning, at...} 升级：\textit{Selvom råderummet er steget, er det stadig sandt, at befolkningen ældes.} $\rightarrow$
\end{enumerate}

\textbf{练习 C：中译丹（综合调用周日社论高阶让步句型与政经词汇）}
\begin{enumerate}
  \item 政府是否能达成一项广泛协议，目前仍处于未知之数。
  \item 无论怎么权衡推敲，社会都必须面对艰难的优先次序抉择。
  \item 诚然，减税可以刺激经济增长，但问题依然在于谁来照料老年人。
  \item 然而，这改变不了劳动力供给正在逐年萎缩的事实。
  \item 短视的竞选空头支票长远来看会掏空民众对社会契约的信任。
  \item 一场成熟的辩论要求政治家们敢于做出清晰取舍，而不是回避问题。
\end{enumerate}

% ---------------------------------------------------------------------
\ncdsubsec{六、PD3 / Studieprøven 实战微任务与参考答案}

\begin{emailbox}{PD3 / Studieprøven 实战写作任务（Avisleder / Debatindlæg · mindst 200 ord）}
\textbf{Opgave}: I den danske valgkamp foreslår nogle partier at bruge det økonomiske råderum på brede skattelettelser, mens andre ønsker at investere hele beløbet i velfærd, skoler og sundhed.\\[3pt]
Skriv en afbalanceret, men markant \textbf{kommentar eller et debatindlæg i lederstil på 220--260 ord} (mindst 200 ord), hvor du tager stilling til dilemmaet mellem velfærd og skattelettelser. Du skal bruge mindst \textbf{tre af lektionens avancerede indrømmelsesstrukturer} (fx \textit{Hvorvidt..., står hen i det uvisse / Bevares, men spørgsmålet forbliver / Uanset hvordan man vender og drejer det / Det ændrer imidlertid ikke ved den kendsgerning, at...}).
\end{emailbox}

\begin{answerbox}
\textbf{Før du læser 参考回答}
\begin{enumerate}
  \item \textit{Fordi det demografiske pres vokser år for år: Når flere ældre lever længere, og færre unge træder ind på arbejdsmarkedet, bliver de ekstra milliarder hurtigt brugt blot på at fastholde det nuværende serviceniveau.}
  \item \textit{Han mener, at lederen opstiller en falsk modsætning mellem velfærd og skattelettelser. Målrettede skattelettelser på arbejde kan øge arbejdsudbuddet, som netop skaber den velstand og de skattekroner, der finansierer fremtidens sygehuse.}
\end{enumerate}

\textbf{Tekst B 阅读核对答案}
\begin{enumerate}
  \item \da{Forkert}（他一开头就指出很难不同意社论关于“需要有勇气排定优先次序”的呼吁）
  \item \da{De får flere til at lægge ekstra timer og blive længere på arbejdsmarkedet, hvilket øger arbejdsudbuddet.}
  \item \da{Kloge reformer, der får kagen til at vokse.}（关于如何把蛋糕做大的明智改革）
\end{enumerate}

\textbf{练习 A 答案}
\begin{enumerate}
  \item \da{Hvorvidt}（正式书面主语从句引导词，“究竟是否……”）
  \item \da{uvisse}（固定社论表达：\textit{at stå hen i det uvisse} 尚无定论、处于未知之中）
  \item \da{Uanset}（固定让步从句开头：\textit{Uanset hvordan man vender og drejer det} 无论怎么看）
  \item \da{drejer}（成语搭配：\textit{at vende og dreje} 反复推敲、从正反面审视）
  \item \da{Bevares}（高阶口语化让步词进入社论体，“诚然、姑且承认”）
  \item \da{forbliver}（核心追问过渡：\textit{men spørgsmålet forbliver, ...} 但问题依然在于）
  \item \da{kendsgerning}（固定强调词：\textit{ved den kendsgerning, at...} ……这一既成事实）
  \item \da{ved}（动词 \textit{at rokke} 与 \textit{at ændre} 在此结构中固定搭配介词 \textit{ved}）
\end{enumerate}

\textbf{练习 B 答案（参考）}
\begin{enumerate}
  \item \textit{Hvorvidt regeringen vil prioritere skattelettelser i år, står endnu hen i det uvisse.}（\textit{Hvorvidt} 引导的间接疑问从句内部须维持“主语 + 动词”不倒装语序）
  \item \textit{Uanset hvordan man vender og drejer det, kan vi ikke undgå svære prioriteringer.}（前置让步从句占据前场，主句须紧跟变位动词 \textit{kan} 实现 V2 倒装）
  \item \textit{Det ændrer imidlertid ikke ved den kendsgerning, at udgifterne til sundhed stiger hvert år.}（固定搭配介词为 \textit{ved}，且 \textit{at} 从句内须恢复“主语 + 动词”不倒装语序）
  \item \textit{Hvorvidt flere arbejdstimer kan betale hele regningen, står fortsat hen i det uvisse.}（用社论悬置句取代平铺直叙的 \textit{Vi ved ikke}，瞬间提升客观权威感）
  \item \textit{Bevares, kortere ventelister er et vigtigt mål. Men spørgsmålet forbliver, hvor pengene skal komme fra.}（用“诚然退让 + 核心追问”拆分长句，增强社论节奏与锋芒）
  \item \textit{Det ændrer imidlertid ikke ved den kendsgerning, at befolkningen bliver ældre, selvom råderummet er steget.}（用不可动摇的事实锚定句强化论证重心）
\end{enumerate}

\textbf{练习 C 答案}
\begin{enumerate}
  \item \textit{Hvorvidt regeringen kan indgå en bred aftale, står stadig hen i det uvisse.}
  \item \textit{Uanset hvordan man vender og drejer det, må samfundet foretage svære prioriteringer.}
  \item \textit{Bevares, skattelettelser kan styrke væksten, men spørgsmålet forbliver, hvem der skal passe de ældre.}
  \item \textit{Det ændrer imidlertid ikke ved den kendsgerning, at arbejdsudbuddet falder år for år.}
  \item \textit{Kortsigtet valgflæsk udhuler på længere sigt borgernes tillid til vores fælles samfundskontrakt.}
  \item \textit{En moden debat kræver, at politikerne tør prioritere frem for at undgå svære valg.}
\end{enumerate}

\textbf{PD3 / Studieprøven 高分范文（220 词，不含标题）}\\
\textit{\textbf{Velfærd og vækst: Tid til at droppe de lette løfter}}

\textit{Når politikerne fordeler næste års finanslov, lyder den gamle debat ofte som en simpel kamp mellem varme hænder og lavere skat. \da{Hvorvidt} næste folketing formår at bryde denne fastlåste konflikt og skabe en holdbar plan for det næste årti, \da{står imidlertid hen i det uvisse}.}

\textit{\da{Bevares}, det er fuldt forståeligt, at mange borgere efterspørger skattelettelser efter år med stigende leveomkostninger. Flere penge mellem hænderne kan give familier større frihed og gøre det mere attraktivt at arbejde ekstra timer. \da{Men spørgsmålet forbliver}, om en bred skatterabat faktisk løser den akutte mangel på personale i vores sygehuse, skoler og daginstitutioner, hvor nedslidning og travlhed allerede presser hverdagen.}

\textit{Omvendt har økonomerne en vigtig pointe, når de advarer mod at betragte statens kasse som bundløs. \dabold{Uanset hvordan man vender og drejer det}, kan et velfærdssamfund ikke fordele større summer, end erhvervslivet og borgerne i fællesskab formår at tjene ind. Et højt serviceniveau forudsætter derfor et stærkt arbejdsudbud og sunde rammer for virksomhederne.}

\textit{Derfor bør vi afvise kortsigtet valgflæsk, der lover begge dele uden finansiering. Nye reformer kan måske skabe et midlertidigt økonomisk råderum, men \dabold{det ændrer ikke ved den kendsgerning, at} en aldrende befolkning kræver benhårde prioriteringer. Kun hvis politikerne tør være ærlige om, hvad der skal vælges til --- og hvad der må vælges fra --- kan vi bevare både velstanden og tilliden til velfærdsstaten.}

\small\textbf{名师点评：}这篇 220 词的社论评论体范文极具丹麦主流媒体（如 \textit{Weekendavisen}）质感：开篇用 \textit{Hvorvidt..., står imidlertid hen i det uvisse} 设定悬念；第二段先以 \textit{Bevares} 共情民众对减税的诉求，随即用 \textit{Men spørgsmålet forbliver, om...} 切入公共服务人手短缺的痛点；第三段公正接纳经济学家视角，以 \textit{Uanset hvordan man vender og drejer det} 亮出财政底线；结尾段则以 \textit{det ændrer ikke ved den kendsgerning, at...} 锚定人口老龄化现实，并呼吁政治诚实。全文语法零瑕疵，立场立体深邃，稳居 C1 高分水准！
\end{answerbox}

\ncdlesson{7}{Ekspertvælde eller sund fornuft?}{专家治国还是常识至上？}
{DU3 Modul 5 (PD3) $\rightarrow$ Modul 6 (Studieprøven, B2 $\rightarrow$ C1)｜对标 NCE 4 \textit{Kildekritik, Referat \& Holdningsspektrum}}
{\textbf{引述动词的态度立场光谱}：客观中立（\textit{konstatere, påpege, anføre}）vs 保持距离/质疑（\textit{hævde, påstå, foregive}）vs 认同强调（\textit{understrege, slå fast, dokumentere}）· 被动引述、倒装插叙与 \textit{ifølge} 避坑}
{在 PD3 与 Studieprøven（特别是多源材料综合写作 \textit{Syntese} 与阅读题态度判定）中，最能体现 C1 学术素养的不是“你说了什么”，而是\textbf{“你如何用动词精确标定他人观点的可信度”}。写 \textit{eksperterne siger} 只能拿 B1 分数；懂得区分 \textit{rapporten dokumenterer}（铁证如山）、\textit{kommissionen konstaterer}（客观陈述）与 \textit{kritikere påstår}（存疑与保持距离），才能在专家报告与民主决策的争论中精准把控立场温度！}

% ---------------------------------------------------------------------
\begin{foerduleaeser}
带着下面两个宏观问题阅读 Tekst A 和 Tekst B：
\begin{enumerate}
  \item \textit{Hvordan fordeles rollerne mellem ekspertkommissioner og folkevalgte politikere ifølge Tekst A, og hvorfor advarer forfatteren mod rent ekspertvælde?}（根据 Tekst A，专家委员会与民选政治人物之间的角色应当如何分工？作者为什么警惕纯粹的“专家治国”？）
  \item \textit{Hvilken indvending har kommunalpolitikeren Sofie i Tekst B mod at lade eksperternes beregninger styre politiske beslutninger alene?}（在 Tekst B 中，地方议员索菲对于单纯由专家测算来主导政治决策提出了什么异议？）
\end{enumerate}
\end{foerduleaeser}

% ---------------------------------------------------------------------
\ncdsubsec{二、Tekster（专家报告评论与电台辩论双语篇）}

\subsection*{Tekst A：\textit{Mellem regneark og folkestyre: Hvem bør bestemme?}（政论与信源分析 · 约 360 词）}

\begin{tipbox}{语篇说明}
本篇探讨现代福利国家中“专家治国（\textit{ekspertvælde}）”与“民主常识（\textit{sund fornuft}）”的张力。阅读时请高度留意引述动词的色彩分层：\textbf{\color{NCDNavy}深蓝色加粗的客观中立与权威认同动词}（\textit{konstatere, påpege, anføre, understrege, slå fast, dokumentere}），以及\textbf{\color{NCDRed}深红色加粗的保持距离与质疑动词}（\textit{hævde, påstå, foregive}）。
\end{tipbox}

Når store velfærdsreformer skal forberedes på Christiansborg, nedsætter regeringen næsten altid en kommission af uafhængige eksperter. Således også sidste måned, da en gruppe førende økonomer, læger og samfundsforskere udgav en omfattende rapport om fremtidens sundhedsvæsen og pensionsalder. Rapporten \da{konstaterer} tørt, at antallet af ældre borgere med kroniske sygdomme vil stige markant frem mod 2040, og i analysen \da{påpeges} det, at det nuværende personalesystem simpelthen ikke kan følge med den voksende efterspørgsel.

Ikke overraskende har udspillet udløst en intens offentlig debat, hvor forskellige aktører læser konklusionerne på hver deres måde. Kommissionens formand \da{understreger} i flere interview, at anbefalingerne bygger på ti års systematisk dataindsamling. En række uafhængige forskere \da{dokumenterer} da også gennem internationale sammenligninger, at en gradvis centralisering af højt specialiserede afdelinger mindsker dødeligheden ved komplicerede operationer og forbedrer behandlingen. På den baggrund \da{slår} rapporten \da{fast}, at strukturelle reformer ikke længere kan udskydes, hvis den høje kvalitet skal bevares.

Fra landets yderområder lyder der imidlertid helt andre toner. Lokale borgmestre \dabold{hævder}, at de økonomiske modeller hviler på teoretiske antagelser fra København, som overser virkeligheden i hverdagen. Flere patientforeninger \dabold{påstår} ligefrem, at kommissionen har valgt sine tal ud for at retfærdiggøre besparelser på mindre sygehuse. I debatten \da{anfører} modstanderne desuden, at længere transporttid til den nærmeste akutmodtagelse skaber utryghed for ældre og børnefamilier --- en menneskelig faktor, som sjældent vejer tungt i et regneark.

Striden blotlægger et grundlæggende demokratisk dilemma mellem ekspertvælde og folkelig legitimitet. På den ene side kan ingen ansvarlig politiker \dabold{foregive} at ignorere faglig viden eller tro, at sund fornuft alene kan erstatte langsigtede prognoser. På den anden side må man anerkende, at ekspertise har en klar grænse: Økonomer kan beregne prisen på en reform og forudsige manglen på hænder, men de kan ikke afgøre, om lighed mellem land og by bør vægtes højere end ren effektivitet.

I et levende demokrati har eksperter altså til opgave at oplyse det faglige grundlag --- ikke at træffe de endelige værdivalg. Når politikere bruger rapporter som et skjold og lader, som om der kun findes én ”nødvendig” løsning, svækkes det folkelige mandat. Kun når faglig dokumentation møder åben politisk debat og borgernes levede erfaring, opstår beslutninger, som både er kloge og demokratisk legitime.

\begin{transbox}
\small
\textbf{\color{NCDTeal}[Tekst A 中文精译]}\\
每当克里斯蒂安堡宫准备推行重大福利改革时，政府几乎总会设立一个由独立专家组成的委员会。上个月也不例外：一组顶尖经济学家、医学专家与社会学者发布了一份关于未来医疗体系与退休年龄的全面报告。报告冷静地指出（konstatere），到 2040 年患有慢性病的老年公民人数将显著上升；分析中还点明（påpege），现有的人员配置体系根本无法跟上不断增长的需求。

不出所料，这份方案引发了激烈的公共辩论，各方行动者以截然不同的方式解读其结论。委员会主席在多次采访中着重强调（understrege），这些建议建立在十年系统性数据积累的基础之上。一批独立学者也确实通过国际对比证实（dokumentere），逐步集中专科科室能够降低复杂手术的死亡率并提升疗效。在此基础上，报告明确断定（slå fast），若要保住高水准医疗质量，结构性改革已刻不容缓。

然而，来自全国偏远地区的声音却截然不同。地方市长们宣称（hævde，暗含未获充分认同），这些经济模型建立在来自哥本哈根的理论假设之上，忽视了日常生活的现实。多个患者协会甚至指责/声称（påstå，带强烈质疑），委员会刻意挑选数据，只为替削减小型医院经费辩护。在辩论中，反对者还列举理由称（anføre），前往最近急诊科的更长通勤时间会给老人和有孩子的家庭带来不安全感——而这一人本因素在电子表格里很少占分量。

这场争论揭示了一个根本性的民主两难：专家治国与民众正当性之间的张力。一方面，任何负责任的政治人物都不能假装（foregive）自己可以无视专业知识，或以为单靠常识就能替代长期预测。另一方面，人们也必须承认专业知识存在清晰的边界：经济学家能够测算出某项改革的成本并预测人手缺口，但他们无法代替公民决定城乡平等是否应当比单纯的效率更受重视。

因此，在一个充满活力的民主制度中，专家的任务在于阐明专业事实基础——而不是代替人民作出最终的价值抉择。当政治人物把专家报告当作挡箭牌，仿佛世上只存在唯一一种“别无选择的必然方案”时，民众授予的政治授权便会被削弱。唯有当专业实证遇上开放的政治辩论以及市民的切身生活经验时，既明智又具民主正当性的决策才会诞生。
\end{transbox}

\subsection*{Tekst B：\textit{Radiodebat: Skal tallene eller vælgerne have det sidste ord?}（电台双人辩论 · 约 145 词）}

\textbf{阅读任务（理解与核对）：}\textit{Læs den korte radiodebat mellem politologen Mikkel og kommunalpolitikeren Sofie, og besvar spørgsmålene nedenfor.}

\begin{emailbox}{Radiodebat｜Mikkel (politolog) \& Sofie (kommunalpolitiker)}
\textbf{Mikkel:} Jeg forstår godt skepsissen ude i de mindre kommuner, Sofie. Men kommissionens tal \da{dokumenterer} jo ret tydeligt, at de små sygehuse allerede i dag mangler erfarne specialister og moderne udstyr. Vi kan ikke bare \dabold{foregive}, at alt kan fortsætte som før.\\
\textbf{Sofie:} Nu \dabold{påstår} du igen, at kritikerne blot vil lukke øjnene for virkeligheden, Mikkel. Jeg \da{anfører} jo kun, at en kold matematisk model ikke kan måle den daglige tryghed ved at have en lokal skadestue i nærheden af borgerne.\\
\textbf{Mikkel:} Det \da{påpeger} rapporten faktisk også i kapitel fire, hvor netop borgernes oplevede tryghed nævnes. Men analysen \da{slår} samtidig \da{fast}, at høj faglig kvalitet og patientsikkerhed kræver større og stærkere enheder.\\
\textbf{Sofie:} Fair nok, men I eksperter \dabold{hævder} alt for ofte, at jeres anbefalinger er rent neutrale. I praksis er enhver samfundsmæssig prioritering et politisk værdivalg, som vælgerne og deres repræsentanter skal tage stilling til.
\end{emailbox}

\textbf{Tekst B 阅读核对（Rigtigt / Forkert / Kort svar，答案见第六部分）：}
\begin{enumerate}
  \item \textit{Mikkel mener, at alle de små sygehuse sagtens kan bevare den høje faglige kvalitet uden ændringer.} \underline{\hspace{2.2cm}}
  \item \textit{Hvad kan en matematisk model ifølge Sofie ikke måle i debatten om de små sygehuse?} \underline{\hspace{2.2cm}}
  \item \textit{Hvad er enhver prioritering i praksis ifølge Sofie, selv når eksperter kalder deres anbefalinger neutrale?} \underline{\hspace{2.2cm}}
\end{enumerate}

\begin{tipbox}{本课发音与语调要点（Udtale \& Prosodi）}
\begin{itemize}
  \item \textbf{引述动词重音与不规则变位读音}：\textit{konstatere} \ipa{kʰʌnsd̥aˈtˢeːˀɐ}（重音在 \textit{-te-}）；\textit{påpege} \ipa{ˈpʰʌˌpʰɑjˀə}（\textit{-ege} 读 \ipa{ɑjˀə}）；\textit{hævde} \ipa{ˈhɛʊ̯də}（\textit{v} 滑音化为 \ipa{ʊ̯}，\textit{d} 仅在过去式或重音下微弱过渡）；\textit{påstå} \ipa{ˈpʰʌˌsd̥ɔːˀ}（长元音带 \textit{stød}）；\textit{foregive} \ipa{ˈfɒːɐˌɡ̊iːˀ}（口语常略去 \textit{-ve} 尾音）。
  \item \textbf{民主与报告高频词}：\textit{et ekspertvælde} \ipa{ɛɡ̊sˈb̥ɛɐ̯d̥ˌvɛlə}（\textit{-vælde} 中的 \textit{d} 不发音！）；\textit{legitimitet} \ipa{leɡ̊itˢimiˈtˢeːˀd̥}（重音在最后一个音节）；\textit{et grundlag} \ipa{ˈɡ̊ʁʌnˌlæːˀ}（\textit{grund} 中的 \textit{d} 不发音！）。
  \item \textbf{质疑语调 vs 实证语调}：读 \textit{påstå} 和 \textit{foregive} 时重音略带拖长和警觉感；读 \textit{dokumenterer} 和 \textit{slår fast} 时声音下沉、干脆利落，体现论断的不可动摇。
\end{itemize}
\end{tipbox}

\begin{transbox}
\small
\textbf{\color{NCDTeal}[Tekst B 中文对照]}\\
\textbf{米克尔（政治学者）：}索菲，我完全理解小市镇里的怀疑态度。但委员会的数据毕竟相当清楚地证实（dokumentere），小型医院如今就已经缺乏经验丰富的专科医生和现代设备。我们不能只是假装（foregive）一切都能照旧继续下去。\\
\textbf{索菲（地方议员）：}米克尔，你现在又声称（påstå，带反感质疑）批评者只是想对现实视而不见。我毕竟只是提出依据（anføre），说明一套冷冰冰的数学模型衡量不出在居民附近拥有一家本地急诊室所带来的日常安全感。\\
\textbf{米克尔：}报告其实在第四章也确实指出了（påpege）这一点，那里专门提到了公民感受到的安全感。但分析同时明确指出（slå fast），高水平的专业医疗质量与患者安全需要规模更大、实力更强的医疗单元。\\
\textbf{索菲：}好吧，但你们专家实在太频繁地宣称（hævde）自己的建议是纯粹中立的。在实践中，每一次社会层面的优先次序取舍都是政治价值抉择，理应由选民及其代表来表态。
\end{transbox}

% ---------------------------------------------------------------------
\ncdsubsec{三、Gloser（核心生词表）}

\begin{tipbox}{词表重点}
本课表中的 8 个引述动词包含多个强变化不规则动词：\textit{at påstå} 随 \textit{stå} 变（\textit{påstod, har påstået}）；\textit{at foregive} 随 \textit{give} 变（\textit{foregav, har foregivet}）；\textit{at slå fast} 随 \textit{slå} 变（\textit{slog fast, har slået fast}）。注意名词 \textit{en antagelse}（假设）与 \textit{en legitimitet}（正当性/合法性）的阴/共性冠词搭配。
\end{tipbox}

\begin{glosertable}
名词 & et ekspertvælde & \textit{(ekspertvældet, kun ental)} & 专家治国，专家专政（含贬义） \\
名词 & en kommission & \textit{(kommissionen, kommissioner, kommissionerne)} & 委员会，专门调查组 \\
名词 & en antagelse & \textit{(antagelsen, antagelser, antagelserne)} & 假设，前提假定 \\
名词 & et værdivalg & \textit{(værdivalget, værdivalg, værdivalgene)} & 价值抉择，价值观选择 \\
名词 & en legitimitet & \textit{(legitimiteten, kun ental)} & 正当性，合法性 \\
名词 & et mandat & \textit{(mandatet, mandater, mandaterne)} & 授权，委任，议席 \\
名词 & et grundlag & \textit{(grundlaget, grundlag, grundlagene)} & 基础，依据，底稿 \\
名词 & en skepsis & \textit{(skepsissen, kun ental)} & 怀疑态度，审慎质疑 \\
\midrule
动词 [har] & at konstatere & \textit{(konstaterer, konstaterede, har konstateret)} & 冷静指出，确认（事实） \\
动词 [har] & at påpege & \textit{(påpeger, påpegede, har påpeget)} & 指出，提醒注意 \\
动词 [har] & at anføre & \textit{(anfører, anførte, har anført)} & 援引，列举（论据/理由） \\
动词 [har] & at hævde & \textit{(hævder, hævdede, har hævdet)} & 宣称，主张（暗含存疑） \\
动词 [har] & at påstå & \textit{(påstår, påstod, har påstået)} & 声称，硬说（强烈质疑） \\
动词 [har] & at foregive & \textit{(foregiver, foregav, har foregivet)} & 假装，佯称，伪装 \\
动词 [har] & at dokumentere & \textit{(dokumenterer, dokumenterede, har dokumenteret)} & 证明，用证据证实 \\
动词 [har] & at slå fast & \textit{(slår fast, slog fast, har slået fast)} & 明确指出，断定，确认 \\
\midrule
形容词 & tvivlsom & \textit{(tvivlsomt, tvivlsomme)} & 可疑的，靠不住的 \\
形容词 & vidtrækkende & \textit{(vidtrækkende, 不变)} & 影响深远的，波及广泛的 \\
\end{glosertable}

% ---------------------------------------------------------------------
\ncdsubsec{四、Grammatik \& Stilistik（引述动词立场光谱与信源语法精讲）}

\begin{grammatikbox}{修辞点 1：引述动词的三阶态度立场光谱（\textit{Holdningsspektrum}）}
在丹麦语学术综述与政论写作中，仅仅把 \textit{sige / mene} 换成高级动词还不够——你必须精准掌控每一个引述动词向读者传递的\textbf{“信任度信号”}：

\vspace{0.3em}
\renewcommand{\arraystretch}{1.24}
{\small
\begin{tabularx}{\linewidth}{@{} >{\bfseries\raggedright\arraybackslash}p{2.7cm} >{\raggedright\arraybackslash}p{4.6cm} >{\raggedright\arraybackslash}X @{}}
\toprule
\textbf{立场光谱段} & \textbf{核心动词与细微语义差} & \textbf{典型语境与立场潜台词} \\
\midrule
\da{1. 客观中立}\newline （Neutralt referat） &
\textit{konstatere}（冷静确认客观现状）\newline
\textit{påpege}（点出容易被忽略的要点）\newline
\textit{anføre}（有理有据地列举理由） &
\textit{Kommissionen \da{konstaterer}, at...} / \textit{Modstandere \da{anfører}, at...}\newline
作者保持记者般的客观中立，既不替对方背书，也不主动拆台。 \\
\midrule
\dabold{2. 保持距离与质疑}\newline （Distance \& skepsis） &
\textit{hævde}（声称、自称，但尚未证实）\newline
\textit{påstå}（强硬宣称、嘴硬，带明显反驳）\newline
\textit{foregive}（假装、佯称，直接戳穿虚伪） &
\textit{Kritikere \dabold{hævder}, at...} / \textit{Han \dabold{påstår} ligefrem, at...}\newline
向读者暗示：“这是他们单方面的说法，作者本人并不买账（或认为证据不足）。” \\
\midrule
\da{3. 认同与实证强调}\newline （Tilslutning \& tyngde） &
\textit{understrege}（着重强调重点）\newline
\textit{slå fast / fastslå}（斩钉截铁断定）\newline
\textit{dokumentere}（用扎实数据/文献证实） &
\textit{Tallene \da{dokumenterer}, at...} / \textit{Analysen \da{slår fast}, at...}\newline
作者强烈认同该信源的权威性与客观真实性，将结论当作稳固基石。 \\
\bottomrule
\end{tabularx}
}
\end{grammatikbox}

\begin{grammatikbox}{语法点 2：高阶引述的四大句法变体（被动引述、插叙倒装与可分动词位置）}
\begin{enumerate}
  \item \textbf{无人称被动引述（\textit{Det + -s passiv / blive-passiv + at...}）}：\\
    当你想聚焦“报告里提出了什么”，而非反复重复主语“专家们”时，使用形式主语 \textit{det} 搭配被动语态，是 C1 政论与报告摘要的黄金句型：\\
    \ding{51}\ \textit{I rapporten \da{påpeges} det, at...} / \textit{Det \dabold{hævdes} ofte fra politisk hold, at...} / \textit{Det \da{slås fast}, at...}
  \item \textbf{尾置/句中插叙引述的 V2 倒装铁律}：\\
    把引语放在前面，引述动词放在逗号后面做插叙时，后面的归属短句必须\textbf{动词在前、主语在后}（V2 倒装！）：\\
    \ding{51}\ \textit{”Tallene taler deres eget tydelige sprog,” \da{understreger} formanden.}（\ding{55}\ \textit{...formanden understreger.}）
  \item \textbf{\da{at slå fast} 的分合位置规律}：\\
    在主动句中，\textit{slå} 与 \textit{fast} 常常分开（中间可插入宾语或状语 \textit{imidlertid, samtidig, klart}）：\textit{Rapporten \da{slår} samtidig \da{fast}, at...}；而如果前面已经出现长宾语名词短语，\textit{fast} 站在宾语后面：\textit{Hun slog sit hovedargument fast.}（不可分写形式 \textit{at fastslå} 则始终连写：\textit{Rapporten fastslår, at...}）。
  \item \textbf{\dabold{at foregive} 后接带 \textit{at} 的不定式或 \textit{at}-从句}：\\
    \ding{51}\ \textit{Ingen kan \dabold{foregive} at sidde inde med alle svar.}（“假装手握所有答案”）；\ding{51}\ \textit{Politikeren \dabold{foregiver}, at løsningen er enkel.}
\end{enumerate}
\end{grammatikbox}

\begin{grammatikbox}{Studieprøven 信源引用避坑清单（Top 3 Pitfalls）}
\begin{itemize}
  \item \textbf{\ding{172} 致命中国式/英语式错误：\textit{Ifølge rapporten viser...}}：\\
    \da{ifølge} 是\textbf{介词}（“根据……”），后面只能接名词短语！整套 \textit{Ifølge + 名词} 充当主句前场（第 1 顺位），后面必须立刻接\textbf{真正的主句变位动词}（V2 倒装），绝对不能再跟一个引述动词！\\
    \ding{55}\ \textit{Ifølge rapporten viser, at antallet stiger.} $\rightarrow$ \ding{51}\ \textit{Ifølge rapporten \da{stiger} antallet.} 或 \ding{51}\ \textit{Rapporten \da{viser / dokumenterer}, at antallet stiger.}
  \item \textbf{\ding{173} 别不小心用 \textit{påstå} 拆自己论据的台}：\\
    很多考生背了高级词就乱换，写出 \ding{55}\ \textit{FN's klimapanel påstår, at temperaturen stiger.}——在丹麦母语者耳中，这句话听起来像你在质疑联合国气候变化委员会在胡说！支持权威数据时务必用 \da{dokumenterer / slår fast / konstaterer}。
  \item \textbf{\ding{174} 本课民主与专家治理高频搭配}：\da{at sidde inde med svarene}（掌握全部答案）；\da{at hvile på et tvivlsomt grundlag}（建立在可疑基础之上）；\da{at bruge noget som et skjold}（拿某物当挡箭牌）；\da{at træffe vidtrækkende beslutninger}（作出影响深远的决策）。
\end{itemize}
\end{grammatikbox}

% ---------------------------------------------------------------------
\ncdsubsec{五、Mønsterøvelser（引述立场精准辨析与信源改写操练）}

\textbf{练习 A：引述动词立场光谱精准填空（根据语境中立、质疑或认同立场，每空填入 1 个最准确的动词变位形式）}
\begin{enumerate}
  \item Ekspertkommissionen \underline{\hspace{2.2cm}} tørt i sin rapport, at udgifterne til ældrepleje stiger frem mod 2035.
  \item Modstanderne \underline{\hspace{2.2cm}} desuden i debatten, at længere afstande skaber utryghed for familier i landdistrikterne.
  \item De nye forskningsdata \underline{\hspace{2.2cm}} utvetydigt, at specialiserede afdelinger giver højere overlevelse.
  \item Regeringens kritikere \underline{\hspace{2.2cm}} ligefrem, at kommissionen bevidst har skjult ubehagelige beregninger.
  \item Ingen seriøs økonom kan \underline{\hspace{2.2cm}} at have fundet en løsning helt uden omkostninger for samfundet.
  \item Rapporten slår i sin konklusion \underline{\hspace{2.2cm}}, at reformer er nødvendige inden for de næste fem år.
  \item Formanden \underline{\hspace{2.2cm}}, at eksperter kun bør levere det faglige grundlag og ikke træffe politiske værdivalg.
  \item Flere borgmestre \underline{\hspace{2.2cm}}, at de centrale modeller ikke tager højde for lokale forhold.
\end{enumerate}

\textbf{练习 B：信源句法纠错与立场切换改写（修正 \textit{ifølge} 语病与语序，或按指定立场重写引述）}
\begin{enumerate}
  \item 【经典介词语病纠错】\textit{Ifølge den nye rapport dokumenterer, at borgernes tillid til eksperter er høj.} $\rightarrow$
  \item 【插叙 V2 倒装纠错】\textit{”Vi må respektere det folkelige mandat,” kommissionens formand understregede i går.} $\rightarrow$
  \item 【立场误用纠错】作者明显认同研究结论，请改掉不妥的质疑动词：\textit{Forskningen påstår tydeligt, at tidlig indsats forebygger alvorlige sygdomme.} $\rightarrow$
  \item 【中立 $\rightarrow$ 质疑保持距离】将平淡句改为表达说话人对反对派怀疑的引述：\textit{Oppositionen siger, at planen intet koster.}（用 \textit{hævde} 或 \textit{påstå} 改写）$\rightarrow$
  \item 【主动 $\rightarrow$ 被动客观摘要体】改写为报告摘要常用的无人称被动结构：\textit{Rapporten påpeger, at mange kommuner mangler speciallæger.}（用 \textit{I rapporten påpeges det, at...} 改写）$\rightarrow$
  \item 【普通引述 $\rightarrow$ 权威实证强调】升级为强有力论断：\textit{Tallene viser meget klart, at uddannelse mindsker uligheden.}（用 \textit{dokumentere} 或 \textit{slå fast} 改写）$\rightarrow$
\end{enumerate}

\textbf{练习 C：中译丹（精准区分客观、质疑与权威引述动词）}
\begin{enumerate}
  \item 报告冷静地指出，到 2040 年人口压力将会加剧。
  \item 批评者宣称，这些模型建立在脱离现实的假设之上。
  \item 研究人员通过最新数据证实，早期干预能够拯救生命。
  \item 没有任何政治人物能够假装拥有所有现成的答案。
  \item 委员会强调，最终的价值抉择必须由选民来决定。
  \item 在辩论中，反对者甚至声称该项改革将摧毁地方社区。
\end{enumerate}

% ---------------------------------------------------------------------
\ncdsubsec{六、PD3 / Studieprøven 实战微任务与参考答案}

\begin{emailbox}{PD3 / Studieprøven 实战写作任务（Referat \& Kommentar · mindst 200 ord）}
\textbf{Baggrund}: En ny ekspertkommission foreslår at hæve pensionsalderen hurtigere og samle landets sygehuse i færre, større enheder. Mens økonomer peger på solide beregninger af fremtidens arbejdsudbud, kalder fagforeninger og borgmestre forslaget for verdensfjernt ekspertvælde.\\[3pt]
\textbf{Opgave}: Skriv en \textbf{sammenhængende kommentar på 220--260 ord} (mindst 200 ord), hvor du først kort \textbf{refererer} de to siders argumenter ved hjælp af varierede holdningsverber og derefter \textbf{vurderer} balancen mellem faglig ekspertise og demokratisk legitimitet. Brug mindst \textbf{fire forskellige referatverber} fra hele spektret (neutral, skeptisk/distancerende og dokumenterende/understregende).
\end{emailbox}

\begin{answerbox}
\textbf{Før du læser 参考回答}
\begin{enumerate}
  \item \textit{Eksperterne skal oplyse det faglige grundlag med fakta og langsigtede prognoser, mens de folkevalgte politikere skal træffe de endelige værdivalg. Forfatteren advarer mod rent ekspertvælde, fordi økonomer nok kan beregne prisen på en reform, men ikke afgøre, hvilket samfund borgerne ønsker at leve i, og fordi politikere ellers svækker det folkelige mandat.}
  \item \textit{Hun anfører, at matematiske modeller ikke kan måle borgernes tryghed ved at have en lokal skadestue, og at prioriteringer i praksis er politiske værdivalg og ikke rent neutrale beregninger.}
\end{enumerate}

\textbf{Tekst B 阅读核对答案}
\begin{enumerate}
  \item \da{Forkert}（米克尔认为小型医院缺乏专科医生，报告也指出高专业质量需要更大规模的医疗单元）
  \item \da{Trygheden ved at have en skadestue i nærheden.}（附近有一家急诊室所带来的安全感）
  \item \da{Et politisk værdivalg, som vælgerne skal tage stilling til.}（本质上都是需要由选民表态的政治价值抉择）
\end{enumerate}

\textbf{练习 A 答案}
\begin{enumerate}
  \item \da{konstaterer}（与客观冷峻副词 \textit{tørt} 搭配，表冷静确认客观预测趋势）
  \item \da{anfører}（中立引述辩论一方所列举的具体理由）
  \item \da{dokumenterer}（与实证主语 \textit{forskningsdata} 及 \textit{utvetydigt} 搭配，表扎实证实）
  \item \da{påstår}（搭配语气强烈的 \textit{ligefrem}，带有明显质疑色彩）
  \item \da{foregive}（与 \textit{at have fundet} 搭配，表“假装/佯称”）
  \item \da{fast}（可分动词固定结构：\textit{at slå fast} 明确断定）
  \item \da{understreger}（搭配权威主体 \textit{Formanden}，表着重强调立场）
  \item \da{hævder}（市长们提出但仍有争议的主张，作者保持中立微距）
\end{enumerate}

\textbf{练习 B 答案（参考）}
\begin{enumerate}
  \item \textit{Ifølge den nye rapport er borgernes tillid til eksperter høj.} / \textit{Den nye rapport dokumenterer, at borgernes tillid til eksperter er høj.}（介词 \textit{Ifølge} 绝不能后接引述动词 \textit{dokumenterer}）
  \item \textit{”Vi må respektere det folkelige mandat,” understregede kommissionens formand i går.}（引语前置时，后面的归属插叙必须遵守 V2 倒装：动词 \textit{understregede} 在主语前）
  \item \textit{Forskningen dokumenterer tydeligt, at tidlig indsats forebygger alvorlige sygdomme.}（说话人明显赞同研究结论时，须用认同实证词 \textit{dokumenterer} 或 \textit{slår fast}，不能用表怀疑的 \textit{påstår}）
  \item \textit{Oppositionen hævder, at planen intet koster.} / \textit{Oppositionen påstår, at planen slet ikke koster noget.}（改用 \textit{hævde / påstå} 标示保持距离或质疑）
  \item \textit{I rapporten påpeges det, at mange kommuner mangler speciallæger.}（利用形式主语 \textit{det} 与 \textit{-s} 被动形成高度紧凑客观的学术转述）
  \item \textit{Tallene dokumenterer utvetydigt, at uddannelse mindsker uligheden.} / \textit{Tallene slår fast, at uddannelse mindsker uligheden.}（升级为实证强调动词）
\end{enumerate}

\textbf{练习 C 答案}
\begin{enumerate}
  \item \textit{Rapporten konstaterer tørt, at det demografiske pres vil vokse frem mod 2040.}
  \item \textit{Kritikerne hævder, at disse modeller hviler på urealistiske antagelser.}
  \item \textit{Forskerne dokumenterer med nye data, at en tidlig indsats redder menneskeliv.}
  \item \textit{Ingen politiker kan foregive at sidde inde med alle de færdige svar.}
  \item \textit{Kommissionen understreger, at de endelige værdivalg skal træffes af vælgerne.}
  \item \textit{I debatten påstår modstanderne ligefrem, at reformen vil ødelægge de lokale samfund.}
\end{enumerate}

\textbf{PD3 / Studieprøven 高分范文（229 词，不含标题）}\\
\textit{\textbf{Eksperter skal oplyse debatten --- ikke afslutte den}}

\textit{Når en ny ekspertkommission fremlægger vidtrækkende forslag om højere pensionsalder og færre, men større sygehuse, blusser konflikten mellem faglighed og folkestyre op på ny. På den ene side \da{konstaterer} kommissionens økonomer i rapporten, at befolkningen lever længere, og de \da{dokumenterer} med solide fremskrivninger, at de nuværende offentlige ordninger kommer under et voldsomt økonomisk pres. På den baggrund \da{slår} eksperterne \da{fast}, at strukturelle ændringer er nødvendige for at fastholde en høj kvalitet i hele sundhedsvæsenet.}

\textit{På den anden side møder planerne skarp modstand fra både lokale borgmestre og fagforeninger. Flere kritikere \dabold{hævder}, at regnearkene fra hovedstaden bygger på virkelighedsfjerne skrivebordsantagelser, og enkelte \dabold{påstår} ligefrem, at kommissionen blot leverer bekvemme argumenter for statslige besparelser. Andre \da{anfører} mere nøgternt, at et fysisk nedslidt menneske eller en borger langt fra det nærmeste hospital betaler en høj pris, som de tørre tal sjældent formår at fange.}

\textit{Efter min opfattelse har vi brug for begge perspektiver, hvis demokratiet skal bevare sin folkelige legitimitet. Ingen ansvarlig regering kan \dabold{foregive} at træffe holdbare beslutninger uden solid forskning i ryggen. Men eksperter må omvendt aldrig overtage det politiske ansvar.}

\textit{Som kommissionens formand selv bør \da{understrege}, er faglig viden et uundværligt grundlag --- ikke en færdig politisk opskrift. Eksperter kan vise konsekvenserne af vores valg, men spørgsmålet om, hvor meget tryghed, geografisk balance og social lighed vi vil prioritere, forbliver et værdivalg, som kun vælgerne og deres folkevalgte kan afgøre.}

\small\textbf{名师点评：}这篇 229 词的综述评论文（\textit{Referat \& Kommentar}）精准击中 Studieprøven 评分最高档：作者不仅以 \textit{konstaterer, dokumenterer, slår fast} 公平呈现委员会的数据实力，还用 \textit{hævder, påstår} 细腻刻画激进批评者的疑虑，同时用中立的 \textit{anfører} 留住“地方普通人真实痛点”的合理声音，随后以 \textit{foregive} 和 \textit{understrege} 完成民主正当性的价值升华。全文立场分明、句句有据，毫无空泛口号！
\end{answerbox}

\ncdlesson{8}{Kulturkamp ved middagsbordet: Mad, moral og klimaskam}{餐桌上的文化之争：饮食、道德与气候羞耻}
{DU3 Modul 5 (PD3) $\rightarrow$ Modul 6 (Studieprøven, B2 $\rightarrow$ C1)｜对标 NCE 4 \textit{Komplekse Sætninger, Indskud \& Prosa-rytme}}
{\textbf{三层以上主从复合句的语序拆解、插入语标点规范（\textit{-- ... --}）与长句节奏控制}（嵌套关系从句与状语从句 · 长前置从句后的 V2 倒装铁律 · \textit{sætningsskema} 纵深平衡与短句点睛）}
{如果说第 5--7 课为你配齐了丹麦语 C1 的修辞锋芒、社论让步与信源立场，那么第 8 课则是整个第二单元的句法压轴课：当你在 Studieprøven 读到或需要亲手写出“嵌套三四层从句、中间还穿插破折号插入语”的报纸长文（\textit{Kronik}）时，如何做到\textbf{读得一眼透亮、写得毫厘不差、节奏张弛有度}？本课借“餐桌上的气候道德与个人自由之争”这一当代丹麦热议文化议题，彻底打通高阶复合长句的拓扑骨架！}

% ---------------------------------------------------------------------
\begin{foerduleaeser}
带着下面两个宏观问题阅读 Tekst A 和 Tekst B：
\begin{enumerate}
  \item \textit{Hvorfor udløser spørgsmålet om kød, plantebaseret mad og klimaskam ofte stærkere følelser i Danmark end mange andre klimapolitiske emner ifølge Tekst A?}（根据 Tekst A，为什么肉类、植物性饮食与气候羞耻问题在丹麦往往比许多其他气候政策议题更容易引发强烈情绪？）
  \item \textit{Hvordan forholder sociologen Freja Bjerre i Tekst B sig til både ideen om ”privat mad” og brugen af moralske pegefingre i klimadebatten?}（在 Tekst B 中，社会学家弗蕾娅·比耶勒如何看待“饮食纯粹属于私事”的观点以及气候辩论中的道德说教？）
\end{enumerate}
\end{foerduleaeser}

% ---------------------------------------------------------------------
\ncdsubsec{二、Tekster（文化长篇随笔与社会学专栏回应双语篇）}

\subsection*{Tekst A：\textit{Når frikadellen bliver et politisk symbol}（深度文化评论 · 约 360 词）}

\begin{tipbox}{语篇说明}
本篇以丹麦主流报刊周末副刊（\textit{Kronik / Essay}）的高阶长短交错笔法撰写，密集呈现三层以上主从复合句与破折号插入语（\textit{-- ... --}）。阅读时请重点观察\textbf{\color{NCDNavy}深蓝色标记的主句骨架与层级引导词}（从句结束后的 V2 变位动词），以及\textbf{\color{NCDRed}深红色标记的 \textit{-- ... --} 插入语及深层从句中场副词}。
\end{tipbox}

Der var engang, hvor et dansk aftensmåltid med frikadeller, kartofler og brun sovs blot handlede om sult, tradition og familiehygge. Den tid er forbi. \da{Når} den daglige indkøbskurv \dabold{-- som tidligere blot afspejlede privat smag og et stramt husholdningsbudget --} pludselig bliver gjort til et offentligt vidnesbyrd om, \da{hvorvidt} man tager klodens fremtid alvorligt, \da{opstår} der en spænding, \da{som} rækker langt ud over selve tallerkenen og dybt ind i vores identitet.

I dag ser man tydeligt, hvordan kosten er blevet en ny og uventet frontlinje i den danske kulturkamp. På den ene side \da{peger} klimaforskere på, \da{at} en markant reduktion i forbruget af okse- og svinekød er helt nødvendig, \da{hvis} Danmark skal nå sine ambitiøse klimamål \dabold{-- et budskab, som sjældent møder modstand i teorien, men som ofte vækker irritation i en travl hverdag}. På den anden side oplever mange borgere i både storbyer og provinsbyer, \da{at} velmente kampagner om ”kødfrie dage” i offentlige kantiner hurtigt glider over i en påtrængende moralisering, \da{der} giver indtryk af, \da{at} den enkelte bør skamme sig over sine traditionelle retter.

Netop begrebet \textit{klimaskam} forklarer, hvorfor tonen ved middagsbordet så let bliver følelsesladet. Mad er nemlig aldrig kun kalorier eller CO\textsubscript{2}-regnskab. Det er erindringer, højtider, værtskab og personlig autonomi. \da{Selvom} de fleste danskere udmærket ved, \da{at} vores vaner må forandres i takt med klimakrisen, \da{nærer} mange en dyb skepsis over for eksperter og medier, \da{der} \dabold{-- ofte med en ubevidst selvretfærdighed --} taler ned til almindelige familier, \da{fordi} de stadig sætter flæskesteg eller en rugbrødsmad med leverpostej på bordet om søndagen. Samtidig frygter kritikere, at al fokus på privatforbrug fjerner blikket fra de store politiske beslutninger om landbrug og energi.

Det afgørende spørgsmål i denne kulturkamp er derfor ikke, hvem der ejer den moralske ret til at løfte pegefingeren over for andre. Spørgsmålet er, hvordan vi som samfund kan \da{skelne} klart mellem nødvendige strukturelle forandringer og unødig udskamning. \da{Hvis} politikerne tør skabe økonomiske og praktiske rammer, \da{der} gør klimavenlige valg både billigere og mere velsmagende \dabold{-- uden at det føles som et formynderisk indgreb i privatsfæren --} \da{kan} samtalen om vores madkultur igen blive saglig. Først da slipper vi for, at hvert eneste måltid forvandles til en anstrengende eksamen i god moral.

\begin{transbox}
\small
\textbf{\color{NCDTeal}[Tekst A 中文精译]}\\
曾几何时，一顿配着丹麦肉丸、土豆和褐色肉汁的晚餐，仅仅关乎饥饿、传统与家庭温馨。那个时代已经过去了。当日常的购物篮——它过去仅仅折射个人的口味与紧张的家庭预算——突然被当作一种公开证词，用来检验一个人是否认真对待地球的未来时，一种远远超出餐盘本身、深入我们身份认同的张力便由此产生。

如今人们可以清晰地看到，饮食如何变成了丹麦文化之争中一条出人意料的新前线。一方面，气候研究人员指出，如果丹麦要实现其雄心勃勃的气候目标，大幅削减牛肉和猪肉的消费是完全必要的——这一信息在理论上很少遭遇反对，但在忙碌的日常生活中却常常引发烦躁情绪。另一方面，无论是在大城市还是外省城镇，许多公民都感受到，公共食堂里本意良好的“无肉日”宣传活动很快会滑向一种咄咄逼人的道德说教，给人一种似乎每个人都应当为自己的传统菜肴感到羞愧的印象。

恰恰是“气候羞耻”这一概念，解释了为何餐桌边的语气如此容易变得充满情绪张力。因为食物从来不只是卡路里或二氧化碳账本。它是记忆、节庆、待客之道与个人自主权。尽管大多数丹麦人十分清楚我们的习惯必须随着气候危机而改变，但许多人对专家和媒体怀有深深的怀疑——他们常常带着一种不自知的道德优越感，居高临下地教训普通家庭，只因为后者周日依然把脆皮烤猪肉或涂着猪肝酱的黑麦面包端上餐桌。与此同时，批评者还担忧，把全部注意力放在私人消费上会转移人们对农业和能源等重大政治决策的视线。

因此，这场文化之争中的决定性问题，并不在于谁拥有居高临下向他人竖起食指说教的道德权利。真正的问题在于，我们作为一个社会如何能够在必要的结构性变革与不必要的羞辱指责之间清晰地划出界线。如果政治人物敢于创造经济与现实条件，让气候友好的选择既更实惠又更美味——而不至于让人觉得这是对私人领域的家长式干预——那么关于我们饮食文化的对话就能重新回归理性务实。唯有到那时，我们才能摆脱把每一顿饭都变成一场疲惫不堪的道德考试的局面。
\end{transbox}

\subsection*{Tekst B：\textit{Læserkommentar: Mad har altid været fælles --- men skam duer ikke}（社会学评论回应 · 约 150 词）}

\textbf{阅读任务（理解与核对）：}\textit{Læs sociologen Freja Bjerres kommentar til Tekst A, og besvar spørgsmålene nedenfor.}

\begin{emailbox}{Læserkommentar i Weekendtillægget｜Fra: Freja Bjerre (sociolog) \quad Emne: Sv: Når frikadellen bliver et politisk symbol}
Mange længes efter en ubekymret tid, hvor aftensmaden kun var et helt privat anliggende. Mad har \dabold{-- uanset hvor meget vi ønsker det anderledes --} imidlertid altid været gennemsyret af sociale normer, klasse og fælles kultur, og vores madvaner har ændret sig markant gennem generationer. Når vores forbrug påvirker det fælles klima, er det derfor urealistisk at tro, \da{at} spisebordet kan holdes helt uden for politik.

Derfor rammer kritikken ved siden af, når den hævder, \da{at} klimadebatten burde holde sig væk fra køkkenet. Men det betyder langtfra, \da{at} den moralske pegefinger virker efter hensigten. \da{Når} klimavenlig omstilling reduceres til individuel klimaskam \dabold{-- som ofte skaber afmagt snarere end varig handlekraft --} \da{svækker} vi netop den brede folkelige opbakning, \da{som} den grønne omstilling ellers kræver.

Hvis vi vil flytte vaner uden at splitte befolkningen i vrede lejre, må vi støtte grønnere landbrugsproduktion og tydelige prissignaler i stedet for at udskamme naboens indkøbskurv.
\end{emailbox}

\textbf{Tekst B 阅读核对（Rigtigt / Forkert / Kort svar，答案见第六部分）：}
\begin{enumerate}
  \item \textit{Freja mener, at aftensmaden i gamle dage var helt fri for normer og samfundskultur.} \underline{\hspace{2.2cm}}
  \item \textit{Hvad skaber individuel klimaskam ofte hos borgerne ifølge Freja i stedet for handlekraft?} \underline{\hspace{2.2cm}}
  \item \textit{Hvad bør man satse på i stedet for at udskamme naboens indkøbskurv ifølge Frejas konklusion?} \\\underline{\hspace{2.2cm}}
\end{enumerate}

\begin{tipbox}{本课发音与语调要点（Udtale \& Prosodi）}
\begin{itemize}
  \item \textbf{文化辩论与饮食伦理词汇}：\textit{en kulturkamp} \ipa{kʰulˈtˢuɐ̯ˌkʰɑmˀb̥}（重音在 \textit{-tur-}）；\textit{en klimaskam} \ipa{ˈkʰliːmaˌsɡ̊ɑmˀ}（复合成词，重音在首音节，\textit{-skam} 带 \textit{stød}）；\textit{en autonomi} \ipa{ɑʊ̯tˢonoˈmiːˀ}（重音在词尾）；\textit{en selvretfærdighed} \ipa{ˈsɛlʁɛd̥ˌfɛɐ̯ˀd̥iheðˀ}（\textit{selv} 尾音 \textit{v} 常弱化脱落，末尾 \textit{-hed} 带软 d 与 \textit{stød}）。
  \item \textbf{极易错读的书面动词与长词}：\textit{at afspejle} \ipa{ˈɑʊ̯ˌsb̥ɑjlə}（前缀 \textit{af-} 读 \ipa{ɑʊ̯}）；\textit{at skelne} \ipa{ˈsɡ̊ɛlnə}（动词 \textit{sk-} 后不送气）；\textit{et vidnesbyrd} \ipa{ˈviðnəsˌb̥yɐ̯ˀð}（两个软 d，中间不插额外元音）；\textit{formynderisk} \ipa{fʌˈmynˀɐʁisɡ̊}（重音在第二音节，\textit{d} 在 \textit{-nd-} 中不发音！）。
  \item \textbf{破折号插入语（\textit{-- ... --}）的朗读“降阶换挡”}：读到进入 \textit{-- ... --} 的瞬间，语速略加快、音高微降半度（仿佛压低嗓音加一句旁白），走出第二个破折号时立刻回升到主句原有音高，这样听众即便遇到 40 词超长句也不会迷路！
\end{itemize}
\end{tipbox}

\begin{transbox}
\small
\textbf{\color{NCDTeal}[Tekst B 中文对照]}\\
\textbf{周末副刊读者来信（弗蕾娅·比耶勒，社会学家｜回应：食物从来不是纯粹的私事——但羞辱解决不了任何问题）：}许多人怀念一个无忧无虑的时代，仿佛晚餐纯粹只是一件完全私人的事务。然而食物——不管我们多么希望它有所不同——从来都浸透着社会规范、阶层属性与公共文化，而且我们的饮食习惯历代以来也已发生了显著变化。当我们的消费影响着共同的气候时，指望把餐桌完全隔绝在政治之外因此是不现实的。

因此，当一些批评声音宣称气候辩论本不该踏进厨房时，这种批评就偏离了靶心。但这远远不意味着道德说教能如愿奏效。当气候友好的转型被简化为个人的气候羞耻——这种羞耻往往制造出无力感而非持久的行动力——我们恰恰是在削弱绿色转型本来最需要的那份广泛民众支持。

如果我们想在不把民众撕裂成愤怒阵营的前提下改变习惯，我们就必须扶持更绿色的农业生产方式和清晰的价格信号，而不是去羞辱邻居的购物篮。
\end{transbox}

% ---------------------------------------------------------------------
\ncdsubsec{三、Gloser（核心生词表）}

\begin{tipbox}{词表重点}
注意几个不规则/易错词形变化：\textit{en pegefinger}（食指/说教手势）复数发生元音脱落变为 \textit{pegefingre}（\ding{55}\ \textit{pegefingerer}）；\textit{et indskud}（插入语）和 \textit{et vidnesbyrd}（证词/见证）均为中性单复同形；形容词 \textit{velment}（本意良好的）在定指/复数时加 \textit{-e}（\textit{velmente kampagner}）。动词 \textit{at skelne} 固定搭配介词 \textit{mellem}（区分 A 与 B）。
\end{tipbox}

\begin{glosertable}
名词 & en kulturkamp & \textit{(kulturkampen, kulturkampe, kulturkampene)} & 文化之争，价值观冲突 \\
名词 & en klimaskam & \textit{(klimaskammen, kun ental)} & 气候羞耻，环保内疚感 \\
名词 & en autonomi & \textit{(autonomien, kun ental)} & 自主权，自治 \\
名词 & en selvretfærdighed & \textit{(selvretfærdigheden, kun ental)} & 道德自义，自以为是 \\
名词 & en pegefinger & \textit{(pegefingeren, pegefingre, pegefingrene)} & 食指；（比喻）道德说教 \\
名词 & en indkøbskurv & \textit{(indkøbskurven, indkøbskurve, indkøbskurvene)} & 购物篮 \\
名词 & et vidnesbyrd & \textit{(vidnesbyrdet, vidnesbyrd, vidnesbyrdene)} & 证词，见证，体现 \\
名词 & en afmagt & \textit{(afmagten, kun ental)} & 无力感，无能为力 \\
名词 & et indskud & \textit{(indskuddet, indskud, indskuddene)} & 插入语；投入（资金） \\
名词 & en opbakning & \textit{(opbakningen, kun ental)} & 支持，声援，后盾 \\
\midrule
动词 [har] & at afspejle & \textit{(afspejler, afspejlede, har afspejlet)} & 反映，折射 \\
动词 [har] & at moralisere & \textit{(moraliserer, moraliserede, har moraliseret)} & 道德说教，板起面孔讲大道理 \\
动词 [har] & at skelne (mellem) & \textit{(skelner, skelnede, har skelnet)} & 区分，辨别 \\
动词 [har] & at udskamme & \textit{(udskammer, udskammede, har udskammet)} & 公开羞辱，令……蒙羞 \\
动词 [har] & at indskrænke & \textit{(indskrænker, indskrænkede, har indskrænket)} & 限制，缩减（自由/范围） \\
\midrule
形容词 & ubekymret & \textit{(ubekymret, ubekymrede)} & 无忧无虑的，毫不在意的 \\
形容词 & følelsesladet & \textit{(følelsesladet, følelsesladede)} & 充满情绪张力的，情绪化的 \\
形容词 & velment & \textit{(velment, velmente)} & 出于好意的，善意的 \\
\end{glosertable}

% ---------------------------------------------------------------------
\ncdsubsec{四、Grammatik \& Stilistik（多层复合句拓扑拆解、插入语标点与长句节奏精讲）}

\begin{grammatikbox}{句法点 1：三层以上主从复合句的“拓扑剥洋葱”拆解术}
C1 级长句看起来令人望而生畏，其实永远逃不出丹麦语拓扑句法（\textit{Sætningsskema}）的铁律。\textbf{不论嵌套多少层，每一层从句内部永远独立重启“连接词 + 主语 + 中场副词 + 动词”}！请看 Tekst A 首段那句经典四层长句的结构解剖：

\vspace{0.3em}
\renewcommand{\arraystretch}{1.24}
{\small
\begin{tabularx}{\linewidth}{@{} >{\bfseries\raggedright\arraybackslash}p{2.7cm} >{\raggedright\arraybackslash}p{5.5cm} >{\raggedright\arraybackslash}X @{}}
\toprule
\textbf{句法层级} & \textbf{丹麦语原句片段} & \textbf{核心语序法则解析} \\
\midrule
\da{第 1 层（前场状语从句）} & \textit{\da{Når} den daglige indkøbskurv ... pludselig bliver gjort til et offentligt vidnesbyrd om,} & \textit{Når} 引导时间条件从句；中场副词 \textit{pludselig} 站在助动词 \textit{bliver} 前面。 \\
\midrule
\dabold{第 2 层（插入语定语从句）} & \textit{\dabold{-- som tidligere blot afspejlede privat smag og husholdningsbudget --}} & 用 \textit{-- ... --} 嵌入修饰 \textit{indkøbskurv}；从句内中场副词 \textit{tidligere blot} 站在动词 \textit{afspejlede} 前。 \\
\midrule
\da{第 2/3 层（介词宾语间接疑问从句）} & \textit{om, \da{hvorvidt} man tager klodens fremtid alvorligt,} & 介词 \textit{om} 后打逗号接 \textit{hvorvidt} 从句，内部保持正常不倒装语序。 \\
\midrule
\da{第 0 层（主句核心 V2 + 尾置定语）} & \textit{\da{opstår} der en spænding, \da{som} rækker langt ud over selve tallerkenen.} & 上述三大从句整体算作主句“前场（Forfelt）”，故逗号后立刻触发 \textbf{V2 倒装}（\textit{opstår der}），再于句末优雅挂出关系从句！ \\
\bottomrule
\end{tabularx}
}
\end{grammatikbox}

\begin{grammatikbox}{标点点 2：插入语标点铁律（\textit{-- ... --}、逗号与括号）与“拇指遮盖检验法”}
\begin{enumerate}
  \item \textbf{破折号插入语（\textit{Tankestreger: -- ... --}）的黄金规范}：\\
    在丹麦语规范排版中，句中插入语使用两个带前后空格的短破折号（\texttt{-- ... --}）。它的修辞力度高于逗号，能把读者注意力引向补充评注、反讽反转或让步条件：\textit{Mad har \dabold{-- uanset hvor meget vi ønsker det anderledes --} altid været politisk.}
  \item \textbf{“拇指遮盖检验法”（Tommelfingerreglen）确保不断句}：\\
    写完带 \textit{-- ... --} 的句子后，用拇指把两个破折号连同中间内容遮住，剩下的主干必须在语法和标点上\textbf{严丝合缝}！注意上面这句遮住插入语后恰好是：\textit{Mad har altid været politisk.}（主句助动词 \textit{har} 在前，\textit{altid været} 接在后，绝不重复 \textit{har}，也不丢失动词！）。同理：\textit{Når klimakampen \dabold{-- som ofte skaber afmagt --} reduceres til skam, \da{svækker} vi opbakningen.} 遮住后为 \textit{Når klimakampen reduceres til skam, svækker vi...}，完全通顺！
  \item \textbf{丹麦语从句逗号（Startkomma）是长句读懂与写对的救命索}：\\
    丹麦语标准传统逗号（\textit{grammatisk komma}）要求在每一个从属连词（\textit{at, som, der, fordi, selvom, hvis, når, hvorvidt}）之前打逗号（无论前面是动词还是介词：\textit{vidnesbyrd om, hvorvidt...} / \textit{indtryk af, at...}）。在三层长句中，这就像路标一样清晰标出每一层从句的起止边界！
  \item \textbf{三种插入语符号的文体分工}：逗号（\textit{, ...,}）最轻柔平实；破折号（\textit{-- ... --}）最具评论张力与节奏顿挫；圆括号（\textit{(...)}）仅用于客观年份、缩写、数字补充，评论写作中切忌滥用圆括号！
\end{enumerate}
\end{grammatikbox}

\begin{grammatikbox}{修辞点 3：长短交错的散文节奏（\textit{Prosa-rytme}）与 Unit 2 终极排雷}
\begin{itemize}
  \item \textbf{黄金节奏公式：长句铺陈 + 短句落锤}：如果通篇全是 40 词嵌套长句，文章会像官僚公文一样令人窒息。请看 Tekst A 首段与第三段的节奏设计——先用两句极短的话劈开空间（\textit{Den tid er forbi.} / \textit{Mad er nemlig aldrig kun kalorier eller CO\textsubscript{2}-regnskab.}），再放出一记结构精密的 35--40 词长句展开纵深论证，最后再用一句 8--10 词短句收束定音！
  \item \textbf{\ding{172} 排雷一：长从句结束后忘了 V2 倒装}：\ding{55}\ \textit{Når indkøbskurven -- som før var privat -- pludselig bliver politisk, en spænding opstår.} $\rightarrow$ 错！必须 \ding{51}\ \textit{...opstår der en spænding.}
  \item \textbf{\ding{173} 排雷二：插入语遮住了眼睛，导致主句缺动词或丢副词}：主句助动词如果放在 \textit{-- ... --} 之前，后面必须紧跟中场副词和实义动词（\textit{Mad \da{har} -- uanset... -- \da{altid været}...}），绝不能再写一遍 \textit{har}！
  \item \textbf{\ding{174} 排雷三：深层从句里中场副词跑到了动词后面}：即便在第三层从句（如 \textit{at den enkelte \dabold{ikke} bør skamme sig} 或 \textit{som \dabold{sjældent} møder modstand}）中，副词也永远站在变位动词前面！
\end{itemize}
\end{grammatikbox}

% ---------------------------------------------------------------------
\ncdsubsec{五、Mønsterøvelser（多层复合句拆解、插入语与节奏操练）}

\textbf{练习 A：三层复合句结构枢纽与标点词汇精准填空（每空填入最准确的 1 个丹麦语词）}
\begin{enumerate}
  \item Når hverdagsmaden pludselig gøres til et bevis på, \underline{\hspace{2.2cm}} man bekymrer sig om klimaet, opstår der let konflikter.
  \item Selvom de fleste borgere udmærket ved, at vores vaner må ændres, \underline{\hspace{2.2cm}} mange en dyb skepsis over for moralske pegefingre.
  \item Hvis politikerne tør skabe rammer, \underline{\hspace{2.2cm}} gør bæredygtige valg billigere, kan samtalen igen blive saglig.
  \item Indkøbskurven -- som tidligere blot afspejlede smag og økonomi -- \underline{\hspace{2.2cm}} i dag ofte gjort til et politisk symbol.
  \item Det afgørende er, at vi lærer at skelne \underline{\hspace{2.2cm}} nødvendige reformer og unødig udskamning af naboer.
  \item Mad har -- uanset hvor meget vi savner en ubekymret fortid -- \underline{\hspace{2.2cm}} været gennemsyret af kulturelle normer.
  \item Klimavenlige kampagner, der \underline{\hspace{2.2cm}} tager hensyn til almindelige familiers hverdag, skaber ofte mere irritation end handling.
  \item En politik, som primært bygger på individuel skyld og skam, svækker netop den brede opbakning, \underline{\hspace{2.2cm}} omstillingen kræver.
\end{enumerate}

\textbf{练习 B：长句拓扑纠错与三层复合句组装升格（修正语序与插入语断裂，或将短句合成为流畅长句）}
\begin{enumerate}
  \item 【长前置从句 V2 纠错】\textit{Når en traditionel ret pludselig bliver gjort til et moralsk spørgsmål, mange borgere føler sig uretfærdigt dømt.} $\rightarrow$
  \item 【插入语后主句语序与重词纠错】\textit{Klimadebatten har -- selvom formålet jo er rigtig godt -- har ofte glemt almindelige familiers vilkår.} $\rightarrow$
  \item 【第三层从句中场副词纠错】\textit{Hun understregede, at man ikke må tro, at forbrugerne ændrer ikke vaner uden politiske rammer.} $\rightarrow$
  \item 【三层嵌套组装升级】将下列零散信息合成为一句带 \textit{-- ... --} 的流畅 C1 长句（用 \textit{Når..., opstår der...} 起句）：\textit{Folk spiser aftensmad. Aftensmaden var tidligere helt privat. Samtidig kritiserer medierne kødforbruget. Derfor opstår der en følelsesladet kulturkamp.} $\rightarrow$
  \item 【插入语嵌入与节奏重组】将让步从句作为 \textit{-- ... --} 插入语嵌入主句中轴：\textit{Klimaskam løser sjældent strukturelle problemer. Mange aktivister mener ellers, at skyldfølelse skaber hurtig handling.} $\rightarrow$
  \item 【长短句节奏微调】先用一个 6 词以内短句定调，再用复合句解释理由：\textit{Maden er aldrig neutral, fordi den afspejler vores historie, vores værdier og den måde, som vi ønsker at indrette samfundet på.} $\rightarrow$
\end{enumerate}

\textbf{练习 C：中译丹（兼顾三层主从语序与地道高阶表达）}
\begin{enumerate}
  \item 当日常饮食变成道德问题时，餐桌旁便会产生强烈的情绪张力。
  \item 尽管大家都知道习惯必须改变，但许多人依然反感道德说教。
  \item 如果政府让绿色食物更便宜，人们就会更容易做出可持续选择。
  \item 我们必须区分必要的社会转型与对普通家庭的不必要羞辱。
  \item 研究表明，制造无力感的气候羞耻往往会削弱公众对改革的支持。
  \item 食物从来不只是卡路里，因为它同时承载着记忆与个人自主权。
\end{enumerate}

% ---------------------------------------------------------------------
\ncdsubsec{六、PD3 / Studieprøven 实战微任务与参考答案}

\begin{emailbox}{PD3 / Studieprøven 实战写作任务（Unit 2 综合 Kronik-indlæg · mindst 200 ord）}
\textbf{Opgave}: Debatten om kødforbrug, kødfrie dage i kantiner og ”klimaskam” skaber ofte heftige diskussioner i Danmark. Nogle mener, at stærke moralske opfordringer er nødvendige for at få folk til at ændre vaner, mens andre oplever klimakampagner som et indgreb i den personlige frihed.\\[3pt]
Skriv en \textbf{velstruktureret kommentar på 230--270 ord} (mindst 200 ord), hvor du forener de vigtigste værktøjer fra Unit 2:
\begin{itemize}
  \item Mindst én \textbf{velafbalanceret tre-lags sætning med indskud (\textit{-- ... --})} fulgt op af en kort, slagkraftig sætning;
  \item Mindst ét \textbf{avanceret referatverbum} (fx \textit{konstatere, hævde, dokumentere}) og én \textbf{indrømmelsesstruktur} (fx \textit{Bevares / Det ændrer ikke ved den kendsgerning, at...});
  \item En afmålt tone, eventuelt med et strejf af \textbf{dansk underdrivelse}.
\end{itemize}
\end{emailbox}

\begin{answerbox}
\textbf{Før du læser 参考回答}
\begin{enumerate}
  \item \textit{Fordi mad ikke blot er kalorier eller et CO\textsubscript{2}-regnskab, men også handler om erindringer, højtider, traditioner og personlig autonomi. Når den daglige indkøbskurv gøres til et moralsk vidnesbyrd, oplever mange borgere kampagnerne som påtrængende selvretfærdighed og udskamning.}
  \item \textit{Hun mener ikke, at aftensmaden kan holdes helt uden for politik, da mad altid har været præget af normer og kultur og vores forbrug påvirker det fælles klima. Samtidig afviser hun moralske pegefingre, fordi individuel klimaskam skaber afmagt og svækker den folkelige opbakning.}
\end{enumerate}

\textbf{Tekst B 阅读核对答案}
\begin{enumerate}
  \item \da{Forkert}（她认为食物从来都浸透着社会规范、阶层属性和公共文化）
  \item \da{Afmagt (snarere end handlekraft).}（无力感，而非行动力）
  \item \da{Grønnere produktion og tydelige prissignaler.}（更绿色的生产方式和清晰的价格信号）
\end{enumerate}

\textbf{练习 A 答案}
\begin{enumerate}
  \item \da{hvorvidt} / \da{om}（接在介词 \textit{på} 之后引导间接疑问从句，书面首选 \textit{hvorvidt}）
  \item \da{nærer}（前置让步从句 \textit{Selvom...} 占据前场，主句须 V2 倒装：动词 \textit{nærer} 居首）
  \item \da{der} / \da{som}（引导定语从句并在从句中作主语，修饰复数名词 \textit{rammer}）
  \item \da{bliver} / \da{er}（注意遮住 \textit{-- ... --} 插入语后，主语 \textit{Indkøbskurven} 后紧接唯一变位动词）
  \item \da{mellem}（固定动词搭配：\textit{at skelne mellem A og B} 区分 A 与 B）
  \item \da{altid}（助动词 \textit{har} 放在 \textit{-- ... --} 之前，实义分词 \textit{været} 之前保留中场副词 \textit{altid}）
  \item \da{ikke} / \da{sjældent}（在关系从句 \textit{der ... tager} 内，否定或频度副词必须位于变位动词之前）
  \item \da{som}（此处关系代词在从句中作动词 \textit{kræver} 的宾语，必须用 \textit{som}，不可用只作主语的 \textit{der}！）
\end{enumerate}

\textbf{练习 B 答案（参考）}
\begin{enumerate}
  \item \textit{Når en traditionel ret pludselig bliver gjort til et moralsk spørgsmål, føler mange borgere sig uretfærdigt dømt.}（长从句位于句首时，主句必须 V2 倒装：\textit{føler mange borgere sig...}）
  \item \textit{Klimadebatten har -- selvom formålet jo er rigtig godt -- ofte glemt almindelige familiers vilkår.}（删除 \textit{-- ... --} 后方重复出现的助动词 \textit{har}，并用“拇指遮盖法”验证主干通顺）
  \item \textit{Hun understregede, at man ikke må tro, at forbrugerne ikke ændrer vaner uden politiske rammer.}（在最内层的 \textit{at} 从句中，\textit{ikke} 依然必须放在变位动词 \textit{ændrer} 之前）
  \item \textit{Når aftensmaden -- som tidligere var et helt privat anliggende -- bliver mødt med kritik af kødforbruget, opstår der let en følelsesladet kulturkamp.}（将补充背景装入 \textit{-- ... --} 定语从句，从句结束后主句干脆利落地实现 \textit{opstår der} 倒装）
  \item \textit{Klimaskam løser -- uanset at mange aktivister tror, at skyldfølelse skaber hurtig handling -- sjældent de strukturelle problemer.}（主句 \textit{Klimaskam løser sjældent...} 中轴跨越插入语稳健连接）
  \item \textit{Maden er aldrig neutral. Den afspejler nemlig både vores historie, vores værdier og den måde, hvorpå vi ønsker at indrette samfundet.}（先立 4 词短句敲钟，再用多层结构句展开阐释，达成高级散文节奏）
\end{enumerate}

\textbf{练习 C 答案}
\begin{enumerate}
  \item \textit{Når hverdagsmaden bliver et moralsk spørgsmål, opstår der stærke følelser ved middagsbordet.}
  \item \textit{Selvom alle ved, at vanerne må ændres, afviser mange stadig moralske pegefingre.}
  \item \textit{Hvis staten gør grøn mad billigere, vælger borgerne lettere bæredygtige alternativer i hverdagen.}
  \item \textit{Vi må skelne mellem nødvendig omstilling og unødig udskamning af almindelige familier.}
  \item \textit{Forskning dokumenterer, at klimaskam, som skaber afmagt, ofte svækker den folkelige opbakning.}
  \item \textit{Mad er aldrig blot kalorier, fordi den også rummer erindringer og personlig autonomi.}
\end{enumerate}

\textbf{PD3 / Studieprøven 高分范文（223 词，不含标题）}\\
\textit{\textbf{Mindre moraliseren, mere handlekraft: Om mad og klima uden skam}}

\textit{Debatten om, hvad danskerne lægger på tallerkenen, er i de senere år blevet \dabold{ikke ligefrem afslappet}. \da{Når} en almindelig hverdagsret \dabold{-- som for mange blot rummer minder om tryghed og tradition --} pludselig gøres til et bevis på, \da{hvorvidt} man tager klimakrisen alvorligt, \da{opstår} der en polarisering, \da{som} sjældent gavner naturen. Måltidet bliver en slagmark.}

\textit{\da{Bevares}, forskningen \da{dokumenterer} tydeligt, at et lavere forbrug af oksekød og mejeriprodukter mindsker vores samlede udledning, og ingen kan \dabold{foregive}, at vores madvaner slet ikke angår fællesskabet. \da{Men spørgsmålet forbliver}, om moralsk udskamning overhovedet er et effektivt redskab. Sociologer \da{påpeger} tværtimod, at en løftet pegefinger ofte vækker trods og afmagt, mens modstandere af klimakampagner \dabold{hævder}, at eksperter ønsker at styre borgernes køkken helt ned i mindste detalje.}

\textit{\dabold{Uanset hvordan man vender og drejer det}, ændrer stærke følelser ikke ved den kendsgerning, at vaner først flytter sig i stor skala, når det grønne valg bliver både nemt og overkommeligt. Hvis kantiner, supermarkeder og politikere gør plantebaseret mad mere indbydende og sænker prisen på klimavenlige råvarer \dabold{-- uden at opføre sig som moralske dommere --} \da{kan} forandringen ske med borgerne i stedet for imod dem.}

\textit{Vi bør derfor skelne skarpere mellem strukturel klimapolitik og privat skyldfølelse. Det ville vist næppe skade debatten, hvis vi lod pegefingeren hvile og i stedet mødtes om et godt måltid --- uanset om der ligger kød eller linser på tallerkenen.}

\small\textbf{名师点评：}这篇 223 词的 Unit 2 综合收官佳作，堪称 Studieprøven 满分示范！开篇先来一记丹麦式低调反语 \textit{ikke ligefrem afslappet}，紧接着用一座精密的 33 词四层复合长句（含 \textit{-- ... --} 插入语与 \textit{hvorvidt} 间接疑问）交代文化撕裂，随即以 4 词短句 \textit{Måltidet bliver en slagmark.} 斩钉截铁落锤；中段完美串联 \textit{Bevares... dokumenterer... foregive... Men spørgsmålet forbliver... påpeger... hævder}，立场层次纤毫毕现；结尾段再用 \textit{Uanset hvordan man vender og drejer det} 与带 \textit{-- ... --} 插入语的长条件句收束对策，最后以一抹克制的丹麦冷幽默 \textit{Det ville vist næppe skade debatten...} 从容谢幕。四课精华，浑然一体！
\end{answerbox}


% =====================================================================
% Unit 3
% =====================================================================
\ncdunit{3}{Etik, Institutioner og Modernitet}{伦理、制度与现代性思辨}{Lektion 9--12：原则与实践辩证论证 · 绩效社会设问与概念拆解 · 丹麦文化不可译核心词辨析 · 生命伦理两难推演与警示句型}

\ncdlesson{9}{Ytringsfrihedens grænser i en polariseret tid}{极化时代的言论自由边界}
{DU3 Modul 5 (PD3) $\rightarrow$ Modul 6 (Studieprøven, B2 $\rightarrow$ C1)｜对标 NCE 4 \& 民主法治与公共伦理辩论}
{\textbf{原则与实践的辩证论证结构}（\textit{principielt set} $\leftrightarrow$ \textit{i praksis} · \textit{på den ene side ukuelig principfasthed...} · \textit{trække en hårfin grænse mellem A og B} · \textit{balancere hensynet til A over for B}）· 丹麦《基本法》第 77 条与刑法仇恨言论条款辨析}
{欢迎来到 Unit 3《伦理、制度与现代性思辨》的第一课！进入 PD3 与 Studieprøven 的高阶议论文，仅仅列出“正反两面”已经不够，考官更看重你能否剖析\textbf{“抽象原则”与“社会实践”之间的深层张力}。在丹麦公共辩论中，言论自由既是神圣的宪法基石，又在算法极化与网络暴力面前不断面临现实考验。掌握本课的 \textbf{\textit{principielt set}} 与 \textbf{\textit{i praksis}} 双轨辩证结构，以及 \textbf{\textit{trække en hårfin grænse mellem}} 和 \textbf{\textit{balancere hensynet til ... over for ...}} 等精准法理搭配，你将具备驾驭任何价值冲突议题的 C1 级思辨力！}

% ---------------------------------------------------------------------
\begin{foerduleaeser}
带着下面两个问题阅读 Tekst A 和 Tekst B：
\begin{enumerate}
  \item \textit{Hvordan stiller Tekst A det forfatningsmæssige princip om ytringsfrihed op over for de praktiske udfordringer med digital chikane og polarisering?}（Tekst A 是如何将言论自由的宪法原则与现实中网络骚扰及社会极化的挑战放在一起对照剖析的？）
  \item \textit{Hvorfor advarer chefredaktøren Henrik i Tekst B mod at trække den juridiske grænse for stramt gennem nye forbud mod krænkende ytringer?}（在 Tekst B 中，主编亨里克为什么警告不要通过新的禁令来对冒犯性言论划定过紧的法律边界？）
\end{enumerate}
\end{foerduleaeser}

% ---------------------------------------------------------------------
\ncdsubsec{二、Tekster（法治伦理专栏与报刊主编反驳评论双语篇）}

\subsection*{Tekst A：\textit{Ytringsfrihedens grænser i en polariseret tid}（原则与实践辩证长文 · 约 370 词）}

\begin{tipbox}{语篇说明}
本篇聚焦“言论自由原则”与“数字时代仇恨言论现实”之间的张力。阅读时请重点标出\textbf{\color{NCDNavy}深蓝色标注的“原则 $\leftrightarrow$ 实践”路标框架与宏观概念}，以及\textbf{\color{NCDRed}深红色标注的法理边界、价值权衡核心搭配}。
\end{tipbox}

I Danmark regnes ytringsfriheden for selve grundstenen i vores folkestyre. Allerede i 1849 gav Grundlovens paragraf 77 enhver borger ret til frit på tryk, i skrift og tale at ytre sine tanker uden forudgående censur --- dog ”under ansvar for domstolene”. I en tidsalder præget af digitale platforme og skarp politisk polarisering stilles dette klassiske ideal imidlertid på en alvorlig prøve, når demokratisk debat støder sammen med hadtale og organiserede trusler.

\da{Principielt set} hviler et modent demokrati på en \dabold{ukuelig principfasthed}, når det gælder retten til at fremsætte upopulære, provokerende eller endda dybt krænkende synspunkter. \da{På den ene side} sikrer denne \dabold{ukuelige principfasthed}, at mindretal, journalister og systemkritikere kan udfordre magthavere uden frygt for statslig vilkårlighed. Hvis borgerne først begynder at lægge bånd på sig selv af angst for straf, mister den offentlige samtale sin kritiske kraft og sin evne til at korrigere fejl.

\da{I praksis} udspiller den offentlige debat sig imidlertid sjældent som en rolig akademisk samtale. Når sociale mediers algoritmer belønner vrede og forargelse, spredes systematisk chikane og dehumaniserende sprog på få sekunder. Mange lokalpolitikere, forskere og repræsentanter for udsatte grupper oplever ikke længere blot modstand mod deres argumenter, men målrettet digital intimidering. En sådan udvikling risikerer at udhule samfundets \da{sammenhængskraft} og skræmme vigtige stemmer helt væk fra den demokratiske arena.

Den afgørende opgave for både lovgivere og domstole består derfor i at \dabold{trække en hårfin grænse mellem} skarp, legitim samfundskritik på den ene side og ulovlige trusler eller forfølgelse på den anden side. Netop dette dilemma afspejles i straffelovens paragraf 266 b, der forbyder offentligt at true, forhåne eller nedværdige befolkningsgrupper. Her må retsstaten hele tiden \dabold{balancere hensynet til} den enkeltes vide ytringsfrihed \dabold{over for} beskyttelsen af medborgernes værdighed og ligebehandling. Trækkes grænsen for snævert, kvæles den nødvendige uenighed; trækkes den for bredt, svigtes de mest sårbare borgere.

\da{Sammenfattende} kan principper og praksis ikke skilles ad i et levende demokrati. \da{Principielt set} bør friheden altid være udgangspunktet, mens indgreb forbliver en snæver undtagelse. \da{I praksis} kræver en holdbar ytringskultur imidlertid både klare juridiske værn mod chikane og en udbredt etisk selvjustits hos den enkelte bruger. Først når vi både forsvarer retten til det ubehagelige ord og værner om samtalens værdighed, kan ytringsfriheden trives i en polariseret tid.

\begin{transbox}
\small
\textbf{\color{NCDTeal}[Tekst A 中文精译]}\\
在丹麦，言论自由常被视为民主制度本身的基石。自 1849 年以来，《基本法》第 77 条便保障公民有权通过印刷、书面和口头方式公开发表思想而不受事先审查——但同时“须在法庭面前承担责任”。然而，在一个充斥着数字平台和尖锐政治极化的时代，当民主辩论与仇恨言论及有组织的威胁迎面相撞时，这一古典理想正面临严峻考验。

从原则上讲，一个成熟的民主社会在涉及发表不受欢迎、挑衅甚至令人深感冒犯的观点的权利时，依赖于一种坚定不移的原则立场。一方面，这种坚定不移的原则立场确保少数群体、记者以及制度批评者能够挑战当权者，而不必担心国家权力的专断任意。倘若公民因害怕受罚而率先自我约束，公共讨论就会丧失其批判力量和纠错能力。

然而在实践中，公共辩论却极少像一场平静的学术对话那样展开。当社交媒体的算法奖励愤怒与激愤情绪时，系统性的骚扰和去人性化语言便会在几秒钟内蔓延。许多地方政治人物、学者以及弱势群体的代表所遭遇的，早已不只是对其论点的反对，而是有针对性的数字恐吓。这样的趋势有可能侵蚀社会的凝聚力，并将重要的声音彻底吓退在民主舞台之外。

因此，摆在立法者与法院面前的决定性任务，就在于如何在尖锐而合法的社会批判与非法的威胁或迫害之间划出一条极其微妙的界线。这一困境正反映在丹麦刑法第 266 b 条中，该条款禁止公开威胁、侮辱或贬损特定人群。在此，法治国家必须始终在保障个人广泛的言论自由与保护同胞公民的尊严及平等待遇这两项考量之间求得平衡。倘若界线划得过窄，必要的意见分歧便会被窒息；倘若划得过宽，最脆弱的公民便会遭到辜负。

综上所述，在一个充满活力的民主制度中，原则与实践无法被割裂。原则上，自由应当始终作为出发点，而干预只能是严格受限的例外。然而在实践中，一种可持续的言论文化既需要针对骚扰的明确法律防线，也需要每位用户普遍具备伦理自律。唯有当我们既捍卫说出刺耳话语的权利，又守护对话本身的尊严时，言论自由才能在极化时代生生不息。
\end{transbox}

\subsection*{Tekst B：\textit{Kommentar: ”Nye forbud løser ikke en polariseret debatkultur”}（报刊主编反驳评论 · 约 155 词）}

\textbf{阅读任务（理解与核对）：}\textit{Læs chefredaktør Henrik Madsens kommentar til Tekst A, og besvar spørgsmålene nedenfor.}

\begin{emailbox}{Avisdebat｜Fra: Henrik Madsen (chefredaktør) \quad Emne: Principfasthed frem for nye gummiparagraffer}
Som redaktør møder jeg dagligt hadske kommentarer i vores digitale kommentarspor, og direkte trusler mod navngivne personer skal naturligvis retsforfølges uden tøven. Alligevel advarer jeg kraftigt imod at tro, at vi kan lovgive os ud af en rå og polariseret tone på nettet.

\da{Principielt set} må staten aldrig fungere som smagsdommer over, hvilke holdninger der er moralsk acceptable i et frit samfund. Hvis Folketinget udvider straffeloven med vage forbud mod krænkende sprog, risikerer vi, at politikere, debattører og forskere af ren forsigtighed tier om ømtålelige emner. \da{I praksis} fører vidtgående forbud desuden sjældent til større respekt: De forvandler derimod provokatører til martyrer og skubber ekstreme synspunkter ned i lukkede netværk, hvor ingen længere kan modsige dem.

Derfor bør vi holde fast i en \dabold{ukuelig principfasthed}, når det gælder lovgivningens rammer, og i stedet \dabold{balancere hensynet til} en anstændig samtale \dabold{over for} ytringsfriheden gennem aktiv redaktionel moderation og modige modargumenter.
\end{emailbox}

\textbf{Tekst B 阅读核对（Rigtigt / Forkert / Kort svar，答案见第六部分）：}
\begin{enumerate}
  \item \textit{Henrik mener, at konkrete trusler mod navngivne personer på nettet bør straffes retsligt.} \underline{\hspace{2.2cm}}
  \item \textit{Hvilke to negative virkninger får vage og vidtgående forbud ifølge Henrik i praksis?} \underline{\hspace{2.2cm}}
  \item \textit{Hvordan bør medier og samfund i stedet sikre en anstændig debat?} \underline{\hspace{2.2cm}}
\end{enumerate}

\begin{tipbox}{本课发音与语调要点（Udtale \& Prosodi）}
\begin{itemize}
  \item \textbf{宪法与抽象核心词减音及喉塞音}：\textit{ytringsfrihed} \ipa{ˈyd̥ʁeŋsˌfʁiːˌheðˀ}（结尾 \textit{d} 读软 \ipa{ð} 且带 \textit{stød}）；\textit{principielt} \ipa{pʰʁensiˈpjɛlˀd}（法语借词重音在后）；\textit{ukuelig} \ipa{uˈkʰuːəli}（词尾 \textit{g} 不发音）；\textit{sammenhængskraft} \ipa{ˈsɑmənˌhɛŋsˌkʰʁɑfd}（\textit{g} 并入鼻音 \ipa{ŋ}）。
  \item \textbf{关键修饰词与外借词读音}：\textit{hårfin} \ipa{ˈhɒːˌfiːˀn}（\textit{hår} 长元音 \ipa{ɒː}，\textit{fin} 带 \textit{stød}，切勿与 \textit{hårfager} 混淆！）；\textit{chikane} \ipa{ʃiˈkʰæːnə}（\textit{ch} 读作 \ipa{ʃ}）；\textit{vilkårlighed} \ipa{velˈkʰɒːliˌheðˀ}（重音在第二音节 \textit{kår}）。
  \item \textbf{辩证语篇的对比重读}：朗读 \textit{Principielt set...} 与 \textit{I praksis...} 时，副词短语首重音分别落在 \textit{principi\textbf{elt}} 与 \textit{\textbf{prak}sis} 上，随后刻意略作停顿再引出 V2 动词，形成鲜明的逻辑对仗节奏。
\end{itemize}
\end{tipbox}

\begin{transbox}
\small
\textbf{\color{NCDTeal}[Tekst B 中文对照]}\\
作为一名报刊编辑，我每天都在数字评论区见到充满敌意的留言；针对具名个人的直接威胁理应毫不犹豫地依法追究。即便如此，我仍强烈警告大家切勿以为仅靠立法就能摆脱网络上粗鄙且极化的言语氛围。

从原则上讲，在一个自由社会里，国家绝不能充当“品味裁判官”，去裁定哪些观点在道德上是可以接受的。如果丹麦议会用含糊不清的禁令去扩大刑法对冒犯性语言的惩罚范围，我们就会面临这样的风险：政治人物、公共评论员和学者出于纯粹谨慎而对敏感议题保持沉默。而在实践中，宽泛的禁令也极少带来更多的尊重：相反，它们会把挑衅者塑造成“殉道者”，并将极端观点推入无人再能够反驳的封闭网络中。

因此，在法律制度的边界上，我们应当坚守坚定不移的原则立场；同时通过积极的编辑部审核管理以及勇敢的针锋相对，在体面的公共对话与言论自由之间取得平衡。
\end{transbox}

% ---------------------------------------------------------------------
\ncdsubsec{三、Gloser（核心生词表）}

\begin{tipbox}{词表重点}
注意法理与伦理高频词的精准构词与复数形态：\textit{en drivkraft} 在后续课中会遇变音复数，本课的 \textit{en sammenhængskraft}（社会凝聚力）与 \textit{en principfasthed}（原则坚定性）均为不可数抽象名词（\textit{kun ental}）。动词 \textit{at trække}（拉、划出界线）属强变化不规则动词（\textit{trækker, trak, har trukket}）；\textit{at kvæle}（扼杀、窒息）也是不规则弱变化（\textit{kvæler, kvalte, har kvalt}）。
\end{tipbox}

\begin{glosertable}
名词 & en ytringsfrihed & \textit{(ytringsfriheden, kun ental)} & 言论自由 \\
名词 & en censur & \textit{(censuren, kun ental)} & 审查制度，新闻检查 \\
名词 & en principfasthed & \textit{(principfastheden, kun ental)} & 原则坚定性，原则立场 \\
名词 & en sammenhængskraft & \textit{(sammenhængskraften, kun ental)} & 社会凝聚力 \\
名词 & en chikane & \textit{(chikanen, chikaner, chikanerne)} & 骚扰，恶意刁难 \\
名词 & en forargelse & \textit{(forargelsen, kun ental)} & 义愤，激愤，众怒 \\
名词 & en vilkårlighed & \textit{(vilkårligheden, vilkårligheder, vilkårlighederne)} & 专断，任意性 \\
名词 & en værdighed & \textit{(værdigheden, kun ental)} & 尊严，体面 \\
名词 & en selvjustits & \textit{(selvjustitsen, kun ental)} & 自律，内部约束（亦指私力救济） \\
名词 & en straffelov & \textit{(straffeloven, straffelove, straffelovene)} & 刑法，刑事法规 \\
名词 & en intimidering & \textit{(intimideringen, intimideringer, intimideringerne)} & 恐吓，威慑 \\
名词 & et hensyn & \textit{(hensynet, hensyn, hensynene)} & 考量，顾及（单复同形） \\
\midrule
动词 [har] & at trække & \textit{(trækker, trak, har trukket)} & 划出（界线），拉 \\
动词 [har] & at balancere & \textit{(balancerer, balancerede, har balanceret)} & 平衡，权衡 \\
动词 [har] & at chikanere & \textit{(chikanerer, chikanerede, har chikaneret)} & 骚扰，纠缠攻击 \\
动词 [har] & at kvæle & \textit{(kvæler, kvalte, har kvalt)} & 窒息，扼杀 \\
\midrule
形容词 & hårfin & \textit{(hårfint, hårfine)} & 极其微妙的，毫厘之间的 \\
形容词 & ukuelig & \textit{(ukueligt, ukuelige)} & 坚定不移的，不屈的 \\
\end{glosertable}

% ---------------------------------------------------------------------
\ncdsubsec{四、Grammatik \& Stilistik（原则与实践辩证架构、法理搭配与 C1 避坑）}

\begin{grammatikbox}{语法点 1：原则与实践的双轨辩证架构（\textit{Principielt set} $\leftrightarrow$ \textit{I praksis}）}
在 PD3 与 Studieprøven 探讨民主权利、司法公正或技术管制时，“先确立抽象价值基准，再检验现实执行困境”是最高级的论证推进方式。请注意句首副词短语引发的 \textbf{V2 倒装铁律}：

\vspace{0.3em}
\renewcommand{\arraystretch}{1.24}
{\small
\begin{tabularx}{\linewidth}{@{} >{\bfseries\raggedright\arraybackslash}p{3.4cm} >{\raggedright\arraybackslash}p{4.3cm} >{\raggedright\arraybackslash}X @{}}
\toprule
\textbf{辩证路标模块} & \textbf{功能与语序说明} & \textbf{Tekst A/B 原句示范} \\
\midrule
\da{Principielt set} + V2 动词 &
从宪法、伦理或制度原则切入；句首占用第 1 位，\textbf{无逗号}，直跟变位动词。 &
\textit{\da{Principielt set} hviler et modent demokrati på en \dabold{ukuelig principfasthed}.} \\
\midrule
\da{I praksis} + V2 动词 \newline (+ \textit{imidlertid}) &
陡然转入现实复杂性；常在主句动词后配 \textit{imidlertid}（然而）或 \textit{desværre}。 &
\textit{\da{I praksis} udspiller den offentlige debat sig \da{imidlertid} sjældent roligt.} \\
\midrule
\da{På den ene side ...} \newline \da{På den anden side ...} &
承接原则与现实两端，用于深入对比两套相互牵制的法理顾虑。 &
\textit{\da{På den ene side} sikrer principfasthed mod vilkårlighed...} \\
\bottomrule
\end{tabularx}
}
\end{grammatikbox}

\begin{grammatikbox}{语法点 2：法理权衡与界线划定三大 C1 核心搭配}
\begin{enumerate}
  \item \textbf{\dabold{trække en hårfin grænse mellem A og B}（在 A 与 B 之间划出微妙界线）}：\newline
    固定动词是 \textit{at trække}（\textit{trak, har trukket}），不能误用中文直译动词 \textit{tegne / lave}。形容词必须精准拼作 \dabold{hårfin}（发丝般纤细的 $\rightarrow$ 毫厘之间、极其微妙的）。例如：\newline
    \ding{51}\ \textit{Domstolene må \da{trække en hårfin grænse mellem} skarp kritik og ulovlig chikane.}
  \item \textbf{\dabold{balancere hensynet til A over for B}（在 A 的考量与 B 之间求得平衡）}：\newline
    注意双介词接口：第一项考量前固定用 \dabold{til}（注意定指 \textit{hensynet til}），对立的另一项原则前用 \dabold{over for}（两个分开写的词！）。例如：\newline
    \ding{51}\ \textit{Staten skal \da{balancere hensynet til} ytringsfriheden \da{over for} borgernes tryghed.}
  \item \textbf{\dabold{en ukuelig principfasthed}（坚定不移的原则立场）}：\newline
    共性名词，配不定冠词 \textit{en}，形容词 \textit{ukuelig} 词尾 \textit{-g} 不发音，用来刻画“绝不因一时舆论风向轻易动摇的宪政坚守”。
\end{enumerate}
\end{grammatikbox}

\begin{grammatikbox}{语法点 3：C1 法理写作高频排雷清单（Top 4 Pitfalls）}
\begin{enumerate}
  \item \textbf{\ding{172} 严防形近错词：\textit{hårfin} $\neq$ \textit{hårfager}！}：\textit{hårfager} 是古雅历史形容词，指“拥有一头秀发的”（如挪威国王 Harald Hårfager）。表达“微乎其微、极其微妙的界线”，\textbf{只能写 \da{hårfin}}：\ding{55}\ \textit{trække en hårfager grænse} $\rightarrow$ \ding{51}\ \textit{trække en \da{hårfin} grænse}。
  \item \textbf{\ding{173} \textit{Principielt set} 与 \textit{I praksis} 后严禁加逗号}：受英语 \textit{In principle, ...} 影响容易误加逗号，但丹麦语中这两组路标短语后直接跟动词：\ding{55}\ \textit{Principielt set, friheden er vigtig.} $\rightarrow$ \ding{51}\ \textit{\da{Principielt set er} friheden vigtig.}
  \item \textbf{\ding{174} 省略 \textit{hvis} 的条件倒装被动句}：Tekst A 中的 \textit{Trækkes grænsen for snævert, kvæles den nødvendige uenighed} 极其凝练——前句用 \textit{-s} 被动动词放句首代替 \textit{Hvis grænsen trækkes...}，后句随即紧跟主句动词 \textit{kvæles} 触发 V2 倒装！
  \item \textbf{\ding{175} \textit{over for} 必须分写，不可误作 \textit{overfor}}：在“权衡 A 与 B”以及表达“面对某对象（\textit{over for domstolene}）”时，规范介词词组写作分开的 \da{over for}（连写的 \textit{overfor} 仅作空间副词“在对面”）。
\end{enumerate}
\end{grammatikbox}

% ---------------------------------------------------------------------
\ncdsubsec{五、Mønsterøvelser（句型与修辞操练）}

\textbf{练习 A：核心法理搭配与辩证连接词精准填空（每空填入最准确的 1 个丹麦语词）}
\begin{enumerate}
  \item \underline{\hspace{2.2cm}} set skal enhver borger have ret til at kritisere regeringen uden frygt for censur.
  \item I \underline{\hspace{2.2cm}} ser debatten på sociale medier dog ofte langt mere barsk og uforsonlig ud.
  \item Domstolene må trække en \underline{\hspace{2.2cm}} grænse mellem lovlig satire og strafbar hadtale.
  \item Staten bør balancere hensynet \underline{\hspace{2.2cm}} den enkeltes frihed over for beskyttelsen af udsatte grupper.
  \item Grundloven kræver på den ene side en ukuelig \underline{\hspace{2.2cm}}, når myndigheder vil indskrænke debatten. (原则坚定性)
  \item Hvis borgere lægger \underline{\hspace{2.2cm}} på sig selv af frygt for chikane, svækkes folkestyret. (克制、束缚)
  \item Systematiske trusler online kan med tiden udhule hele samfundets \underline{\hspace{2.2cm}}. (社会凝聚力)
  \item Lovgivere skal passe på, at brede forbud ikke \underline{\hspace{2.2cm}} den nødvendige uenighed i samfundet. (kvæle，现在时)
\end{enumerate}

\textbf{练习 B：C1 句法纠错与辩证语体升级（修正错词/语序，或按提示升级句子）}
\begin{enumerate}
  \item 【形近词与搭配纠错】\textit{Lovgiverne prøver at lave en hårfager grænse mellem kritik og chikane.} $\rightarrow$
  \item 【句首副词 V2 倒装纠错】\textit{Principielt set, alle holdninger bør kunne ytres frit, men i praksis folk føler sig truet.} $\rightarrow$
  \item 【介词接口纠错】\textit{Kommunen skal balancere hensynet om borgernes tryghed overfor ytringsfriheden.} $\rightarrow$
  \item 【省略 \textit{hvis} 被动条件升级】用省略 \textit{hvis} 的对称句式改写：\textit{Hvis grænsen bliver trukket for bredt, bliver sårbare grupper svigtet.} $\rightarrow$
  \item 【名词化与权威表达升级】用 \textit{på en ukuelig principfasthed} 升级：\textit{Vi skal stå meget fast på vores principper, når vi forsvarer friheden.} $\rightarrow$
  \item 【双轨辩证连接升级】将两句合并为 \textit{Principielt set ... i praksis ... imidlertid ...} 句式：\textit{Ytringsfriheden gælder for alle i teorien. Mange forskere tier i virkeligheden på grund af digital chikane.} $\rightarrow$
\end{enumerate}

\textbf{练习 C：中译丹实战（灵活运用原则与实践辩证结构，控制句长精准表达）}
\begin{enumerate}
  \item 原则上，公民应当有权在不受审查的情况下表达不受欢迎的观点。
  \item 然而在实践中，社交平台上的匿名骚扰正威胁着社会凝聚力。
  \item 法治国家必须在尖锐批评与非法威胁之间划出极其微妙的界线。
  \item 我们应当在个人言论自由的考量与弱势群体的尊严之间求得平衡。
  \item 一方面需要坚定不移的原则立场，另一方面也需要数字自律。
  \item 倘若界线划得过窄，公开辩论中的必要分歧就会遭到扼杀。
\end{enumerate}

% ---------------------------------------------------------------------
\ncdsubsec{六、PD3 / Studieprøven 实战微任务与参考答案}

\begin{emailbox}{PD3 / Studieprøven 写作微任务（Debatartikel · 200--250+ ord）}
\textbf{Situation}: På et dansk universitet foreslår en gruppe studerende et nyt ordensreglement, der forbyder ”stødende og krænkende politiske udtalelser” i alle campusbygninger og digitale studiegrupper. Forslaget møder både støtte og skarp kritik.\\[3pt]
\textbf{Opgave}: Skriv en argumenterende kommentar på \textbf{200--250 ord}, hvor du diskuterer forslaget ud fra spændingen mellem princip og praksis. Du skal mindst inddrage \textbf{fire} af lektionens kerneudtryk (\textit{principielt set, i praksis, ukuelig principfasthed, trække en hårfin grænse mellem, balancere hensynet til ... over for ..., sammenhængskraft}) og afslutte med en velbegrundet konklusion.
\end{emailbox}

\begin{answerbox}
\textbf{Før du læser 参考回答}
\begin{enumerate}
  \item \textit{Ifølge Tekst A garanterer Grundlovens § 77 ytringsfrihed uden forudgående censur under ansvar for domstolene. Principielt set kræver demokratiet ukuelig principfasthed for upopulære holdninger, men i praksis udfordres debatten af digital chikane, der truer samfundets sammenhængskraft.}
  \item \textit{Henrik advarer mod vage forbud mod krænkende sprog, fordi staten ikke må være smagsdommer, og fordi brede forbud i praksis kan få forskere til at tie af forsigtighed samt gøre provokatører til martyrer i lukkede netværk.}
\end{enumerate}

\textbf{Tekst B 阅读核对答案}
\begin{enumerate}
  \item \da{Rigtigt}（亨里克明确指出针对具名个人的直接威胁应毫不犹豫依法追诉：\textit{skal naturligvis retsforfølges uden tøven}）
  \item \da{Forskere/politikere tier af forsigtighed, og provokatører bliver martyrer i lukkede netværk.}（专业人士出于谨慎沉默，挑衅者成为殉道者并转入封闭网络）
  \item \da{Gennem aktiv redaktionel moderation og modige modargumenter.}（通过积极的编辑部审核与勇敢的反驳论证）
\end{enumerate}

\textbf{练习 A 答案}
\begin{enumerate}
  \item \da{Principielt}（固定辩证开头：\textit{Principielt set} 从原则上讲）
  \item \da{praksis}（对应转折路标：\textit{I praksis} 在实践中）
  \item \da{hårfin}（固定搭配：\textit{trække en hårfin grænse mellem...}，不可误作 \textit{hårfager}）
  \item \da{til}（双介词接口：\textit{balancere hensynet til A over for B}）
  \item \da{principfasthed}（固定搭配：\textit{en ukuelig principfasthed} 坚定不移的原则立场）
  \item \da{bånd}（习惯搭配：\textit{lægge bånd på sig selv} 自我约束）
  \item \da{sammenhængskraft}（高阶抽象名词：\textit{samfundets sammenhængskraft} 社会凝聚力）
  \item \da{kvæler}（动词 \textit{at kvæle} 的现在时，注意条件句结果用一般现在时）
\end{enumerate}

\textbf{练习 B 答案（参考）}
\begin{enumerate}
  \item \textit{Lovgiverne prøver at \da{trække en hårfin grænse} mellem kritik og chikane.}（改用正确搭配动词 \textit{trække}，并将错词 \textit{hårfager} 修正为 \textit{hårfin}）
  \item \textit{\da{Principielt set bør} alle holdninger kunne ytres frit, men \da{i praksis føler} folk sig truet.}（删去句首副词后多余逗号，并在两部分主句中严格执行 V2 动词提前倒装）
  \item \textit{Kommunen skal \da{balancere hensynet til} borgernes tryghed \da{over for} ytringsfriheden.}（\textit{hensynet} 固定接 \textit{til}，对比项前用分写的 \textit{over for}）
  \item \textit{\da{Trækkes} grænsen for bredt, \da{svigtes} sårbare grupper.}（用省略 \textit{hvis} 的 \textit{-s} 被动倒装，形成对称凝练的法理表达）
  \item \textit{Forsvaret for friheden hviler \da{på en ukuelig principfasthed}.}（以抽象名词化短语替换平淡动词短语，显著提升 C1 书面语感）
  \item \textit{\da{Principielt set} gælder ytringsfriheden for alle, men \da{i praksis} tier mange forskere \da{imidlertid} på grund af digital chikane.}（运用双轨辩证副词与 \textit{imidlertid} 构建紧扣张力的复合长句）
\end{enumerate}

\textbf{练习 C 答案}
\begin{enumerate}
  \item \textit{Principielt set bør borgerne have ret til at ytre upopulære holdninger uden censur.}
  \item \textit{I praksis truer anonym chikane på sociale medier imidlertid samfundets sammenhængskraft.}
  \item \textit{Retsstaten må trække en hårfin grænse mellem skarp kritik og ulovlige trusler.}
  \item \textit{Vi bør balancere hensynet til individets ytringsfrihed over for sårbare gruppers værdighed.}
  \item \textit{På den ene side kræves der ukuelig principfasthed, på den anden side digital selvjustits.}
  \item \textit{Trækkes grænsen for snævert, kvæles den nødvendige uenighed i den offentlige debat.}
\end{enumerate}

\textbf{PD3 / Studieprøven 高分范文（约 230 词，不含标题）}\\
\textit{\textbf{Universitetet bør beskytte både ubehagelige tanker og trygge rammer}}

\textit{Forslaget om at forbyde ”stødende og krænkende” politiske udtalelser på campus rejser et klassisk og afgørende spørgsmål om ytringsfrihedens grænser. \da{Principielt set} bør et universitet være det sidste sted i et demokratisk samfund, der indfører snævre regler for, hvilke meninger der må komme til orde i undervisningen. \da{På den ene side} kræver fri forskning og kritisk tænkning en \dabold{ukuelig principfasthed}: Studerende skal netop udfordres af ubehagelige og uvante perspektiver for at skærpe deres egne argumenter. Hvis man straffer upræcist definerede krænkelser, risikerer man hurtigt udbredt selvcensur i både auditorier og studiegrupper.}

\textit{\da{I praksis} kan vi imidlertid ikke ignorere, at digitale studiegrupper og sociale platforme ofte forstærker personangreb og hårde fronter. Når minoritetsstuderende eller forskere udsættes for systematisk chikane og nedværdigelse, handler det ikke længere om fri idéudveksling, men om målrettet eksklusion, der svækker institutionens \da{sammenhængskraft} og tillid. Derfor må ledelsen \dabold{trække en hårfin grænse mellem} skarp politisk uenighed om idéer på den ene side og personlig intimidering på den anden side.}

\textit{Løsningen er derfor ikke et bredt forbud mod ”stødende” holdninger, som let kan misbruges til at lukke munden på upopulære mindretal eller politiske modstandere. I stedet bør universitetet \dabold{balancere hensynet til} den åbne akademiske debat \dabold{over for} studerendes ret til et trygt og chikanefrit læringsmiljø. Konkret bør eksisterende regler mod direkte trusler og personforfølgelse håndhæves konsekvent, mens kontroversielle politiske synspunkter fortsat mødes med åbne og saglige modargumenter.}

\small\textbf{名师点评}：正文（不含标题）约 230 词，完美契合 PD3 / Studieprøven 辩论短文篇幅与 C1 评分标准。文章以“高校言论边界”切入，第一段借助 \textit{Principielt set} 与 \textit{På den ene side ... ukuelig principfasthed} 立起学术自由基准；第二段以 \textit{I praksis ... imidlertid} 切入现实排斥危机，顺势引出 \textit{trække en hårfin grænse mellem}；第三段摒弃一刀切禁令，利用 \textit{balancere hensynet til ... over for ...} 提出“严惩人身恐吓、宽容观点交锋”的可行方案，逻辑首尾咬合，法理分寸恰到好处！
\end{answerbox}
\ncdlesson{10}{Præstationssamfundet og jagten på det perfekte liv}{绩效社会与对完美生活的追逐}
{DU3 Modul 5 (PD3) $\rightarrow$ Modul 6 (Studieprøven, B2 $\rightarrow$ C1)｜对标 NCE 4 \& 现代社会心理学与文化哲学评论}
{\textbf{哲学与心理学思辨评论}（连续设问推进 \textit{Hvornår blev...? Er det i grunden...?} · 概念拆解句型 \textit{Begrebet dækker ikke blot over..., men...} · 悖论揭示 \textit{Paradokset består i, at...}）· 韩炳哲（Byung-Chul Han）《倦怠社会》语汇体系}
{欢迎来到第 10 课！在高福利的丹麦，为什么越来越多年轻人仍感到喘不过气，陷入“12 分女孩（\textit{12-talspiger}）”、完美履历与职业倦怠（\textit{udbrændthed}）的焦虑漩涡？在 PD3 与 Studieprøven 中，哲学与心理反思类评论最忌空洞说教，而最推崇\textbf{以层层递进的连续设问刺破表象}，再用\textbf{严密的概念拆解与悖论揭示}展开深度推演。本课借助韩炳哲（Byung-Chul Han）的\textbf{“绩效社会”（\textit{præstationssamfund}）}理论，全面升级你的 C1 思辨笔法！}

% ---------------------------------------------------------------------
\begin{foerduleaeser}
带着下面两个问题阅读 Tekst A 和 Tekst B：
\begin{enumerate}
  \item \textit{Hvordan adskiller filosoffen Byung-Chul Han ifølge Tekst A det moderne præstationssamfund fra det traditionelle disciplinsamfund, og hvorfor skaber selvoptimering en følelse af utilstrækkelighed?}（根据 Tekst A，哲学家韩炳哲如何区分现代“绩效社会”与传统“规训社会”？为什么“自我优化”反而催生了持续的匮乏感与不足感？）
  \item \textit{Hvad mener psykologen Kasper i Tekst B med ”evalueringsfrie rum”, og hvorfor påpeger Maja, at individet ikke kan løse præstationspresset alene?}（在 Tekst B 中，心理学家卡斯珀所说的“免于评价的空间”指什么？为什么玛雅指出仅靠个人无法独自化解绩效压力？）
\end{enumerate}
\end{foerduleaeser}

% ---------------------------------------------------------------------
\ncdsubsec{二、Tekster（哲学思辨专栏与校园心理对话双语篇）}

\subsection*{Tekst A：\textit{Når frihed bliver et nyt fængsel}（哲学与心理学思辨长文 · 约 370 词）}

\begin{tipbox}{语篇说明}
本篇是一篇极具思辨锋芒的文化哲学评论。阅读时请重点标出\textbf{\color{NCDNavy}深蓝色标注的连续设问推进与概念界定框架}，以及\textbf{\color{NCDRed}深红色标注的绩效社会核心概念与悖论转折词}。
\end{tipbox}

\da{Hvornår blev} hverdagen et evigt projekt, der hele tiden skal forbedres? \da{Hvornår blev} fritiden reduceret til endnu en arena for optimering af krop, netværk og CV? \da{Er det i grunden} ægte frihed, vi oplever i dag, eller blot en mere subtil form for indre tvang? I et velfærdssamfund som det danske, hvor de fleste materielle grundbehov er dækket, burde roen egentlig sænke sig. Alligevel vokser antallet af unge og voksne, der rammes af stress, søvnløshed og dyb \dabold{udbrændthed}.

For at forstå denne gåde peger den sydkoreansk-tyske filosof Byung-Chul Han på et afgørende historisk skift. I det klassiske \dabold{disciplinsamfund} blev mennesket styret udefra gennem forbud, autoriteter og klare pligter: ”Du må ikke” og ”Du skal”. Det moderne \dabold{præstationssamfund} taler derimod med en tilsyneladende frigørende stemme: ”Du kan alt!” \da{Begrebet dækker ikke blot over} stigende krav på arbejdsmarkedet eller jagten på topkarakterer i gymnasiet, \da{men over} en gennemgribende mentalitet, hvor individet gør sig selv til både entreprenør og arbejder i sit eget liv.

\da{Paradokset består i, at} ydre undertrykkelse dermed afløses af frivillig \dabold{selvoptimering}. Når vi konstant måler vores værd i tolv-taller, løbetider, sunde madpakker og polerede familiebilleder på sociale medier, opstår der en snigende \dabold{selvudnyttelse}. \da{Er det i grunden} muligt at gøre oprør mod en chef, når den hårdeste pisk ligger inde i ens egen bevidsthed? \da{Når man skræller overfladen væk, viser det sig, at} den uendelige jagt på det perfekte liv sjældent bringer varig ro --- tværtimod efterlader den et kronisk tomrum af \da{utilstrækkelighed}.

Denne psykologiske dynamik rammer især den unge generation. Hvor tidligere tider skabte klare grænser mellem arbejde og hvile, udviskes skellet i dag af digitale platforme og uophørlig sammenligning. Selv hvilen risikerer at blive instrumentaliseret: Man går ikke længere en tur i skoven for turens egen skyld, men for at ”lade batterierne op” til næste dags præstation. Når ethvert øjeblik skal retfærdiggøres som en investering i fremtiden, mister mennesket evnen til blot at være til stede i nuet.

Vejen ud af denne udmattelse begynder næppe med endnu en effektiv app til tidsstyring. Vi må i fællesskab genvinde modet til at sætte grænser, acceptere ufuldkommenhed og vælge langsommeligheden til. Først når værdighed ikke længere måles i resultater, kan friheden igen blive et åndehul frem for et usynligt fængsel.

\begin{transbox}
\small
\textbf{\color{NCDTeal}[Tekst A 中文精译]}\\
日常生活究竟是从何时起，变成了一个必须不断自我改良的永恒工程？闲暇时光又是从何时起，沦为了又一个用来优化身材、人脉和简历的竞技场？我们今天体验到的，归根结底究竟是真正的自由，还是只是一种更加隐蔽的内在强迫？在像丹麦这样物质基本需求大多已得到满足的福利社会里，平静按理说本该降临。然而，遭受压力、失眠与深层职业倦怠困扰的年轻人与成年人却反而日益增多。

为了理解这一谜题，韩裔德国哲学家韩炳哲指出了一个决定性的历史转向。在传统的规训社会中，人受到外部禁令、权威和明确义务的支配：“你不准”以及“你必须”。相比之下，现代绩效社会却用一种看似解放人性的声音说道：“你无所不能！”这一概念所涵盖的，并不仅仅是劳动力市场上日益严苛的要求或高中阶段对高分成绩的追逐，更是一种全面渗透的心态——在其中，个体将自己同时塑造成自己人生中的创业者与打工人。

悖论恰恰在于，外部压迫由此被自愿的“自我优化”所取代。当我们不断用满分 12 分、跑步配速、健康便当以及社交媒体上光鲜修饰的家庭照来衡量自身价值时，一种悄然蔓延的“自我剥削”便产生了。当最严厉的鞭子就握在自己脑海里时，人从根本上还能反抗老板吗？剥开表层之后便会发现，对完美生活无休止的追逐极少带来持久的安宁——恰恰相反，它留下的是一种慢性的“自我不足感”与空虚。

这种心理机制对年轻一代的冲击尤为剧烈。以往的时代曾在工作与休息之间划出清晰边界，而如今这道界线却被数字平台与无休止的横向攀比模糊了。就连休息本身也面临被工具化的危险：人们去森林散步不再是为了散步本身，而是为了给明天的绩效“充电”。当每一个瞬间都必须被证明是对未来的投资时，人就失去了纯粹安住当下的能力。

走出这种疲惫的出路，几乎不可能始于又一款高效的时间管理应用。我们必须共同重拾勇气，为生活划定边界、接纳不完美并主动选择慢下来。唯有当人的尊严不再用可量化的成绩来衡量时，自由才能重新成为喘息的空间，而非一座无形的牢笼。
\end{transbox}

\subsection*{Tekst B：\textit{Samtale i studierådgivningen: ”Man kan ikke hvile på bestilling”}（校园心理咨询对话 · 约 155 词）}

\textbf{阅读任务（理解与核对）：}\textit{Læs samtalen mellem studenterpsykolog Kasper og den universitetsstuderende Maja, og besvar spørgsmålene nedenfor.}

\begin{emailbox}{Samtale på campus｜Kasper (studenterpsykolog) \& Maja (universitetsstuderende)}
\textbf{Maja:} Kasper, jeg læste om præstationssamfundet i går, og jeg genkendte næsten alt. Selv når jeg holder fri i weekenden, får jeg dårlig samvittighed over ikke at opdatere mit CV, træne igennem eller forberede næste seminar.

\textbf{Kasper:} Det hører jeg fra studerende næsten hver eneste dag, Maja. Mange tror, at problemet blot handler om for mange arbejdstimer, men i grunden skyldes presset, at næsten alle aktiviteter bliver vurderet. Derfor har unge desperat brug for evalueringsfrie rum, hvor man kan være sammen uden at præstere.

\textbf{Maja:} Men er det ikke lidt naivt at lægge hele ansvaret over på den enkelte? Så længe uddannelsessystemet og arbejdsmarkedet belønner konstant \dabold{selvoptimering}, bliver det jo dyrt for den studerende, der vælger langsommeligheden til.

\textbf{Kasper:} Du rammer hovedet på sømmet. Individet kan ikke bryde dette mønster alene; netop derfor må både universiteter og arbejdspladser skabe en ny kultur, hvor fejl, tøven og uforpligtende pauser betragtes som naturlige.
\end{emailbox}

\textbf{Tekst B 阅读核对（Rigtigt / Forkert / Kort svar，答案见第六部分）：}
\begin{enumerate}
  \item \textit{Maja oplever, at hun uden problemer kan slappe af i weekenden uden dårlig samvittighed.} \underline{\hspace{2.2cm}}
  \item \textit{Hvad skyldes presset i grunden ifølge Kasper, og hvad har unge derfor brug for?} \underline{\hspace{2.2cm}}
  \item \textit{Hvorfor mener Maja, at den enkelte studerende ikke bare kan vælge langsommeligheden til på egen hånd?} \underline{\hspace{2.2cm}}
\end{enumerate}

\begin{tipbox}{本课发音与语调要点（Udtale \& Prosodi）}
\begin{itemize}
  \item \textbf{拉丁借词与复合抽象名词重音}：\textit{præstationssamfund} \ipa{pʰʁɛsd̥aˈɕoːˀnsˌsɑmˀfʌnˀ}（\textit{-tion-} 读 \ipa{ˈɕoːˀn}，且 \textit{samfund} 尾音 \textit{d} 不发音带 \textit{stød}）；\textit{selvoptimering} \ipa{ˈsɛlˀʌbd̥iˌmeːˀɐeŋ}；\textit{udbrændthed} \ipa{ˈuðˌbʁɛnˀdheðˀ}（软 \textit{d} 加双处 \textit{stød}）。
  \item \textbf{高频思辨词组语音细节}：\textit{i grunden} \ipa{i ˈɡʁʌnˀən}（\textit{d} 不发音！读作鼻音加 \textit{stød}）；\textit{utilstrækkelighed} \ipa{utˢilˈsd̥ʁɛɡəliˌheðˀ}（重音在 \textit{-stræk-}，结尾 \textit{g} 不发音）；\textit{paradoks} \ipa{pʰɑʁɑˈdʌɡs}（重音在最后音节 \textit{-doks}）。
  \item \textbf{连续设问的升调与节奏}：朗读以 \textit{Hvornår blev...} 开头的特殊疑问句时，丹麦语\textbf{句尾不扬反降}（核心重音落在 \textit{Hvor\textbf{når}} 和句末新信息词上）；而以 \textit{Er det i grunden...} 开头的是非问句句尾略微上扬，与前者形成极富哲学张力的语调落差。
\end{itemize}
\end{tipbox}

\begin{transbox}
\small
\textbf{\color{NCDTeal}[Tekst B 中文对照]}\\
\textbf{玛雅：}卡斯珀，我昨天读了一篇关于绩效社会的文章，里面写的一切我几乎都深有同感。哪怕在周末休息的时候，我也会因为没去更新简历、没有努力锻炼或者没提前准备下一次研讨课而感到内疚。

\textbf{卡斯珀：}玛雅，这话我几乎每天都能从学生那里听到。很多人以为问题仅仅出在学习或工作时长太长，但从根本上说，这种压力源于几乎所有活动都在被评判打分。因此，年轻人迫切需要“免于评价的空间”，在那里大家可以仅仅在一起相处而无需去表现绩效。

\textbf{玛雅：}可是把全部责任都推给个人，难道不是有点天真吗？只要教育体制和劳动力市场还在奖赏持续不断的自我优化，选择“慢下来”的学生就必然要付出高昂的代价呀。

\textbf{卡斯珀：}你一针见血地说到了关键。个体确实无法独自打破这种模式；正因如此，大学和用人单位都必须共同营造一种新文化，把犯错、犹豫以及不带功利目的的停顿视为自然之事。
\end{transbox}

% ---------------------------------------------------------------------
\ncdsubsec{三、Gloser（核心生词表）}

\begin{tipbox}{词表重点}
注意两处核心变化规律：\textit{en drivkraft}（驱动力）在变复数时元音发生变音（\textit{drivkræfter, drivkræfterne}），与 \textit{kraft $\rightarrow$ kræfter} 完全一致；强变化动词 \textit{at opstå}（产生、出现）在完成时中表示状态变化，\textbf{必须搭配助动词 \textit{er}}（\textit{er opstået}）！而 \textit{en udbrændthed}（职业倦怠）和 \textit{en utilstrækkelighed}（不足感）均为不可数抽象名词。
\end{tipbox}

\begin{glosertable}
名词 & et præstationssamfund & \textit{(præstationssamfundet, -samfund, -ene)} & 绩效社会，功绩社会 \\
名词 & et disciplinsamfund & \textit{(disciplinsamfundet, -samfund, -ene)} & 规训社会 \\
名词 & en selvoptimering & \textit{(selvoptimeringen, kun ental)} & 自我优化 \\
名词 & en selvudnyttelse & \textit{(selvudnyttelsen, kun ental)} & 自我剥削 \\
名词 & en udbrændthed & \textit{(udbrændtheden, kun ental)} & 精疲力竭，职业倦怠 \\
名词 & en utilstrækkelighed & \textit{(utilstrækkeligheden, kun ental)} & 不足感，力不从心感 \\
名词 & en tvang & \textit{(tvangen, kun ental)} & 强迫，强制力 \\
名词 & en drivkraft & \textit{(drivkraften, drivkræfter, drivkræfterne)} & 驱动力，内在动力 \\
名词 & en langsommelighed & \textit{(langsommeligheden, kun ental)} & 缓慢，慢节奏 \\
名词 & et paradoks & \textit{(paradokset, paradokser, paradokserne)} & 悖论，矛盾现象 \\
名词 & et åndehul & \textit{(åndehullet, åndehuller, åndehullerne)} & 喘息空间，避风港 \\
\midrule
动词 [har] & at præstere & \textit{(præsterer, præsterede, har præsteret)} & 表现，完成绩效 \\
动词 [har] & at optimere & \textit{(optimerer, optimerede, har optimeret)} & 优化，追求极致效率 \\
动词 [har] & at udnytte & \textit{(udnytter, udnyttede, har udnyttet)} & 剥削，利用 \\
动词 [har] & at retfærdiggøre & \textit{(retfærdiggør, -gjorde, har retfærdiggjort)} & 为……辩护，使合理化 \\
动词 [er] & at opstå & \textit{(opstår, opstod, er opstået)} & 产生，涌现，形成 \\
\midrule
形容词 & subtil & \textit{(subtilt, subtile)} & 微妙的，隐蔽难察的 \\
形容词 & kronisk & \textit{(kronisk, kroniske)} & 慢性的，长期持续的 \\
\end{glosertable}

% ---------------------------------------------------------------------
\ncdsubsec{四、Grammatik \& Stilistik（连续设问推进、概念拆解与 C1 哲学修辞）}

\begin{grammatikbox}{语法点 1：连续设问推动深度论述（\textit{Retoriske spørgsmål som argumentationsmotor}）}
在哲学与社会心理学评论中，平铺直叙容易流于干瘪；通过\textbf{排比式连续设问}（2 至 3 句紧扣的追问）直接拷问习以为常的生活前提，能瞬间确立思想深度：

\vspace{0.3em}
\renewcommand{\arraystretch}{1.24}
{\small
\begin{tabularx}{\linewidth}{@{} >{\bfseries\raggedright\arraybackslash}p{3.4cm} >{\raggedright\arraybackslash}p{4.3cm} >{\raggedright\arraybackslash}X @{}}
\toprule
\textbf{设问修辞层} & \textbf{句式功能与语序结构} & \textbf{Tekst A 高分示范} \\
\midrule
\textbf{1. 时间溯源设问}\newline \da{Hvornår blev...?} &
连用两个 \textit{Hvornår blev X til Y?}，揭示“习以为常”其实只是晚近历史产物。 &
\textit{\da{Hvornår blev} hverdagen et evigt projekt? \da{Hvornår blev} fritiden reduceret til...?} \\
\midrule
\textbf{2. 本质追问设问}\newline \da{Er det i grunden...?} &
直接质疑“自由/进步”等漂亮口号是否名实相符，准备引出概念翻转。 &
\textit{\da{Er det i grunden} ægte frihed, vi oplever i dag, eller blot en mere subtil form for tvang?} \\
\midrule
\textbf{3. 反诘悖论设问}\newline \textit{Er det... muligt, når...?} &
用自身矛盾的处境反问，让读者自己推导出“系统内无处可逃”的结论。 &
\textit{\da{Er det i grunden} muligt at gøre oprør mod en chef, når den hårdeste pisk ligger inde i én?} \\
\bottomrule
\end{tabularx}
}
\end{grammatikbox}

\begin{grammatikbox}{语法点 2：概念拆解与悖论揭示四大核心句型}
\begin{enumerate}
  \item \textbf{否定表象 + 递进深挖：\da{Begrebet dækker ikke blot over A, men over B}}：\newline
    用 \textit{dække over}（涵盖、指代）做学术界定；注意对称连接词后若接名词短语，介词 \dabold{over} \textbf{必须重复出现}：\newline
    \ding{51}\ \textit{Begrebet dækker ikke blot \dabold{over} lange arbejdsdage, men \dabold{over} en indre drivkraft.}
  \item \textbf{揭示反常机制：\da{Paradokset består i, at...}}（悖论在于……）：\newline
    动词 \textit{bestå} 必须搭配介词 \dabold{i} 再接 \textit{at}-从句（内部不倒装）：\newline
    \ding{51}\ \textit{Paradokset består \dabold{i, at} ydre undertrykkelse afløses af frivillig selvoptimering.}
  \item \textbf{剥离表象隐喻：\da{Når man skræller overfladen væk, viser det sig, at...}}（剥开表层就会发现……）：\newline
    句首长条件时间从句占据第 1 位，因此主句必须紧跟变位动词 \dabold{viser}（V2 倒装）：\newline
    \ding{51}\ \textit{Når man skræller overfladen væk, \dabold{viser det sig}, at jagten skaber utilstrækkelighed.}
  \item \textbf{历史对照界定：\da{I det klassiske A ... Det moderne B taler derimod...}}：\newline
    通过 \textit{derimod}（相比之下，置于主句动词后）把传统规训与现代绩效并置对比。
\end{enumerate}
\end{grammatikbox}

\begin{grammatikbox}{语法点 3：哲学评论常见中丹语序与搭配误区}
\begin{enumerate}
  \item \textbf{\ding{172} \textit{i grunden} 放主句与放从句的位置差异}：在主句是非问句里放在主语后（\textit{Er det \da{i grunden} frihed?}）；在普通主句中置于动词后（\textit{Det handler \da{i grunden} om vurdering.}）；在 \textit{at}-从句中置于动词前（\textit{...at friheden \da{i grunden} er en ny tvang.}）。
  \item \textbf{\ding{173} \textit{bestå i} vs \textit{bestå af}}：\textit{bestå af} 指“由具体零部件/成员组成”；阐述抽象道理、核心原因或悖论所在时，\textbf{只能用 \da{bestå i}}（不能误作 \ding{55}\ \textit{Paradokset består af, at...}）。
  \item \textbf{\ding{174} \textit{at opstå} 与 \textit{at forsvinde} 必配助动词 \textit{er}}：凡表示状态“产生/消失”的不可及物变化动词，完成时均用 \textit{er}，不可误作 \ding{55}\ \textit{der har opstået en ny træthed} $\rightarrow$ \ding{51}\ \textit{der \da{er opstået} en ny træthed}。
  \item \textbf{\ding{175} 分清 \textit{i stedet for} 与 \textit{frem for}}：C1 抽象价值抉择中，表达“宁可要 A 而不要 B（……而非……）”常用 \da{frem for}：\textit{et åndehul \da{frem for} et usynligt fængsel}。
\end{enumerate}
\end{grammatikbox}

% ---------------------------------------------------------------------
\ncdsubsec{五、Mønsterøvelser（句型与修辞操练）}

\textbf{练习 A：连续设问、概念拆解与绩效社会核心词精准填空（每空填入最准确的 1 个丹麦语词）}
\begin{enumerate}
  \item \underline{\hspace{2.2cm}} blev hverdagen et projekt, der hele tiden skal forbedres og optimeres?
  \item Er det i \underline{\hspace{2.2cm}} ægte frihed, når man aldrig tør holde fri uden dårlig samvittighed?
  \item Begrebet dækker ikke blot over lange arbejdstimer, men \underline{\hspace{2.2cm}} en dyb indre drivkraft.
  \item Paradokset består \underline{\hspace{2.2cm}}, at borgerne frivilligt overvåger og presser sig selv.
  \item Når man skræller overfladen væk, \underline{\hspace{2.2cm}} det sig, at jagten på perfektion skaber tomhed.
  \item I det moderne præstationssamfund er der \underline{\hspace{2.2cm}} en udbredt følelse af utilstrækkelighed. (opstå，完成分词)
  \item Konstant \underline{\hspace{2.2cm}} på sociale medier fører let til stress og udbrændthed hos unge. (自我优化)
  \item Unge har brug for et mentalt \underline{\hspace{2.2cm}}, hvor de ikke hele tiden bliver målt og vurderet. (喘息空间)
\end{enumerate}

\textbf{练习 B：哲学思辨句法纠错与修辞升级（修正搭配/语序错误，或按提示升级）}
\begin{enumerate}
  \item 【悖论介词纠错】\textit{Paradokset består af, at vi har mere fritid, men føler os mere stressede end før.} $\rightarrow$
  \item 【助动词与从句语序纠错】\textit{Artiklen forklarer, at der har opstået en ny form for selvudnyttelse blandt studerende.} $\rightarrow$
  \item 【对称介词与平行纠错】\textit{Udbrændthed dækker ikke blot over træthed i kroppen, men også en følelse af manglende værdi.} $\rightarrow$
  \item 【前置时间从句 V2 纠错】\textit{Når man skræller de polerede billeder væk, det viser sig, at mange kæmper med angst.} $\rightarrow$
  \item 【连续设问修辞升级】将平淡陈述改写为“\textit{Hvornår blev...}”修辞设问：\textit{Vores fritid er desværre blevet til en slags arbejde, hvor man skal optimere sin sundhed.} $\rightarrow$
  \item 【概念对照升级】用 \textit{I modsætning til ... er ... præget af ...} 整合句意：\textit{Disciplinsamfundet brugte ydre forbud. Præstationssamfundet bruger indre selvoptimering.} $\rightarrow$
\end{enumerate}

\textbf{练习 C：中译丹实战（运用连续设问与概念拆解表达，精准锤炼 C1 语感）}
\begin{enumerate}
  \item 日常生活是从何时起变成一项必须不断优化的工程的？
  \item 我们今天体验到的，归根结底究竟是自由还是一种隐蔽的强迫？
  \item 这一概念不仅涵盖高分要求，也涵盖持续不断的自我剥削。
  \item 悖论在于，外部指令已被现代人的自愿自我优化所取代。
  \item 剥开表层就会发现，对完美生活的追逐反而滋生出慢性不足感。
  \item 当任何休息都必须被证明是合理的，人们便失去了当下的喘息空间。
\end{enumerate}

% ---------------------------------------------------------------------
\ncdsubsec{六、PD3 / Studieprøven 实战微任务与参考答案}

\begin{emailbox}{PD3 / Studieprøven 写作微任务（Reflekterende essay · 200--250+ ord）}
\textbf{Situation}: En ny undersøgelse viser, at næsten hver tredje unge dansker oplever dagligt præstationspres i forbindelse med studier, fritid og sociale medier. På en ungdomsuddannelse diskuterer man nu, om man skal afskaffe karakterer i det første år for at dæmpe presset.\\[3pt]
\textbf{Opgave}: Skriv et reflekterende debatindlæg på \textbf{200--250 ord}, hvor du analyserer årsagerne til præstationspresset og tager stilling til forslaget. Brug mindst \textbf{to retoriske spørgsmål} til at drive argumentationen frem samt \textbf{tre} begrebsanalytiske udtryk fra lektionen (\textit{præstationssamfund, selvoptimering, paradokset består i, at..., i grunden, utilstrækkelighed, udbrændthed}).
\end{emailbox}

\begin{answerbox}
\textbf{Før du læser 参考回答}
\begin{enumerate}
  \item \textit{Ifølge Byung-Chul Han blev mennesket i disciplinsamfundet styret af ydre forbud og pligter, mens præstationssamfundet lokker med ”Du kan alt”, så individet selv bliver både chef og arbejder. Når værdien hele tiden måles gennem selvoptimering, opstår en kronisk følelse af utilstrækkelighed og udbrændthed.}
  \item \textit{Kasper mener, at unge har brug for evalueringsfrie rum, hvor de blot kan være sammen uden at blive vurderet. Maja påpeger imidlertid, at den enkelte studerende betaler en høj pris for at sige nej, så længe uddannelsessystemet og arbejdsmarkedet fortsat belønner konstant selvoptimering.}
\end{enumerate}

\textbf{Tekst B 阅读核对答案}
\begin{enumerate}
  \item \da{Forkert}（玛雅一到周末反而因为没有完善简历或预习课程而感到内疚：\textit{får jeg dårlig samvittighed}）
  \item \da{At næsten alle aktiviteter bliver vurderet; evalueringsfrie rum.}（根本原因在于几乎一切活动都在被评价打分，因此需要免于评价的空间）
  \item \da{Fordi uddannelsessystemet og arbejdsmarkedet belønner konstant selvoptimering.}（因为教育体制和劳动力市场仍在奖赏持续的自我优化）
\end{enumerate}

\textbf{练习 A 答案}
\begin{enumerate}
  \item \da{Hvornår}（时间溯源连续设问核心词：\textit{Hvornår blev...}）
  \item \da{grunden}（本质追问副词短语：\textit{i grunden} 归根结底、从根本上说）
  \item \da{over}（对称复现介词：\textit{dækker ikke blot over A, men over B}）
  \item \da{i}（固定搭配：\textit{bestå i, at...} 在于……，切勿误用 \textit{af}）
  \item \da{viser}（前置时间从句后触发主句 V2 倒装：\textit{viser det sig}）
  \item \da{opstået}（状态变化动词 \textit{at opstå} 与助动词 \textit{er} 构成现在完成时）
  \item \da{selvoptimering}（课内核心哲学心理概念名词：自我优化）
  \item \da{åndehul}（名词：喘息空间，避风港）
\end{enumerate}

\textbf{练习 B 答案（参考）}
\begin{enumerate}
  \item \textit{Paradokset består \da{i}, at vi har mere fritid, men føler os mere stressede end før.}（阐述“悖论在于……”须用 \textit{bestå i}，不可误用 \textit{bestå af}）
  \item \textit{Artiklen forklarer, at der \da{er opstået} en ny form for selvudnyttelse blandt studerende.}（\textit{opstå} 属状态变化动词，完成时必须搭配助动词 \textit{er}）
  \item \textit{Udbrændthed dækker ikke blot \da{over} træthed i kroppen, men også \da{over} en følelse af manglende værdi.}（注意 \textit{ikke blot ... men også ...} 并列两个名词短语时须对称复现介词 \textit{over}）
  \item \textit{Når man skræller de polerede billeder væk, \da{viser det sig}, at mange kæmper med angst.}（前置从句占第 1 位，主句变位动词 \textit{viser} 必须紧跟实现 V2 倒装）
  \item \textit{\da{Hvornår blev} vores fritid reduceret til et arbejde, hvor man skal optimere sin sundhed?}（借助连续设问将平淡叙述升级为尖锐哲学追问）
  \item \textit{\da{I modsætning til} disciplinsamfundet, der brugte ydre forbud, er præstationssamfundet \da{præget af} indre selvoptimering.}（运用 \textit{I modsætning til} 与关系从句完成概念对照拆解）
\end{enumerate}

\textbf{练习 C 答案}
\begin{enumerate}
  \item \textit{Hvornår blev hverdagen til et projekt, som man hele tiden skal optimere?}
  \item \textit{Er det i grunden frihed eller en subtil form for tvang, vi oplever?}
  \item \textit{Begrebet dækker ikke blot over høje karakterkrav, men også over konstant selvudnyttelse.}
  \item \textit{Paradokset består i, at ydre pligter er blevet erstattet af frivillig selvoptimering.}
  \item \textit{Når man skræller overfladen væk, nærer jagten på perfektion en kronisk utilstrækkelighed.}
  \item \textit{Hvis enhver pause skal retfærdiggøres, mister vi vores nødvendige åndehul i nuet.}
\end{enumerate}

\textbf{PD3 / Studieprøven 高分范文（约 235 词，不含标题）}\\
\textit{\textbf{Karakterfrihed er kun første skridt ud af præstationsfælden}}

\textit{\da{Hvornår blev} ungdomsårene forvandlet fra en tid med nysgerrighed, vildveje og fri eksperimenteren til et stramt optimeringsprojekt? \da{Er det i grunden} mangel på viljestyrke, der plager dagens unge, eller reagerer de blot fornuftigt på et grænseløst ydre og indre pres? Når næsten hver tredje unge dansker oplever dagligt stress i \dabold{præstationssamfundet}, peger forslaget om et karakterfrit første studieår i den rigtige retning --- men det løser næppe hele problemet alene.}

\textit{\da{Begrebet præstationspres dækker nemlig ikke blot over} frygten for en lav karakter på eksamensbeviset, \da{men over} en gennemgribende kultur, hvor selv fritid, venskaber og krop underkastes konstant \dabold{selvoptimering}. \da{Paradokset består i, at} jo mere frihed de unge formelt har fået til at forme deres egen fremtid, desto hårdere dømmer og overvåger de sig selv gennem sociale medier og urealistiske idealer. \da{Når man skræller overfladen væk, viser det sig, at} den udbredte \dabold{udbrændthed} og træthed ofte udspringer af en kronisk følelse af \da{utilstrækkelighed}.}

\textit{Et karakterfrit første år kan skabe et værdifuldt og nødvendigt åndehul, hvor fejl ikke straks straffes, og hvor samtale, leg og faglig fordybelse igen får plads uden målestok. Hvis vi imidlertid fortsat måler menneskelig værdi på hastighed, effektivitet og fejlfrie resultater i resten af uddannelsessystemet og arbejdslivet, flytter presset blot videre til andet år. Ægte forandring kræver derfor både konkrete reformer i skolen, evalueringsfrie rum i hverdagen og et modigt opgør med selve præstationskulturen uden for klasseværelset.}

\small\textbf{名师点评}：正文（不含标题）约 235 词，堪称哲学心理学思辨佳作。开篇两连设问（\textit{Hvornår blev...} 与 \textit{Er det i grunden...}）直击问题本质；中段以 \textit{Begrebet ... dækker ikke blot over..., men over...} 拆解“绩效压力”的深层结构，再用 \textit{Paradokset består i, at...} 与 \textit{Når man skræller overfladen væk, viser det sig, at...} 揭示年轻一代越自由越疲惫的深层悖论；末段辩证评估“高一取消打分制”的限度，既有理论高度又有教育现实感！
\end{answerbox}
\ncdlesson{11}{Folkeskolen mellem faglighed og dannelse}{公立学校在“学科能力”与“全人教养”之间}
{DU3 Modul 5 (PD3) $\rightarrow$ Modul 6 (Studieprøven, B2 $\rightarrow$ C1)｜对标 NCE 4 \& 丹麦教育制度与国民文化内核}
{\textbf{丹麦文化核心不可译词深度辨析}（\textit{dannelse} vs \textit{uddannelse} · \textit{tryghed} vs \textit{sikkerhed} · \textit{fællesskab} vs \textit{samfund}）· 格伦特维（Grundtvig）与哈尔·科赫（Hal Koch）民主对话传统 · 对比并列句式 \textit{Hvor A retter sig mod..., handler B om...}}
{为什么英文和中文都很难用一个词精准对译丹麦语中的 \textit{dannelse}、\textit{tryghed} 与 \textit{fællesskab}？这三组词恰恰构成了丹麦公立基础学校（\textit{folkeskolen}）乃至整个北欧福利社会的精神密码！在 PD3 与 Studieprøven 涉及教育改革、儿童成长或公民价值的命题中，能否精准区分 \textbf{\textit{dannelse}} 与 \textbf{\textit{uddannelse}}、\textbf{\textit{tryghed}} 与 \textbf{\textit{sikkerhed}}、\textbf{\textit{fællesskab}} 与 \textbf{\textit{samfund}}，并援引格伦特维（Grundtvig）与哈尔·科赫（Hal Koch）的民主传统，是考官评判 C1 文化理解深度的核心标尺。}

% ---------------------------------------------------------------------
\begin{foerduleaeser}
带着下面两个问题阅读 Tekst A 和 Tekst B：
\begin{enumerate}
  \item \textit{Hvordan definerer Tekst A forskellen mellem de tre danske begrebspar, og hvordan knytter Grundtvigs og Hal Kochs tanker an til folkeskolens opgave?}（Tekst A 是如何界定这三组丹麦语核心概念之差异的？格伦特维与哈尔·科赫的思想又如何与公立学校的使命相联系？）
  \item \textit{Hvorfor understreger læreren Mette og skolelederen Thomas i Tekst B, at kontrol og høje karakterer ikke i sig selv kan skabe tryghed og demokratisk dannelse?}（在 Tekst B 中，教师梅特和校长托马斯为什么强调管控手段与高分本身并不能自动造就心理安全感与民主教养？）
\end{enumerate}
\end{foerduleaeser}

% ---------------------------------------------------------------------
\ncdsubsec{二、Tekster（文化概念精析长文与教研室对话双语篇）}

\subsection*{Tekst A：\textit{Skolen for livet --- eller skolen for arbejdsmarkedet?}（教育文化与概念辨析长文 · 约 380 词）}

\begin{tipbox}{语篇说明}
本篇深入剖析丹麦公立教育百年争论背后的三组“不可译文化词”。阅读时请重点留意\textbf{\color{NCDNavy}深蓝色标注的三组核心文化概念对}，以及\textbf{\color{NCDRed}深红色标注的概念界定与对比论证结构}。
\end{tipbox}

I årtier har politikere, lærere og forældre diskuteret, hvad den danske folkeskole først og fremmest skal måles på. Bør undervisningen fokusere benhårdt på nationale test, Pisa-placeringer og snæver \da{faglighed}, eller skal skolen i lige så høj grad forme hele og kritiske medborgere? For at forstå kernen i denne debat må man undersøge tre danske nøglebegreber, som næsten er umulige at oversætte med et enkelt ord til andre sprog.

Det første og mest grundlæggende begrebspar er \da{uddannelse} over for \da{dannelse}. \dabold{Hvor} \da{uddannelse} \dabold{retter sig mod} målbare kvalifikationer, eksamener og færdigheder til arbejdsmarkedet, \dabold{handler} \da{dannelse} \dabold{om} menneskelig modning, dømmekraft og evnen til at orientere sig etisk i verden. Allerede i 1800-tallet skelnede N.F.S. Grundtvig skarpt mellem ”skolen for livet” og ”skolen for døden” og fremhævede levende folkelig oplysning frem for mekanisk udenadslære. Efter Anden Verdenskrig videreførte teologen Hal Koch denne arv med sit berømte budskab om, at demokrati ikke blot er en styreform med stemmesedler, men en daglig samtale og livsform, som børn må øve sig i allerede i klasseværelset.

Det andet par, man nødvendigvis må \dabold{skelne mellem}, er \da{sikkerhed} og \da{tryghed}. På mange fremmedsprog dækker ét og samme ord begge betydninger, men i dansk skolekultur er forskellen afgørende. \da{Sikkerhed} drejer sig om ydre og målbare tiltag --- såsom brandalarmer, adgangskontrol eller klare ordensregler i skolegården. \da{Tryghed} derimod betegner en indre oplevelse af tillid, nærvær og visheden om, at man kan række hånden op og begå fejl uden at blive mødt med hån. Uden denne dybe tryghed tør elever sjældent stille de skæve spørgsmål, hvor nysgerrighed og reel læring ellers opstår.

Endelig rummer dansk en vigtig skillelinje mellem \da{samfund} og \da{fællesskab}. \da{Samfundet} udgør den overordnede juridiske og økonomiske ramme af love, skatter og rettigheder, vi alle indgår i som borgere. Et \da{fællesskab} kan derimod ikke vedtages af Folketinget; det opstår gennem frivilligt engagement, fælles erfaringer og gensidig forpligtelse mellem konkrete mennesker --- for eksempel i en blandet folkeskoleklasse, hvor børn fra forskellige sociale lag lærer at lytte til hinanden.

Når skoledebatten spidser til, bliver faglighed og dannelse ofte \dabold{stillet op som modsætninger}. I Grundtvigs og Hal Kochs ånd \dabold{forudsætter de imidlertid hinanden}: En solid uddannelse giver eleverne redskaberne til at handle i samfundet, mens dannelse, tryghed og et levende fællesskab giver dem visdommen og modet til at bruge redskaberne klogt.

\begin{transbox}
\small
\textbf{\color{NCDTeal}[Tekst A 中文精译]}\\
几十年来，政治人物、教师与家长一直在争论：丹麦公立基础学校（Folkeskole）究竟首先应当以什么来衡量？教学应当强硬地聚焦于全国统测、PISA 国际排名和狭义的“学科专业能力”，还是应当在同等程度上塑造完整而具备批判精神的公民？要理解这场辩论的核心，就必须剖析三个几乎无法用单一外语词汇直接翻译的丹麦语关键词。

第一组也是最根本的概念对，是“职业/学业教育（\textit{uddannelse}）”与“全人教养（\textit{dannelse}）”。如果说 \textit{uddannelse} 指向可衡量的资质、考试以及面向劳动力市场的技能，那么 \textit{dannelse} 关乎的则是人的心智成熟、独立判断力以及在世界中进行伦理定位的能力。早在 19 世纪，N.F.S. 格伦特维便尖锐区分了“为了生命的学校”与“死气沉沉的书斋学校”，推崇充满生机的民众启蒙，而非机械的死记硬背。二战结束后，神学家兼思想家哈尔·科赫延续了这一传统，提出了那句著名论断：民主不仅仅是一种凭选票运作的政体，更是一种日常的对话与生活方式，孩子们必须早在教室里就开始练习它。

人们必须严格区分的第二组词，是“外在安全（\textit{sikkerhed}）”与“内在安心感（\textit{tryghed}）”。在许多外语中，这两个意思往往只对应同一个词，但在丹麦校园文化里，二者的差异具有决定性意义。\textit{Sikkerhed} 指的是外在的、技术性的且可衡量的防范措施——例如火警装置、门禁管控或操场上的明确纪律规定。而 \textit{tryghed} 则指一种内在的信任感、真诚在场与确信：即使自己举手答错问题，也不会遭到嘲笑。若缺少这种深层安心感，学生就很少敢于提出那些奇特而发散的问题——而真正的好奇心与深度学习恰恰诞生于此。

最后，丹麦语在“国家社会（\textit{samfund}）”与“价值共同体（\textit{fællesskab}）”之间也保留着重要分界线。\textit{Samfundet} 构成了我们作为公民置身其中的宏观法律与经济框架，由法律、税收与权利组成。相比之下，\textit{fællesskab} 却无法由议会立法直接通过；它产生于具体的人与人之间自愿的投入、共同的经历以及相互的道义承担——例如在一个来自不同社会阶层的混编公立学校班级里，孩子们学会彼此倾听。

当教育辩论趋于白热化时，学科能力与全人教养常常被摆成一对水火不容的对立面。然而，在格伦特维与哈尔·科赫的精神传统中，二者其实互为前提：扎实的学业教育赋予学生进入社会行动的工具，而全人教养、深层安心感与充满生机的共同体，则赋予他们明智运用这些工具的智慧与勇气。
\end{transbox}

\subsection*{Tekst B：\textit{Lærerværelset efter skolebestyrelsesmødet}（校董会后教研室对话 · 约 155 词）}

\textbf{阅读任务（理解与核对）：}\textit{Læs samtalen mellem folkeskolelærer Mette og skoleleder Thomas efter et møde med skolebestyrelsen, og besvar spørgsmålene nedenfor.}

\begin{emailbox}{Skoledebat｜Mette (folkeskolelærer) \& Thomas (skoleleder)}
\textbf{Mette:} Thomas, lagde du mærke til forældrenes uro over vores nye trivselsforløb? Flere spurgte direkte, om vi forsømmer den hårde faglighed, når klassen bruger en hel formiddag på samtaleøvelser og projekter.

\textbf{Thomas:} Ja, det er en klassisk bekymring på mange skoler. Desværre blander mange \da{uddannelse} og \da{dannelse} sammen. Gode karakterer i dansk og matematik er selvfølgelig vigtige, men uden evnen til at lytte, tænke selvstændigt og diskutere uenigheder svigter vi folkeskolens opgave i Hal Kochs forstand.

\textbf{Mette:} Præcis. En far foreslog endda strammere overvågning og faste sanktioner mod uro i timerne. Sådanne tiltag skaber måske ydre \da{sikkerhed}, men de skaber ikke i sig selv den \da{tryghed}, som gør eleverne modige i undervisningen.

\textbf{Thomas:} Helt enig. \da{Samfundet} kan udstede nok så mange læreplaner og love, men det forpligtende \da{fællesskab} må lærere og elever sammen bygge op i klasselokalet hver eneste dag.
\end{emailbox}

\textbf{Tekst B 阅读核对（Rigtigt / Forkert / Kort svar，答案见第六部分）：}
\begin{enumerate}
  \item \textit{Både Mette og Thomas mener, at karakterer i dansk og matematik slet ikke har nogen betydning.} \underline{\hspace{2.2cm}}
  \item \textit{Hvad er problemet med streng overvågning og sanktioner ifølge Mette?} \underline{\hspace{2.2cm}}
  \item \textit{Hvilken opgave kan samfundet ifølge Thomas ikke løse gennem love og læreplaner alene?} \underline{\hspace{2.2cm}}
\end{enumerate}

\begin{tipbox}{本课发音与语调要点（Udtale \& Prosodi）}
\begin{itemize}
  \item \textbf{丹麦文化核心词与人物专名读音}：\textit{dannelse} \ipa{ˈdɑnəlsə}（注意 \textit{nn} 发短元音 \ipa{ɑ}）；\textit{tryghed} \ipa{ˈtˢʁyɡˌheðˀ}（软 \textit{d} 加 \textit{stød}）；\textit{fællesskab} \ipa{ˈfɛləsˌsɡæːˀb}（重音在最前面 \textit{fæl-}）；\textit{N.F.S. Grundtvig} \ipa{ˈɡʁɔnd̥vi}（\textit{dt} 与结尾 \textit{g} 均不发音！）；\textit{Hal Koch} \ipa{hɑl ˈkʰʌɡ}。
  \item \textbf{概念辨析关键动词与名词}：\textit{at skelne} \ipa{ˈsɡɛlnə}；\textit{at forudsætte} \ipa{ˈfɒːuðˌsɛd̥ə}（重音在前缀 \textit{for-}，\textit{d} 读软 \textit{d}）；\textit{dømmekraft} \ipa{ˈdømˀəˌkʰʁɑfd}（\textit{døm-} 带 \textit{stød}）；\textit{udenadslære} \ipa{ˈuːðənɑðsˌlɛːɐ}（\textit{uden-} 与 \textit{-ad-} 均含软 \textit{d}）。
  \item \textbf{对比句式重读节奏}：朗读 \textit{Hvor \textbf{ud}dannelse retter sig mod..., handler \textbf{dan}nelse om...} 时，必须将清晰对比重音分别放在两对反义核心词的首音节上（\textit{\textbf{ud}-} vs \textit{\textbf{dan}-}、\textit{\textbf{sik}-} vs \textit{\textbf{tryg}-}、\textit{\textbf{sam}-} vs \textit{\textbf{fæl}-}）。
\end{itemize}
\end{tipbox}

\begin{transbox}
\small
\textbf{\color{NCDTeal}[Tekst B 中文对照]}\\
\textbf{梅特：}托马斯，你注意到家长们对我们新推行的班级心理福祉课程感到不安了吗？好几位家长直接发问，当全班花整整一上午做对话练习和专题项目时，我们是不是荒废了扎实的学科学习。

\textbf{托马斯：}是的，这是许多学校里都非常经典的顾虑。遗憾的是，很多人把“学业培训（\textit{uddannelse}）”和“全人教养（\textit{dannelse}）”混为一谈了。丹麦语和数学拿到好成绩当然很重要，但如果缺少倾听、独立思考以及讨论分歧的能力，我们就辜负了哈尔·科赫意义上的公立学校使命。

\textbf{梅特：}说得太对了。有位父亲甚至建议在课堂上实行更严格的监控和固定惩戒，好让教室保持安静。这类措施或许能制造出外在的“安全秩序（\textit{sikkerhed}）”，但它们本身并不会自动生发出让学生在课堂上勇敢提问的“内在安心感（\textit{tryghed}）”。

\textbf{托马斯：}完全赞同。“国家社会（\textit{samfundet}）”可以颁布再多的教学大纲和法律条文，但那种彼此承担道义的“共同体（\textit{fællesskab}）”，必须靠老师和学生每一天在教室里共同搭建起来。
\end{transbox}

% ---------------------------------------------------------------------
\ncdsubsec{三、Gloser（核心生词表）}

\begin{tipbox}{词表重点}
本课三大文化词对的“可数性”极易出错：\textit{en dannelse}（全人教养）、\textit{en faglighed}（学科专业性）、\textit{en tryghed}（内在安心感）以及 \textit{en sikkerhed}（外在安全秩序，一般语义下）均属不可数共性名词（\textit{kun ental}）；而 \textit{en uddannelse}（学历/教育项目）可数，复数为 \textit{uddannelser}。另外注意不规则动词：\textit{at forudsætte}（以……为前提）过去式与完成时分别作 \textit{forudsatte, har forudsat}；\textit{at opstå}（出现/产生）完成时用 \dabold{er}。
\end{tipbox}

\begin{glosertable}
名词 & en dannelse & \textit{(dannelsen, kun ental)} & 全人教养，人文陶冶 \\
名词 & en uddannelse & \textit{(uddannelsen, uddannelser, uddannelserne)} & 学业教育，职业培训，学历 \\
名词 & en faglighed & \textit{(fagligheden, kun ental)} & 学科能力，专业水准 \\
名词 & en tryghed & \textit{(trygheden, kun ental)} & 内在安心感，心理安全与信任 \\
名词 & en sikkerhed & \textit{(sikkerheden, kun ental)} & 外在安全，治安防卫保障 \\
名词 & et fællesskab & \textit{(fællesskabet, fællesskaber, fællesskaberne)} & 共同体，凝聚的群体关系 \\
名词 & et samfund & \textit{(samfundet, samfund, samfundene)} & 社会，国家制度共同体 \\
名词 & en dømmekraft & \textit{(dømmekraften, kun ental)} & 判断力，辨别力 \\
名词 & en livsform & \textit{(livsformen, livsformer, livsformerne)} & 生活方式，生命形态 \\
名词 & en udenadslære & \textit{(udenadslæren, kun ental)} & 死记硬背 \\
名词 & en foranstaltning & \textit{(foranstaltningen, -er, foranstaltningerne)} & 措施，防范举措 \\
名词 & et nærvær & \textit{(nærværet, kun ental)} & 真诚在场，陪伴感 \\
\midrule
动词 [har] & at skelne (mellem) & \textit{(skelner, skelnede, har skelnet)} & 区分，辨别 \\
动词 [har] & at forudsætte & \textit{(forudsætter, forudsatte, har forudsat)} & 以……为前提，预设 \\
动词 [har] & at svigte & \textit{(svigter, svigtede, har svigtet)} & 辜负，失职，背弃 \\
动词 [har] & at forsømme & \textit{(forsømmer, forsømte, har forsømt)} & 荒废，疏忽 \\
动词 [er] & at opstå & \textit{(opstår, opstod, er opstået)} & 产生，出现，萌发 \\
\midrule
形容词 & forpligtende & \textit{(forpligtende, 不变)} & 彼此承担责任的，有约束力的 \\
\end{glosertable}

% ---------------------------------------------------------------------
\ncdsubsec{四、Grammatik \& Stilistik（不可译文化词辨析、对比句法与 C1 避坑）}

\begin{grammatikbox}{语法点 1：丹麦文化三大不可译核心词深度辨析矩阵}
\vspace{0.2em}
\renewcommand{\arraystretch}{1.24}
{\small
\begin{tabularx}{\linewidth}{@{} >{\bfseries\raggedright\arraybackslash}p{2.9cm} >{\raggedright\arraybackslash}p{4.6cm} >{\raggedright\arraybackslash}X @{}}
\toprule
\textbf{概念对照对} & \textbf{丹麦文化与制度内核差异} & \textbf{地道语境辨析例句} \\
\midrule
\da{dannelse}\newline vs \da{uddannelse} &
\textit{uddannelse} = 可测量的学历、职业技能与资格；\newline \textit{dannelse} = 人格成熟、判断力、伦理与公民修养（Grundtvig 传统）。 &
\textit{\da{Uddannelse} giver et eksamensbevis, mens \da{dannelse} lærer os at bruge vores viden etisk.} \\
\midrule
\da{tryghed}\newline vs \da{sikkerhed} &
\textit{sikkerhed} = 门锁、监控、警报、制度规则等客观防护；\newline \textit{tryghed} = 信任、接纳、不怕犯错的心理安全感。 &
\textit{Kameraer øger måske den ydre \da{sikkerhed}, men skaber ikke automatisk indre \da{tryghed}.} \\
\midrule
\da{fællesskab}\newline vs \da{samfund} &
\textit{samfund} = 法律、税收与公共福利构成的宏观制度框架；\newline \textit{fællesskab} = 基于参与和互助的情感与责任共同体。 &
\textit{\da{Samfundet} vedtager love, men et forpligtende \da{fællesskab} bygges nedefra i hverdagen.} \\
\bottomrule
\end{tabularx}
}
\end{grammatikbox}

\begin{grammatikbox}{语法点 2：概念界定与对比修辞四大 C1 句型}
\begin{enumerate}
  \item \textbf{对比界定主句铁律：\dabold{Hvor A retter sig mod X, handler B om Y}}：\newline
    注意：\textit{Hvor} 在此处不是空间副词“在哪里”，而是高阶对比连词“如果说 A 指向 X，那么 B 关乎 Y”！前面整个 \textit{Hvor}-从句内部不倒装，但占据主句第 1 位，因此逗号后面主句\textbf{必须 V2 倒装}：\newline
    \ding{51}\ \textit{\da{Hvor} uddannelse retter sig mod eksamener, \dabold{handler dannelse} om dømmekraft.}
  \item \textbf{精准区分：\dabold{at skelne mellem A og B}}：\newline
    后面固定接介词 \dabold{mellem}（不可用 \textit{fra}），对应英文 \textit{distinguish between}：\newline
    \ding{51}\ \textit{Man må nødvendigvis \da{skelne mellem} ydre sikkerhed og indre tryghed.}
  \item \textbf{破除二元对立：\dabold{stille A og B op som modsætninger} vs \dabold{forudsætte hinanden}}：\newline
    先指出舆论的误区（把两者对立起来），再用“互为前提”完成辩证升华：\newline
    \ding{51}\ \textit{Ofte \da{stilles} faglighed og dannelse \da{op som modsætninger}, selvom de i virkeligheden \da{forudsætter hinanden}.}
  \item \textbf{偏好与取舍：\dabold{X frem for Y}}（重 X 轻 Y / 宁要 X 而非 Y）：\newline
    \ding{51}\ \textit{Grundtvig fremhævede levende oplysning \da{frem for} mekanisk udenadslære.}
\end{enumerate}
\end{grammatikbox}

\begin{grammatikbox}{语法点 3：文化概念误用与语序排雷清单}
\begin{enumerate}
  \item \textbf{\ding{172} 切勿用 \textit{sikkerhed} 翻译“班级心理安全感”}：中文的“安全感”若指孩子敢于发言、不怕被霸凌，丹麦语必须用 \da{tryghed}！说 \ding{55}\ \textit{eleverne mangler sikkerhed i klassen} 会让人以为教室门锁坏了或有消防隐患。
  \item \textbf{\ding{173} \textit{Hvor}-对比从句后的 V2 倒装不可漏}：常见错误是受英语 \textit{Whereas A focuses on X, B deals with Y} 影响忘了倒装：\ding{55}\ \textit{Hvor samfundet handler om love, fællesskabet handler om tillid.} $\rightarrow$ \ding{51}\ \textit{Hvor samfundet handler om love, \da{handler fællesskabet} om tillid.}
  \item \textbf{\ding{174} 区分 \textit{at blande A og B sammen} 与 \textit{at skelne mellem A og B}}：“混淆 A 和 B”用可分动词短语 \textit{blande A og B sammen}；“区分 A 和 B”用 \textit{skelne mellem A og B}，切勿写成 \ding{55}\ \textit{skelne A fra B} 或 \ding{55}\ \textit{blande A med B}。
  \item \textbf{\ding{175} 不规则过去分词 \textit{har forudsat}}：\textit{forudsætte} 变化跟 \textit{sætte} 一样（\textit{forudsatte, har forudsat}），注意末尾不加 \textit{-et}，不可误写为 \ding{55}\ \textit{har forudsattet}。
\end{enumerate}
\end{grammatikbox}

% ---------------------------------------------------------------------
\ncdsubsec{五、Mønsterøvelser（句型与修辞操练）}

\textbf{练习 A：丹麦文化核心词与概念对比句式精准填空（每空填入最准确的 1 个丹麦语词）}
\begin{enumerate}
  \item Hvor \underline{\hspace{2.2cm}} giver eleverne konkrete kvalifikationer, former dannelse deres menneskelige dømmekraft.
  \item I debatten om folkeskolen må vi skelne \underline{\hspace{2.2cm}} ydre kontrol og ægte tillid i klassen.
  \item Låste døre og faste ordensregler giver måske teknisk \underline{\hspace{2.2cm}}, men skaber ikke automatisk varme relationer.
  \item Børn tør først række hånden op og begå fejl, når der hersker reel \underline{\hspace{2.2cm}} i klasseværelset. (内在安心感)
  \item Et forpligtende \underline{\hspace{2.2cm}} kan ikke vedtages ved lov, men må gro mellem konkrete mennesker. (共同体)
  \item Grundtvig fremhævede levende samtale frem \underline{\hspace{2.2cm}} mekanisk udenadslære.
  \item Faglighed og dannelse er ikke fjender; i virkeligheden \underline{\hspace{2.2cm}} de hinanden. (forudsætte，现在时)
  \item Hvis skolen kun måler testresultater, \underline{\hspace{2.2cm}} den sin demokratiske opgave i Hal Kochs ånd. (svigte，现在时)
\end{enumerate}

\textbf{练习 B：文化概念辨析纠错与 C1 句法升级（修正误用/语序，或按提示合并）}
\begin{enumerate}
  \item 【文化概念误用纠错】\textit{Når eleverne ikke bliver latterliggjort for fejl, oplever de stor sikkerhed i undervisningen.} $\rightarrow$
  \item 【对比从句 V2 倒装纠错】\textit{Hvor samfundet udgør den juridiske ramme, et fællesskab bygger på frivillig tillid.} $\rightarrow$
  \item 【动词搭配与介词纠错】\textit{Vi bør ikke skelne faglighed fra dannelse, og heller ikke blande dem med hinanden.} $\rightarrow$
  \item 【不规则过去完成时纠错】\textit{Hal Kochs demokratiopfattelse har altid forudsattet en aktiv samtale mellem borgerne.} $\rightarrow$
  \item 【对比句型合并升级】用 \textit{Hvor ... retter sig mod ..., handler ... om ...} 合并：\textit{Uddannelse retter sig mod arbejdsmarkedet. Dannelse handler om etisk dømmekraft.} $\rightarrow$
  \item 【破除二元对立升级】用 \textit{stille ... op som modsætninger} 与 \textit{forudsætte hinanden} 改写：\textit{Politikere siger tit, at test og dannelse ikke passer sammen, men de har brug for hinanden.} $\rightarrow$
\end{enumerate}

\textbf{练习 C：中译丹实战（精准驾驭三大文化词对与对比句型）}
\begin{enumerate}
  \item 在关于公立学校的讨论中，我们必须严格区分学业培训与全人教养。
  \item 如果说学业教育指向具体技能，全人教养则关乎人的道德判断力。
  \item 外在的安全规则本身并不会自动在课堂上创造内在的安心感。
  \item 国家社会可以通过法律，但有凝聚力的共同体却产生于日常相处之中。
  \item 学科能力与全人教养不应被摆成对立面，因为二者其实互为前提。
  \item 在哈尔·科赫的精神中，民主首先是一种必须从小练习的日常对话。
\end{enumerate}

% ---------------------------------------------------------------------
\ncdsubsec{六、PD3 / Studieprøven 实战微任务与参考答案}

\begin{emailbox}{PD3 / Studieprøven 写作微任务（Kultur- og debatkommentar · 200--250+ ord）}
\textbf{Situation}: En politisk tænketank foreslår, at den danske folkeskole fremover skal reducere antallet af kreative fag og klassens time for i stedet at afholde flere nationale test i læsning og matematik samt indføre strengere disciplinære regler.\\[3pt]
\textbf{Opgave}: Skriv en nuanceret debatartikel på \textbf{200--250 ord}, hvor du vurderer forslaget. Brug lektionens tre danske kernebegrebspar (\textit{uddannelse / dannelse, sikkerhed / tryghed, samfund / fællesskab}) aktivt i din argumentation og inddrag kort enten Grundtvigs eller Hal Kochs syn på skolens demokratiske rolle.
\end{emailbox}

\begin{answerbox}
\textbf{Før du læser 参考回答}
\begin{enumerate}
  \item \textit{Tekst A skelner mellem uddannelse (målbare kvalifikationer) og dannelse (modning og dømmekraft), mellem ydre sikkerhed og indre tryghed samt mellem samfundets juridiske ramme og menneskelige fællesskaber. Grundtvig talte om ”skolen for livet”, og Hal Koch så demokrati som en daglig samtale, der skal øves i skolen.}
  \item \textit{De påpeger, at høje karakterer og faglighed ikke kan stå alene uden demokratisk samtaleevne, og at streng overvågning højst skaber ydre sikkerhed, mens reel tryghed og et forpligtende fællesskab må opbygges gennem daglig tillid mellem lærere og elever.}
\end{enumerate}

\textbf{Tekst B 阅读核对答案}
\begin{enumerate}
  \item \da{Forkert}（托马斯明确承认丹麦语和数学的成绩当然也很重要：\textit{Gode karakterer ... er selvfølgelig vigtige}）
  \item \da{Det skaber måske ydre sikkerhed, men ikke i sig selv den tryghed, der gør elever modige.}（严格监控顶多带来外在秩序，却不会自动产生让学生敢于表达的心理安心感）
  \item \da{At bygge det forpligtende fællesskab op i klasselokalet.}（在教室里建构彼此担责的共同体）
\end{enumerate}

\textbf{练习 A 答案}
\begin{enumerate}
  \item \da{uddannelse}（与后文 \textit{dannelse} 构成“学历技能 vs 全人教养”对比）
  \item \da{mellem}（固定介词搭配：\textit{skelne mellem A og B}）
  \item \da{sikkerhed}（对应门禁、规章等客观外在安全措施）
  \item \da{tryghed}（指不怕犯错、彼此信任的内在心理安全感）
  \item \da{fællesskab}（指无法靠立法强制生成、在人与人之间自然生长的共同体）
  \item \da{for}（固定对比取舍短语：\textit{frem for} ……而非……）
  \item \da{forudsætter}（动词 \textit{at forudsætte} 现在时：互为前提）
  \item \da{svigter}（动词 \textit{at svigte} 现在时：辜负、失职）
\end{enumerate}

\textbf{练习 B 答案（参考）}
\begin{enumerate}
  \item \textit{Når eleverne ikke bliver latterliggjort for fejl, oplever de stor \da{tryghed} i undervisningen.}（心理层面安心不怕出错必须用 \textit{tryghed}，不可误用 \textit{sikkerhed}）
  \item \textit{Hvor samfundet udgør den juridiske ramme, \da{bygger et fællesskab} på frivillig tillid.}（\textit{Hvor}-对比从句置于句首时，主句变位动词 \textit{bygger} 须触发 V2 倒装）
  \item \textit{Vi bør \da{skelne mellem} faglighed og dannelse, men heller ikke \da{blande} dem \da{sammen}.}（“区分”固定搭配 \textit{skelne mellem A og B}；“混淆”用 \textit{blande A og B sammen}）
  \item \textit{Hal Kochs demokratiopfattelse har altid \da{forudsat} en aktiv samtale mellem borgerne.}（不规则过去分词作 \textit{forudsat}，不加 \textit{-tet}）
  \item \textit{\da{Hvor} uddannelse retter sig mod arbejdsmarkedet, \da{handler} dannelse om etisk dømmekraft.}（借助 \textit{Hvor ..., handler ...} 构建严密概念界定句并正确倒装）
  \item \textit{Politikere \da{stiller} ofte test og dannelse \da{op som modsætninger}, selvom de faktisk \da{forudsætter hinanden}.}（以固定短语替代口语化表述，突出“互为前提”的辩证力）
\end{enumerate}

\textbf{练习 C 答案}
\begin{enumerate}
  \item \textit{I debatten om folkeskolen må vi skelne klart mellem uddannelse og dannelse.}
  \item \textit{Hvor uddannelse retter sig mod konkrete færdigheder, handler dannelse om etisk dømmekraft.}
  \item \textit{Ydre regler om sikkerhed skaber ikke i sig selv indre tryghed i klassen.}
  \item \textit{Samfundet vedtager love, men et stærkt fællesskab opstår gennem dagligt samvær.}
  \item \textit{Faglighed og dannelse bør ikke stilles op som modsætninger, da de forudsætter hinanden.}
  \item \textit{I Hal Kochs ånd er demokrati først og fremmest en samtale, som skal øves.}
\end{enumerate}

\textbf{PD3 / Studieprøven 高分范文（约 235 词，不含标题）}\\
\textit{\textbf{En stærk folkeskole kræver både faglighed og levende dannelse}}

\textit{Forslaget om flere nationale test og strengere disciplin på bekostning af kreative fag og klassens time rammer skævt, fordi det bygger på en falsk modsætning. Ganske vist er gode læse- og regnefærdigheder uundværlige for alle unge, men en folkeskole reduceret til ren testtræning overser selve kernen i den danske skoletradition.}

\textit{Her er det helt afgørende at skelne klart mellem \da{uddannelse} og \da{dannelse}. \dabold{Hvor uddannelse retter sig mod} målbare kundskaber og fremtidig beskæftigelse, \dabold{handler dannelse om} dømmekraft, nysgerrighed og evnen til at indgå konstruktivt i demokratisk uenighed. Allerede i 1945 mindede Hal Koch os om, at demokrati først og fremmest er en livsform og en daglig samtale, som børn må øve sig i gennem undervisningen --- ikke en mekanisk eksamensdisciplin, man kan lære ved udenadslære.}

\textit{Samtidig forveksler forslaget ydre \da{sikkerhed} med indre \da{tryghed}. Strenge disciplinære regler og øget kontrol kan måske dæmpe uroen kortvarigt, men reel læring og mod opstår kun, når eleverne føler sig så trygge i klassens \da{fællesskab}, at de tør begå fejl og stille skæve spørgsmål uden frygt for hån. \da{Samfundet} kan udstede nok så mange test, læreplaner og sanktioner, men det kan aldrig vedtage engagement og gensidig tillid ved lov.}

\textit{Derfor bør faglighed og dannelse ikke stilles op som modsætninger i folkeskolen. De forudsætter tværtimod hinanden: Kun når børn trives i trygge, forpligtende fællesskaber med plads til samtale og kreativ udfoldelse, får de også det mentale overskud til at løfte deres faglige uddannelse.}

\small\textbf{名师点评}：正文（不含标题）约 235 词，文化纵深感极强！文章开篇否定智库提案的虚假二元对立，随后自然流畅地展开三组丹麦不可译文化词的精确辨析：\textit{uddannelse} vs \textit{dannelse}（援引哈尔·科赫 1945 年民主对话观）、\textit{sikkerhed} vs \textit{tryghed}（揭示纪律控制与班级安全感之别）、\textit{samfund} vs \textit{fællesskab}（点明立法与共同体的边界），最后收束于“互为前提”，是 Studieprøven 教育文化题的满分样板！
\end{answerbox}
\ncdlesson{12}{Bioetik og fremtidens medicin: Hvor langt bør vi gå?}{生命伦理与未来医学：我们应当走多远？}
{DU3 Modul 5 (PD3) $\rightarrow$ Modul 6 (Studieprøven, B2 $\rightarrow$ C1)｜对标 NCE 4 \& 科技伦理与生命医学深度思辨}
{\textbf{伦理两难推演与警示句型}（\textit{rejse fundamentale etiske dilemmaer} · \textit{skabe præcedens for} · \textit{uoverskuelige konsekvenser} · \textit{en glidebane} · \textit{hvor sætter vi grænsen?}）· \textit{Jo..., desto...} 递进警示句式与医学伦理核心词汇}
{欢迎来到 Unit 3《伦理、制度与现代性思辨》的收官课！当 CRISPR 基因编辑能够抹去致命遗传病、AI 诊断系统能够在症状显现前数年预测疾病风险时，人类正站在“治愈疾病”与“人为设计生命”的临界点上。在 PD3 与 Studieprøven 面对生物科技或人工智能伦理题时，绝不能停留在笼统的“有好有坏”，而应熟练调用\textbf{\textit{rejse fundamentale etiske dilemmaer}}、\textbf{\textit{skabe præcedens for}}、\textbf{\textit{en glidebane}} 与 \textbf{\textit{uoverskuelige konsekvenser}} 进行层层推演，最终发出极具警醒力的追问：\textbf{\textit{Hvor sætter vi grænsen?}}}

% ---------------------------------------------------------------------
\begin{foerduleaeser}
带着下面两个问题阅读 Tekst A 和 Tekst B：
\begin{enumerate}
  \item \textit{Hvilke etiske dilemmaer opstår ifølge Tekst A, når genredigering bevæger sig fra sygdomsbehandling til forbedring af menneskelige egenskaber, og hvordan kan AI-diagnostik true patientens selvbestemmelse?}（根据 Tekst A，当基因编辑从治疗疾病滑向“改良人类性状”时会引发哪些伦理困境？AI 诊断又如何威胁患者的自主决定权？）
  \item \textit{Hvorfor advarer overlæge Sofie Nielsen i Tekst B mod at lade frygten for en glidebane føre til et totalt forbud mod ny medicinsk teknologi?}（在 Tekst B 中，主任医师索菲·尼尔森为什么警告不要让对“滑坡效应”的恐惧导致对新医疗技术的全面禁止？）
\end{enumerate}
\end{foerduleaeser}

% ---------------------------------------------------------------------
\ncdsubsec{二、Tekster（生命伦理深度专栏与临床学者回应双语篇）}

\subsection*{Tekst A：\textit{Fra helbredelse til optimering --- bioetikkens grænser}（生命伦理两难推演长文 · 约 365 词）}

\begin{tipbox}{语篇说明}
本篇以 CRISPR 基因编辑与 AI 预测诊断为切入口，展开严密的伦理风险推演。阅读时请重点留意\textbf{\color{NCDNavy}深蓝色标注的医疗伦理核心概念与法理框架}，以及\textbf{\color{NCDRed}深红色标注的伦理两难、先例滑坡与终极警示句型}。
\end{tipbox}

Fremtidens medicin rummer løfter, der for få årtier siden lignede ren science fiction. Med præcis genredigering som CRISPR kan forskere i dag gribe direkte ind i menneskets DNA for at fjerne alvorlige arvelige sygdomme, mens kunstig intelligens i diagnostikken kan afsløre skjulte mønstre i scanninger og blodprøver længe før de første symptomer. Men jo hurtigere teknologien udvikler sig, desto skarpere \dabold{rejser den fundamentale etiske dilemmaer} for hele samfundet.

Så længe nye bioteknologiske metoder udelukkende anvendes \da{i behandlingens tjeneste} til at helbrede livstruende lidelser, nyder forskningen bred moralsk opbakning. Problemet opstår i det øjeblik, skellet mellem sygdomsforebyggelse og ikke-medicinsk forbedring begynder at skride. Hvad udgør egentlig en ”defekt”, som bør rettes, og hvem skal definere det normale menneske? Hvis vi først tillader små genetiske justeringer for at ”optimere” sundhed, intelligens eller udseende, kan selv velmente indgreb \dabold{skabe præcedens for}, at menneskelige egenskaber forvandles til en handelsvare.

Bioetikere og medlemmer af Det Etiske Råd advarer derfor om, at samfundet let kan bevæge sig ud ad \dabold{en farlig glidebane}. I modsætning til traditionel medicin begrænser ændringer i kønsceller eller tidlige fostre sig nemlig ikke til én patient; de går i arv til kommende generationer og kan dermed få \dabold{uoverskuelige konsekvenser} for både den biologiske mangfoldighed og den sociale lighed. I værste fald risikerer vi et nyt genetisk klasseskel, hvor økonomisk privilegerede familier køber sig adgang til bedre biologiske forudsætninger.

Samme spænding mellem nytte og risiko viser sig i brugen af kunstig intelligens i sundhedsvæsenet. En klog og trænet algoritme kan redde liv gennem tidlig opsporing af kræft eller demens, men forudsiger samtidig statistiske sygdomsrisici, som patienten måske slet ikke ønsker at kende. Hvad sker der med \da{patientens selvbestemmelse} og retten til ikke at vide? Og hvordan forhindrer vi, at forsikringsselskaber eller arbejdsgivere får adgang til følsomme genetiske og digitale risikoprofiler, så \da{solidariteten} i det offentlige sundhedsvæsen smuldrer?

Teknologiske muligheder har en tendens til at blive til samfundsmæssige forventninger, medmindre vi griber ind i tide. Derfor må lovgivere, forskere og borgere stille det uafviselige spørgsmål: \dabold{Hvor sætter vi grænsen?} Kun gennem bindende internationale regler, gennemsigtig demokratisk debat og modet til at sige stop i tide kan vi sikre, at medicinen fortsat tjener menneskets værdighed frem for blind teknologisk optimisme.

\begin{transbox}
\small
\textbf{\color{NCDTeal}[Tekst A 中文精译]}\\
未来医学蕴含着几十年前还宛如纯粹科幻小说的承诺。借助 CRISPR 等精准基因编辑技术，研究者如今能够直接介入人类 DNA 以消除严重的遗传性疾病；与此同时，诊断领域的人工智能也能早在首批症状出现之前，就从影像扫描和血液样本中识别出隐藏的病变模式。然而，技术发展得越快，它为整个社会所引发的根本性伦理困境也就越加尖锐。

只要新的生物技术方法完全用于服务治病救人、治愈危及生命的疾苦，科学研究便能获得广泛的道德支持。问题恰恰产生于“疾病预防”与“非医疗目的的人体改良”之间的界线开始松动滑移的那一刻。究竟什么才算一种应当被修正的“缺陷”？又该由谁来定义“正常的人”？如果我们一旦允许为了“优化”健康水平、智力或外貌而进行微小基因调整，那么哪怕是出于善意的干预，也可能为“人类性状沦为商品”开创危险先例。

因此，生命伦理学家及丹麦伦理委员会成员警告称，社会极易走上一条危险的滑坡。与传统药物不同，对生殖细胞或早期胚胎所做的改动并不局限于某一位患者；它们会遗传给后代，从而对生物多样性与社会平等造成难以预料的后果。在最坏的情况下，我们甚至面临催生新型“基因阶层分化”的风险——经济优越的家庭得以花钱为子女购买更优的先天生理条件。

益处与风险之间的同样张力，也体现在医疗系统对人工智能的运用中。基于算法的诊断能够通过早期筛查癌症或痴呆症挽救生命，但与此同时，它也会预测出患者或许根本不愿知晓的概率性患病风险。届时，患者的自主决定权与“不知情权”将何去何从？我们又该如何防止保险公司或雇主获取敏感的基因与数字风险档案，从而避免公共医疗体系内的团结互助精神土崩瓦解？

除非我们及时出手干预，否则技术上的可能性往往有演变为社会普遍期待的趋势。正因如此，立法者、学者和公民必须提出那个无可回避的问题：我们究竟把边界划在哪里？唯有依靠具有约束力的国际法规、透明的民主辩论以及在关键时刻敢于叫停的勇气，我们才能确保医学继续服务于人的尊严，而非盲目的技术乐观主义。
\end{transbox}

\subsection*{Tekst B：\textit{Kronik-replik: ”Passivitet har også en etisk pris”}（临床学者回应评论 · 约 155 词）}

\textbf{阅读任务（理解与核对）：}\textit{Læs overlæge og genforsker Sofie Nielsens modsvar til debatten om bioetik, og besvar spørgsmålene nedenfor.}

\begin{emailbox}{Kronik-replik｜Fra: Sofie Nielsen (overlæge og genforsker) \quad Emne: Etik handler også om at handle i tide}
Advokater og bioetikere gør helt ret i at advare mod \dabold{en glidebane}, hvor menneskets gener bliver et kommercielt marked. Ingen seriøse forskere i Danmark ønsker designerbørn eller et samfund uden solidaritet og lige adgang.

Alligevel savner jeg ofte én afgørende vinkel i den offentlige debat: Også tøven og forbud har dybe menneskelige omkostninger. Når lovgivningen af ren frygt for \dabold{uoverskuelige konsekvenser} bremser sikre, målrettede genterapier eller lovende AI-screening, svigter vi de patienter og familier, der i dag lever med smertefulde og uhelbredelige sygdomme. At nægte at hjælpe, selvom de medicinske redskaber faktisk findes, er også et tungt moralsk valg.

Derfor bør vi ikke besvare spørgsmålet ”\dabold{hvor sætter vi grænsen?}” med et blankt nej til al medicinsk udvikling. I stedet må vi trække skarpe juridiske rammer, der klart adskiller behandling fra forbedring, kræve fuldt informeret \da{samtykke} og sikre lige adgang gennem det offentlige sundhedsvæsen.
\end{emailbox}

\textbf{Tekst B 阅读核对（Rigtigt / Forkert / Kort svar，答案见第六部分）：}
\begin{enumerate}
  \item \textit{Sofie Nielsen mener, at danske forskere bør have fri adgang til at skabe designerbørn på et kommercielt marked.} \underline{\hspace{2.2cm}}
  \item \textit{Hvilken overset etisk omkostning peger Sofie på, hvis lovgivningen bremser sikre genterapier og AI-screening?} \underline{\hspace{2.2cm}}
  \item \textit{Hvilke tre løsninger foreslår Sofie i stedet for et blankt nej til teknologisk udvikling?} \underline{\hspace{2.2cm}}
\end{enumerate}

\begin{tipbox}{本课发音与语调要点（Udtale \& Prosodi）}
\begin{itemize}
  \item \textbf{生命伦理与法律高阶借词重音}：\textit{bioetik} \ipa{biːoeˈtˢiɡ}（重音在最后音节 \textit{-tik}）；\textit{et dilemma} \ipa{diˈlɛmɑ}；\textit{en præcedens} \ipa{pʰʁɛseˈdɛnˀs}（法/拉丁借词！\textit{c} 读作 \ipa{s}，重音在最后且带 \textit{stød}）；\textit{en konsekvens} \ipa{kʰʌnsəˈkʰvɛnˀs}（重音在尾音节，带 \textit{stød}）。
  \item \textbf{复合并列与滑音细节}：\textit{en glidebane} \ipa{ˈɡliːðəˌbæːnə}（软 \textit{d}）；\textit{uoverskuelig} \ipa{ˈuːʌʊ̯ɐˌsɡuːəli}（否定前缀 \textit{u-} 常重读！\textit{over-} 读 \ipa{ʌʊ̯ɐ}，词尾 \textit{g} 不发音）；\textit{et samtykke} \ipa{ˈsɑmˌtˢyɡə}；\textit{arvemasse} \ipa{ˈɑːvəˌmɑsə}。
  \item \textbf{递进句与伦理设问语调}：朗读 \textit{Jo hurtigere..., desto skarpere...} 时，两个比较级副词（\textit{hurtigere}, \textit{skarpere}）须读出陡升强调；最后一段的 \textit{Hvor sætter vi grænsen?} 用冷静而坚定的降调收束在 \textit{\textbf{græn}sen} 上，极具警示权威感。
\end{itemize}
\end{tipbox}

\begin{transbox}
\small
\textbf{\color{NCDTeal}[Tekst B 中文对照]}\\
法学家与生命伦理学家完全有理由对人类基因沦为商业市场的“滑坡效应”提出警告。在丹麦，没有任何严肃的研究者希望出现“定制婴儿”或丧失团结互助精神与平等可及性的社会。

然而，我常常感到公共讨论中缺少一个关键视角：犹豫不决与一味禁止同样伴随着沉重的人道代价。当立法因单纯害怕难以预料的后果而叫停安全、靶向的基因疗法或前景广阔的 AI 筛查时，我们便辜负了今天正在忍受痛苦与不治之症折磨的患者和家庭。明明手中已经握有有效的医疗工具却选择拒绝施救，这本身也是一种沉重的道德抉择。

因此，我们不应以对一切医学进步一刀切说“不”，来回答“我们把边界划在哪里”这一问题。相反，我们应当划定严格的法律框架将“治病”与“增强”清晰区分开来，坚持充分的知情同意，并通过公共医疗体系保障所有人的平等可及性。
\end{transbox}

% ---------------------------------------------------------------------
\ncdsubsec{三、Gloser（核心生词表）}

\begin{tipbox}{词表重点}
注意两处极易写错的双写辅音与词形：\textit{en bioetik}（生命伦理学）和 \textit{en diagnostik}（诊断学）定指形式均须\textbf{双写 \textit{k}}（\textit{bioetikken}, \textit{diagnostikken}）；\textit{et klasseskel}（阶级鸿沟）定指与复数定指均须\textbf{双写 \textit{l}}（\textit{klasseskellet, klasseskellene}）。强变化动词 \textit{at skride}（滑动、动摇）和 \textit{at gribe ind}（出手干预）均为核心不规则变化，且 \textit{skride} 完成时使用助动词 \dabold{er}！
\end{tipbox}

\begin{glosertable}
名词 & en bioetik & \textit{(bioetikken, kun ental)} & 生命伦理学 \\
名词 & et dilemma & \textit{(dilemmaet, dilemmaer, dilemmaerne)} & 两难困境，进退维谷 \\
名词 & en præcedens & \textit{(præcedensen, kun ental)} & 先例，判例先导 \\
名词 & en glidebane & \textit{(glidebanen, glidebaner, glidebanerne)} & 滑坡（效应），危险的逐步滑落 \\
名词 & en konsekvens & \textit{(konsekvensen, konsekvenser, -erne)} & 后果，影响 \\
名词 & en arvemasse & \textit{(arvemassen, arvemasser, arvemasserne)} & 基因组，遗传物质 \\
名词 & en genredigering & \textit{(genredigeringen, -er, -erne)} & 基因编辑 \\
名词 & en diagnostik & \textit{(diagnostikken, kun ental)} & 诊断技术，诊断学 \\
名词 & en selvbestemmelse & \textit{(selvbestemmelsen, kun ental)} & 自主决定权 \\
名词 & et klasseskel & \textit{(klasseskellet, klasseskel, klasseskellene)} & 阶层鸿沟，阶级分化 \\
名词 & en solidaritet & \textit{(solidariteten, kun ental)} & 团结互助精神 \\
名词 & et samtykke & \textit{(samtykket, samtykker, samtykkerne)} & 同意，知情同意书 \\
\midrule
动词 [har] & at helbrede & \textit{(helbreder, helbredte, har helbredt)} & 治愈，治好 \\
动词 [er] & at skride & \textit{(skrider, skred, er skredet)} & 滑动，松动，偏移 \\
动词 [har] & at gribe (ind) & \textit{(griber, greb, har grebet ind)} & 干预，出手制止 \\
动词 [har] & at bremse & \textit{(bremser, bremsede, har bremset)} & 阻滞，煞停，遏制 \\
\midrule
形容词 & uoverskuelig & \textit{(uoverskueligt, uoverskuelige)} & 不可预见的，难以估量的 \\
形容词 & bindende & \textit{(bindende, 不变)} & 具有法律约束力的 \\
\end{glosertable}

% ---------------------------------------------------------------------
\ncdsubsec{四、Grammatik \& Stilistik（伦理推演警示句型、递进句式与 C1 避坑）}

\begin{grammatikbox}{语法点 1：生命伦理“两难—先例—滑坡—底线”四步推演模板}
在 PD3 与 Studieprøven 处理基因编辑、人工智能、安乐死或数据隐私等前沿伦理议题时，用下面四组核心警示搭配串联论证，极具学者风范：

\vspace{0.3em}
\renewcommand{\arraystretch}{1.24}
{\small
\begin{tabularx}{\linewidth}{@{} >{\bfseries\raggedright\arraybackslash}p{3.6cm} >{\raggedright\arraybackslash}p{4.1cm} >{\raggedright\arraybackslash}X @{}}
\toprule
\textbf{推演环节与核心搭配} & \textbf{语法接口与语义功能} & \textbf{Tekst A/B 经典示范} \\
\midrule
\textbf{1. 引发伦理困境}\newline \da{rejse fundamentale etiske dilemmaer} &
必须用动词 \textit{at rejse}（提出/引发），不可用 \textit{lave}；常配复数 \textit{etiske dilemmaer}。 &
\textit{Teknologien \da{rejser fundamentale etiske dilemmaer} for hele samfundet.} \\
\midrule
\textbf{2. 开创危险先例}\newline \da{skabe præcedens for} (+ 名词 / \textit{at}-从句) &
\textit{præcedens} 零冠词单数，固定接介词 \dabold{for}，推演“一旦开此先例将如何”。 &
\textit{Selv velmente indgreb kan \da{skabe præcedens for}, at gener bliver en vare.} \\
\midrule
\textbf{3. 滑坡与失控后果}\newline \da{en glidebane} \& \da{uoverskuelige konsekvenser} &
常用动词 \textit{bevæge sig ud ad en glidebane}；后果配 \textit{få uoverskuelige konsekvenser for...}。 &
\textit{Samfundet kan bevæge sig ud ad \da{en farlig glidebane} med \da{uoverskuelige konsekvenser}.} \\
\midrule
\textbf{4. 发出底线追问}\newline \da{Hvor sætter vi grænsen?} &
动词用 \textit{sætte}（划定边界），常以独立短问句置于结尾段首，力透纸背。 &
\textit{Lovgivere og borgere må stille spørgsmålet: \da{Hvor sætter vi grænsen?}} \\
\bottomrule
\end{tabularx}
}
\end{grammatikbox}

\begin{grammatikbox}{语法点 2：“越……越……”递进警示句式（\textit{Jo..., desto...}）与紧迫临界点}
\begin{enumerate}
  \item \textbf{\dabold{Jo + 比较级 + 主语 + 动词, desto + 比较级 + 动词 + 主语}}：\newline
    这是丹麦语高阶议论文最易丢分的语序考点！牢记“\textbf{前从后主，前不倒装、后必倒装}”：
    \begin{itemize}
      \item 前半截（\textit{Jo}-从句）：比较级放最前面，后面保持正常从句语序（\textbf{主语在前，动词在后}）：\textit{Jo hurtigere \da{teknologien udvikler} sig...}
      \item 后半截（\textit{desto}-主句）：比较级放最前面，但整个 \textit{desto + 比较级} 后\textbf{紧跟变位动词}（V2 倒装）：\textit{desto skarpere \dabold{rejser den} fundamentale dilemmaer.}
    \end{itemize}
  \item \textbf{临界点警示从句：\dabold{i det øjeblik, ...}} 与 \textbf{\dabold{medmindre ... i tide}}：\newline
    用 \textit{Problemet opstår i det øjeblik, skellet begynder at skride}（问题产生于界线开始松动的那一刻）和 \textit{medmindre vi griber ind i tide}（除非我们及时干预）制造强烈的紧迫感与行动力呼吁。
\end{enumerate}
\end{grammatikbox}

\begin{grammatikbox}{语法点 3：C1 生命伦理写作四大高频避坑点}
\begin{enumerate}
  \item \textbf{\ding{172} 切勿把 \textit{Jo..., desto...} 后半句写错语序！}：\newline
    \ding{55}\ \textit{Jo mere vi optimerer generne, desto større risikoen bliver.} $\rightarrow$ \ding{51}\ \textit{Jo mere vi optimerer generne, desto større \da{bliver risikoen}.}（注意在后半句，若比较级修饰名词 \textit{større risiko}，完整短语 \textit{desto større risiko} 在前也可能触发 \textit{opstår der} 等倒装！）
  \item \textbf{\ding{173} \textit{præcedens} 既不可数也不能加不定冠词 \textit{en}}：规范表达为 \da{skabe præcedens for...}（不写 \ding{55}\ \textit{skabe en præcedens}）。
  \item \textbf{\ding{174} 区分 \textit{at helbrede}（彻底治愈）与 \textit{at behandle}（治疗/护理）}：能根治遗传病用 \textit{helbrede}；泛指医疗过程与疗法用 \textit{behandle / behandling}。
  \item \textbf{\ding{175} \textit{at skride}（滑动/松动）完成时用 \textit{er}}：写“界线已经开始松动”时必须写 \textit{skellet \da{er begyndt} at skride} 或 \textit{grænsen \da{er skredet}}，切勿用 \textit{har skredet}。
\end{enumerate}
\end{grammatikbox}

% ---------------------------------------------------------------------
\ncdsubsec{五、Mønsterøvelser（句型与修辞操练）}

\textbf{练习 A：生命伦理核心搭配与警示句型精准填空（每空填入最准确的 1 个丹麦语词）}
\begin{enumerate}
  \item Kunstig intelligens i sundhedsvæsenet \underline{\hspace{2.2cm}} fundamentale etiske dilemmaer om patientens privatliv. (rejse，现在时)
  \item At redigere gener af kosmetiske årsager kan skabe \underline{\hspace{2.2cm}} for et nyt genetisk klasseskel. (先例)
  \item Hvis vi tillader genetisk optimering, bevæger vi os hurtigt ud ad en farlig \underline{\hspace{2.2cm}}. (滑坡)
  \item Ændringer i menneskets arvemasse kan få \underline{\hspace{2.2cm}} konsekvenser for kommende generationer. (难以预料的)
  \item Lovgivere og borgere må stille det svære spørgsmål: Hvor sætter vi \underline{\hspace{2.2cm}}? (边界，带定指词尾)
  \item Jo mere præcis genredigeringen bliver, \underline{\hspace{2.2cm}} skarpere bliver behovet for klare regler.
  \item Problemet opstår i det øjeblik, skellet mellem behandling og forbedring er \underline{\hspace{2.2cm}}. (skride，完成分词)
  \item Samfundet mister kontrollen med udviklingen, \underline{\hspace{2.2cm}} vi griber ind med bindende love i tide. (除非)
\end{enumerate}

\textbf{练习 B：C1 伦理句法纠错与推演升级（修正语序/搭配错误，或按提示升级）}
\begin{enumerate}
  \item 【\textit{Jo..., desto...} 倒装纠错】\textit{Jo hurtigere teknologien vokser, desto sværere det bliver at regulere den.} $\rightarrow$
  \item 【先例搭配与冠词纠错】\textit{Dette forslag kan lave en præcedens til kommerciel brug af fostres DNA.} $\rightarrow$
  \item 【助动词与过去分词纠错】\textit{Mange etiske eksperter frygter, at grænsen mellem helbredelse og optimering allerede har skredet.} $\rightarrow$
  \item 【核心动词搭配纠错】\textit{Ny medicinsk teknologi laver mange etiske dilemmaer og får uoverskuelig konsekvenser.} $\rightarrow$
  \item 【两难推演升级】用 \textit{i behandlingens tjeneste ... men skabe præcedens for ...} 改写：\textit{Genredigering er god til at kurere sygdom, men kan også føre til designerbørn.} $\rightarrow$
  \item 【递进警示升级】用 \textit{Jo ..., desto ...} 合并：\textit{Algoritmerne bliver klogere. Det bliver også mere presserende at beskytte patientens selvbestemmelse.} $\rightarrow$
\end{enumerate}

\textbf{练习 C：中译丹实战（精准运用伦理两难推演与警示句型）}
\begin{enumerate}
  \item 人工智能在诊断中的应用引发了关于隐私保护的根本性伦理困境。
  \item 哪怕出于善意的基因干预，也可能为将人类性状商品化开创先例。
  \item 许多伦理学家警告，社会正在走上一条带来难以预料后果的滑坡。
  \item 技术发展得越快，立法者及时出手干预的需求就越迫切。
  \item 除非我们尽早划定具有约束力的法律界线，否则医疗团结互助将瓦解。
  \item 面对未来的基因医学，我们必须追问自己究竟把边界划在哪里。
\end{enumerate}

% ---------------------------------------------------------------------
\ncdsubsec{六、PD3 / Studieprøven 实战微任务与参考答案}

\begin{emailbox}{PD3 / Studieprøven 写作微任务（Etisk debatindlæg · 200--250+ ord）}
\textbf{Situation}: Et privat medicinalfirma ønsker at lancere en ny genteknologi i Danmark, der både kan fjerne alvorlige arvelige sygdomme og mod ekstra betaling reducere risikoen for almindelige livsstilssygdomme samt styrke børns fysiske udholdenhed allerede før fødslen.\\[3pt]
\textbf{Opgave}: Skriv et argumenterende debatindlæg på \textbf{200--250 ord}, hvor du analyserer de etiske muligheder og risici ved forslaget og tager klar stilling. Inddrag mindst \textbf{fire} af lektionens nøgleudtryk (\textit{rejse fundamentale etiske dilemmaer, skabe præcedens for, en glidebane, uoverskuelige konsekvenser, hvor sætter vi grænsen?, i behandlingens tjeneste}) samt en \textit{Jo..., desto...}-sætning.
\end{emailbox}

\begin{answerbox}
\textbf{Før du læser 参考回答}
\begin{enumerate}
  \item \textit{Ifølge Tekst A opstår etiske dilemmaer, når skellet mellem helbredelse og menneskelig forbedring skrider, fordi indgreb i arvemassen kan skabe præcedens for et genetisk klasseskel med uoverskuelige konsekvenser. AI-diagnostik kan desuden true patientens selvbestemmelse, retten til ikke at vide og solidariteten, hvis forsikringsselskaber får adgang til risikoprofiler.}
  \item \textit{Hun advarer mod et totalt forbud, fordi passivitet og tøven også har høje etiske omkostninger for patienter, der lider af uhelbredelige sygdomme, når sikre behandlinger faktisk findes.}
\end{enumerate}

\textbf{Tekst B 阅读核对答案}
\begin{enumerate}
  \item \da{Forkert}（她明确表示丹麦没有任何严肃研究者希望出现定制婴儿或商业基因市场：\textit{Ingen seriøse forskere i Danmark ønsker designerbørn}）
  \item \da{At man svigter patienter, der i dag lever med smertefulde og uhelbredelige sygdomme.}（辜负了正在忍受痛苦与不治之症的当代患者）
  \item \da{Skarpe juridiske rammer mellem behandling og forbedring, fuldt informeret samtykke og lige adgang i det offentlige sundhedsvæsen.}（划定区分治疗与改良的严格法律框架、坚持知情同意并确保公共医疗体系内的平等可及）
\end{enumerate}

\textbf{练习 A 答案}
\begin{enumerate}
  \item \da{rejser}（固定搭配：\textit{rejse fundamentale etiske dilemmaer} 引发根本性伦理困境）
  \item \da{præcedens}（零冠词固定搭配：\textit{skabe præcedens for} 为……开创先例）
  \item \da{glidebane}（警示隐喻：\textit{bevæge sig ud ad en farlig glidebane} 走上危险滑坡）
  \item \da{uoverskuelige}（复数形容词，修饰复数 \textit{konsekvenser}：难以预料的后果）
  \item \da{grænsen}（核心警示设问：\textit{Hvor sætter vi grænsen?}）
  \item \da{desto}（递进比较连词结构：\textit{Jo ..., desto ...}）
  \item \da{skredet}（变化动词 \textit{at skride} 过去分词，配助动词 \textit{er}）
  \item \da{medmindre}（限制条件连词：“除非……”，后接肯定从句）
\end{enumerate}

\textbf{练习 B 答案（参考）}
\begin{enumerate}
  \item \textit{Jo hurtigere teknologien vokser, desto sværere \da{bliver det} at regulere den.}（\textit{desto}-主句部分在比较级短语后必须触发 V2 倒装，将动词 \textit{bliver} 置于主语 \textit{det} 之前）
  \item \textit{Dette forslag kan \da{skabe præcedens for} kommerciel brug af fostres DNA.}（用动词 \textit{skabe} 替代 \textit{lave}，去掉多余不定冠词 \textit{en}，并将介词改用 \textit{for}）
  \item \textit{Mange etiske eksperter frygter, at grænsen mellem helbredelse og optimering allerede \da{er skredet}.}（状态变化动词 \textit{at skride} 的完成时必须使用助动词 \textit{er}）
  \item \textit{Ny medicinsk teknologi \da{rejser} mange etiske dilemmaer og får \da{uoverskuelige} konsekvenser.}（改用固定搭配动词 \textit{rejser}，并将复数形容词 \textit{uoverskuelige} 补全词尾 \textit{-e}）
  \item \textit{Genredigering bør kun bruges \da{i behandlingens tjeneste}, da kommercielle tilbud ellers kan \da{skabe præcedens for} designerbørn.}（运用课内高级名词短语与先例搭配完成论证升级）
  \item \textit{\da{Jo} klogere algoritmerne bliver, \da{desto} mere presserende \da{bliver det} at beskytte patientens selvbestemmelse.}（用 \textit{Jo..., desto...} 严格按“前不倒装、后必倒装”规范合并）
\end{enumerate}

\textbf{练习 C 答案}
\begin{enumerate}
  \item \textit{Brugen af kunstig intelligens i diagnostik rejser fundamentale etiske dilemmaer om privatlivet.}
  \item \textit{Selv velmente genetiske indgreb kan skabe præcedens for at gøre menneskelige egenskaber til varer.}
  \item \textit{Mange etikere advarer om, at samfundet bevæger sig ud ad en farlig glidebane.}
  \item \textit{Jo hurtigere teknologien udvikler sig, desto stærkere bliver behovet for lovgivning i tide.}
  \item \textit{Medmindre vi sætter bindende juridiske grænser snart, smuldrer solidariteten i sundhedsvæsenet.}
  \item \textit{Over for fremtidens genmedicin må vi spørge os selv, hvor vi sætter grænsen.}
\end{enumerate}

\textbf{PD3 / Studieprøven 高分范文（约 230 词，不含标题）}\\
\textit{\textbf{Genmedicin skal helbrede sygdom --- ikke sælge menneskelig optimering}}

\textit{Det private medicinalfirmas forslag om at kombinere sygdomsbehandling med betalt genetisk forbedring før fødslen \dabold{rejser fundamentale etiske dilemmaer} for hele samfundet. Ganske vist er det både nødvendigt og moralsk forsvarligt at bekæmpe alvorlige arvelige lidelser, når forskningen endelig stiller sikre og præcise redskaber til rådighed. Så længe genredigering udelukkende anvendes \da{i behandlingens tjeneste} til at fjerne livstruende sygdomme, tjener den direkte patienternes værdighed og livskvalitet.}

\textit{Problemet opstår imidlertid i det øjeblik, en kommerciel aktør begynder at sælge fysisk udholdenhed og forebyggende optimering mod ekstra betaling til raske familier. Her risikerer vi nemlig at \dabold{skabe præcedens for}, at menneskets arvemasse gøres til en almindelig handelsvare. Hvis velhavende forældre kan købe sig til bedre biologiske forudsætninger for deres ufødte børn, bevæger samfundet sig hurtigt ud ad \dabold{en farlig glidebane}, der kan få \dabold{uoverskuelige konsekvenser} for den sociale lighed og \da{solidariteten} i velfærdssamfundet. \da{Jo} mere markedskræfterne får lov til at definere det ”perfekte” barn, \da{desto} skarpere tegner konturerne af et nyt genetisk klasseskel sig over generationer.}

\textit{Spørgsmålet er derfor uafviseligt: \dabold{Hvor sætter vi grænsen?} Svaret bør naturligvis ikke være et totalt forbud mod genforskning, for det ville svigte patienter med alvorlige arvelige sygdomme og uhelbredelige lidelser. I stedet bør Folketinget trække en knivskarp juridisk streg: Dokumenteret terapeutisk behandling skal godkendes og finansieres i det offentlige sundhedsvæsen med fuldt samtykke, mens kommerciel genetisk forbedring af raske fostre bør forbydes ved lov.}

\small\textbf{名师点评}：正文（不含标题）约 230 词，紧扣商业化基因增强提案展开犀利而稳健的伦理辨析。首段用 \textit{Ganske vist} 与 \textit{i behandlingens tjeneste} 肯定治病价值，确立审慎立场；第二段借 \textit{i det øjeblik...}、\textit{skabe præcedens for}、\textit{en farlig glidebane}、\textit{uoverskuelige konsekvenser} 以及严丝合缝的 \textit{Jo..., desto...} 倒装句完成风险推演；末段以 \textit{Hvor sætter vi grænsen?} 设问收束，提出“公立医疗资助治病、立法禁止商业优化”的政策区分，兼具人文温度与法理严密度！
\end{answerbox}

% =====================================================================
% Unit 4
% =====================================================================
\ncdunit{4}{Litteratur, Historie og Akademisk Perspektiv}{文学、历史与高阶学术（衔接 C1）}{Lektion 13--16：安徒生童话双关与讽刺解读 · 高阶书面连词与代词 · 多源综合论述（Syntese）· 双语人生哲学散文总收官}

\ncdlesson{13}{H.C. Andersens skygge: Eventyr for børn eller satire for voksne?}{安徒生的影子：儿童童话还是成人讽刺？}
{DU3 Modul 5 (PD3) $\rightarrow$ Modul 6 (Studieprøven, B2 $\rightarrow$ C1)｜对标 NCE 4 \& \textit{Litterær Kritik, Satire \& Undertekst}}
{\textbf{文学批评与文本潜台词分析语体}（双关多义 \textit{tvetydighed} · 文本双层意涵 \textit{dobbeltbund} · 潜台词 \textit{undertekst} · 阶级讽刺 \textit{klassesatire} · 隐喻寓言 \textit{allegori}）· 文学评论核心动词搭配（\textit{afspejle, udstille, gennemskue, fortolke som}）与表里对比句式}
{离开丹麦后，人们提起汉斯·克里斯蒂安·安徒生（H.C. Andersen），脑海中浮现的往往是温馨的睡前童话或动画改编。然而在丹麦文学批评（\textit{litteraturkritik}）与 Studieprøven 文化分析写作中，读懂安徒生恰恰意味着揭开“天真表层”之下的\textbf{文本双层意涵（\textit{dobbeltbund}）与成人讽刺}。本课带你精读《皇帝的新衣》《影子》与《丑小鸭》的深层潜台词，掌握客观、敏锐的丹麦语文学评论语体。}

% ---------------------------------------------------------------------
\begin{foerduleaeser}
带着下面两个问题阅读 Tekst A 和 Tekst B：
\begin{enumerate}
  \item \textit{Hvorfor mener litteraturforskere i Tekst A, at H.C. Andersens eventyr rummer en ”dobbeltbund”, og hvad udstiller både ”Kejserens nye klæder” og ”Skyggen” under overfladen?}（根据 Tekst A，为什么文学学者认为安徒生的童话包含“双层意涵”？《皇帝的新衣》和《影子》在表层之下分别揭露了什么？）
  \item \textit{Hvilken uenighed opstår mellem Henrik og Sofie i Tekst B, når de diskuterer, hvordan ”Den grimme ælling” bør læses?}（在 Tekst B 中，当亨里克与索菲娅讨论应如何解读《丑小鸭》时，两人产生了什么分歧？）
\end{enumerate}
\end{foerduleaeser}

% ---------------------------------------------------------------------
\ncdsubsec{二、Tekster（文学批评专论与播客对谈双语篇）}

\subsection*{Tekst A：\textit{Bag den uskyldige overflade: Den skarpe satire hos H.C. Andersen}（文学评论专论 · 约 360 词）}

\begin{tipbox}{语篇说明}
本篇是一篇标准的 C1 文学评论短论，剖析安徒生童话如何同时面向儿童与成人读者。请重点关注\textbf{\color{NCDNavy}深蓝色标记的文学批评术语与分析动词}（\textit{dobbeltbund, undertekst, klassesatire, allegori, afspejle, udstille, gennemskue}），以及\textbf{\color{NCDRed}深红色标记的表层与深层对照句式}。
\end{tipbox}

Uden for Danmark bliver H.C. Andersen ofte opfattet som en harmløs fortæller af søde godnathistorier for børn, hvor alting til sidst falder trygt på plads. Blandt danske litteraturforskere hersker der imidlertid bred enighed om, at hans eventyr er skrevet med en bevidst \da{dobbeltbund}. Allerede i sin egen levetid understregede Andersen, at han fortalte for både børn og voksne, og \dabold{under den tilsyneladende uskyldige overflade gemmer der sig} en skarp \da{tvetydighed}, der kræver, at læseren evner at læse mellem linjerne.

Tag for eksempel \textit{Kejserens nye klæder} fra 1837. Ved første øjekast ligner fortællingen en humoristisk fabel om en narragtig hersker, der lader sig narre af to snedige vævere. \dabold{Ved nærmere eftersyn fungerer eventyret snarere som} en bidende \da{klassesatire} over konformitet, smiger og intellektuel fejhed i magtens korridorer. Ingen ved hoffet tør indrømme, at de ikke kan se det usynlige stof, af frygt for at blive stemplet som dumme eller uegnede til deres høje embede. Dermed \da{udstiller} teksten et kollektivt \da{selvbedrag}, hvor kejserens \da{forfængelighed} og omgivelsernes tavshed gensidigt forstærker hinanden, indtil et barn bryder illusionen med en enkel, ufiltreret sandhed.

Endnu mørkere bliver denne \da{undertekst} i \textit{Skyggen} fra 1847, der ofte regnes for et af Andersens mest \da{uhyggelige} og \da{flertydige} værker. Her mister en godtroende lærd sin egen skygge under en rejse i syden, hvorefter skyggen år senere vender tilbage som en velhavende og kynisk verdensmand. Skridt for skridt bytter de to skikkelser roller: den tomme form overtager magten, gifter sig med en prinsesse og lader til sidst sin tidligere herre henrette for at skjule sandheden om sin oprindelse. \dabold{Snarere end blot at være} en fantastisk spøgelseshistorie læses \textit{Skyggen} i dag som en eksistentiel \da{allegori} over modernitetens overfladiskhed, hvor kynisme og opportunistisk imagepleje kan fortrænge ægte integritet.

Selv et elsket værk som \textit{Den grimme ælling} \da{afspejler} en social smerte, som rækker langt ud over børneværelset. Fortællingen kan ganske vist læses som en trøstende historie om at finde sit sande jeg, men den rummer samtidig en bitter kritik af 1800-tallets stive klasseskel og snæversynede andegård. Andersens kunstneriske storhed ligger netop i denne uforløste spænding mellem barnlig naivitet og voksen \da{spidsfindighed}: den, der kun ser det trygge eventyr, overser halvdelen af værket.

\begin{transbox}
\small
\textbf{\color{NCDTeal}[Tekst A 中文精译]}\\
在丹麦以外的地方，安徒生常常被看作一位人畜无害、专门为儿童讲述甜蜜睡前故事的作家——在他的故事里，一切最终都会安稳妥帖地各归其位。然而在丹麦文学学者之间却存在广泛共识：他的童话是带着自觉的“双层意涵”写就的。早在生前，安徒生就曾强调自己既为孩子讲故事，也为成年人写作；而在看似天真无邪的表层之下，实则隐藏着敏锐的双关多义，要求读者具备“读出字里行间弦外之音”的能力。

以 1837 年的《皇帝的新衣》为例。乍看之下，这个故事像是一则幽默寓言，讲述一位愚蠢的统治者如何被两个狡猾的织工欺骗。但经仔细审视，这篇童话更像是一部辛辣的阶级讽刺之作，直指权力走廊中的从众心理、阿谀奉承与智识怯懦。宫廷里没有人敢承认自己看不见那块隐形布料，因为人人都害怕被贴上愚蠢或不称职于自己高位的标签。文本由此暴露出一种集体自我欺骗：皇帝的虚荣心与周围人的沉默彼此强化，直到一个孩子用一句简单而未经修饰的真话打破幻象。

这种潜台词在 1847 年的《影子》中变得更加幽暗，该作常被视为安徒生最令人毛骨悚然且意涵繁复的作品之一。故事里，一位轻信他人的学者在南方旅行时失去了自己的影子，多年后影子却以一位富有而犬儒的世故名流身份归来。一步步地，这两个形象互换了角色：空洞的躯壳篡夺了权力，迎娶了公主，最终甚至下令处决昔日的主人，以掩盖自己出身的真相。《影子》与其说仅仅是一个奇幻鬼故事，不如说在今天被解读为一则关于现代性浅薄本质的存在主义寓言——在这个世界里，犬儒主义与投机取巧的形象包装能够排挤并取代真正的品格操守。

就连《丑小鸭》这样广受喜爱的作品，也折射出一种远远超出儿童房范畴的社会痛感。这个故事固然可以被读作寻找真实自我的心灵慰藉，但与此同时，它也包含着对十九世纪僵化阶级壁垒以及狭隘“鸭群世界”的苦涩批判。安徒生的艺术伟大之处，恰恰在于这种孩童般的纯真与成人般的精微敏锐之间那未曾消解的张力：只看得见令人安心的童话表层的人，便错过了这部作品的一半。
\end{transbox}

\subsection*{Tekst B：\textit{Debat i litteraturpodcasten: Bør ”Den grimme ælling” læses som trøst eller samfundskritik?}（文学播客对谈 · 约 140 词）}

\textbf{阅读任务（理解与核对）：}\textit{Læs samtalen mellem litteraturformidleren Henrik og gymnasielæreren Sofie, og besvar spørgsmålene nedenfor.}

\begin{emailbox}{播客对谈｜Henrik（文学节目主持）\& Sofie（高中丹麦语教师）}
\textbf{Henrik:} Sofie, mange voksne husker stadig \textit{Den grimme ælling} som en varm fortælling om håb og tålmodighed. Men når vi læser den igen i dag, virker den næsten lidt ubarmhjertig, ikke?\\
\textbf{Sofie:} Jo, det gør den faktisk. \dabold{Hvor} mange opfatter historien som en simpel trøst, \dabold{rummer} teksten jo en tydelig \da{klassesatire}. Ællingen bliver slet ikke accepteret for sin karakter eller venlighed, men først respekteret, da de andre opdager, at den tilhører en finere art.\\
\textbf{Henrik:} Præcis --- hele moralen hviler næsten på, at man er født som svane fra begyndelsen. Betyder det så, at eventyret i virkeligheden er elitært?\\
\textbf{Sofie:} Det kan godt virke sådan ved første øjekast. Alligevel bør man \da{fortolke} eventyret \da{som} både et selvforsvar fra en fattig skomagersøn og en skarp kritik af snobberiet i datidens København. Hvis vi \da{gennemskuer} den \da{tvetydighed}, bliver Andersens \da{undertekst} langt mere spændende end nogen enkel moral.
\end{emailbox}

\textbf{Tekst B 阅读核对（Rigtigt / Forkert / Kort svar，答案见第六部分）：}
\begin{enumerate}
  \item \textit{Henrik og Sofie er enige om, at ”Den grimme ælling” udelukkende er en varm og simpel børnehistorie.} \underline{\hspace{2.2cm}}
  \item \textit{Hvornår bliver ællingen ifølge Sofie endelig respekteret af de andre i andegården?} \underline{\hspace{2.2cm}}
  \item \textit{Hvilke to ting kan man ifølge Sofie fortolke eventyret som?} \underline{\hspace{2.2cm}}
\end{enumerate}

\begin{tipbox}{本课发音与语调要点（Udtale \& Prosodi）}
\begin{itemize}
  \item \textbf{文学批评核心词的重音与静音字母}：\textit{en tvetydighed} \ipa{tˢveˈtˢyːðiˌheðˀ}（重音在第二音节 \textit{-ty-}，词尾带弱化 \textit{stød}）；\textit{tilsyneladende} \ipa{tˢelˈsyːnəˌlæːˀðənə}；\textit{en forfængelighed} \ipa{fʌˈfɛŋˀəliˌheðˀ}（\textit{-lig-} 中的 \textit{g} 不发音）；\textit{uhyggelig} \ipa{uˈhyɡəli}（词尾 \textit{g} 不发音！）。
  \item \textbf{外来文学学术词重音后移}：\textit{en allegori} \ipa{aleɡoˈʁiːˀ}（重音与 \textit{stød} 落在最末长元音 \textit{-ri} 上！）；\textit{en satire} \ipa{saˈtˢiːɐ}（重音在 \textit{-ti-}）；\textit{spidsfindig} \ipa{sb̥esˈfenˀi}（\textit{spids} 中的 \textit{d} 不发音，词尾 \textit{g} 不发音）。
  \item \textbf{对照句型的语流重音}：朗读 \textit{Ved første øjekast... Ved nærmere eftersyn...} 时，对比重音必须清晰落在 \textit{\textbf{før}ste} 与 \textit{\textbf{nær}mere} 上，借音调反差体现“表层误解 vs 深度洞见”的批评转折。
\end{itemize}
\end{tipbox}

\begin{transbox}
\small
\textbf{\color{NCDTeal}[Tekst B 中文对照]}\\
\textbf{亨里克：}索菲娅，许多成年人至今仍把《丑小鸭》记作一个关于希望与耐心的温暖故事。但当我们今天重读它时，它给人的感觉几乎有点冷酷无情，不是吗？\\
\textbf{索菲娅：}是的，确实如此。人们往往把这则故事看作简单的安慰，但文本中其实包含着明显的阶级讽刺。小鸭子根本不是因为自己的品格或善良而被接纳，而是直到其他动物发现它本就属于更高贵的物种时，才开始受到尊重。\\
\textbf{亨里克：}没错——整篇故事的寓意几乎建立在“你从一开始就天生是天鹅”这一点上。那这是不是意味着，这篇童话实际上带有精英主义色彩？\\
\textbf{索菲娅：}乍看之下确实可能给人这种印象。尽管如此，我们还是应当把这篇童话既解读为一位贫苦鞋匠之子的自我辩护，也解读为对当时哥本哈根势利风气的尖锐批判。只要我们看穿这种双关多义，安徒生的潜台词就会比任何单一的道德说教都精彩得多。
\end{transbox}

% ---------------------------------------------------------------------
\ncdsubsec{三、Gloser（核心生词表）}

\begin{tipbox}{词表重点}
文学分析极其讲究“动作指向与搭配”：表示“把某物解读为……”固定使用 \textit{at fortolke noget \dabold{som} noget}（而非 \textit{til}）；表示“对某现象的寓言”用 \textit{en allegori \dabold{over / på} noget}；表示“揭露/曝露虚伪”则使用动词 \textit{at udstille} 或 \textit{at gennemskue}。此外，\textit{tvetydighed} 与 \textit{flertydig} 中词根 \textit{tyd-} 均来自 \textit{betyde}（意指），指“多重含义”，切勿与表示“怀疑”的 \textit{tvivl} 混淆。
\end{tipbox}

\begin{glosertable}
名词 & en skygge & \textit{(skyggen, skygger, skyggerne)} & 影子，阴影 \\
名词 & et eventyr & \textit{(eventyret, eventyr, eventyrene)} & 童话，奇遇故事 \\
名词 & en tvetydighed & \textit{(tvetydigheden, tvetydigheder, tvetydighederne)} & 双关多义，模棱两可 \\
名词 & en undertekst & \textit{(underteksten, undertekster, underteksterne)} & 潜台词，言外之意 \\
名词 & en dobbeltbund & \textit{(dobbeltbunden, kun ental)} & 文本双层意涵，深层底蕴 \\
名词 & en allegori & \textit{(allegorien, allegorier, allegorierne)} & 寓言，隐喻象征作品 \\
名词 & en klassesatire & \textit{(klassesatiren, klassesatirer, klassesatirerne)} & 阶级讽刺（作品） \\
名词 & en forfængelighed & \textit{(forfængeligheden, kun ental)} & 虚荣，爱慕虚荣 \\
名词 & et selvbedrag & \textit{(selvbedraget, selvbedrag, selvbedragene)} & 自我欺骗，自欺 \\
名词 & en ælling & \textit{(ællingen, ællinger, ællingerne)} & 小鸭，雏鸭 \\
名词 & en spidsfindighed & \textit{(spidsfindigheden, spidsfindigheder, spidsfindighederne)} & 精微机敏，尖锐讽刺感 \\
\midrule
动词 [har] & at afspejle & \textit{(afspejler, afspejlede, har afspejlet)} & 折射，反映出 \\
动词 [har] & at udstille & \textit{(udstiller, udstillede, har udstillet)} & 暴露，揭露（缺点/虚伪） \\
动词 [har] & at fortolke & \textit{(fortolker, fortolkede, har fortolket)} & 解读，阐释 \\
动词 [har] & at gennemskue & \textit{(gennemskuer, gennemskuede, har gennemskuet)} & 看穿，识破 \\
\midrule
形容词 & tilsyneladende & \textit{(tilsyneladende, 不变)} & 表面上的，看似的 \\
形容词 & flertydig & \textit{(flertydigt, flertydige)} & 多义的，含意丰富的 \\
形容词 & uhyggelig & \textit{(uhyggeligt, uhyggelige)} & 阴森的，令人不寒而栗的 \\
\end{glosertable}

% ---------------------------------------------------------------------
\ncdsubsec{四、Grammatik \& Stilistik（文学批评修辞与潜台词分析精讲）}

\begin{grammatikbox}{修辞支柱 1：由表及里的文学批评“双层对照”句型框架}
写 C1 级别文学或文化评论时，最核心的推进逻辑是\textbf{先交代大众对文本的“表层误读”，再翻转切入“深层潜台词分析”}。请熟记以下四组高阶句型公式：

\vspace{0.3em}
\renewcommand{\arraystretch}{1.24}
{\small
\begin{tabularx}{\linewidth}{@{} >{\bfseries\raggedright\arraybackslash}p{3.0cm} >{\raggedright\arraybackslash}p{5.2cm} >{\raggedright\arraybackslash}X @{}}
\toprule
\textbf{分析维度} & \textbf{核心句型公式（注意 V2 倒装）} & \textbf{Tekst A/B 示范句} \\
\midrule
\da{1. 表层印象 vs 深度洞见} &
\textit{\dabold{Ved første øjekast} ligner X...}\newline
\textit{\dabold{Ved nærmere eftersyn} fungerer X snarere som...} &
\textit{Ved nærmere eftersyn fungerer eventyret snarere som en bidende klassesatire.} \\
\midrule
\da{2. 否定单一读法并升级} &
\textit{\dabold{Snarere end blot at} være A, læses værket som B.} &
\textit{Snarere end blot at være en spøgelseshistorie læses Skyggen som en allegori.} \\
\midrule
\da{3. 揭示潜台词与张力} &
\textit{Under den tilsyneladende... gemmer der sig...}\newline
\textit{Teksten rummer en bevidst dobbeltbund.} &
\textit{Under den uskyldige overflade gemmer der sig en skarp tvetydighed.} \\
\midrule
\da{4. 对比大众与专业视角} &
\textit{\dabold{Hvor} mange opfatter X som..., \dabold{peger} forskere på...} &
\textit{Hvor mange opfatter historien som en trøst, rummer teksten jo en klassesatire.} \\
\bottomrule
\end{tabularx}
}
\end{grammatikbox}

\begin{grammatikbox}{语法点 2：文学批评的核心句法结构与长句凝练技巧}
\begin{enumerate}
  \item \textbf{介词短语置于句首时的 V2 倒装铁律}：\\
    \textit{Ved første øjekast}, \textit{Ved nærmere eftersyn}, \textit{Under overfladen} 以及 \textit{Snarere end blot at + 不定式} 放在句首充当状语成分时，后面主句的第一位必须是\textbf{变位动词}（丹麦语不在这些句首状语后打逗号）：\\
    \ding{51}\ \textit{Under overfladen \dabold{gemmer} der sig en skarp kritik.}（\ding{55}\ \textit{*Under overfladen, der gemmer sig...}）\\
    \ding{51}\ \textit{Snarere end blot at underholde \dabold{udfordrer} eventyret læserens syn på magt.}
  \item \textbf{客观引述与被动评价句型（避免滥用“我认为”）}：\\
    文学学术文章常借助非人称或被动表达提升客观权威感：\\
    \ding{172} \textit{Værket \da{regnes ofte for} at være...}（该作常被视为……）\\
    \ding{173} \textit{Fortællingen \da{kan læses som} en allegori over...}（故事可被解读为关于……的寓言）\\
    \ding{174} \textit{Blandt litteraturforskere \da{hersker der enighed om, at}...}（学者之间普遍认同……）
  \item \textbf{形式主语 \textit{der} + 存现/隐藏动词}：表达“字里行间隐藏着某种深意”时，用 \textit{gemmer der sig + 不定名词} 或 \textit{ligger der + 不定名词}，将信息焦点自然留到句末。
\end{enumerate}
\end{grammatikbox}

\begin{grammatikbox}{文学批评避坑清单（Top 4 Pitfalls）}
\begin{enumerate}
  \item \textbf{\ding{172} \textit{fortolke} 与 \textit{læse} 表“解读为”时只能接 \dabold{som}}：\ding{55}\ \textit{*Man kan fortolke Skyggen til en kritik af magt.} $\rightarrow$ \ding{51}\ \textit{Man kan fortolke Skyggen \da{som} en kritik af magt.}
  \item \textbf{\ding{173} \textit{en allegori} 后接 \dabold{over} 或 \dabold{på}，切勿直译英文 \textit{of}}：\ding{51}\ \textit{en allegori \da{over / på} modernitetens overfladiskhed}（而非 \ding{55}\ \textit{*en allegori af}）。
  \item \textbf{\ding{174} 区分 \textit{en skæbne}（命运）与 \textit{et eventyr}（童话/奇遇）的词性}：\textit{eventyr} 是中性名词，单复同形（\textit{et eventyr, alle hans eventyr}），切勿受“很多童话”影响误加复数词尾 \ding{55}\ \textit{*eventyrer}（\textit{en eventyrer} 意思是“冒险家”！）。
  \item \textbf{\ding{175} 区分 \textit{snarere end} 与 \textit{hellere end}}：\textit{hellere end} 表达个人偏好（宁愿）；文学批评中表达“与其说是 A，不如说是 B”的客观重新定性，必须用 \da{snarere end}。
\end{enumerate}
\end{grammatikbox}

% ---------------------------------------------------------------------
\ncdsubsec{五、Mønsterøvelser（文学分析句型操练）}

\textbf{练习 A：文学批评核心词汇与介词搭配精准填空（每空填入最准确的 1 个丹麦语词）}
\begin{enumerate}
  \item Under den uskyldige overflade rummer Andersens eventyr en bevidst \underline{\hspace{2.2cm}}.
  \item Ved nærmere eftersyn fungerer \textit{Kejserens nye klæder} som en bidende \underline{\hspace{2.2cm}} over magtens smiger.
  \item Teksten \underline{\hspace{2.2cm}} et kollektivt selvbedrag, hvor ingen ved hoffet tør sige sandheden højt.
  \item Man kan fortolke \textit{Skyggen} \underline{\hspace{2.2cm}} en eksistentiel advarsel mod tom opportunisme.
  \item Eventyret \underline{\hspace{2.2cm}} på én gang tidens stive klasseskel og forfatterens egen erfaring med social udstødelse.
  \item Snarere \underline{\hspace{2.2cm}} blot at trøste børn stiller fortællingen ubehagelige spørgsmål til de voksne.
  \item Hvis man kun læser handlingen bogstaveligt, overser man let den kritiske \underline{\hspace{2.2cm}} mellem linjerne.
  \item Først når læseren \underline{\hspace{2.2cm}} ironien i slutningen, åbner værket sig i hele sin dybde.
\end{enumerate}

\textbf{练习 B：文学评论句式纠错与 C1 凝练升级（修正语法与搭配错误，或按要求改写）}
\begin{enumerate}
  \item 【语序纠错】\textit{Ved nærmere eftersyn, eventyret er ikke blot en sød fortælling for børn.} $\rightarrow$
  \item 【搭配纠错】\textit{Litteraturforskere fortolker ofte værket til en skarp allegori af det borgerlige samfund.} $\rightarrow$
  \item 【复数词形纠错】\textit{H.C. Andersen skrev mange berømte eventyrer, som stadig diskuteres i dag.} $\rightarrow$
  \item 【表里对比升级】用 \textit{Ved første øjekast... Ved nærmere eftersyn...} 改写：\textit{Historien virker først meget enkel, men senere opdager man, at den kritiserer samfundet.} $\rightarrow$
  \item 【否定升级】用 \textit{Snarere end blot at være... læses værket som...} 改写：\textit{Skyggen er ikke kun en uhyggelig historie; den handler i høj grad om tomme facader.} $\rightarrow$
  \item 【潜台词提炼】用 \textit{Under den tilsyneladende... gemmer der sig...} 改写：\textit{Teksten ser harmløs ud, men den indeholder faktisk en bitter kritik af snobberi.} $\rightarrow$
\end{enumerate}

\textbf{练习 C：中译丹（运用文学批评词汇与标准句式翻译下列评论句）}
\begin{enumerate}
  \item 在看似天真的表层之下，这些童话隐藏着尖锐的潜台词。
  \item 经仔细审视，该文本实际上是一出关于虚荣心的辛辣阶级讽刺。
  \item 这部作品不应只被读作儿童故事，而应被视为一则现代寓言。
  \item 安徒生的文本暴露了人们在面对权力时的集体自欺。
  \item 他的许多故事都折射出十九世纪僵化的社会等级。
  \item 只有能够看穿双关多义的读者，才能领会作者的艺术深度。
\end{enumerate}

% ---------------------------------------------------------------------
\ncdsubsec{六、PD3 / Studieprøven 实战微任务与参考答案}

\begin{emailbox}{Studieprøven / PD3 文学文化评论写作任务（220--260 ord）}
\textbf{Opgave}: Mange udlændinge kender kun H.C. Andersens eventyr fra barndommen eller fra romantiske filmudgaver. Skriv en \textbf{sammenhængende litterær og kulturel kommentar på 220--260 ord}, hvor du forklarer, hvorfor hans værker i Danmark ofte læses med en helt anden alvor. Inddrag eksempler fra mindst to af de nævnte eventyr (\textit{Kejserens nye klæder}, \textit{Skyggen} eller \textit{Den grimme ælling}), og brug relevante litterære analysebegreber samt kontraststrukturer fra lektionen.
\end{emailbox}

\begin{answerbox}
\textbf{Før du læser 参考回答}
\begin{enumerate}
  \item \textit{Litteraturforskere peger på, at Andersens eventyr har en bevidst dobbeltbund, der både henvender sig til børn og voksne. ”Kejserens nye klæder” udstiller konformitet, smiger og et kollektivt selvbedrag ved hoffet, mens ”Skyggen” fungerer som en mørk allegori over overfladiskhed og opportunisme, der fortrænger integritet.}
  \item \textit{Henrik finder eventyret næsten ubarmhjertigt, mens Sofie påpeger en klar klassesatire i fortællingen, fordi ællingen først respekteres, når omgivelserne ser, at den er en svane. Samtidig mener hun, at man bør læse det som både selvforsvar og kritik af tidens snobberi.}
\end{enumerate}

\textbf{Tekst B 阅读核对答案}
\begin{enumerate}
  \item \da{Forkert}（两人均认为现代重读时会感到某种冷酷，索菲娅更指出其中包含阶级讽刺与多重潜台词）
  \item \da{Først når de andre opdager, at den tilhører en finere art (er en svane).}
  \item \da{Både som et selvforsvar fra en fattig skomagersøn og som en kritik af snobberiet i datidens København.}
\end{enumerate}

\textbf{练习 A 答案}
\begin{enumerate}
  \item \da{dobbeltbund}（文学分析核心词：文本双层意涵，配不定冠词 \textit{en}）
  \item \da{klassesatire}（阶级讽刺作品）
  \item \da{udstiller}（动词：暴露、揭露虚伪或缺点）
  \item \da{som}（固定搭配：\textit{fortolke noget som noget} 将……解读为……）
  \item \da{afspejler}（动词：折射、反映出某种社会现实）
  \item \da{end}（固定升级结构：\textit{Snarere end blot at...} 与其仅仅……不如……）
  \item \da{undertekst}（潜台词、言外之意）
  \item \da{gennemskuer}（动词：看穿、识破反讽或表象）
\end{enumerate}

\textbf{练习 B 答案（参考）}
\begin{enumerate}
  \item \textit{Ved nærmere eftersyn er eventyret ikke blot en sød fortælling for børn.}（去掉英语式逗号，且句首状语触发主句 V2 倒装：\textit{er eventyret}）
  \item \textit{Litteraturforskere fortolker ofte værket som en skarp allegori over det borgerlige samfund.}（搭配修正：\textit{fortolke som} 与 \textit{allegori over / på}）
  \item \textit{H.C. Andersen skrev mange berømte eventyr, som stadig diskuteres i dag.}（中性名词 \textit{et eventyr} 单复同形，\textit{eventyrer} 意为“冒险家”）
  \item \textit{Ved første øjekast virker historien meget enkel, men ved nærmere eftersyn kritiserer den samfundet.}（运用“乍看……细看……”标准对照结构提升思辨层次）
  \item \textit{Snarere end blot at være en uhyggelig historie læses Skyggen som en allegori over tomme facader.}（用 \textit{Snarere end blot at...} 实现表里定性升级）
  \item \textit{Under den tilsyneladende harmløse form gemmer der sig en bitter kritik af snobberi.}（用形式主语存现结构突出潜台词焦点）
\end{enumerate}

\textbf{练习 C 答案}
\begin{enumerate}
  \item \textit{Under den tilsyneladende uskyldige overflade gemmer disse eventyr en skarp undertekst.}
  \item \textit{Ved nærmere eftersyn er teksten faktisk en bidende klassesatire over forfængelighed.}
  \item \textit{Dette værk bør ikke kun læses som børnelitteratur, men som en moderne allegori.}
  \item \textit{Andersens tekst udstiller menneskers kollektive selvbedrag i mødet med magten.}
  \item \textit{Mange af hans historier afspejler de stive sociale skel i 1800-tallet.}
  \item \textit{Kun den læser, der kan gennemskue tvetydigheden, forstår forfatterens kunstneriske dybde.}
\end{enumerate}

\textbf{PD3 / Studieprøven 高分范文（约 250 词）}\\
\textit{\textbf{Mellem barnets blik og voksensatire: H.C. Andersens skjulte verden}}

\textit{For mange læsere uden for Danmark står H.C. Andersen først og fremmest som en mild eventyrdigter, hvis historier hører hjemme i barndommens trygge rammer. \dabold{Ved nærmere eftersyn rummer} hans forfatterskab imidlertid en bevidst \da{dobbeltbund}, der retter sig direkte mod voksne og kritiske læsere. Under den tilsyneladende enkle og legende form gemmer der sig nemlig en skarp samfundskritik og en filosofisk \da{tvetydighed}, som gør teksterne vedvarende aktuelle også i det moderne samfund.}

\textit{Dette dobbelte blik træder særligt tydeligt frem i \textit{Kejserens nye klæder}. \dabold{Ved første øjekast} virker fortællingen som en munter anekdote om en forfængelig monark, men i virkeligheden fungerer den som en præcis \da{klassesatire}. Ved at beskrive et hof, hvor ingen tør indrømme det åbenlyse af frygt for at tabe ansigt eller stilling, \da{udstiller} Andersen et kollektivt \da{selvbedrag} skabt af konformitet og magtangst. Det er ikke blot kejserens \da{forfængelighed}, der afsløres, men hele elitens mangel på moralsk mod.}

\textit{Endnu mere foruroligende er \textit{Skyggen}. \dabold{Snarere end blot at være} en dyster gyserhistorie læses værket i dag som en tidløs \da{allegori} over, hvordan ydre overflade og kynisk imagepleje kan fortrænge sandhed og integritet. Hvor skyggen ubesværet stiger i graderne gennem tom charme, går den lærde til grunde med sine idealer intakte. På tilsvarende vis \da{afspejler} selv \textit{Den grimme ælling} en dyb og smertelig erfaring med social udstødelse og snobberi.}

\textit{Netop derfor bør Andersens værker aldrig reduceres til ufarlig nostalgi. Først når man \da{gennemskuer} den \da{flertydige} \da{undertekst}, opdager man, hvor ubarmhjertig --- og samtidig befriende ærlig --- hans litterære kunst egentlig er.}

\small\textbf{名师点评：}这篇约 250 词的范文严格契合 C1 文学文化评论写作范式：开篇点出跨文化印象与丹麦本土学术视角之别；中段紧扣《皇帝的新衣》《影子》与《丑小鸭》三篇经典，自然调度 \textit{dobbeltbund, klassesatire, selvbedrag, allegori, undertekst} 等核心批评词汇；全篇句型节奏张弛有度，精准呈现表里对照（\textit{Ved første øjekast / Ved nærmere eftersyn / Snarere end blot at...}）与主句 V2 倒装，结尾升华点题，毫无生硬拼凑感。
\end{answerbox}

\ncdlesson{14}{Danmark i Europa: Småstatsmentalitet og globalt udsyn}{欧洲中的丹麦：小国心态与全球视野}
{DU3 Modul 5 (PD3) $\rightarrow$ Modul 6 (Studieprøven, B2 $\rightarrow$ C1)｜对标 NCE 4 \& \textit{Historisk-Politisk Essayistik}}
{\textbf{高阶书面连词与关系/副词代词}（原因 \textit{eftersom, idet} · 疑问限定 \textit{hvorvidt, i hvilken udstrækning} · 让步转折 \textit{desuagtet} · 归属列举 \textit{herunder} · 方式结果 \textit{hvormed}）· 历史政论文跨时态衔接与从句语序铁律}
{丹麦在世界舞台上的处世哲学，始终处于一对耐人寻味的张力之中：一方面是 1864 年战败后深植民族记忆的\textbf{“小国心态”（\textit{småstatsmentalitet}）}与对主权让渡的审慎保留；另一方面则是依赖海洋贸易、合作社运动（\textit{andelsbevægelsen}）与绿色外交的\textbf{“全球视野”（\textit{globalt udsyn}）}。本课通过这一历史政论主题，攻克 C1 书面语中最具权威感的七大高阶连词与副词代词。}

% ---------------------------------------------------------------------
\begin{foerduleaeser}
带着下面两个问题阅读 Tekst A 和 Tekst B：
\begin{enumerate}
  \item \textit{Hvordan påvirkede nederlaget i 1864 ifølge Tekst A Danmarks selvforståelse, og hvorfor kan småstatsmentalitet og international åbenhed alligevel eksistere side om side?}（根据 Tekst A，1864 年的战败如何影响了丹麦的自我定位？为什么小国心态与国际开放依然能够并存？）
  \item \textit{Hvilket paradoks peger historikeren i Tekst B på i forbindelse med de danske EU-forbehold og afstemningen i 2022?}（在 Tekst B 中，历史学家围绕丹麦的欧盟保留条款与 2022 年公投指出了哪一悖论？）
\end{enumerate}
\end{foerduleaeser}

% ---------------------------------------------------------------------
\ncdsubsec{二、Tekster（历史政论专篇与学者短评双语篇）}

\subsection*{Tekst A：\textit{Mellem beskyttelse og åbenhed: Danmarks lange vej i Europa}（历史政论专篇 · 约 355 词）}

\begin{tipbox}{语篇说明}
本篇以严谨的历史政论文笔梳理丹麦自 1864 年至今的欧洲定位演变。阅读时请特别关注\textbf{\color{NCDRed}深红色标记的高阶书面连词与关系/副词代词}（\textit{eftersom, idet, hvorvidt, desuagtet, herunder, hvormed, i hvilken udstrækning}）各自引导的语序变化，以及\textbf{\color{NCDNavy}深蓝色标记的政治与历史核心概念}。
\end{tipbox}

For at forstå Danmarks nuværende forhold til Europa må man gå tilbage til det skæbnesvangre \da{nederlag} i 1864, hvor landet mistede omkring en tredjedel af sit territorium til Preussen og Østrig. Med ét slag blev det gamle rige reduceret til en sårbar \da{småstat}, \dabold{hvormed} en dyb forsigtighed over for stormagter slog rod i den politiske kultur. Under slagordet ”Hvad udadtil tabes, skal indadtil vindes” vendte samfundet blikket indad og byggede både folkehøjskoler og en stærk \da{andelsbevægelse}. Denne \da{småstatsmentalitet} \da{prægede} generationer af danskere, \dabold{eftersom} national overlevelse blev kædet tæt sammen med indre sammenhængskraft og folkelig dannelse.

Samtidig har Danmark aldrig haft råd til at isolere sig økonomisk fra omverdenen. Som en lille nation uden store råstoffer har velstanden tværtimod altid hvilet på søfart, landbrugseksport og teknologisk specialisering, \dabold{herunder} medicinalindustri og vedvarende energi. Efter Anden Verdenskrig og især med indtrædelsen i EF i 1973 trådte landet målrettet ind i det europæiske marked, \dabold{idet} en åben og regelbaseret handel skabte trygge rammer for eksporten. Dansk \da{udenrigspolitik} har derfor i årtier balanceret mellem ønsket om fri adgang til globale markeder og frygten for at miste demokratisk kontrol.

Denne indre spænding blev særlig tydelig under folkeafstemningen i 1992, da et snævert flertal af vælgerne sagde nej til Maastrichttraktaten og derved førte til de fire danske \da{EU-forbehold}. Befolkningen diskuterede intenst gennem hele valgkampen, \dabold{hvorvidt} et tættere samarbejde ville udhule lokal \da{selvbestemmelse}, og \dabold{i hvilken udstrækning} en småstat bør \da{afgive} \da{suverænitet} til overnationale organer i Bruxelles. \dabold{Desuagtet} har Danmark i praksis ofte hørt til de medlemslande, der hurtigst og mest loyalt har gennemført fælles europæiske regler på klima-, arbejdsmarkeds- og miljøområdet.

I de seneste år har krigen i Ukraine, ændrede stormagtsrelationer og den globale klimakrise imidlertid tvunget landet til at genoverveje sin traditionelle tilbageholdenhed. Da et klart flertal ved folkeafstemningen i 2022 stemte for at \da{ophæve} forsvarsforbeholdet, markerede det et historisk skift i dansk sikkerhedspolitik. Samtidig søger Danmark i dag aktiv indflydelse gennem grønt diplomati og multilaterale alliancer, \dabold{idet} små stater sjældent kan forme deres egen fremtid alene. Erfaringen viser således, at småstatsmentalitet og \da{globalt udsyn} ikke nødvendigvis er modsætninger, men to uundværlige sider af samme overlevelsesstrategi.

\begin{transbox}
\small
\textbf{\color{NCDTeal}[Tekst A 中文精译]}\\
要理解丹麦如今与欧洲的关系，必须回溯到 1864 年那场改变国运的战败——当时丹麦将大约三分之一的领土割让给了普鲁士与奥地利。转瞬之间，这个古老的王国缩水为一个脆弱的小国，由此一种对大国保持深层警惕的谨慎心态在政治文化中扎下了根。在“凡对外失去的，必于对内赢得”的口号下，丹麦社会将目光转向内部，兴办民众高等学院并建立起强大的合作社运动。这种小国心态深刻塑造了数代丹麦人，因为民族的生存自此与内部凝聚力及民众教养紧密捆绑在了一起。

与此同时，丹麦在经济上从来负担不起与外部世界自我孤立的代价。作为一个缺乏丰富矿产资源的弹丸之地，其繁荣反而始终建立在航海贸易、农产品出口与技术专业化之上——其中包括制药业与可再生能源。二战之后，尤其是 1973 年加入欧洲共同体（EF）以来，丹麦目标明确地融入欧洲市场，因为开放且以规则为基础的贸易为出口创造了安全框架。因此，丹麦外交政策数十年来一直在“渴望自由进入全球市场”与“担忧丧失民主控制权”之间寻求平衡。

这种内在张力在 1992 年公投中表现得尤为明显：当时微弱多数选民对《马斯特里赫特条约》投下反对票，由此催生了丹麦的四项欧盟保留条款。民众在整个竞选期间激烈辩论：更紧密的合作究竟是否会掏空地方自决权，以及一个欧洲小国应当在多大程度上向布鲁塞尔的超国家机构让渡主权。尽管如此，在实践中丹麦却往往跻身于最迅速、最忠实地在气候、劳动力市场与环境领域落实欧盟共同法规的成员国之列。

然而近年来，乌克兰战争、大国关系的变化与全球气候危机迫使这个国家重新审视其传统的克制姿态。当 2022 年公投中明确多数选民投票赞成废除防务保留条款时，这标志着丹麦安全政策的一次历史性转折。与此同时，今天的丹麦正通过绿色外交与多边联盟积极争取影响力，因为小国往往难以凭一己之力塑造自己的未来。历史经验由此表明，小国心态与全球视野未必非此即彼，而是同一套生存战略不可或缺的一体两面。
\end{transbox}

\subsection*{Tekst B：\textit{Kommentar: Når skepsis møder indflydelse i Bruxelles}（历史政治短评 · 约 145 词）}

\textbf{阅读任务（理解与核对）：}\textit{Læs historikeren Thomas Linds kommentar om Danmarks rolle i EU, og besvar spørgsmålene nedenfor.}

\begin{emailbox}{报刊短评｜Fra: Thomas Lind (lektor i europæisk historie) \quad Emne: Småstatens dilemma}
I den danske EU-debat fremstilles \da{suverænitet} ofte som en fast størrelse, man enten beholder udelt derhjemme eller taber i Bruxelles. Denne forsimplede opfattelse overser imidlertid det klassiske dilemma for enhver \da{småstat}: Formelt at stå udenfor garanterer ikke reel \da{selvbestemmelse}, \dabold{eftersom} afgørende beslutninger om sikkerhed, frihandel og grønne produktstandarder alligevel træffes af de store nabolande.

I tre årtier holdt Danmark fast i sine \da{EU-forbehold}, \dabold{herunder} både euro- og forsvarsforbeholdet, ud fra et forståeligt ønske om at beskytte nationale kerneområder. \dabold{Desuagtet} har skiftende danske regeringer måttet følge de fleste europæiske direktiver tæt for ikke at isolere landets eksport og erhvervsliv. Afskaffelsen af forsvarsforbeholdet i 2022 viser tydeligt, \dabold{hvorvidt} udefrakommende kriser kan ændre dybe historiske vaner. Spørgsmålet for fremtiden er derfor ikke længere, om Danmark skal samarbejde med Europa, men \dabold{i hvilken udstrækning} vi selv vil sidde med ved bordet, når de fælles regler skrives.
\end{emailbox}

\textbf{Tekst B 阅读核对（Rigtigt / Forkert / Kort svar，答案见第六部分）：}
\begin{enumerate}
  \item \textit{Ifølge Thomas Lind sikrer det automatisk et lands reelle selvbestemmelse, hvis landet formelt står uden for EU-samarbejdet.} \underline{\hspace{2.2cm}}
  \item \textit{Hvorfor har skiftende danske regeringer alligevel måttet følge de fleste europæiske direktiver trods forbeholdene?} \underline{\hspace{2.2cm}}
  \item \textit{Hvilket afgørende spørgsmål peger Thomas Lind på for Danmarks fremtid?} \underline{\hspace{2.2cm}}
\end{enumerate}

\begin{tipbox}{本课发音与语调要点（Udtale \& Prosodi）}
\begin{itemize}
  \item \textbf{高阶书面连词与代词读音}：\textit{eftersom} \ipa{ˈɛfdɐˌsʌm}（重音在首音节 \textit{ef-}）；\textit{idet} \ipa{iˈdeð} 或弱读 \ipa{iˈdɛ}（注意末尾 \textit{t} 不发硬爆破音，常近于弱化或静音！）；\textit{hvorvidt} \ipa{vɒˈvid}（\textit{h} 与 \textit{d} 均不发音）；\textit{desuagtet} \ipa{dɛsuˈɑɡd̥əð}（重音在 \textit{-ag-}）；\textit{hvormed} \ipa{vɒˈmɛð}（词首 \textit{h} 不发音）。
  \item \textbf{历史政治高频词}：\textit{en suverænitet} \ipa{suvəʁɛniˈtˢeːˀd}（重音在末音节 \textit{-tet}）；\textit{en udstrækning} \ipa{ˈuðˌsdʁɛɡneŋ}（前缀 \textit{ud-} 为软 \textit{d}）；\textit{et forbehold} \ipa{ˈfɒːbeˌhʌlˀ}（后段 \textit{-hold} 带 \textit{stød}，\textit{d} 不发音！）；\textit{et nederlag} \ipa{ˈneːðɐˌlɑʊ̯ˀ}（\textit{-lag} 读滑音 \ipa{lɑʊ̯ˀ}）。
  \item \textbf{长从句政论语调}：遇到 \textit{...eftersom...} 或 \textit{...hvormed...} 引导的后置长从句时，逗号前微升保连贯，切忌在连词后突兀停顿；而在句首 \textit{Desuagtet} 后需紧接动词，一气呵成读出 V2 节奏。
\end{itemize}
\end{tipbox}

\begin{transbox}
\small
\textbf{\color{NCDTeal}[Tekst B 中文对照]}\\
\textbf{学者短评（托马斯·林德，欧洲史副教授｜小国的两难处境）：}在丹麦的欧盟辩论中，主权常常被描绘成一种固定不变的东西——要么原封不动留在国内，要么在布鲁塞尔失去。然而，这种简化的看法忽略了任何一个小国所面临的经典两难：形式上置身事外并不能保证真正的自决权，因为有关安全、自由贸易与绿色产品标准的重大决策终究是由周边大国与邻国做出的。

三十年来，丹麦一直坚持其欧盟保留条款——其中包括欧元与防务保留条款——这是出于一种可以理解的愿望，即保护国家核心领域。尽管如此，历届丹麦政府仍然不得不紧跟大多数欧洲指令，以免孤立本国的出口与工商业界。2022 年废除防务保留条款之举清楚表明，外部危机究竟是否以及如何能够改变根深蒂固的历史习惯。因此，未来的问题不再是丹麦是否应当与欧洲合作，而是在共同规则制定之时，我们在多大程度上愿意亲自坐在谈判桌旁。
\end{transbox}

% ---------------------------------------------------------------------
\ncdsubsec{三、Gloser（核心生词表）}

\begin{tipbox}{词表重点}
本课词汇表囊括了历史政论文与高阶逻辑连接的两大核心集群：政治概念词（如 \textit{suverænitet, selvbestemmelse, forbehold, nederlag}）及七大 C1 书面连接词。特别留意：\textit{et nederlag}（战败/失利）与 \textit{et forbehold}（保留条款）均为单复同形的中性名词；而 \textit{udenrigspolitik} 在加定冠词后缀时须双写 \textit{k}（\textit{udenrigspolitikken}）。
\end{tipbox}

\begin{glosertable}
名词 & en småstat & \textit{(småstaten, småstater, småstaterne)} & 小国，小规模国家 \\
名词 & en småstatsmentalitet & \textit{(småstatsmentaliteten, kun ental)} & 小国心态（谨慎与危机意识） \\
名词 & et udsyn & \textit{(udsynet, udsyn, udsynene)} & 视野，国际格局 \\
名词 & et nederlag & \textit{(nederlaget, nederlag, nederlagene)} & 战败，挫折 \\
名词 & en andelsbevægelse & \textit{(andelsbevægelsen, kun ental)} & 合作社运动（丹麦历史制度） \\
名词 & en udenrigspolitik & \textit{(udenrigspolitikken, -politikker, -erne)} & 外交政策 \\
名词 & et forbehold & \textit{(forbeholdet, forbehold, forbeholdene)} & 保留条款，保留意见 \\
名词 & en suverænitet & \textit{(suveræniteten, kun ental)} & 主权 \\
名词 & en selvbestemmelse & \textit{(selvbestemmelsen, kun ental)} & 自决权，自主决定权 \\
名词 & en udstrækning & \textit{(udstrækningen, udstrækninger, -erne)} & 程度，范围，跨度 \\
\midrule
动词 [har] & at præge & \textit{(præger, prægede, har præget)} & 深刻影响，塑造，打上烙印 \\
动词 [har] & at afgive & \textit{(afgiver, afgav, har afgivet)} & 让渡，交出（权力/主权） \\
动词 [har] & at ophæve & \textit{(ophæver, ophævede, har ophævet)} & 废除，取消（条款/禁令） \\
\midrule
连词 & eftersom & \textit{(从属连词，不变)} & 因为，鉴于（正式因果） \\
连词 & idet & \textit{(从属连词，不变)} & 因为，在于，同时（正式说明） \\
副词 & desuagtet & \textit{(句间副词，不变)} & 尽管如此，即便这样（触发 V2） \\
副词 & herunder & \textit{(副词，不变)} & 其中包括，含以下各项 \\
连词 & hvormed & \textit{(关系副词连词，不变)} & 借此，由此，凭借该举措 \\
\end{glosertable}

% ---------------------------------------------------------------------
\ncdsubsec{四、Grammatik \& Stilistik（高阶书面连词、副词代词与政论时态精讲）}

\begin{grammatikbox}{语法点 1：七大 C1 书面连接词与关系/副词代词全景功能矩阵}
掌握下面这七个高阶书面连接工具，关键在于分清它们的\textbf{句法身份}——究竟是\textbf{从属连词}（引导从句，内部“主语 + 状语 + 动词”不倒装），还是\textbf{句间副词}（位于句首时主句必须发生 V2 倒装）：

\vspace{0.3em}
\renewcommand{\arraystretch}{1.24}
{\small
\begin{tabularx}{\linewidth}{@{} >{\bfseries\raggedright\arraybackslash}p{2.9cm} >{\raggedright\arraybackslash}p{4.6cm} >{\raggedright\arraybackslash}X @{}}
\toprule
\textbf{C1 连接词 / 短语} & \textbf{词性与语序铁律} & \textbf{典型政论示范句} \\
\midrule
\dabold{eftersom}\newline （鉴于，因为） &
\textbf{因果从属连词}：引导客观已知原因；从句内部不倒装，但若放句首，后面主句必倒装。 &
\textit{Mentaliteten prægede landet, \dabold{eftersom} overlevelse krævede sammenhold.} \\
\midrule
\dabold{idet}\newline （由于，在于，同时） &
\textbf{书面说明/同时连词}：常接在主句后补充理据或同步机制；内部绝不倒装。 &
\textit{Landet trådte ind i EF, \dabold{idet} åben handel skabte tryghed.} \\
\midrule
\dabold{hvorvidt}\newline （究竟是否） &
\textbf{间接疑问从属连词}：用于探讨或争论两可议题；内部绝不倒装（禁代以 \textit{hvis}）。 &
\textit{Vælgerne diskuterede, \dabold{hvorvidt} samarbejdet ville udhule demokratiet.} \\
\midrule
\dabold{i hvilken udstrækning}\newline （在多大程度上） &
\textbf{程度疑问短语}：引出定量或界限追问；引导间接问句，内部维持主语 + 动词语序。 &
\textit{Debatten drejer sig om, \dabold{i hvilken udstrækning} man bør afgive suverænitet.} \\
\midrule
\dabold{hvormed}\newline （借此，由此） &
\textbf{关系副词连词}：回指前句整体行动或手段；引导非限定性从句，内部不倒装。 &
\textit{Danmark blev en småstat, \dabold{hvormed} en dyb forsigtighed slog rod.} \\
\midrule
\dabold{desuagtet}\newline （尽管如此） &
\textbf{高阶让步句间副词}（= 更正式的 \textit{alligevel}）：放句首时\textbf{立即触发 V2 倒装}！ &
\textit{\dabold{Desuagtet} har Danmark ofte gennemført fælles regler hurtigt.} \\
\midrule
\dabold{herunder}\newline （其中包括） &
\textbf{归属列举副词}：紧跟名词短语后补充具体子类，不用动词；放句首亦触发 V2。 &
\textit{Eksporten hviler på flere sektorer, \dabold{herunder} vedvarende energi.} \\
\bottomrule
\end{tabularx}
}
\end{grammatikbox}

\begin{grammatikbox}{语法点 2：历史政论文的“跨时态精密切换”与复杂从句语序}
\begin{enumerate}
  \item \textbf{三大时态在历史政论中的精准分工}：\\
    \ding{172} \textbf{一般过去时（Datid）}：定位历史上已发生且结束的节点事件（如 \textit{i 1864 mistede Danmark..., i 1992 sagde vælgerne nej..., i 2022 stemte et flertal...}）。\\
    \ding{173} \textbf{现在完成时（Førnutid）}：连接过去决策与当代持续影响（如 \textit{Danmark har aldrig haft råd til..., I de senere år har krisen tvunget landet til...}）。\\
    \ding{174} \textbf{一般现在时（Nutid）}：阐述当今现实格局、结构性规律与未来追问（如 \textit{Spørgsmålet er..., Erfaringen viser, at...}）。
  \item \textbf{从句内部的否定词位置永远前移}：即使在 \textit{eftersom, idet, hvorvidt, hvormed} 引导的极正式从句中，否定词 \textit{ikke} 或副词 \textit{ofte / sjældent} 依然必须严格站在变位动词前面：\\
    \ding{51}\ \textit{...idet små stater \dabold{sjældent kan} forme deres egen sikkerhed alene.}（\ding{55}\ \textit{*...idet små stater kan sjældent forme...}）
\end{enumerate}
\end{grammatikbox}

\begin{grammatikbox}{高阶连词四大常见失误排雷（Top 4 Pitfalls）}
\begin{enumerate}
  \item \textbf{\ding{172} 误把 \textit{hvormed / idet} 当成 \textit{derfor} 并实施倒装}：\ding{55}\ \textit{*Landet mistede territorium, hvormed opstod en ny mentalitet.} $\rightarrow$ \ding{51}\ \textit{Landet mistede territorium, \dabold{hvormed} en ny mentalitet opstod.}（记住：\textit{hvormed, hvorved, idet, eftersom} 全是从属连词，从句内主语先于动词！）
  \item \textbf{\ding{173} 句首 \textit{Desuagtet} 忘记 V2 倒装或乱加逗号}：\ding{55}\ \textit{*Desuagtet, regeringen fulgte direktiverne.} $\rightarrow$ \ding{51}\ \textit{\dabold{Desuagtet} fulgte regeringen direktiverne.}
  \item \textbf{\ding{174} 混淆 \textit{i hvilken udstrækning}（在多大程度上）中的格与语序}：由于 \textit{udstrækning} 为共性单数名词（\textit{en udstrækning}），疑问限定词必须用 \dabold{hvilken}（\ding{55}\ \textit{*i hvilket udstrækning}），且后面从句不能倒装：\ding{51}\ \textit{i \dabold{hvilken} udstrækning vi vil sidde med ved bordet}。
  \item \textbf{\ding{175} 在 \textit{herunder} 后面画蛇添足加动词或 \textit{at}}：\textit{herunder} 用于直接在逗号后列举同位名词成分（\textit{flere forbehold, herunder euroforbeholdet}），不能接完整句子除非放在句首并倒装（\textit{Herunder spiller...}）。
\end{enumerate}
\end{grammatikbox}

% ---------------------------------------------------------------------
\ncdsubsec{五、Mønsterøvelser（政论高阶句法操练）}

\textbf{练习 A：C1 书面连词与历史政治核心词精准填空（每空填入最准确的 1 个丹麦语词）}
\begin{enumerate}
  \item Danmark blev i 1864 en sårbar småstat, \underline{\hspace{2.2cm}} en dyb forsigtighed slog rod i kulturen.
  \item Sammenholdet blev livsvigtigt, \underline{\hspace{2.2cm}} nationens overlevelse afhang af indre styrke.
  \item Landet valgte at deltage i det europæiske marked, \underline{\hspace{2.2cm}} åben handel sikrede velstanden.
  \item Eksporten hviler på flere stærke brancher, \underline{\hspace{2.2cm}} søfart og grøn energiteknologi.
  \item Vælgerne var uenige om, \underline{\hspace{2.2cm}} et tættere EU-samarbejde ville svække demokratiet.
  \item Debatten handler nu om, i hvilken \underline{\hspace{2.2cm}} et lille land bør dele suverænitet med andre.
  \item Befolkningen stemte nej i 1992; \underline{\hspace{2.2cm}} har Danmark ofte indført EU-regler meget hurtigt.
  \item Et flertal af vælgerne valgte i 2022 at \underline{\hspace{2.2cm}} det danske forsvarsforbehold.
\end{enumerate}

\newpage
\textbf{练习 B：高阶书面连接句纠错与升级（修正从句语序、V2 或用词错误，或按提示整合升级）}
\begin{enumerate}
  \item 【从句语序纠错】\textit{Småstater må ofte indgå alliancer, idet kan de sjældent klare sikkerheden alene.} $\rightarrow$
  \item 【句首倒装纠错】\textit{Desuagtet, Danmark har længe været en loyal partner i det europæiske miljøsamarbejde.} $\rightarrow$
  \item 【词形与语序纠错】\textit{Vi må undersøge, i hvilket udstrækning påvirker europæiske regler dansk erhvervsliv.} $\rightarrow$
  \item 【因果升级】用 \textit{eftersom} 合并两句：\textit{Danmark har få naturlige råstoffer. Landet er dybt afhængigt af global handel.} $\rightarrow$
  \item 【手段结果升级】用 \textit{, hvormed} 合并两句：\textit{Vælgerne fjernede forsvarsforbeholdet i 2022. Dermed fik Danmark en ny rolle i Europa.} $\rightarrow$
  \item 【归属列举升级】用 \textit{, herunder} 浓缩句子：\textit{Politikerne diskuterede flere danske forbehold, og de talte blandt andet om euroen og retsområdet.} $\rightarrow$
\end{enumerate}

\textbf{练习 C：中译丹（灵活调用本课高阶书面连词与历史政治词汇）}
\begin{enumerate}
  \item 丹麦依赖开放贸易，因为这个小国几乎没有自然资源。
  \item 选民们争论过，这是否会削弱国家的民主自决权。
  \item 关键问题在于，我们在多大程度上愿意让渡主权。
  \item 历史带来了疑虑，尽管如此，丹麦在实践中仍然忠实遵守欧洲规则。
  \item 该国出口多种商品，其中包括绿色能源和医药产品。
  \item 政府通过了新协议，从而加强了自身的国际影响力。
\end{enumerate}

% ---------------------------------------------------------------------
\ncdsubsec{六、PD3 / Studieprøven 实战微任务与参考答案}

\begin{emailbox}{Studieprøven / PD3 历史政论写作任务（220--260 ord）}
\textbf{Opgave}: Danmark beskrives ofte som et land, der både værner om sin nationale selvstændighed og samtidig er dybt afhængigt af omverdenen. Skriv en \textbf{sammenhængende historisk-politisk kommentar på 220--260 ord} om spændingsforholdet mellem småstatsmentalitet og globalt udsyn. Brug mindst \textbf{fem} af de syv avancerede forbindere fra lektionen (\textit{eftersom, idet, hvorvidt, desuagtet, herunder, hvormed, i hvilken udstrækning}) med korrekt ordstilling.
\end{emailbox}

\begin{answerbox}
\textbf{Før du læser 参考回答}
\begin{enumerate}
  \item \textit{Nederlaget i 1864 gjorde Danmark til en sårbar småstat, hvormed en dyb forsigtighed over for stormagter slog rod. Alligevel kan småstatsmentalitet og åbenhed eksistere side om side, fordi et lille land uden råstoffer har brug for fri handel og multilaterale alliancer for at sikre både velstand og sikkerhed.}
  \item \textit{Thomas Lind peger på det paradoks, at formelt at stå uden for EU-samarbejdet ikke sikrer reel selvbestemmelse, eftersom store beslutninger om handel, sikkerhed og grønne standarder alligevel træffes af nabolandene, så Danmark i praksis må følge direktiverne trods forbeholdene.}
\end{enumerate}

\textbf{Tekst B 阅读核对答案}
\begin{enumerate}
  \item \da{Forkert}（他指出形式上置身事外并不能保证真正的自决权，因为邻国仍会做出关键决定）
  \item \da{For ikke at isolere dansk erhvervsliv.}（为了不孤立丹麦的工商业界）
  \item \da{I hvilken udstrækning Danmark vil sidde med ved bordet, når reglerne skrives.}
\end{enumerate}

\textbf{练习 A 答案}
\begin{enumerate}
  \item \da{hvormed}（关系副词连词：由此，借此，回指前句变成小国的事实）
  \item \da{eftersom}（正式因果连词：因为，鉴于）
  \item \da{idet}（书面解释说明连词：因为，在于）
  \item \da{herunder}（副词：其中包括，后直接接名词短语）
  \item \da{hvorvidt}（正式间接疑问连词：究竟是否）
  \item \da{udstrækning}（固定搭配：\textit{i hvilken udstrækning} 在多大程度上）
  \item \da{desuagtet}（句间让步副词：尽管如此，分号后放句首触发 V2 倒装 \textit{har Danmark}）
  \item \da{ophæve}（动词原形：废除，取消，搭配保留条款）
\end{enumerate}

\textbf{练习 B 答案（参考）}
\begin{enumerate}
  \item \textit{Småstater må ofte indgå alliancer, idet de sjældent kan klare sikkerheden alene.}（\textit{idet} 引导从句，主语 \textit{de} 在前，副词 \textit{sjældent} 站在动词 \textit{kan} 前）
  \item \textit{Desuagtet har Danmark længe været en loyal partner i det europæiske miljøsamarbejde.}（去掉逗号，\textit{Desuagtet} 置句首立即触发 V2 倒装 \textit{har Danmark}）
  \item \textit{Vi må undersøge, i hvilken udstrækning europæiske regler påvirker dansk erhvervsliv.}（共性名词用 \textit{hvilken}，且间接疑问从句内部维持主语 + 动词不倒装）
  \item \textit{Danmark er dybt afhængigt af global handel, eftersom landet har få naturlige råstoffer.}（用正式因果连词 \textit{eftersom} 整合因果逻辑）
  \item \textit{Vælgerne fjernede forsvarsforbeholdet i 2022, hvormed Danmark fik en ny rolle i Europa.}（用 \textit{, hvormed} 顺滑承接结果，从句不倒装）
  \item \textit{Politikerne diskuterede flere danske forbehold, herunder euroen og retsområdet.}（用 \textit{herunder} 精炼替换冗长并列句）
\end{enumerate}

\textbf{练习 C 答案}
\begin{enumerate}
  \item \textit{Danmark er afhængigt af åben handel, eftersom det lille land næsten ingen råstoffer har.}
  \item \textit{Vælgerne diskuterede, hvorvidt dette ville svække landets demokratiske selvbestemmelse.}
  \item \textit{Det centrale spørgsmål er, i hvilken udstrækning vi vil afgive suverænitet.}
  \item \textit{Historien skabte skepsis; desuagtet følger Danmark i praksis loyalt de europæiske regler.}
  \item \textit{Landet eksporterer mange varer, herunder grøn energi og medicinske produkter.}
  \item \textit{Regeringen vedtog en ny aftale, hvormed den styrkede sin internationale indflydelse.}
\end{enumerate}

\textbf{PD3 / Studieprøven 高分范文（约 235 词）}\\
\textit{\textbf{Mellem forsigtighed og globalt ansvar: Småstatens europæiske balance}}

\textit{Gennem mere end halvandet århundrede har Danmarks forhold til Europa været præget af et dobbelt hensyn. Efter nederlaget i 1864 opstod en dyb \da{småstatsmentalitet}, \dabold{hvormed} forsigtighed over for stormagter og værn om national sammenhængskraft blev politiske kerneværdier. Landet lærte at opbygge stærke indre fællesskaber gennem folkehøjskoler og andelsbevægelse, \dabold{eftersom} en sårbar nation uden stor befolkning ikke kunne forlade sig på militær styrke.}

\textit{Alligevel har denne naturlige forsigtighed aldrig ført til varig økonomisk isolation. Tværtimod har dansk velstand siden efterkrigstiden hvilet på international åbenhed og frihandel, \dabold{idet} et lille hjemmemarked uden store råstoffer ikke kan bære en moderne velfærdsøkonomi alene. I dag bygger landets globale profil på specialiserede sektorer, \dabold{herunder} søfart, medicinalindustri og grøn energiteknologi, som alle forudsætter forudsigelige og fælles internationale spilleregler.}

\textit{Netop dette krydspres mellem national tryghed og global afhængighed forklarer den mangeårige debat om europæisk integration. Siden nej-stemmen i 1992 har danske vælgere gentagne gange spurgt, \dabold{hvorvidt} tættere samarbejde truer den lokale selvbestemmelse, og \dabold{i hvilken udstrækning} man bør \da{afgive} formel \da{suverænitet} til Bruxelles. \dabold{Desuagtet} viser nyere historie, at Danmark i praksis sjældent står udenfor, når afgørende europæiske standarder fastlægges. Ophævelsen af forsvarsforbeholdet i 2022 understregede yderligere, at reel sikkerhed og demokratisk indflydelse for et lille land ikke skabes i ensomhed.}

\textit{Fremtidens vigtigste opgave for dansk \da{udenrigspolitik} er derfor ikke at vælge mellem nære rødder og \da{globalt udsyn}, men at forene begge dele gennem aktivt, pragmatisk og konstruktivt europæisk engagement.}

\small\textbf{名师点评：}这篇约 235 词的范文堪称历史政论文的语法教科书示范：全部七大目标书面连接工具（\textit{hvormed, eftersom, idet, herunder, hvorvidt, i hvilken udstrækning, desuagtet}）尽数自然嵌入论证肌理，从句内语序与句首倒装无一疏漏；同时精准示范了 1864 史实过去时、持续影响的现在完成时与当下结论的现在时三层切换。
\end{answerbox}

\ncdlesson{15}{Det gode liv: Lykkeforskning og det skandinaviske paradoks}{美好生活：幸福感研究与斯堪的纳维亚悖论}
{DU3 Modul 6 (Studieprøven, C1) 核心写作课｜对标 NCE 4 \& \textit{Skriftlig Fremstilling: Syntese af To Kilder}}
{\textbf{多源观点综合论述（\textit{Syntese}）}（交叉对比 \textit{Hvor Kilde A fremhæver..., anlægger Kilde B et markant mere kritisk perspektiv...} · 共识归并 \textit{På trods af denne uenighed deler de to tekster den opfattelse, at...} · 态度与方法论评判）}
{在丹麦最高级别语言考试 Studieprøven（C1）的写作核心题型 \textbf{\textit{Syntese}} 中，考生绝不能把两份背景材料机械地“先摘要 A、再摘要 B”，而必须围绕共同主题把截然相反的证据交织对比、提炼隐性共识并做出独立评判。本课以“丹麦幸福感研究”与“斯堪的纳维亚悖论”为双源语料，手把手带你写出一篇真正达到 C1 学术水准的综合论述长文！}

% ---------------------------------------------------------------------
\begin{foerduleaeser}
带着下面两个宏观问题对比阅读 Kilde A（Tekst A）与 Kilde B（Tekst B）：
\begin{enumerate}
  \item \textit{Hvordan forklarer Institut for Nordisk Lykkeforskning i Tekst A, at Danmark konsekvent ligger højt i internationale målinger af livstilfredshed?}（根据 Tekst A，北欧幸福感研究所如何解释丹麦在国际生活满意度测评中持续名列前茅？）
  \item \textit{Hvilke modstridende tendenser og kulturelle skyggesider fremhæver samfundskritikeren i Tekst B som tegn på ”det skandinaviske paradoks”?}（在 Tekst B 中，社会评论家指出哪些相互矛盾的趋势与文化阴暗面，以印证所谓的“斯堪的纳维亚悖论”？）
\end{enumerate}
\end{foerduleaeser}

% ---------------------------------------------------------------------
\ncdsubsec{二、Tekster（双源背景材料：机构研究报告 vs 社会批判评论）}

\subsection*{Tekst A：\textit{Kilde A: Tillid og velfærd som fundament for dansk livstilfredshed}（幸福研究机构背景材料 · 约 335 词）}

\begin{tipbox}{语篇说明}
Tekst A 模拟丹麦《幸福研究所》风格的制度与实证研究文本，聚焦“结构性安全感与社会信任”。阅读时请注意\textbf{\color{NCDNavy}深蓝色标记的幸福感研究与福利制度核心术语}，以及\textbf{\color{NCDRed}深红色标记的因果归因与科学界定表达}，为下文展开 \textit{Syntese} 对比做好标记。
\end{tipbox}

År efter år placerer Danmark sig helt i toppen af internationale undersøgelser som FN's World Happiness Report, hvor borgere på tværs af mere end hundrede lande bedømmer deres generelle tilværelse på en ti-trins skala. Når man inden for moderne \da{lykkeforskning} forsøger at forklare dette bemærkelsesværdige mønster, skelner forskerne skarpt mellem flygtig begejstring i øjeblikket og varig \da{livstilfredshed}. Det handler med andre ord sjældent om konstant eufori eller godt humør, men om hvorvidt borgerne oplever personlig frihed, forudsigelig stabilitet og reel kontrol over deres eget liv på længere sigt.

Ifølge Institut for Nordisk Lykkeforskning hviler de høje danske resultater først og fremmest på en usædvanlig høj grad af \da{institutionel tillid} kombineret med stærk mellemmenneskelig tillid. Danskerne stoler som udgangspunkt ikke kun på familie og venner, men også på fremmede medborgere, domstole og offentlige myndigheder, \dabold{hvilket} skaber lave transaktionsomkostninger og markant mindre hverdagsskepsis i samfundet. Når korruptionen er minimal, og det progressive skattesystem opleves som forholdsvis retfærdigt, vokser følelsen af, at samfundets fælles institutioner arbejder for almenvellet frem for snævre særinteresser.

Derudover spiller det udbyggede velfærdssamfund en helt afgørende rolle som socialt \da{sikkerhedsnet}. Gratis adgang til videregående uddannelse, offentlig sundhedspleje og økonomisk støtte i form af dagpenge ved arbejdsløshed fjerner meget af den eksistentielle utryghed, der i mange andre lande tynger almindelige familier gennem hele livet. Man behøver ikke frygte social deroute eller økonomisk ruin ved pludselig sygdom eller jobskifte, \dabold{og netop denne grundlæggende tryghed frigør} mental energi til videreuddannelse, frivilligt foreningsliv og langsigtet personlig udfoldelse.

Endelig \da{fremhæver} rapporten den danske balance mellem arbejde og fritid som en central drivkraft bag den høje tilfredshed. Kortere gennemsnitlige arbejdsuger, fleksible arbejdstider, generøse barselsordninger og en flad ledelseskultur giver reel plads til både nært familieliv og sociale fællesskaber uden for arbejdspladsen. Ganske vist peger forskerne selv på, at Danmark fortsat har alvorlige udfordringer med arbejdsrelateret stress og psykisk sårbarhed i visse befolkningsgrupper. Alligevel \dabold{konkluderer rapporten, at} kombinationen af økonomisk tryghed, lighed og social tillid danner et ualmindeligt robust fundament for et godt og meningsfuldt liv.

\begin{transbox}
\small
\textbf{\color{NCDTeal}[Tekst A 中文精译]}\\
年复一年，丹麦在联合国《世界幸福报告》等国际调查中始终稳居前列——在这类调查中，超过一百个国家的公民会在十分制量表上评价自己的整体生活状态。当现代幸福感研究试图解释这一引人注目的规律时，学者们严格区分了“转瞬即逝的兴奋”与“持久的生活满意度”。换言之，问题很少在于人是否时刻保持狂喜或好心情，而在于公民从长远来看是否感受到个人自由、可预期的稳定以及对自身生活的真正掌控感。

据北欧幸福感研究所指出，丹麦的优异成绩首先建立在极高程度的制度信任以及强韧的人际信任之上。丹麦人原则上不仅信任亲友，也信任陌生同胞、司法机构与公共行政部门，这在社会中降低了交易成本，并显著减少了日常猜忌。当腐败微乎其微、累进税制也被视为相对公正时，人们便愈发感到社会的公共机构是在为整体福祉而非狭隘的特殊利益运转。

此外，完善的福利社会作为一张社会安全网发挥着极为关键的作用。免费接受高等教育、公共医疗保障以及失业时的津贴救济，消除了在许多其他国家终生压在普通家庭头上的诸多生存焦虑。人们无需担心因突发疾病或更换工作而陷入社会滑坡或经济崩溃，而正是这种基本的安稳感释放了心理能量，使人得以投入继续教育、志愿社团生活以及长期的个人自我实现。

最后，报告强调丹麦在工作与闲暇之间的平衡是推动高满意度的重要动力。较短的平均工作周、灵活的工作时间、优厚的育儿假制度以及扁平化的管理文化，为工作场所之外的亲密家庭生活与社交共同体留出了切实空间。诚然，研究者自身也指出丹麦在特定人群中依然面临与工作相关的压力以及心理脆弱等严峻挑战。尽管如此，报告依然得出结论：经济安全感、平等与社会信任的结合，为美好而有意义的生活构筑了格外坚实的基础。
\end{transbox}

\subsection*{Tekst B：\textit{Kilde B: Lykkepresset --- ensomhed, Jantelov og medicinforbrug bag facaden}（社会批判对比材料 · 约 175 词）}

\textbf{阅读任务（理解与核对）：}\textit{Læs kultursociologen Henrik Madsens kritiske kommentar om det skandinaviske paradoks, og besvar spørgsmålene nedenfor.}

\begin{emailbox}{Kilde B｜Uddrag af kronik af Henrik Madsen (kultursociolog)}
Bag de rosenrøde internationale målinger gemmer der sig imidlertid et foruroligende paradoks. Når jeg betragter det moderne danske velfærdssamfund tæt på, ser jeg nemlig en dyb kløft mellem den officielle statistik og mange borgeres faktiske virkelighed. I et land, der konstant kåres som et af verdens lykkeligste, oplever stadig flere unge og enlige husstande en voksende \da{ensomhed}, samtidig med at \da{forbruget} af \da{antidepressiv} medicin og antallet af langvarige stresssygemeldinger hører til blandt de højeste i Europa.

Denne kontrast skyldes ikke blot tilfældigheder, men et usynligt kulturelt \da{pres}. På den ene side skaber den stærke norm om lighed og Jantelovens ”du skal ikke tro, du er noget” en høj grad af \da{konformitet}, hvor det kan føles tabubelagt at skille sig ud eller udtrykke dyb utilfredshed. På den anden side medfører fortællingen om det perfekte velfærdsliv et massivt \da{lykkepres}: Når alle ydre rammer tilsyneladende er på plads, bliver personlig mistrivsel let oplevet som en skamfuld individuel fiasko. Derfor bør man \da{anfægte}, om spørgeskemaer om ”lykke” i virkeligheden snarere måler taknemmelig pligtfølelse end ægte menneskelig \da{trivsel}.
\end{emailbox}

\textbf{Tekst B 阅读核对（Rigtigt / Forkert / Kort svar，答案见第六部分）：}
\begin{enumerate}
  \item \textit{Henrik Madsen mener, at de internationale lykkemålinger giver et fuldstændigt billede af danskernes faktiske trivsel.} \underline{\hspace{2.2cm}}
  \item \textit{Hvilke to konkrete tegn på mistrivsel bag statistikken nævner Henrik Madsen i første afsnit?} \underline{\hspace{2.2cm}}
  \item \textit{Hvorfor oplever borgere ifølge Henrik Madsen ofte personlig mistrivsel som en individuel fiasko i et velfærdssamfund?} \underline{\hspace{2.2cm}}
\end{enumerate}

\begin{tipbox}{本课发音与语调要点（Udtale \& Prosodi）}
\begin{itemize}
  \item \textbf{研究与心理核心长词}：\textit{en livstilfredshed} \ipa{ˈliʊ̯sˌtˢefʁɛsˌheðˀ}（\textit{tilfreds} 中的 \textit{d} 不发音！）；\textit{eksistentiel} \ipa{ɛɡsisdɛnˈɕɛlˀ}（重音在尾音节，带 \textit{stød}）；\textit{antidepressiv} \ipa{ˈantˢiˌdebʁəˌsiʊ̯ˀ}（词尾 \textit{v} 滑音化为 \ipa{ʊ̯ˀ}）；\textit{en trivsel} \ipa{ˈtˢʁiʊ̯ˀsəl}（\textit{v} 滑音带 \textit{stød}）。
  \item \textbf{学术综合分析高频词}：\textit{en syntese} \ipa{synˈtˢeːsə}（重音在第二音节长元音 \textit{-te-}）；\textit{et paradoks} \ipa{pʰɑʁɑˈdʌɡs}（重音在末音节）；\textit{at anfægte} \ipa{ˈanˌfɛɡd̥ə}（前缀重读）；\textit{en uenighed} \ipa{uˈeːniˌheðˀ}（\textit{-ig-} 中 \textit{g} 不发音）。
  \item \textbf{\textit{Syntese} 交叉对比朗读节奏}：朗读 \textit{Hvor Kilde A fremhæver..., anlægger Kilde B et markant mere kritisk perspektiv...} 时，两个分句主语 \textit{Kilde \textbf{A}} 与 \textit{Kilde \textbf{B}} 必须给予清晰的对比重音，展现辩证比较的逻辑张力。
\end{itemize}
\end{tipbox}

\begin{transbox}
\small
\textbf{\color{NCDTeal}[Tekst B 中文对照]}\\
\textbf{文化社会学家亨里克·马森专栏节选（外表光鲜背后的幸福压力）：}然而，在这些玫瑰色的国际测评背后，潜藏着一个令人不安的悖论。当我近距离审视现代丹麦福利社会时，我看到了官方统计数据与许多公民实际现实之间的巨大鸿沟。在这个屡屡被评为世界上最幸福的国家之一的地方，越来越多的年轻人和单身家庭正经历着日益加剧的孤独感，与此同时，抗抑郁药物的消耗量与长期因压力请病假的人数却位居欧洲最高行列。

这种反差并非纯粹出于巧合，而是源于一种无形的文化压力。一方面，“绝对平等”的强大规范以及《詹代法则》中“不要以为你有多了不起”的准则营造出高度的从众氛围，使人感到标新立异或表露深层不满仿佛成了一种禁忌。另一方面，关于“完美福利生活”的叙事带来了沉重的幸福压力：既然一切外部条件看似都已齐备，个人的不顺心便容易被视为一种令人羞愧的自身失败。因此，我们应当质疑：那些关于“幸福”的问卷调查，测量的究竟是心怀感恩的义务感，还是真正的人类身心福祉。
\end{transbox}

% ---------------------------------------------------------------------
\ncdsubsec{三、Gloser（核心生词表）}

\begin{tipbox}{词表重点}
\textit{Syntese} 写作极为依赖“表达观点立场与对比”的精密动词搭配：表示“采取某种视角”固定用 \textit{at anlægge et perspektiv}；表示“质疑某种前提或结论”用 \textit{at anfægte}；表示“把两份材料放在一起比对”用 \textit{at sammenholde A med B}。注意中性名词 \textit{et sikkerhedsnet}（安全网）与 \textit{et pres}（压力）在定指形式中分别须双写 \textit{t}（\textit{sikkerhedsnettet}）和 \textit{s}（\textit{presset}）。
\end{tipbox}

\begin{glosertable}
名词 & en lykkeforskning & \textit{(lykkeforskningen, kun ental)} & 幸福感研究 \\
名词 & en livstilfredshed & \textit{(livstilfredsheden, kun ental)} & 生活满意度 \\
名词 & et sikkerhedsnet & \textit{(sikkerhedsnettet, sikkerhedsnet, -tene)} & 社会安全网 \\
名词 & en ensomhed & \textit{(ensomheden, kun ental)} & 孤独，孤寂感 \\
名词 & et paradoks & \textit{(paradokset, paradokser, paradokserne)} & 悖论，自相矛盾之处 \\
名词 & en facade & \textit{(facaden, facader, facaderne)} & 门面，光鲜表象 \\
名词 & et pres & \textit{(presset, kun ental)} & 压力，迫力 \\
名词 & en trivsel & \textit{(trivslen, kun ental)} & 身心福祉，良好状态 \\
名词 & en uenighed & \textit{(uenigheden, uenigheder, uenighederne)} & 分歧，意见不一致 \\
名词 & en syntese & \textit{(syntesen, synteser, synteserne)} & 多源综合论述，综合 \\
名词 & et forbrug & \textit{(forbruget, kun ental)} & 消耗量，消费 \\
\midrule
动词 [har] & at fremhæve & \textit{(fremhæver, fremhævede, har fremhævet)} & 强调，突出指出 \\
动词 [har] & at anlægge & \textit{(anlægger, anlagde, har anlagt)} & 采取，持有（某种视角） \\
动词 [har] & at anfægte & \textit{(anfægter, anfægtede, har anfægtet)} & 质疑，挑战（某观点/数据） \\
动词 [har] & at sammenholde & \textit{(sammenholder, sammenholdt, har sammenholdt)} & 对比，对照检验 \\
动词 [har] & at nedtone & \textit{(nedtoner, nedtonede, har nedtonet)} & 淡化，低调处理 \\
\midrule
形容词 & markant & \textit{(markant, markante)} & 显著的，鲜明的 \\
形容词 & antidepressiv & \textit{(antidepressivt, antidepressive)} & 抗抑郁的 \\
\end{glosertable}

% ---------------------------------------------------------------------
\ncdsubsec{四、Grammatik \& Stilistik（Studieprøven C1 多源综合论述 \textit{Syntese} 精讲）}

\begin{grammatikbox}{C1 核心考型突破：多源综合论述（\textit{Syntese}）的四大黄金句式库}
在 Studieprøven 书面写作中，\textbf{绝对不能}先用整段复述 Kilde A、再用整段复述 Kilde B（这被称为“平行摘要”，得分无法上 C1）。真正的 \textit{Syntese} 必须按\textbf{核心议题}分段，并在每一段内将两篇文本交叉对话。请牢记以下四大句式引擎：

\vspace{0.3em}
\renewcommand{\arraystretch}{1.24}
{\small
\begin{tabularx}{\linewidth}{@{} >{\bfseries\raggedright\arraybackslash}p{3.0cm} >{\raggedright\arraybackslash}p{5.5cm} >{\raggedright\arraybackslash}X @{}}
\toprule
\textbf{Syntese 逻辑功能} & \textbf{核心高分句型公式} & \textbf{紧扣本课双源示范句} \\
\midrule
\da{1. 交叉对比立场}\newline （核心高频武器） &
\textit{\dabold{Hvor} Kilde A fremhæver..., \dabold{anlægger} Kilde B et markant mere kritisk perspektiv, idet...} &
\textit{Hvor Kilde A fremhæver tryghed, anlægger Kilde B et markant mere kritisk perspektiv.} \\
\midrule
\da{2. 针锋相对与方法质疑} &
\textit{I modsætning til Kilde A's fokus på..., \dabold{anfægter} Kilde B denne tolkning ved at...}\newline
\textit{Mens Kilde A forklarer... med..., peger Kilde B på, at...} &
\textit{I modsætning til Kilde A's fokus på tillid anfægter Kilde B selve målemetoden.} \\
\midrule
\dabold{3. 提炼隐性共识}\newline （C1 区分度关键！） &
\textit{\dabold{På trods af denne uenighed deler} de to tekster den opfattelse, at...}\newline
\textit{Begge kilder synes at være enige om, at..., om end de vurderer årsagerne forskelligt.} &
\textit{På trods af denne uenighed deler de to tekster den opfattelse, at velfærdssystemet ikke fjerner al sårbarhed.} \\
\midrule
\da{4. 综合权衡收束} &
\textit{\dabold{Sammenholdes} de to perspektiver, tegner der sig et mere nuanceret billede, hvor...} &
\textit{Sammenholdes de to kilder, tegner der sig et komplekst billede af dansk trivsel.} \\
\bottomrule
\end{tabularx}
}
\end{grammatikbox}

\begin{grammatikbox}{语法点 2：对比连词 \textit{Hvor...} 与让步短语 \textit{På trods af...} 的语序铁律}
\begin{enumerate}
  \item \textbf{句首对比从句 \dabold{Hvor...} 触发主句 V2 倒装}：\\
    在英语里“Whereas Source A stresses..., Source B takes...”后面不倒装，但丹麦语把从属连词 \da{Hvor} 引导的整条对比从句放在句首充当第 1 成分时，\textbf{主句必须立刻以变位动词开头}！\\
    \ding{51}\ \textit{\da{Hvor} Kilde A fremhæver de ydre rammer, \dabold{fokuserer} Kilde B på det indre pres.}（主句动词 \textit{fokuserer} 紧贴逗号后！）\\
    \ding{55}\ \textit{*Hvor Kilde A fremhæver de ydre rammer, Kilde B fokuserer på det indre pres.}
  \item \textbf{介词让步短语 \dabold{På trods af...} 触发主句 V2 倒装，且后只能接名词短语}：\\
    \ding{51}\ \textit{\da{På trods af} denne uenighed \dabold{deler} de to tekster en vigtig indsigt.}（若要接完整句子，须改用 \textit{På trods af \dabold{at} de er uenige, \dabold{deler}...}！）
  \item \textbf{\textit{Sammenholde} 的条件被动倒装开篇}：\\
    \textit{\dabold{Sammenholdes} de to tekster, bliver det tydeligt, at...}（相当于省略 \textit{Hvis} 的条件句：\textit{Hvis de to tekster sammenholdes, bliver det tydeligt...}，主句同样发生倒装！）。
\end{enumerate}
\end{grammatikbox}

\begin{grammatikbox}{Syntese 考场四大失误排雷（Top 4 Pitfalls）}
\begin{enumerate}
  \item \textbf{\ding{172} 陷入“两段式流水账”}：只分别写“Tekst A 说…… Tekst B 说……”而没有在同段内用 \textit{Hvor...} 或 \textit{I modsætning til...} 建立直接对话。
  \item \textbf{\ding{173} 遗漏隐性共识}：高分 \textit{Syntese} 不仅要找“吵架点”，更要像侦探一样发现两者潜在共识（如本题中 Kilde A 也承认了压力与心理脆弱的存在）。
  \item \textbf{\ding{174} 搭配错误}：汉语说“采取批判视角”，丹麦语固定是 \textit{at \da{anlægge} et kritisk perspektiv}，绝不能写成 \ding{55}\ \textit{*at lave/tage et kritisk perspektiv}。
  \item \textbf{\ding{175} 混淆 \textit{enighed} 与 \textit{uenighed} 后面的介词}：表示“就某事达成一致/有分歧”，固定用介词 \dabold{om}：\ding{51}\ \textit{uenighed \da{om} årsagerne} / \textit{være enige \da{om}, at...}（非 \textit{på}）。
\end{enumerate}
\end{grammatikbox}

% ---------------------------------------------------------------------
\ncdsubsec{五、Mønsterøvelser（多源综合句型操练）}

\textbf{练习 A：Syntese 核心句式与对比词精准填空（每空填入最准确的 1 个丹麦语词）}
\begin{enumerate}
  \item \underline{\hspace{2.2cm}} Kilde A fremhæver trygheden i velfærdsstaten, fokuserer Kilde B på det skjulte lykkepres.
  \item Kultursociologen \underline{\hspace{2.2cm}} et markant mere kritisk perspektiv på de internationale målinger.
  \item På trods af denne \underline{\hspace{2.2cm}} deler de to tekster den opfattelse, at psykisk sårbarhed er en reel udfordring.
  \item I \underline{\hspace{2.2cm}} til rapportens fokus på tillid peger Henrik Madsen på stigende ensomhed blandt unge.
  \item \underline{\hspace{2.2cm}} de to perspektiver, tegner der sig et langt mere nuanceret billede af danskernes liv.
  \item Kilde B vælger at \underline{\hspace{2.2cm}}, hvorvidt spørgeskemaer overhovedet kan måle ægte menneskelig trivsel.
  \item Begge kilder er enige om, at materielle rammer spiller en rolle, \underline{\hspace{2.2cm}} end de vægter faktorerne forskelligt.
  \item Rapporten forsøger ikke at \underline{\hspace{2.2cm}}, at stress fortsat rammer bestemte grupper i samfundet.
\end{enumerate}

\textbf{练习 B：Syntese 语序纠错与双源句型合并手术（修正错误，或将 A、B 两方观点合成为一句 C1 综合对比句）}
\begin{enumerate}
  \item 【V2 倒装纠错】\textit{Hvor Kilde A understreger institutionel tillid, Kilde B peger på kulturel konformitet og Jantelov.} $\rightarrow$
  \item 【搭配与倒装纠错】\textit{På trods af deres uenighed om målingerne, de to forfattere tager et alvorligt blik på stress.} $\rightarrow$
  \item 【交叉对比合成】用 \textit{Hvor Kilde A fremhæver..., anlægger Kilde B et markant mere kritisk perspektiv, idet...} 合并：Kilde A 强调免费医疗与教育消除焦虑；Kilde B 指出完美条件反而让失败被视为个人过错。$\rightarrow$
  \item 【反驳质疑合成】用 \textit{I modsætning til Kilde A's tolkning anfægter Kilde B..., idet...} 合并：Kilde A 认为高分反映持久生活满意度；Kilde B 认为问卷测出的只是感恩的义务感。$\rightarrow$
  \item 【隐性共识提炼】用 \textit{På trods af denne uenighed deler de to tekster den opfattelse, at...} 合成两篇均承认“部分群体仍面临压力与心理脆弱”。$\rightarrow$
  \item 【综合权衡收束】用 \textit{Sammenholdes de to kilder, bliver det tydeligt, at...} 归结物质保障无法自动消除孤独感。$\rightarrow$
\end{enumerate}

\textbf{练习 C：中译丹（操练 Syntese 交叉对比与幸福悖论词汇）}
\begin{enumerate}
  \item A 源材料强调高水平的制度信任，而 B 源材料采取更加批判的视角。
  \item 尽管存在这种分歧，两篇文本都认为心理脆弱是一个真实问题。
  \item 与 A 源材料相反，评论家质疑调查问卷是否能真正衡量幸福。
  \item 社会安全网消除了许多焦虑，却无法自动消除孤独感。
  \item 若将两份材料对照来看，便会浮现出一幅更加细腻的图景。
  \item 完美的福利表象背后往往隐藏着无形的文化从众压力。
\end{enumerate}

% ---------------------------------------------------------------------
\ncdsubsec{六、PD3 / Studieprøven 实战微任务与参考答案}

\begin{emailbox}{Studieprøven C1 Syntese 核心写作任务（250--300 ord）}
\textbf{Opgave}: Skriv en \textbf{sammenhængende akademisk analyse på 250--300 ord}, hvor du syntetiserer de to kilder (Tekst A og Tekst B) om dansk lykkeforskning og det skandinaviske paradoks. Du må ikke blot referere teksterne hver for sig, men skal strukturere din tekst tematisk omkring både \textbf{de modsatrettede forklaringer} og \textbf{de fælles berøringsflader}, samt afslutningsvis vurdere, hvad kilderne tilsammen fortæller om det gode liv i Danmark.
\end{emailbox}

\begin{answerbox}
\textbf{Før du læser 参考回答}
\begin{enumerate}
  \item \textit{Instituttet forklarer Danmarks høje placering med en usædvanlig høj institutionel tillid, et stærkt velfærdsmæssigt sikkerhedsnet, der mindsker eksistentiel utryghed, samt en god balance mellem arbejde og fritid. Forskerne lægger vægt på varig livstilfredshed frem for kortvarig eufori.}
  \item \textit{Henrik Madsen peger på voksende ensomhed blandt unge og enlige samt et højt forbrug af antidepressiv medicin og mange stresssygemeldinger. Han fremhæver desuden, at Janteloven skaber konformitet, og at velfærdsfortællingen skaber et lykkepres, hvor mistrivsel opleves som en personlig fiasko.}
\end{enumerate}

\textbf{Tekst B 阅读核对答案}
\begin{enumerate}
  \item \da{Forkert}（他认为官方统计数据与许多公民的实际体验之间存在巨大鸿沟）
  \item \da{Voksende ensomhed (blandt unge og enlige) samt et højt forbrug af antidepressiv medicin og stresssygemeldinger.}
  \item \da{Fordi alle ydre rammer allerede er på plads, så personlig mistrivsel let opleves som en individuel fiasko.}
\end{enumerate}

\textbf{练习 A 答案}
\begin{enumerate}
  \item \da{Hvor}（对比从属连词，放句首引导对立面，主句随即倒装）
  \item \da{anlægger}（固定搭配：\textit{anlægge et kritisk perspektiv} 采取批判视角）
  \item \da{uenighed}（固定衔接：\textit{På trods af denne uenighed...} 尽管存在这一分歧）
  \item \da{modsætning}（固定搭配：\textit{I modsætning til...} 与……相反）
  \item \da{Sammenholdes}（省略 \textit{hvis} 的条件被动倒装：若将两者放在一起对比）
  \item \da{anfægte}（动词原形：质疑、挑战某观点或测算方式）
  \item \da{om}（让步连词短语：\textit{om end} 尽管、虽然）
  \item \da{nedtone}（动词原形：淡化、轻描淡写）
\end{enumerate}

\textbf{练习 B 答案（参考）}
\begin{enumerate}
  \item \textit{Hvor Kilde A understreger institutionel tillid, peger Kilde B på kulturel konformitet og Jantelov.}（句首 \textit{Hvor} 从句整体占第 1 位，主句动词 \textit{peger} 必须 V2 倒装）
  \item \textit{På trods af deres uenighed om målingerne deler de to forfattere en bekymring for stress.}（去掉逗号并实行 V2 倒装；同时注意搭配为 \textit{dele en bekymring / anlægge et kritisk blik}）
  \item \textit{Hvor Kilde A fremhæver, at gratis velfærd fjerner utryghed, anlægger Kilde B et markant mere kritisk perspektiv, idet perfekte rammer får mistrivsel til at føles som en individuel fiasko.}（交叉对比句型精确合成）
  \item \textit{I modsætning til Kilde A's tolkning anfægter Kilde B selve målemetoden, idet spørgeskemaer måske snarere måler pligtfølelse end ægte trivsel.}（聚焦方法论分歧）
  \item \textit{På trods af denne uenighed deler de to tekster den opfattelse, at psykisk sårbarhed fortsat udgør en alvorlig udfordring.}（精准提炼隐性共识）
  \item \textit{Sammenholdes de to kilder, bliver det tydeligt, at materiel tryghed ikke automatisk fjerner menneskelig ensomhed.}（对比收束句型）
\end{enumerate}

\textbf{练习 C 答案}
\begin{enumerate}
  \item \textit{Hvor Kilde A fremhæver høj institutionel tillid, anlægger Kilde B et mere kritisk perspektiv.}
  \item \textit{På trods af denne uenighed deler de to tekster den opfattelse, at psykisk sårbarhed er reel.}
  \item \textit{I modsætning til Kilde A anfægter kritikeren, hvorvidt spørgeskemaer virkelig kan måle ægte lykke.}
  \item \textit{Det sociale sikkerhedsnet fjerner meget utryghed, men kan ikke automatisk fjerne ensomheden.}
  \item \textit{Sammenholdes de to kilder, tegner der sig et langt mere nuanceret billede.}
  \item \textit{Bag den perfekte velfærdsfacade gemmer der sig ofte et usynligt kulturelt konformitetspres.}
\end{enumerate}

\textbf{PD3 / Studieprøven 高分范文（约 290 词）}\\
\textit{\textbf{Mellem trygge rammer og indre pres: En syntese af det danske lykkeparadoks}}

\textit{Danmarks mangeårige placering i toppen af internationale lykkemålinger rejser et centralt spørgsmål om, hvad der egentlig skaber et godt liv i et moderne samfund. \dabold{Hvor} Kilde A (Institut for Nordisk Lykkeforskning) \dabold{fremhæver} de strukturelle styrker bag de høje tal, \dabold{anlægger} kultursociologen Henrik Madsen i Kilde B \dabold{et markant mere kritisk perspektiv}, idet han fokuserer på den voksende kløft mellem officiel statistik og menneskers levede virkelighed.}

\textit{Kilde A forklarer den høje \da{livstilfredshed} ud fra høj \da{institutionel tillid}, et finmasket socialt \da{sikkerhedsnet} og en sund balance mellem arbejde og fritid, som tilsammen fjerner meget af den eksistentielle utryghed i hverdagen. I modsætning hertil peger Kilde B på det skandinaviske \da{paradoks}: Trods denne materielle tryghed kæmper stadig flere unge og enlige danskere med \da{ensomhed}, langvarige stresssygemeldinger og et bemærkelsesværdigt højt \da{forbrug} af \da{antidepressiv} medicin. Ifølge Kilde B skaber lighedsidealet og Janteloven en kulturel \da{konformitet}, der forstærker et usynligt \da{lykkepres}, så individuel mistrivsel hurtigt opleves som et skamfuldt personligt nederlag.}

\textit{Uenigheden gælder også selve datagrundlaget. Mens Kilde A skelner nøje mellem kortvarig eufori og stabil tilfredshed for at forsvare målingernes værdi, \da{anfægter} Kilde B metodens gyldighed ved at spørge, om respondenter i virkeligheden blot udtrykker taknemmelig pligtfølelse frem for ægte velbefindende. \dabold{På trods af denne uenighed deler de to tekster den opfattelse, at} det danske samfund langt fra er frit for psykisk sårbarhed; selv Kilde A \da{nedtoner} ikke, at stress fortsat rammer bestemte befolkningsgrupper.}

\textit{\dabold{Sammenholdes} de to perspektiver, tegner der sig et langt mere nuanceret billede. Et robust velfærdssystem kan skabe uvurderlig ydre tryghed og social frihed, men det garanterer ikke automatisk dyb menneskelig forbundethed eller mental \da{trivsel}. Netop derfor bør fremtidens \da{lykkeforskning} ikke nøjes med at måle institutionelle rammer, men også turde undersøge de kulturelle forventninger og den stille ensomhed bag facaden.}

\small\textbf{名师点评：}这篇约 290 词的 C1 \textit{Syntese} 范文是 Studieprøven 考场的绝佳样本：全篇完全摒弃“A 一段、B 一段”的低阶摘要法，而是围绕“结构成因 vs 心理压力”、“数据方法论争议”与“双方共同承认的脆弱性”三大主题线让两源材料同框对话。精准调用了 \textit{Hvor Kilde A fremhæver..., anlægger Kilde B...}、\textit{I modsætning hertil...}、\textit{På trods af denne uenighed deler de to tekster den opfattelse, at...} 及 \textit{Sammenholdes de to perspektiver...} 等高阶综合公式，收尾独立评判老练深刻，具备冲击顶档满分的说服力。
\end{answerbox}

\ncdlesson{16}{Sproget som hjem: Tanker om at leve på to sprog}{以语言为家：关于双语生活的哲思 [终章]}
{全五册（104 课）总收官 · DU3 Modul 5 (PD3) $\rightarrow$ Modul 6 (Studieprøven, B2 $\rightarrow$ C1)｜对标 NCE 4 终章 \textit{Refleksivt Essay}}
{\textbf{哲理散文体议论文修辞与全五册知识谱系总览}（双重视角 \textit{tvesynet} · 可译性边界 \textit{oversættelighed} · 母语与心语 \textit{modersmål vs hjertesprog} · 细微之境 \textit{nuancernes rige}）· 全五册 104 课语音—句法—修辞螺旋总收官}
{欢迎抵达《新概念丹麦语 1--5》全套 104 课的终章！从第一册第一课小心翼翼问出 \textit{Undskyld, er det plads 14?}，到今天能够自如研读安徒生潜台词、历史政论与学术综述，丹麦语对你而言已不再只是一张需要攻克的考卷，而逐渐化作第二重思考与感知世界的居所。本课通过一篇充满文学温度的双语哲思散文与读者回信，带你走入\textbf{“细微之境（\textit{nuancernes rige}）”}，完成全书总梳理与终极长文收官！}

% ---------------------------------------------------------------------
\begin{foerduleaeser}
带着下面两个宏观问题阅读 Tekst A 和 Tekst B：
\begin{enumerate}
  \item \textit{Hvordan beskriver forfatteren i Tekst A overgangen fra at bruge et nyt sprog som et ”redskab” til at opleve det som et ”hjem”, og hvad ligger der i begrebet ”tvesynet”?}（在 Tekst A 中，作者如何描述一门新语言从“工具”转变为“家园”的过程？“双重视角（tvesynet）”又蕴含着什么意义？）
  \item \textit{Hvilken personlig erfaring med ”mellemrummet” mellem to sprog deler Maja i Tekst B, og hvornår føles dansk ifølge hende som et ægte ”hjertesprog”?}（在 Tekst B 中，玛雅分享了关于两种语言之间“间隙”的什么个人体会？根据她的说法，丹麦语在什么时候会真正变成一种“心语”？）
\end{enumerate}
\end{foerduleaeser}

% ---------------------------------------------------------------------
\ncdsubsec{二、Tekster（哲理散文随笔与双语者回信双语篇）}

\subsection*{Tekst A：\textit{Sproget som hjem: Når to verdener mødes i én stemme}（哲理散文议论文 · 约 360 词）}

\begin{tipbox}{语篇说明}
本篇以兼顾散文诗意与哲学思辨的语调收官全书。请细细品味\textbf{\color{NCDNavy}深蓝色标记的核心隐喻与语言哲学词汇}（\textit{redskab, modersmål, hjertesprog, tvesyn, oversættelighed, nuancernes rige}），以及\textbf{\color{NCDRed}深红色标记的排比对照与哲思收束句型}。
\end{tipbox}

I begyndelsen er et nyt sprog sjældent andet end et praktisk \da{redskab} til at klare hverdagens gøremål. Man lærer at bestille en kop kaffe, spørge om vej på stationen og kæmpe sig igennem sine første \da{famlende} sætninger med en blanding af stolthed og sårbarhed. Længe føles de fremmede ord som lånt tøj, der sidder lidt skævt om skuldrene. Men en dag sker der noget næsten umærkeligt: Man opdager, at man ikke længere blot oversætter sine tanker udefra, men begynder at tøve, drømme og finde præcise billeder direkte på det nye sprog. I det øjeblik forvandles grammatikken langsomt fra en tør regelbog til et levende indre hjem.

At leve på to sprog betyder dog ikke, at det ene skal fortrænge det andet. Snarere opstår der en særlig balance mellem ens \da{modersmål} og det tillærte \da{hjertesprog}. Hvor modersmålet bærer barndommens umiddelbare sansninger og dybeste instinkter, åbner det andet sprog et rum for fornyelse, hvor man kan formulere sig med ny frihed og skærpet opmærksomhed. Denne dobbelte erfaring skaber det, man kunne kalde \da{tvesynet}: evnen til at se den samme virkelighed gennem to sproglige linser på én gang, så hvert sprog kaster lys over det, som det andet lader stå \da{uudtalt}.

Netop i spændingsfeltet mellem to kulturer opdager man spørgsmålet om \da{oversættelighed}. Hver gang man støder på et dansk begreb som \textit{tryghed}, \textit{dannelse} eller \textit{overskud} --- eller et kinesisk ordsprog, hvis klang modstår enhver ordret gengivelse --- træder man ind i \da{nuancernes rige}. \dabold{Snarere end at se} disse sproglige grænser som et nederlag bør man betragte dem som en invitation til at lytte dybere og tålmodigere. Det \da{uoversættelige} minder os nemlig om, at forskellige samfund har lagt mærke til forskellige fine skift i menneskelivet og bygget ord netop dér, hvor et andet sprog tier.

\dabold{At finde hjem i et nyt sprog kræver} derfor ikke, at man visker sin accent eller sin fortid ud. Tværtimod \dabold{beriges sproget netop af} de mange erfaringer og minder, man bærer med sig over grænser. Når to verdener mødes i én stemme, bliver hvert ord en bro over afstand og stille misforståelse --- og læreren, der engang stod som fremmed på perronen, står nu tilbage som borger i begge sprogs \da{rige}.

\begin{transbox}
\small
\textbf{\color{NCDTeal}[Tekst A 中文精译]}\\
起初，一门新语言往往不过是一件用来应付日常事务的实用工具。人们学习如何点一杯咖啡、在车站问路，并带着一丝自豪与脆弱，磕磕绊绊地吐出最初那些尚不熟练的句子。在很长一段时间里，这些陌生的词句就像借来的衣服，穿在肩上总显得有些不合身。然而有一天，某种几乎难以察觉的变化发生了：你发现自己不再只是从外部把念头生硬地翻译过来，而是开始直接用这门新语言去犹疑、做梦，并找到精确的意象。就在那一刻，语法慢慢从一本枯燥的规则手册，蜕变成了一座鲜活的内在家园。

然而，在两种语言中生活并不意味着一种语言必须排挤另一种语言。相反，人们在自己的母语与后天习得的心语之间，会建立起一种独特的平衡。如果说母语承载着童年最直接的感官记忆与最深层的直觉，那么第二语言则开辟了一片重塑自我的空间，让人得以带着新的自由与敏锐的自觉去表达。这种双重经验孕育出了所谓的“双重视角”：能够同时透过两面语言的透镜观照同一片现实，从而让一种语言照亮另一种语言未曾言说的幽微之处。

恰恰是在两种文化的张力地带，人们遇见了“可译性”的问题。每当遇到诸如 \textit{tryghed}（安稳信任）、\textit{dannelse}（全人教养）或 \textit{overskud}（从容余裕）这样的丹麦语核心词——或者某句音韵抗拒任何字面直译的中文谚语时——你便踏入了“细微之境”。与其把这些语言边界视作一种挫败，不如将它们看作更深沉、更耐心地倾听的邀约。因为那些“不可译之物”正提醒着我们：不同的社会各自留意到了人类生命中不同的微妙转向，并恰恰在另一门语言保持沉默的地方筑起了语词。

因此，在一门新语言中找到家的感觉，绝不要求人抹去自己的口音或过去。恰恰相反，语言正是因人们跨越边界带来的丰厚经验与记忆而变得更加丰盈。当两个世界在一个声音中相遇，每一个词都成为跨越距离与静默误解的桥梁——而那位曾经独自站在站台上的异乡旅人，如今已成了两座语言王国共同的公民。
\end{transbox}

\subsection*{Tekst B：\textit{Svarbrev: ”Jeg bor stadig i mellemrummet --- og det er der, stemmen lever”}（读者回信 · 约 155 词）}

\textbf{阅读任务（理解与核对）：}\textit{Læs Maja Linds personlige svarbrev til essayet om at leve på to sprog, og besvar spørgsmålene nedenfor.}

\begin{emailbox}{读者回信｜Fra: Maja Lind (tidligere kursist og litterær oversætter) \quad Emne: Sv: Sproget som hjem}
Tak for et dybt rørende essay, som jeg fuldt ud genkender fra mit eget liv efter ti år i Danmark. I begyndelsen troede jeg, at målet var at tale helt fejlfrit, så ingen kunne høre, hvor jeg kom fra. Alligevel opdagede jeg hurtigt, at livet mellem to sprog også \da{rummer} et stille \da{mellemrum} --- en daglig træthed over altid at skulle oversætte og forklare sig selv fra én kulturel logik ind i en anden.

Netop derfor sætter jeg så stor pris på tanken om \da{tvesynet}. Når jeg arbejder med tekster mellem dansk og mit \da{modersmål}, møder jeg ofte sætninger, der nægter at lade sig \da{indfange} i en ordbog. Men det er faktisk i disse ufuldkomne øjeblikke, at mit danske \da{særpræg} bliver en styrke. Dansk blev først et ægte \da{hjertesprog} for mig den dag, jeg holdt op med at jage perfektion og turde bruge sprogets små nuancer til at udtrykke tvivl, humor og ægte nærvær.
\end{emailbox}

\textbf{Tekst B 阅读核对（Rigtigt / Forkert / Kort svar，答案见第六部分）：}
\begin{enumerate}
  \item \textit{Maja mener fra starten til slutningen, at det vigtigste mål er at tale dansk helt uden accent.} \underline{\hspace{2.2cm}}
  \item \textit{Hvilken udfordring beskriver Maja i det daglige ”mellemrum” mellem to sprog?} \underline{\hspace{2.2cm}}
  \item \textit{Hvornår blev dansk ifølge Maja først et ægte ”hjertesprog” for hende?} \underline{\hspace{2.2cm}}
\end{enumerate}

\begin{tipbox}{本课发音与语调要点（Udtale \& Prosodi）}
\begin{itemize}
  \item \textbf{语言哲学核心词与长复合词}：\textit{et tvesyn} \ipa{ˈtˢveːˌsyːˀn}（后半截 \textit{-syn} 带 \textit{stød}）；\textit{oversættelighed} \ipa{ˌʌʊ̯ɐˈsɛdəliˌheðˀ}（重音在 \textit{-sæt-}，\textit{lig} 中 \textit{g} 不发音）；\textit{et hjertesprog} \ipa{ˈjɛɐ̯d̥əˌsbʁʌʊ̯ˀ}（词首 \textit{hj-} 中 \textit{h} 不发音！\textit{sprog} 带软滑音与 \textit{stød}）；\textit{et modersmål} \ipa{ˈmoːðɐsˌmɔːˀl}（软 \textit{d} + 尾音 \textit{-mål} 带 \textit{stød}）。
  \item \textbf{文学散文细微词汇}：\textit{en nuance} \ipa{nyˈɑŋsə}（法语借词读音，重音在 \textit{-an-}）；\textit{uudtalt} \ipa{ˈuːˌuðtˢæːˀld}（双重 \textit{u} 连写分读清晰，\textit{ud-} 为软 \textit{d}）；\textit{et særpræg} \ipa{ˈsɛɐ̯ˌpʰʁɛːˀj}（\textit{-præg} 尾音滑化带 \textit{stød}）；\textit{famlende} \ipa{ˈfɑmlənə}。
  \item \textbf{全五册语音总复盘提示}：朗读哲思散文时，不再刻意重读每个语法词，而是让轻词（\textit{et, en, at, som, der}）自然归附于实词重音节之上，并用适度的长句顿挫体现丹麦语独有的沉静语流。
\end{itemize}
\end{tipbox}

\begin{transbox}
\small
\textbf{\color{NCDTeal}[Tekst B 中文对照]}\\
\textbf{读者回信（玛雅·林德，前语言班学员暨文学译者｜回复：以语言为家）：}感谢这篇深深打动人心的随笔，我在丹麦生活十年后的亲身经历对此完全感同身受。起初我以为，学习的目标是说出毫无瑕疵的丹麦语，让任何人都听不出我来自哪里。然而我很快发现，在两种语言之间生活也包含着一片静默的“间隙”——一种日复一日、总要将自己从一套文化逻辑翻译并解释进另一套文化逻辑的疲惫感。

正因如此，我格外珍视“双重视角”这个想法。当我在丹麦语与我的母语之间处理文本时，常常遇到那些拒绝被装进词典的句子。但恰恰是在这些并不完美的时刻里，我独特的语言印记反而成了一种力量。直到我不再追逐完美、敢于运用这门语言的细微色调去表达疑虑、幽默与真切的共在的那一天，丹麦语才真正成了我的“心语”。
\end{transbox}

% ---------------------------------------------------------------------
\ncdsubsec{三、Gloser（核心生词表）}

\begin{tipbox}{词表重点}
注意本课中大量的\textbf{中性复合抽象名词}：\textit{et tvesyn, et modersmål, et hjertesprog, et redskab, et mellemrum, et særpræg} 全部使用冠词 \textit{et}（注意 \textit{et mellemrum} 复数仍为 \textit{mellemrum}，定指复数双写 \textit{m} 为 \textit{mellemrummene}）。另外，形容词 \textit{uudtalt}（未言明的）与动词 \textit{at berige}（使丰富、滋养）是极富文学质感的 C1 散文高分用词。
\end{tipbox}

\begin{glosertable}
名词 & et tvesyn & \textit{(tvesynet, kun ental)} & 双重视角，双语透镜 \\
名词 & en oversættelighed & \textit{(oversætteligheden, kun ental)} & 可译性，可翻译程度 \\
名词 & et modersmål & \textit{(modersmålet, modersmål, modersmålene)} & 母语 \\
名词 & et hjertesprog & \textit{(hjertesproget, hjertesprog, hjertesprogene)} & 心语（承载深层情感的语言） \\
名词 & en nuance & \textit{(nuancen, nuancer, nuancerne)} & 细微差别，色调，神韵 \\
名词 & et rige & \textit{(riget, riger, rigerne)} & 王国，领域，境界 \\
名词 & et redskab & \textit{(redskabet, redskaber, redskaberne)} & 工具，手段 \\
名词 & et mellemrum & \textit{(mellemrummet, mellemrum, mellemrummene)} & 间隙，过渡地带 \\
名词 & et særpræg & \textit{(særpræget, særpræg, særprægene)} & 独特印记，鲜明特征 \\
名词 & en identitet & \textit{(identiteten, identiteter, identiteterne)} & 身份认同 \\
\midrule
动词 [har] & at rumme & \textit{(rummer, rummede, har rummet)} & 容纳，包含，蕴藏 \\
动词 [har] & at indfange & \textit{(indfanger, indfangede, har indfanget)} & 捕捉，精准表达 \\
动词 [har] & at overskride & \textit{(overskrider, overskred, har overskredet)} & 跨越，超越（边界） \\
动词 [har] & at berige & \textit{(beriger, berigede, har beriget)} & 使丰盈，滋养，丰富 \\
\midrule
形容词 & famlende & \textit{(famlende, 不变)} & 摸索的，磕磕绊绊的 \\
形容词 & uudtalt & \textit{(uudtalt, uudtalte)} & 未曾言明的，默契留白的 \\
形容词 & uoversættelig & \textit{(uoversætteligt, uoversættelige)} & 不可译的 \\
形容词 & tværsproglig & \textit{(tværsprogligt, tværsproglige)} & 跨语言的 \\
\end{glosertable}

% ---------------------------------------------------------------------
\ncdsubsec{四、Grammatik \& Stilistik（哲理散文修辞与 NCD 全五册 104 课知识谱系收官全景）}

\begin{grammatikbox}{修辞支柱 1：散文体议论文（\textit{Det refleksive essay}）的四大C1美学技巧}
C1 级别的哲理散文议论文不同于刻板的汇报，它需要在严密逻辑之外拥有文字的温度、画面感与概念张力：
\begin{enumerate}
  \item \textbf{具体意象引出抽象命题}：用 \textit{lånt tøj, et praktisk redskab, et indre hjem, en bro over afstand} 等具象隐喻锚定“双语认同”这一抽象主题，切忌空喊口号。
  \item \textbf{概念二分与对偶对比（\textit{Hvor A..., B...} / \textit{fra A til B}）}：将 \textit{redskab}（工具）对照 \textit{hjem}（家园），将 \textit{modersmål} 对照 \textit{hjertesprog}，在字里行间形成对称张力（如 \textit{Hvor modersmålet bærer..., åbner det andet sprog...}）。
  \item \textbf{“与其……不如……”重构视角}：用 \textit{\dabold{Snarere end at} se disse grænser som et nederlag bør man betragte dem som...} 巧妙翻转看似消极的体验，赋予文章哲学深度。
  \item \textbf{不定式名词化主语增强金句质感}：如 \textit{\dabold{At finde hjem i et nyt sprog kræver} ikke, at...} 或 \textit{\dabold{At leve på to sprog betyder} ikke, at...}，让哲理结论显得凝练而从容。
\end{enumerate}
\end{grammatikbox}

\begin{grammatikbox}{终章总览：NCD 全五册 104 课知识谱系收官全景（Fra \textit{Plads 14} til \textit{Nuancernes Rige}）}
恭喜你走完从《新概念丹麦语 1》到《新概念丹麦语 5》的全部 \textbf{104 课}！下表为你浓缩整套五册螺旋式课程从 A1 起步至 C1 精研的完整能力跃迁图谱：

\vspace{0.3em}
\renewcommand{\arraystretch}{1.24}
{\small
\begin{tabularx}{\linewidth}{@{} >{\bfseries\raggedright\arraybackslash}p{2.3cm} >{\raggedright\arraybackslash}p{2.9cm} >{\raggedright\arraybackslash}p{5.1cm} >{\raggedright\arraybackslash}X @{}}
\toprule
\textbf{分册与篇幅} & \textbf{对标级别与课数} & \textbf{核心语法与语篇引擎} & \textbf{代表性里程碑突破} \\
\midrule
\da{NCD 1}\newline 起步篇 & A1 $\rightarrow$ A2-\newline （DU3 Modul 1 · 24 课） & 软 d / 吞音与 \textit{stød} 感知；\textit{en/et}；主句 V2 倒装；中心状语位置；基础时态与问答 & 从 \textit{Undskyld, er det plads 14?} 起步，掌握丹麦语最底层的语音肌肉与句法骨架。 \\
\midrule
\da{NCD 2}\newline 进阶篇 & A2\newline （DU3 Modul 2 · 24 课） & 从句语序（\textit{ikke} 前移）；完成时 vs 过去时切换；\textit{da/når, som/der, sin/sit/sine} & 实现从孤立造句到 70--180 词连贯邮件叙事、信息差口语互动与多人物对照阅读。 \\
\midrule
\da{NCD 3}\newline 社会篇 & B1\newline （DU3 Modul 3 · 20 课） & \textit{-s} 被动 vs \textit{blive} 被动；形式宾语 \textit{det}；省略 \textit{hvis} 条件倒装；长复合词构词 & 读懂丹麦职场扁平文化、福利制度、\textit{Jantelov} 与数字化社会，产出 250--300 词说明文。 \\
\midrule
\da{NCD 4}\newline 思辨篇 & B1+ $\rightarrow$ B2-\newline （DU3 Modul 4 · 20 课） & 灵魂情态小品词（\textit{jo, da, nok, vist}）；让步与虚拟推演；动词名词化；设问与分裂强调 & 练就礼貌而不失锋芒的公共辩论力，写出 350--400 词报刊读者来信（\textit{Debatindlæg}）。 \\
\midrule
\dabold{NCD 5}\newline 精研篇 & B2 $\rightarrow$ C1\newline （PD3 $\rightarrow$ Studieprøven · 16 课） & 数据图表三步评论；低调反讽（\textit{Underdrivelse}）；引述光谱；文化不可译词；\textit{Syntese} 多源综合与文学散文 & 彻底跨越 C1 门槛，在数据的理性、反讽的分寸、伦理的深度与双语哲思之间游刃有余。 \\
\bottomrule
\end{tabularx}
}
\end{grammatikbox}

\begin{grammatikbox}{从 B2 迈向 C1 终身双语阅读与表达的四个“不完美自由”（Stilistik）}
\begin{enumerate}
  \item \textbf{\ding{172} 学会保护\textit{nuancernes rige}，不再迷信逐词硬译}：遇到 \textit{tryghed, dannelse, overskud, tvesyn} 时，用丹麦语语境去体会其边界，而不是用中文单字强行锁死含义。
  \item \textbf{\ding{173} 区分散文体与报刊体，善用节制留白}：散文随笔允许使用 \textit{snarere, næsten umærkeligt, stille mellemrum} 等富于呼吸感的修饰，但主干仍需坚守 V2 倒装与清晰主谓。
  \item \textbf{\ding{174} 区分 \textit{forandre}（改变）、\textit{overskride}（跨越）与 \textit{berige}（使丰富）的及物搭配}：表示“因某物而变得丰盈”在被动语态中固定用 \textit{beriges \da{af}}（而非 \textit{med}）。
  \item \textbf{\ding{175} 接纳双语者独有的声调与印记（\textit{sprogligt særpræg}）}：正如 Tekst B 所言，真正的语言自信始于停止机械模仿零口音，转而用精准的词汇和真诚的思考发声。
\end{enumerate}
\end{grammatikbox}

% ---------------------------------------------------------------------
\ncdsubsec{五、Mønsterøvelser（散文哲思句型与全册语汇操练）}

\textbf{练习 A：双语哲思与散文修辞核心词汇精准填空（每空填入最准确的 1 个丹麦语词）}
\begin{enumerate}
  \item I begyndelsen fungerer et nyt sprog ofte blot som et praktisk \underline{\hspace{2.2cm}} i hverdagen.
  \item Når man ikke længere oversætter udefra, bliver det tillærte sprog et egentligt \underline{\hspace{2.2cm}}.
  \item Begrebet \underline{\hspace{2.2cm}} beskriver evnen til at se verden gennem to sproglige linser på én gang.
  \item Når ord modstår enhver ordret gengivelse, træder man ind i \underline{\hspace{2.2cm}} rige.
  \item Det uoversættelige minder os om det, som et andet sprog lader stå \underline{\hspace{2.2cm}}.
  \item Det nye sprog \underline{\hspace{2.2cm}} netop af de mange erfaringer, man bærer med sig over grænser.
  \item Hvert ord bliver en \underline{\hspace{2.2cm}} over kulturel afstand og stille misforståelser.
  \item I svære sætninger findes der nuancer, som nægter at lade sig \underline{\hspace{2.2cm}} i en almindelig ordbog.
\end{enumerate}

\textbf{练习 B：散文体升级手术与语法搭配纠错（修正表达或按要求升级为 C1 散文句）}
\begin{enumerate}
  \item 【搭配与语序纠错】\textit{Snarere end at se grænser som fejl, man bør betragte dem som muligheder.} $\rightarrow$
  \item 【被动介词纠错】\textit{Det danske sprog bliver beriget med mange tværsproglige stemmer og erfaringer.} $\rightarrow$
  \item 【散文隐喻升级】用 \textit{forvandles fra ... til ...} 改写：\textit{Før var dansk kun et værktøj for mig, men nu føles det som mit eget hjem.} $\rightarrow$
  \item 【对偶对照升级】用 \textit{Hvor modersmålet bærer..., åbner det nye sprog...} 合并：母语保存童年记忆；新语言带来重新表达自我的自由。$\rightarrow$
  \item 【不定式主语凝练】用 \textit{At finde hjem i et nyt sprog kræver ikke, at...} 改写：\textit{Hvis du vil føle dig hjemme på dansk, behøver du ikke glemme din accent.} $\rightarrow$
  \item 【哲思反转升级】用 \textit{Snarere end at... bør man...} 改写：\textit{Vi skal ikke være kede af uoversættelige ord, men lytte bedre til kulturen.} $\rightarrow$
\end{enumerate}

\textbf{练习 C：中译丹（综合调用全五册核心语序与双语哲思词汇）}
\begin{enumerate}
  \item 起初，一门新语言往往不过是一件实用工具。
  \item 生活在两种语言之间并不意味着一种语言必须取代另一种语言。
  \item 母语承载着童年记忆，而心语则开辟了一片新的自由空间。
  \item 不可译的词语邀请我们更加深入地倾听另一种文化。
  \item 在一门新语言中找到家园并不要求我们抹去自己的过去。
  \item 当两个世界在一个声音中相遇时，每一个词都成为一座桥梁。
\end{enumerate}

% ---------------------------------------------------------------------
\ncdsubsec{六、PD3 / Studieprøven 实战微任务与参考答案}

\begin{emailbox}{全五册收官暨 Studieprøven 终极长散文写作任务（250--300 ord）}
\textbf{Afsluttende opgave}: Tillykke med at have nået lektion 104 --- den allersidste lektion i \textit{Ny Koncept Dansk 1--5}! Skriv et \textbf{refleksivt essay eller et afslutningsbrev til dig selv på 250--300 ord} om rejsen gennem to sprog. Se tilbage på vejen fra de første famlende sætninger til mødet med nuancernes rige, reflekter over balancen mellem \textit{modersmål} og \textit{hjertesprog}, og beskriv, hvad \textit{tvesynet} betyder for dit fremtidige liv på dansk.
\end{emailbox}

\begin{answerbox}
\textbf{Før du læser 参考回答}
\begin{enumerate}
  \item \textit{I begyndelsen er et nyt sprog blot et praktisk redskab til hverdagen, men når man begynder at tøve, drømme og tænke i præcise billeder direkte på sproget, forvandles det til et indre hjem. ”Tvesynet” betyder evnen til at se virkeligheden gennem to sproglige linser på én gang, så hvert sprog belyser det, det andet lader stå uudtalt.}
  \item \textit{Maja oplever et stille mellemrum, hvor man dagligt bliver træt af hele tiden at skulle forklare sig selv ind i en ny kulturel logik. Dansk blev først et ægte hjertesprog for hende, da hun holdt op med at jage perfektion og turde bruge sprogets små nuancer til at udtrykke tvivl, humor og nærvær.}
\end{enumerate}

\textbf{Tekst B 阅读核对答案}
\begin{enumerate}
  \item \da{Forkert}（她最初曾这样以为，但后来意识到正是放弃对完美的追逐、接纳自己的语言特点，才让丹麦语真正成为心语）
  \item \da{En daglig træthed over altid at skulle forklare sig selv fra én kulturel logik ind i en anden.}
  \item \da{Den dag, hun holdt op med at jage perfektion og turde bruge sprogets små nuancer til at udtrykke tvivl, humor og nærvær.}
\end{enumerate}

\textbf{练习 A 答案}
\begin{enumerate}
  \item \da{redskab}（中性名词：实用工具，搭配不定冠词 \textit{et}）
  \item \da{hjertesprog}（心语，与母语相对的第二情感语言）
  \item \da{tvesynet}（特指概念：双重视角）
  \item \da{nuancernes}（所有格复数定指：\textit{nuancernes rige} 细微之境）
  \item \da{uudtalt}（未曾言明的，中性/单数基本形）
  \item \da{beriges}（被动语态：\textit{beriges af} 因……而变得丰富）
  \item \da{bro}（隐喻名词：\textit{en bro over...} 跨越……的桥梁）
  \item \da{indfange}（动词原形：捕捉、精准传达）
\end{enumerate}

\textbf{练习 B 答案（参考）}
\begin{enumerate}
  \item \textit{Snarere end at se grænser som fejl bør man betragte dem som muligheder.}（去掉逗号，句首短语触发主句 V2 倒装 \textit{bør man}）
  \item \textit{Det danske sprog bliver beriget af mange tværsproglige stemmer og erfaringer.}（被动句施事者/来源介词固定用 \textit{af}）
  \item \textit{Sproget har forvandlet sig fra et praktisk redskab til et indre hjem for mig.}（用 \textit{forvandles fra... til...} 提升散文画面感）
  \item \textit{Hvor modersmålet bærer barndommens minder, åbner det nye sprog et rum for frihed.}（对称结构 + 句首从句触发 V2 倒装）
  \item \textit{At finde hjem i et nyt sprog kræver ikke, at man glemmer sin egen accent.}（不定式短语作主语，句式凝练庄重）
  \item \textit{Snarere end at ærgre sig over uoversættelige ord bør man lytte dybere til kulturen.}（运用 C1 视角翻转框架）
\end{enumerate}

\textbf{练习 C 答案}
\begin{enumerate}
  \item \textit{I begyndelsen er et nyt sprog sjældent andet end et praktisk redskab.}
  \item \textit{At leve på to sprog betyder ikke, at det ene skal erstatte det andet.}
  \item \textit{Hvor modersmålet bærer barndommens minder, åbner hjertesproget et nyt rum for frihed.}
  \item \textit{De uoversættelige ord inviterer os til at lytte dybere til en anden kultur.}
  \item \textit{At finde hjem i et nyt sprog kræver ikke, at man sletter sin fortid.}
  \item \textit{Når to verdener mødes i én stemme, bliver hvert eneste ord en bro.}
\end{enumerate}

\textbf{PD3 / Studieprøven 高分范文（约 260 词）}\\
\textit{\textbf{Brev til den, der turde begynde forfra i nuancernes rige}}

\textit{Kære medrejsende --- og kære fremtidige jeg. Når jeg ser tilbage på rejsen fra den allerførste sætning i en dansk togkupé til dette sidste kapitel, står det klart, hvor meget et sprog kan \da{forandre} et menneske indefra. I begyndelsen føltes dansk blot som et stift \da{redskab}: En lang række bløde d-lyde, stramme V2-regler og \da{famlende} forsøg på at forklare helt enkle tanker. Dengang troede jeg fejlagtigt, at målet var en fejlfri overflade.}

\textit{Undervejs lærte jeg imidlertid, at vejen fra overlevelse til fordybelse går gennem en stille forvandling. \dabold{Hvor} mit \da{modersmål} stadig bærer barndommens dufte og de mest umiddelbare følelser, \dabold{er} dansk gradvist vokset til et andet \da{hjertesprog}. Det skete ikke den dag, jeg lærte alle uregelmæssige verber udenad, men den dag jeg forstod ironien bag en underspillet dansk bemærkning, mærkede betydningen af \textit{tryghed} og \textit{dannelse}, og opdagede, at man kan være uenig med både varme og præcision.}

\textit{At leve i \da{mellemrummet} mellem to kulturer er ikke altid let, og nogle erfaringer modstår stadig enhver \da{oversættelighed}. \dabold{Snarere end at} se disse grænser som personlige mangler vælger jeg nu at forstå dem som \da{tvesynets} gave. Netop fordi jeg kender to måder at navngive verden på, kan jeg høre det \da{uudtalte} i begge sprog. Min accent og min historie er ikke fejl, der skal slettes, men det \da{særpræg}, som gør min stemme ægte og genkendelig.}

\textit{Nu slutter de 104 lektioner, men selve samtalen begynder først for alvor. Med ro i stemmen, mod i hjertet og nysgerrighed i blikket træder jeg videre ind i \da{nuancernes rige} --- hjemme i begge verdener.}

\small\textbf{名师点评：}这篇约 260 词的收官长散文，将你从 NCD 1 第一课一路走来的学习足迹凝练为一封写给自己的深情致敬信。全篇兼具散文意境与 C1 句法掌控：长短句错落如呼吸，精准调用 \textit{Hvor... er dansk...}、\textit{Snarere end at...}、抽象不定式主语与哲理意象；既呼应了全套五册从软 d、V2 到低调反讽与文化不可译词的脉络，也宣告了真正的毕业——从“学习一门外语”，走进“用两种语言思考与栖居”的广阔天地。祝贺你圆满完成《新概念丹麦语 1--5》全部 104 课的学习！
\end{answerbox}


\end{document}
